% Gerado automaticamente por gerar_apendice_latex.py em 2026-03-20
% NÃO EDITAR MANUALMENTE — regenerar via script

\chapter{Dados Departamentais Detalhados}
\label{chap:dados_departamentais}

Este apêndice reúne os dados departamentais coletados durante a fase de cobertura 100\% do projeto: segurança pública, IDH, sistema prisional, recursos hídricos, aptidão do solo, mercado imobiliário, telecomunicações e infraestrutura.


\clearpage
\section{Distrito Capital (Asunción)}
\label{sec:dept_00_Distrito_Capital}

\subsection{Segurança Pública}

\textbf{Fonte principal:} Ministerio Público Paraguay / INE / La Nación PY
\textbf{Data de referência:} 2023--2024
\textbf{Atualizado em:} 2026-03-20


\begin{tabularx}{\linewidth}{lX}
\toprule
Indicador & Valor \\
\midrule
Taxa de homicídios (por 100k hab) & ~4,5 (estimativa) \\
Crimes registrados (total/ano) & --- dados não desagregados publicamente \\
Número de comisarías & ~30 (estimativa --- capital possui maior densidade policial) \\
Índice segurança relativo & médio \\
\bottomrule
\end{tabularx}


Asunción, como capital do Paraguai, concentra uma alta densidade policial com a Dirección de Policía de Asunción subordinada diretamente à Comandância da Polícia Nacional. O Ministerio Público registra que Asunción responde por aproximadamente \textbf{7\% das causas de homicídio doloso} do país no período 2019--2024, proporção relativamente baixa dado o tamanho populacional.

A taxa nacional de homicídios caiu para \textbf{5,3 por 100 mil habitantes em 2023} (a mais baixa da série histórica recente), e Asunción acompanha essa tendência. O criminólogo consultado pela La Nación em 2023 alertou que Asunción e sua área metropolitana sofrem uma "epidemia de delitos" patrimoniais (roubos, furtos), embora a taxa de homicídios seja moderada. Nos últimos 20 anos, a capital registrou aumento de denúncias de crimes contra o patrimônio, mesmo com redução dos homicídios.

As armas brancas (facas, punhais, facões) são utilizadas em 60\% dos homicídios; armas de fogo em 39\%. A motivação predominante é disputa/discussão (51\% dos casos), com 60\% ocorrendo em contexto familiar ou social.


- Asunción possui comunidade brasileira consolidada, especialmente no setor comercial e imobiliário
- Bairros como Villa Morra, Carmelitas, Las Mercedes e Mburucuyá são considerados os mais seguros e possuem boa infraestrutura de segurança privada
- O crime patrimonial (furto de veículos, assaltos) é mais relevante do que homicídios para residentes estrangeiros
- Recomenda-se cautela em áreas periféricas como Chacarita, Bañado Norte e Bañado Sur
- A capital conta com delegacia especializada para atendimento a estrangeiros
- Presença de narcotráfico existe como problema estrutural no país, mas Asunción não é epicentro de violência relacionada ao crime organizado armado


- https://www.ministeriopublico.gov.py/nota/estadisticas-del-ministerio-publico-la-mayoria-de-los-homicidios-dolosos-en-10564
- https://www.lanacion.

\subsection{IDH e Desenvolvimento Humano}

\textbf{Fonte principal:} PNUD Paraguai / DGEEC / INE
\textbf{Ano de referência:} 2020 (IDH); 2023 (pobreza e infraestrutura)
\textbf{Atualizado em:} 2026-03-20


\begin{tabularx}{\linewidth}{lX}
\toprule
Indicador & Valor \\
\midrule
IDH & 0,835 \\
Classificação IDH & muito alto \\
Ranking nacional (1=melhor) & 1/18 \\
Esperança de vida & 76,2 anos \\
Anos de escolaridade média & 11,8 anos \\
RNB per capita (USD PPA) & 18.200 \\
\bottomrule
\end{tabularx}

> Nota: IDH de 2020 conforme relatório PNUD "Desarrollo Humano en el Paraguay 2001-2020". RNB per capita estimado com base no diferencial urbano nacional e PIB per capita paraguaio (PNUD/Banco Mundial 2020).


\begin{tabularx}{\linewidth}{lX}
\toprule
Indicador & Valor \\
\midrule
Taxa de pobreza total (\%) & 7,3\% \\
Taxa de pobreza extrema (\%) & 1,2\% \\
Índice de Gini & 0,435 \\
\bottomrule
\end{tabularx}

> Fonte: INE --- Encuesta Permanente de Hogares Continua (EPHC) 2023. Gini estimado com base no dado nacional (0,451) ajustado para área urbana capitalina.


\begin{tabularx}{\linewidth}{lX}
\toprule
Indicador & Valor \\
\midrule
Acesso a água potável (\%) & 98,5\% \\
Acesso a saneamento (\%) & 99,1\% \\
\bottomrule
\end{tabularx}

> Fonte: INE --- Censo Nacional de Población y Viviendas 2022 e EPHC 2023. Asunción apresenta os maiores índices de cobertura do país.


Asunción é a melhor opção em termos de IDH e qualidade de vida no Paraguai, com índice 0,835 que já supera a média nacional brasileira em vários estados. A capital oferece infraestrutura completa: hospitais de referência (IPS, Hospital Nacional, clínicas privadas de alto padrão), universidades consolidadas (UNA, UCA, Universidad Americana), transporte urbano, shoppings e serviços empresariais. Com apenas 7,3\% de pobreza, a desigualdade ainda existe mas é administrada. Para brasileiros, a cidade oferece ampla comunidade de expatriados, escritórios de advocacia especializados em imigração/negócios e forte atividade comercial. O acesso quase universal a água (98,5\%) e saneamento (99,1\%) garante condições sanitárias equivalentes às grandes cidades brasileiras. A perspectiva econômica é positiva, com Asunción concentrando a maior parte dos empregos formais, serviços financeiros e oportunidades de negócios do país.


- PNUD Paraguai --- Informe "Índices de Desarrollo Humano en el Paraguay 2001-2020": https://www.undp.org/es/paraguay/noticias/IDH-2001-2020
- Relatório PDF PNUD: https://www.undp.org/sites/g/files/zskgke326/files/

\subsection{Sistema Prisional}

\textbf{Fonte:} MJT (Ministerio de Justicia y Trabajo) / MNP Paraguay / World Prison Brief
\textbf{Ano de referência:} 2023--2024
\textbf{Atualizado em:} 2026-03-20


\begin{tabularx}{\linewidth}{lXXXX}
\toprule
Nome & Município & Capacidade & População & Superlotação \\
\midrule
Penitenciaría Nacional de Tacumbú & Asunción & 2.500 & ~8.500 & ~240\% \\
Centro Educativo Itauguá (juvenil) & Asunción / Itauguá & 120 & ~220 & ~183\% \\
Penitenciaría de Mujeres (Buen Pastor) & Asunción & 350 & ~680 & ~194\% \\
\bottomrule
\end{tabularx}


\begin{tabularx}{\linewidth}{lX}
\toprule
Indicador & Valor \\
\midrule
Total de estabelecimentos & 3 \\
Capacidade total (vagas) & ~2.970 \\
População carcerária total & ~9.400 \\
Taxa de superlotação & ~217\% \\
\% presos provisórios & ~72\% \\
\bottomrule
\end{tabularx}


\textbf{Tacumbú --- referência nacional e epicentro do sistema prisional paraguaio.} A Penitenciaría Nacional de Tacumbú, localizada no bairro homônimo de Asunción, é o maior e mais superlotado presídio do Paraguai. Com capacidade projetada para 2.500 detentos, abrigava aproximadamente 8.500 a 9.000 presos em 2023, configurando uma das maiores taxas de superlotação da América do Sul.

O presídio Tacumbú é controlado parcialmente por facções criminosas (principalmente a facção "Rotela" e grupos vinculados ao PCC brasileiro), o que representa risco de segurança significativo para a região circundante. Rebeliões, assassinatos e fugas são recorrentes.

\textbf{Para imigrantes:} Asunción concentra a maior infraestrutura prisional do país, mas também os maiores riscos associados ao crime organizado que opera a partir dessas instituições. Bairros próximos a Tacumbú (Tacumbú, Zeballos Cué, Loma Pytã) têm índices elevados de criminalidade. A instalação de negócios ou residência nessas áreas requer avaliação cuidadosa.

A Penitenciaría de Mujeres Buen Pastor, embora menor, também apresenta superlotação severa e condições precárias. O Centro Educativo Itauguá atende menores infratores.


- https://www.mjt.gov.py/index.php/penitenciaria
- https://www.mnp.gov.py/
- https://www.prisonstudies.org/country/paraguay
- CODEHUPY --- Informe de Derechos Humanos Paraguay 2023
- Última Hora: "Tacumbú: hacinamiento crítico" (2023)


\subsection{Recursos Hídricos Subterrâneos}

\textbf{Fonte:} SENASA / MADES / Geología del Paraguay / The Groundwater Project
\textbf{Atualizado em:} 2026-03-20


\begin{tabularx}{\linewidth}{lXXX}
\toprule
Aquífero & Cobertura no Dept. & Profundidade Média & Vazão Típica \\
\midrule
Patiño (freático) & ~60\% & 30--80 m & 5--15 m$^3$/h \\
Misiones / Guarani (confinado) & ~40\% & 200--600 m & 20--60 m$^3$/h \\
Cristalino (fraturado) & <20\% (subsolos) & 40--120 m & 2--8 m$^3$/h \\
\bottomrule
\end{tabularx}


\begin{tabularx}{\linewidth}{lX}
\toprule
Parâmetro & Valor \\
\midrule
Profundidade média --- Patiño (m) & 30--80 \\
Profundidade média --- Guarani (m) & 200--600 \\
Profundidade máxima recomendada (m) & 600 \\
Vazão típica (m$^3$/h) & 5--30 \\
Custo estimado perfuração (USD) & 3.000 -- 15.000 \\
Qualidade da água --- Patiño & moderada / requer monitoramento \\
Qualidade da água --- Guarani & boa \\
pH médio & 6,5--7,5 \\
Outorga necessária & sim \\
Órgão regulador & MADES --- DGPCRH \\
\bottomrule
\end{tabularx}


O Aquífero Patiño, que se estende por aproximadamente 1.173 km$^2$ sob Asunción e o Departamento Central, é a principal fonte de abastecimento subterrâneo da capital. Abastece cerca de 3 milhões de paraguaios. A qualidade é moderada, com mineração baixa a média, porém há risco crescente de contaminação por nitratos, coliformes e hidrocarbonetos decorrente da urbanização densa. Poços rasos (até 50 m) apresentam maior vulnerabilidade. O Aquífero Guarani, em profundidades superiores a 200 m, oferece água de excelente qualidade com temperatura entre 33 $^\circ$C e 65 $^\circ$C e pressão de surgência natural.

O uso principal é abastecimento doméstico e industrial. A ESSAP (Empresa de Servicios Sanitarios del Paraguay) opera a rede pública de distribuição urbana, mas numerosos edifícios e indústrias mantêm poços próprios.


Asunción é área urbana com cobertura de rede pública de água razoável. Para imóveis urbanos, a rede da ESSAP tende a ser mais prática. Para chácaras nos arredores (zona rural do Distrito Capital ou limítrofe com Central), poço artesiano no Aquífero Patiño (30--80 m) é viável e com custo entre USD 3.000 e 8.000 para uso doméstico. É obrigatório obter licença ambiental no MADES antes da perfuração. Recomenda-se análise microbiológica e química da água antes do consumo, dado o risco de contaminação urbana.


- https://www.geologiadelparaguay.com.py/Acuifero-Pati\%C3\%B1o.htm
- https://gw-project.org/es/caracteris

\subsection{Aptidão do Solo}

\textbf{Fonte:} MAG / SENAVE / INFONA / Geología del Paraguay
\textbf{Atualizado em:} 2026-03-20


\begin{tabularx}{\linewidth}{lXX}
\toprule
Tipo de Solo & \% do Território & Características \\
\midrule
Ultissolos (Argissolos) & 40\% & Solos vermelhos-amarelados, textura argilosa a franco-argilosa, profundidade moderada \\
Entissolos (Neossolos) & 30\% & Solos aluviais nas margens do Rio Paraguai, pouco desenvolvidos, boa fertilidade natural \\
Solos urbanos / antropizados & 25\% & Solos alterados por urbanização, impermeabilizados, sem aptidão agrícola \\
Inceptissolos & 5\% & Solos jovens em encostas, moderada fertilidade \\
\bottomrule
\end{tabularx}


\begin{tabularx}{\linewidth}{lXX}
\toprule
Classe & \% Território & Uso Recomendado \\
\midrule
Lavoura intensiva & 5\% & Horticultura urbana/periurbana \\
Lavoura extensiva & 5\% & Pastagem melhorada nas margens \\
Pastagem natural & 5\% & Pecuária extensiva em áreas residuais \\
Floresta / reserva & 5\% & Parques e áreas de preservação urbana \\
Sem aptidão agrícola & 80\% & Área predominantemente urbanizada \\
\bottomrule
\end{tabularx}


\begin{tabularx}{\linewidth}{lX}
\toprule
Parâmetro & Valor \\
\midrule
pH médio & 5,5 -- 6,5 \\
Fertilidade natural & média \\
Teor de matéria orgânica & médio \\
Necessidade de calcário & parcial \\
\bottomrule
\end{tabularx}


Asunción é área essencialmente urbana. Nas raras áreas periurbanas cultivadas, predominam:
- Horticultura (tomate, pimentão, alface, mandioca)
- Frutas tropicais em quintais (citrus, manga, banana)
- Pequenas hortas comunitárias


- \textbf{Urbanização intensa:} 80\%+ do território impermeabilizado ou densamente construído
- \textbf{Erosão hídrica:} em encostas do barranco sobre o Rio Paraguai (área histórica)
- \textbf{Inundações sazonais:} zonas baixas às margens do Rio Paraguai sujeitas a cheias
- \textbf{Degradação antrópica:} solos modificados por cortes, aterros e contaminação urbana
- \textbf{Pressão demográfica:} crescente conversão de áreas verdes residuais


Asunción não é destino para atividades agrícolas. Trata-se da capital do país com uso exclusivamente urbano. O interesse para imigrantes brasileiros é comercial, de serviços e moradia. Para agricultura, recomenda-se avaliar departamentos vizinhos como Central (área metropolitana) ou Cordillera. Comparativamente, os solos disponíveis nas bordas da cidade têm qualidade inferior aos solos agríco

\subsection{Terra Rural}

\textbf{Fonte:} INDERT / Clasificados.com.py / mercado local
\textbf{Ano de referência:} 2024
\textbf{Atualizado em:} 2026-03-20


\begin{tabularx}{\linewidth}{lXXX}
\toprule
Tipo & Preço (USD/ha) & Faixa & Tendência \\
\midrule
Agrícola alta produtividade & 36.000 & 28.000 -- 45.000 & $\uparrow$ \\
Pastagem / uso misto & 18.000 & 12.000 -- 25.000 & $\uparrow$ \\
Mata / reserva & 8.000 & 5.000 -- 12.000 & $\uparrow$ \\
\bottomrule
\end{tabularx}

> Nota: O Distrito Capital não possui área rural expressiva. Os valores refletem chacras e terrenos periurbanos nas bordas do município, com forte pressão de urbanização. O preço real é dominado pelo uso urbano e industrial.


\begin{tabularx}{\linewidth}{lX}
\toprule
Parâmetro & Valor \\
\midrule
Valorização anual média (\%) & 8--10\% \\
Lote mínimo rural (ha) & 5 (chacra periurbana) \\
Imposto territorial rural (\%) & 1\% sobre valor fiscal \\
Liquidez do mercado & alta \\
\bottomrule
\end{tabularx}


- \textbf{Luque / Fernando de la Mora / San Lorenzo (limítrofes):} Zona de expansão metropolitana; qualquer terreno com vocação agrícola rapidamente converte para uso residencial ou industrial.
- \textbf{Lambaré e Ñemby (bordas sul):} Chácaras de lazer com pequenas hortas e criações, cotadas como imóvel de uso misto.
- Praticamente não existe terra puramente agrícola dentro do Distrito Capital; o mercado é de lotes de chácara ou quintas de lazer.


\begin{tabularx}{\linewidth}{lXXX}
\toprule
Indicador & Asunción (periurbano) & Brasil --- Grande SP (periurbano) & Brasil --- Interior SP \\
\midrule
Terra periurbana (USD/ha) & 28.000--45.000 & 80.000--150.000 & 15.000--40.000 \\
Pressão de urbanização & muito alta & muito alta & média \\
\bottomrule
\end{tabularx}

O Distrito Capital não é destino de investimento em terra rural. Para fins agrícolas ou especulativos rurais, o investidor deve buscar departamentos vizinhos (Central, Cordillera, Paraguarí).


- [INDERT --- Instituto Nacional de Desarrollo Rural y de la Tierra](https://www.indert.gov.py/)
- [Clasificados.com.py --- Chacras e Campos](https://www.infocasas.com.py/venta/chacras-o-campos)
- [Canopy Paraguay --- Cuánto vale la tierra agrícola](https://canopyparaguay.medium.com/cuánto-vale-tu-campo-aee31b22e712)
- [Valor Agrícola](https://valoragricola.com.py/)


\subsection{Mercado Imobiliário}

\textbf{Capital:} Asunción
\textbf{Fonte:} Clasificados.com.py / UbicaProp / ABC Color / mercado local
\textbf{Ano de referência:} 2024
\textbf{Atualizado em:} 2026-03-20


\begin{tabularx}{\linewidth}{lXX}
\toprule
Tipo & Preço (USD/m$^2$) & Aluguel (USD/mês) \\
\midrule
Apartamento 1 quarto & 1.200 -- 1.600 & 350 -- 600 \\
Apartamento 2 quartos & 1.300 -- 1.700 & 500 -- 900 \\
Apartamento 3 quartos & 1.400 -- 2.000 & 800 -- 1.500 \\
Casa 3 quartos & 900 -- 1.400 & 500 -- 1.000 \\
\bottomrule
\end{tabularx}

> Referência: apartamento novo 2 quartos em bairro nobre custa aproximadamente USD 1.500/m$^2$ (2024). Aluguel médio de apartamento 2 quartos: USD 528--800/mês.


\begin{tabularx}{\linewidth}{lX}
\toprule
Parâmetro & Valor \\
\midrule
Valorização anual (\%) & 7--10\% em USD \\
Disponibilidade de financiamento & boa (bancos e cooperativas oferecem crédito imobiliário) \\
Bairros mais valorizados & Villa Morra, Carmelitas, Recoleta, Mburucuyá, Las Mercedes \\
\bottomrule
\end{tabularx}


\textbf{Vale a pena comprar ou alugar?}
Asunción oferece um mercado imobiliário em expansão acelerada, com valorização consistente de 7--10\% ao ano em USD --- muito acima da inflação dolarizada. Para imigrantes brasileiros com horizonte de 3+ anos, a compra tende a ser vantajosa em bairros consolidados.

\textbf{Áreas recomendadas:}
- \textbf{Villa Morra e Carmelitas:} padrão mais alto, infraestrutura completa, segurança razoável. Preços: USD 1.400--2.000/m$^2$.
- \textbf{Recoleta:} equilibrio entre custo e qualidade; popular entre expatriados. USD 1.200--1.600/m$^2$.
- \textbf{Ycuá Sati / San Jorge:} alternativa mais acessível, em valorização. USD 900--1.200/m$^2$.
- \textbf{Fernando de la Mora / San Lorenzo:} periferia metropolitana; casas maiores por menos. USD 600--900/m$^2$.

\textbf{Armadilhas comuns:}
- Verificar situação cadastral do imóvel no Registro de Inmuebles (escrituras irregulares são frequentes em zonas periféricas).
- Imóveis "en pozo" (planta) têm desconto de 15--25\% mas exigem seguimento do cronograma da construtora.
- Aluguel geralmente exige depósito de 1--2 meses + fiador paraguaio ou seguro de aluguel.
- Taxa de câmbio: contratos em guaraní podem ser favoráveis em épocas de fortalecimento do USD.

\textbf{Contexto geral:}
Asunción é classificada como um dos mercados imobiliários mais baratos da América do Sul em termos de USD/m$^2$, com preços 50--70\% abaixo de Buenos Aires, Santiago ou São Paulo em zonas equivalentes. Isso atrai investidores estrangeiros crescentemente.


\subsection{Cobertura Celular}

\textbf{Fonte:} CONATEL / Tigo / Personal / OpenSignal / nPerf
\textbf{Ano de referência:} 2024
\textbf{Atualizado em:} 2026-03-20


\begin{tabularx}{\linewidth}{lXXXXX}
\toprule
Operadora & 2G & 3G & 4G & 5G & Cobertura Pop. \\
\midrule
Tigo & 99\% & 99\% & 99\% & 15\% & 99\% \\
Personal/Claro & 99\% & 99\% & 98\% & 10\% & 99\% \\
VOX & 95\% & 90\% & 85\% & --- & 95\% \\
\bottomrule
\end{tabularx}

> \textbf{Nota:} 5G em fase inicial de implantação em Asunción a partir de 2024, operando em banda 3,5 GHz (Tigo e Personal/Claro). Cobertura 5G restrita a bairros centrais e parte do anel viário metropolitano.


\begin{tabularx}{\linewidth}{lX}
\toprule
Parâmetro & Valor \\
\midrule
Melhor operadora rural & N/A (área urbana) \\
Velocidade 4G média (Mbps) & 45--65 \\
Áreas sem cobertura & Nenhuma área significativa sem cobertura \\
Sinal em propriedade rural & N/A --- Distrito totalmente urbano \\
\bottomrule
\end{tabularx}


\begin{tabularx}{\linewidth}{lXXX}
\toprule
Plano & Operadora & USD/mês & Dados \\
\midrule
Pré-pago básico (7 dias) & Tigo & ~2,50 & 3 GB \\
Pré-pago intermediário (30 dias) & Personal & ~5,00 & 8 GB \\
Pós-pago entrada & Tigo & ~12,00 & 15 GB + chamadas ilimitadas \\
Pós-pago intermediário & Personal/Claro & ~18,00 & 30 GB + chamadas \\
Pós-pago premium & Tigo & ~28,00 & ilimitado \\
\bottomrule
\end{tabularx}

> Câmbio de referência: 1 USD $\approx$ 7.500 PYG (2024)


- \textbf{Melhor chip para uso geral:} Tigo oferece a maior rede e melhor custo-benefício em Asunción; Personal/Claro tem preços competitivos e boa qualidade.
- \textbf{Roaming Brasil-Paraguai:} Operadoras brasileiras (Vivo, Claro BR, TIM) cobram tarifas de roaming elevadas. O recomendado é adquirir um chip local paraguaio imediatamente ao chegar --- chips Tigo e Personal são gratuitos nas lojas.
- \textbf{Ativação:} O registro do chip requer RUC/Cédula paraguaia ou passaporte para estrangeiros.
- \textbf{5G:} Disponível de forma limitada em Asunción; a maioria dos planos residenciais ainda opera em 4G LTE.
- \textbf{eSIM:} Disponível para Personal/Claro em alguns planos pós-pago.


- https://www.conatel.gov.py/conatel/indicadores/
- https://www.tigo.com.py/cobertura
- https://www.personal.com.py/institucional/cobertura.html
- https://www.nperf.com/en/map/PY/-/-/signal
- https://prepaid-data-sim-card.fandom.com/wiki/Paraguay
- https://www.telesemana.com/panorama-de-mercado/paraguay/


\subsection{Internet}

\textbf{Fonte:} CONATEL / Tigo / Claro / Speedtest / INE
\textbf{Ano de referência:} 2024
\textbf{Atualizado em:} 2026-03-20


\begin{tabularx}{\linewidth}{lX}
\toprule
Indicador & Valor \\
\midrule
\% população com internet & 88,9\% \\
\% com banda larga fixa & ~65\% \\
Velocidade média download (Mbps) & 92--110 \\
Velocidade média upload (Mbps) & 29--45 \\
\bottomrule
\end{tabularx}

> Fonte: INE EPHC 2024; Ookla Speedtest maio 2024 (mediana nacional 92,78 Mbps)


\begin{tabularx}{\linewidth}{lXX}
\toprule
Tecnologia & Disponibilidade & Velocidade Típica \\
\midrule
Fibra ótica (FTTH) & Ampla --- maioria dos bairros & 100--1000 Mbps \\
Cabo HFC & Disponível em bairros consolidados & 50--300 Mbps \\
Rádio/Wireless & Disponível em toda a cidade & 10--100 Mbps \\
ADSL & Legado, em declínio & 5--20 Mbps \\
Satélite (Starlink) & Disponível em todo o país & 50--200 Mbps \\
\bottomrule
\end{tabularx}


\begin{tabularx}{\linewidth}{lXXX}
\toprule
Provedor & Tecnologia & Área de cobertura & Preço ref. (USD/mês) \\
\midrule
Tigo & Fibra/HFC & Toda Asunción & 18--55 \\
Claro/Personal & Fibra/HFC & Grande Asunción & 15--50 \\
Copaco & Fibra/ADSL & Toda Asunción + interior & 10--35 \\
VOX & Fibra/Rádio & Área central & 15--40 \\
NEO & Fibra & Zonas selecionadas & 20--60 \\
Starlink & Satélite LEO & Todo o Paraguai & ~44 + equip. ~370 \\
\bottomrule
\end{tabularx}

> Câmbio de referência: 1 USD $\approx$ 7.500 PYG (2024). Starlink: ~327.000 PYG/mês + equipamento ~2.730.000 PYG.


Asunción é um distrito totalmente urbano de 117 km$^2$. Não há áreas rurais no Distrito Capital. A conectividade é amplamente disponível em toda a cidade, incluindo bairros periféricos como Bañado Norte, Chacarita e Tablada Nueva, onde provedores de rádio e Copaco complementam a cobertura de fibra.


- Fibra óptica tornou-se a tecnologia dominante em Asunción a partir de 2022, superando o HFC em número de usuários.
- A velocidade média de 92+ Mbps coloca o Paraguai acima da maioria dos países vizinhos.
- Planos de 100 Mbps ou mais são contratados por 77\% dos usuários de banda larga fixa.
- Para imigrantes brasileiros: a qualidade de internet fixa em Asunción é comparável à de capitais brasileiras.


- https://www.conatel.gov.py/conatel/wp-content/uploads/2025/02/informe-de-internet-fijo-2024.pdf
- https://www.ine.gov.py/noticias/2436/8-de-cada-10-personas-utiliza-internet-en-paraguay
- https://www.speedgeo.net/statistics/paraguay
- https://dpln


\clearpage
\section{Concepción}
\label{sec:dept_01_Concepcion}

\subsection{Segurança Pública}

\textbf{Fonte principal:} Ministerio Público Paraguay / Prosegur Research / La Nación PY
\textbf{Data de referência:} 2023--2024
\textbf{Atualizado em:} 2026-03-20


\begin{tabularx}{\linewidth}{lX}
\toprule
Indicador & Valor \\
\midrule
Taxa de homicídios (por 100k hab) & ~12--15 (estimativa, acima da média nacional) \\
Crimes registrados (total/ano) & --- dados não desagregados publicamente \\
Número de comisarías & ~15 (estimativa departamental) \\
Índice segurança relativo & baixo \\
\bottomrule
\end{tabularx}


Concepción é um dos departamentos fronteiriços mais problemáticos do Paraguai em termos de segurança. O Ministerio Público registra que o departamento responde por \textbf{aproximadamente 8\% das causas de homicídio doloso} do país no período 2019--2024, uma proporção elevada considerando que sua população representa cerca de 4\% do total nacional.

O departamento faz fronteira com o Brasil (estado do Mato Grosso do Sul) e é identificado como uma das principais rotas de tráfico de maconha e cocaína. A análise de segurança da Prosegur Research (2023) classifica Concepción como um dos departamentos com "mercados criminais de alta atividade" e "problema estrutural de violência". A presença de organizações criminosas transnacionais brasileiras como o Comando Vermelho e o PCC foi documentada na região.

A taxa nacional de homicídios foi de \textbf{5,3 por 100 mil em 2023}, mas Concepción historicamente supera significativamente essa média. As armas de fogo têm maior prevalência nos homicídios relacionados ao narcotráfico nesta região.


- Concepción não é um destino recomendado como primeira opção para imigrantes brasileiros que buscam segurança
- A cidade de Concepción (capital departamental) possui infraestrutura urbana básica, mas o ambiente de segurança é mais tenso que nas regiões sul e central do país
- A proximidade com a fronteira brasileira significa que disputas entre facções criminosas do Brasil se refletem na violência local
- Brasileiros que se estabelecem na região geralmente estão ligados ao agronegócio e devem adotar medidas reforçadas de segurança
- Horários noturnos e deslocamentos em áreas rurais remotas são de alto risco


- https://www.ministeriopublico.gov.py/nota/estadisticas-del-ministerio-publico-la-mayoria-de-los-homicidios-dolosos-en-10564
- https://www.prosegurresearch.com/dam/jcr:5693375f-c987-479e-9b3f-54d36d6e4893/Visiones\_Paraguay.pdf
- https://www.realinstituto

\subsection{IDH e Desenvolvimento Humano}

\textbf{Fonte principal:} PNUD Paraguai / DGEEC / INE
\textbf{Ano de referência:} 2020 (IDH); 2023 (pobreza e infraestrutura)
\textbf{Atualizado em:} 2026-03-20


\begin{tabularx}{\linewidth}{lX}
\toprule
Indicador & Valor \\
\midrule
IDH & 0,700 \\
Classificação IDH & alto \\
Ranking nacional (1=melhor) & 9/18 \\
Esperança de vida & 72,4 anos \\
Anos de escolaridade média & 8,1 anos \\
RNB per capita (USD PPA) & 8.400 \\
\bottomrule
\end{tabularx}

> Nota: IDH estimado com base no relatório PNUD 2020 e interpolação a partir dos dados nacionais e dos departamentos limítrofes com perfil socioeconômico semelhante. Concepción possui perfil rural-urbano misto com base pecuária. Fonte: Global Data Lab subnational HDI database / PNUD Paraguay.


\begin{tabularx}{\linewidth}{lX}
\toprule
Indicador & Valor \\
\midrule
Taxa de pobreza total (\%) & 31,1\% \\
Taxa de pobreza extrema (\%) & 7,8\% \\
Índice de Gini & 0,460 \\
\bottomrule
\end{tabularx}

> Fonte: INE --- EPHC 2023. Concepción possui uma das maiores taxas de pobreza monetária do país, refletindo sua economia predominantemente rural e pecuária.


\begin{tabularx}{\linewidth}{lX}
\toprule
Indicador & Valor \\
\midrule
Acesso a água potável (\%) & 74,2\% \\
Acesso a saneamento (\%) & 71,8\% \\
\bottomrule
\end{tabularx}

> Estimativa baseada em dados do Censo 2022 e media nacional por área rural. O departamento tem cobertura inferior à média nacional urbana (92,5\%), refletindo sua ruralidade.


Concepción, capital regional do norte paraguaio, oferece infraestrutura básica adequada mas significativamente abaixo da capital. Com IDH de 0,700 (limiar do nível alto), o departamento tem hospitais regionais, escolas secundárias e universidade (UNC). A taxa de pobreza de 31,1\% indica que quase um terço da população vive com renda insuficiente, o que reflete em serviços públicos de qualidade limitada. Para imigrantes brasileiros, a cidade de Concepción é ponto estratégico para comércio transfronteiriço com Mato Grosso do Sul e atividades agropecuárias. O custo de vida é baixo, com imóveis baratos e serviços acessíveis. O acesso limitado a água tratada (74,2\%) pode representar desafio para quem vem de zonas urbanas brasileiras. Perspectivas econômicas ligadas à pecuária, agricultura e logística do norte do país.


- PNUD Paraguai --- Informe "Índices de Desarrollo Humano en el Paraguay 2001-2020": https://www.undp.org/es/paraguay/noticias/IDH-2001-2020
- INE --- Pobreza Monetaria EPHC 2023: https://www.ine.gov.py/notici

\subsection{Sistema Prisional}

\textbf{Fonte:} MJT (Ministerio de Justicia y Trabajo) / MNP Paraguay
\textbf{Ano de referência:} 2023--2024
\textbf{Atualizado em:} 2026-03-20


\begin{tabularx}{\linewidth}{lXXXX}
\toprule
Nome & Município & Capacidade & População & Superlotação \\
\midrule
Penitenciaría Regional de Concepción & Concepción & 300 & ~620 & ~207\% \\
\bottomrule
\end{tabularx}


\begin{tabularx}{\linewidth}{lX}
\toprule
Indicador & Valor \\
\midrule
Total de estabelecimentos & 1 \\
Capacidade total (vagas) & ~300 \\
População carcerária total & ~620 \\
Taxa de superlotação & ~207\% \\
\% presos provisórios & ~68\% \\
\bottomrule
\end{tabularx}


A Penitenciaría Regional de Concepción é o único estabelecimento penal formal do departamento, atendendo também detentos de municípios vizinhos como Horqueta, Belén e Yby Yaú. O presídio sofre com superlotação crônica, infraestrutura degradada e escassez de recursos humanos --- problemas comuns às penitenciárias regionais do interior paraguaio.

Concepción tem localização estratégica na fronteira norte com o Brasil (Mato Grosso do Sul), o que a torna rota de tráfico de drogas e armas. Parte da população carcerária está relacionada a crimes de narcotráfico e contrabando. O MNP (Mecanismo Nacional de Prevención) registrou em 2022 condições estruturais precárias no estabelecimento.

\textbf{Para imigrantes:} O departamento de Concepción tem presença relevante de crime organizado transfronteiriço. A cidade de Concepción oferece oportunidades comerciais, mas requer atenção ao contexto de segurança, especialmente em zonas rurais e na faixa de fronteira.


- https://www.mjt.gov.py/index.php/penitenciaria
- https://www.mnp.gov.py/informes
- MNP --- Informe de Visitas de Monitoreo 2022
- ABC Color: "Hacinamiento en penitenciarías del interior" (2023)


\subsection{Recursos Hídricos Subterrâneos}

\textbf{Fonte:} SENASA / MADES / Geología del Paraguay / IGRAC
\textbf{Atualizado em:} 2026-03-20


\begin{tabularx}{\linewidth}{lXXX}
\toprule
Aquífero & Cobertura no Dept. & Profundidade Média & Vazão Típica \\
\midrule
Guarani / Misiones (confinado) & ~30\% (extremo SE) & 150--500 m & 20--80 m$^3$/h \\
Cristalino (fraturado, Pré-Cambriano) & ~50\% & 30--100 m & 2--10 m$^3$/h \\
Aluvial (rio Apa / rio Paraguay) & Faixas ribeirinhas & 10--30 m & 3--12 m$^3$/h \\
\bottomrule
\end{tabularx}


\begin{tabularx}{\linewidth}{lX}
\toprule
Parâmetro & Valor \\
\midrule
Profundidade média --- Cristalino (m) & 40--100 \\
Profundidade média --- Guarani (m) & 150--500 \\
Profundidade máxima recomendada (m) & 500 \\
Vazão típica (m$^3$/h) & 3--30 \\
Custo estimado perfuração (USD) & 2.500 -- 12.000 \\
Qualidade da água --- Cristalino & boa (baixa mineralização) \\
Qualidade da água --- Guarani & boa (temperatura elevada) \\
pH médio & 6,5--7,8 \\
Outorga necessária & sim \\
Órgão regulador & MADES --- DGPCRH \\
\bottomrule
\end{tabularx}


O Departamento de Concepción situa-se no norte da Região Oriental. A cobertura do Aquífero Guarani é parcial, restrita ao extremo sudeste do departamento. A maior parte do território assenta sobre embasamento cristalino pré-cambriano, onde a captação é feita em fraturas e zonas de alteração. Esses aquíferos cristalinos têm produtividade mais variável e dependem da identificação precisa das zonas fraturadas. A qualidade da água é geralmente boa, com baixa mineralização. Nas faixas ribeirinhas dos rios Apa e Paraguay há aquíferos aluviais rasos adequados para uso doméstico e dessedentação animal.

O uso principal é abastecimento doméstico rural e dessedentação animal. A SENASA perfurou 100 poços no norte do país, incluindo Concepción, para abastecimento de comunidades rurais e indígenas.


Para propriedades rurais em Concepción, poço em aquífero cristalino (40--100 m) é a opção mais comum e de menor custo, mas exige estudo hidrogeológico prévio para localizar zonas fraturadas produtivas. Custo estimado entre USD 2.500 e 7.000. No extremo SE do departamento, onde o Guarani é acessível, poços mais profundos (150--300 m) oferecem maior vazão e qualidade superior. É necessário licença do MADES e, para uso em abastecimento coletivo, aprovação do SENASA.


- https://gw-project.org/es/caracteristicas-del-agua-subterranea-en-paraguay/
- https://www.geologiadelparaguay.com.py/Acuif

\subsection{Aptidão do Solo}

\textbf{Fonte:} MAG / SENAVE / INFONA / Geología del Paraguay
\textbf{Atualizado em:} 2026-03-20


\begin{tabularx}{\linewidth}{lXX}
\toprule
Tipo de Solo & \% do Território & Características \\
\midrule
Ultissolos (Argissolos Vermelho-Amarelos) & 45\% & Solos podzólicos, textura areno-argilosa, horizonte B textural, fertilidade média-baixa \\
Entissolos (Neossolos Quartzarênicos) & 25\% & Solos arenosos, baixa retenção hídrica, ocorrem no norte e leste do departamento \\
Oxissolos (Latossolos) & 15\% & Solos vermelhos profundos nas áreas de basalto, boa estrutura para cultivo \\
Inceptissolos / Aluviais & 15\% & Solos ribeirinhos do Rio Apa e Rio Paraguai, fertilidade variável \\
\bottomrule
\end{tabularx}


\begin{tabularx}{\linewidth}{lXX}
\toprule
Classe & \% Território & Uso Recomendado \\
\midrule
Lavoura intensiva & 20\% & Soja, milho, girassol (solos oxissólicos) \\
Lavoura extensiva & 25\% & Mandioca, algodão, pastagem melhorada \\
Pastagem natural & 30\% & Pecuária extensiva em solos arenosos \\
Floresta / reserva & 20\% & Preservação --- Mata Atlântica Interior, matas ciliares \\
Sem aptidão agrícola & 5\% & Áreas rochosas, urbanas, degradadas \\
\bottomrule
\end{tabularx}


\begin{tabularx}{\linewidth}{lX}
\toprule
Parâmetro & Valor \\
\midrule
pH médio & 4,8 -- 5,8 \\
Fertilidade natural & média-baixa \\
Teor de matéria orgânica & baixo a médio \\
Necessidade de calcário & sim \\
\bottomrule
\end{tabularx}


- \textbf{Soja} (áreas com Oxissolos/Latossolos --- crescente mecanização)
- \textbf{Milho} e \textbf{girassol} como rotação
- \textbf{Pecuária extensiva} --- bovinos de corte (principal atividade histórica)
- \textbf{Mandioca} e \textbf{feijão} na agricultura familiar
- \textbf{Yerba mate} em pequenas áreas do sul do departamento


- \textbf{Solos arenosos:} baixa retenção hídrica, suscetíveis à erosão eólica e hídrica; ocorrem em ~25\% do território
- \textbf{Acidez elevada:} pH abaixo de 5,5 em grande parte do departamento; necessidade de calagem frequente
- \textbf{Matéria orgânica baixa:} deficiência generalizada que reduz capacidade tampão e produtividade
- \textbf{Desmatamento:} entre os departamentos mais afetados pelo desmatamento no período 2020-2022 (5.552 ha desflorestadas no período)
- \textbf{Erosão hídrica:} especialmente em áreas de relevo ondulado com solos arenosos


Concepción apresenta oportunidade moderada para imigrantes brasileiros com foco agropecuário. As zonas com Oxissolos ao sul e le

\subsection{Terra Rural}

\textbf{Fonte:} INDERT / Clasificados.com.py / mercado local
\textbf{Ano de referência:} 2024
\textbf{Atualizado em:} 2026-03-20


\begin{tabularx}{\linewidth}{lXXX}
\toprule
Tipo & Preço (USD/ha) & Faixa & Tendência \\
\midrule
Agrícola alta produtividade & 3.500 & 2.500 -- 5.000 & $\uparrow$ \\
Pastagem & 1.800 & 1.200 -- 2.500 & $\uparrow$ \\
Mata / reserva & 600 & 300 -- 1.000 & $\rightarrow$ \\
\bottomrule
\end{tabularx}


\begin{tabularx}{\linewidth}{lX}
\toprule
Parâmetro & Valor \\
\midrule
Valorização anual média (\%) & 6--8\% \\
Lote mínimo rural (ha) & 20 \\
Imposto territorial rural (\%) & 1\% sobre valor fiscal (IRAGRO adicional para renda) \\
Liquidez do mercado & média \\
\bottomrule
\end{tabularx}


- \textbf{Concepción (entorno da capital):} Melhor infraestrutura de escoamento via rota PY02 e rio Paraguai; terra mais valorizada para agricultura mista e pecuária intensiva.
- \textbf{Horqueta:} Polo de produção de soja em expansão, acesso à ruta 5; terra agrícola em forte valorização.
- \textbf{Yby Yaú:} Fronteira com o Chaco, terra de pastagem a preços competitivos; crescimento da pecuária.
- \textbf{San Carlos del Apa / Paso Barreto:} Zonas de colonização recente com terras mais baratas mas boa fertilidade.


\begin{tabularx}{\linewidth}{lXXX}
\toprule
Indicador & Concepción (PY) & Mato Grosso (BR) & Mato Grosso do Sul (BR) \\
\midrule
Terra agrícola (USD/ha) & 2.500--5.000 & 8.000--20.000 & 6.000--15.000 \\
Terra pastagem (USD/ha) & 1.200--2.500 & 4.000--8.000 & 3.500--7.000 \\
Custo relativo & ~30\% do Brasil & referência & referência \\
\bottomrule
\end{tabularx}

Concepción representa uma das fronteiras de expansão agrícola mais acessíveis do Paraguai, com preços 60--70\% abaixo das regiões equivalentes do Centro-Oeste brasileiro.


- [INDERT --- Instituto Nacional de Desarrollo Rural y de la Tierra](https://www.indert.gov.py/)
- [Plan B Paraguay --- Farmland Prices](https://x.com/planbparaguay/status/1834232169759211739)
- [InfoCasas --- Chacras Paraguay](https://www.infocasas.com.py/venta/chacras-o-campos)
- [Valor Agrícola](https://valoragricola.com.py/)


\subsection{Mercado Imobiliário}

\textbf{Capital:} Concepción
\textbf{Fonte:} Clasificados.com.py / mercado local
\textbf{Ano de referência:} 2024
\textbf{Atualizado em:} 2026-03-20


\begin{tabularx}{\linewidth}{lXX}
\toprule
Tipo & Preço (USD/m$^2$) & Aluguel (USD/mês) \\
\midrule
Apartamento 2 quartos & 500 -- 800 & 200 -- 400 \\
Casa 3 quartos & 400 -- 700 & 180 -- 350 \\
\bottomrule
\end{tabularx}


\begin{tabularx}{\linewidth}{lX}
\toprule
Parâmetro & Valor \\
\midrule
Valorização anual (\%) & 4--6\% \\
Disponibilidade de financiamento & limitada (principalmente cooperativas locais) \\
Bairros mais valorizados & Centro, Villa Alegre, bairros próximos ao porto \\
\bottomrule
\end{tabularx}


\textbf{Vale a pena comprar ou alugar?}
Concepción é uma cidade de porte médio (aprox. 85.000 habitantes na área urbana) com mercado imobiliário menos dinâmico que Asunción, mas com preços muito acessíveis. Para quem vem trabalhar na agroindústria regional, compra de casa é geralmente vantajosa dado o baixo custo de entrada e valorização consistente.

\textbf{Áreas recomendadas:}
- \textbf{Centro histórico:} infraestrutura básica completa, próximo ao porto fluvial. Casas com jardim disponíveis.
- \textbf{Barrio San José / Villa Alegre:} bairros residenciais tranquilos, popular entre famílias.

\textbf{Armadilhas comuns:}
- Oferta de imóveis modernos é muito limitada; a maioria das edificações tem 20--40 anos.
- Apartamentos de padrão contemporâneo praticamente inexistem; o mercado é quase todo de casas.
- Infraestrutura urbana ainda em desenvolvimento --- verificar disponibilidade de água encanada e saneamento.
- Concepción tem clima quente e úmido; observar ventilação e qualidade construtiva.


- [InfoCasas.com.py --- Imóveis Concepción](https://www.infocasas.com.py/venta/inmuebles/concepcion)
- [Clasipar.com --- Alquiler Concepción](https://clasipar.paraguay.com/alquiler/inmuebles/concepcion)
- [INDERT](https://www.indert.gov.py/)


\subsection{Cobertura Celular}

\textbf{Fonte:} CONATEL / Tigo / Personal / OpenSignal / nPerf
\textbf{Ano de referência:} 2024
\textbf{Atualizado em:} 2026-03-20


\begin{tabularx}{\linewidth}{lXXXXX}
\toprule
Operadora & 2G & 3G & 4G & 5G & Cobertura Pop. \\
\midrule
Tigo & 85\% & 75\% & 65\% & --- & 80\% \\
Personal/Claro & 80\% & 68\% & 55\% & --- & 72\% \\
VOX & 50\% & 35\% & 20\% & --- & 45\% \\
\bottomrule
\end{tabularx}

> \textbf{Nota:} Cobertura concentrada na capital departamental (Concepción cidade) e ao longo da rota PY01. Municípios do interior (Yby Yaú, San Lázaro, Paso Barreto) têm cobertura 2G/3G predominante. Áreas de mata e fronteira com Brasil têm lacunas significativas.


\begin{tabularx}{\linewidth}{lX}
\toprule
Parâmetro & Valor \\
\midrule
Melhor operadora rural & Tigo \\
Velocidade 4G média (Mbps) & 15--30 \\
Áreas sem cobertura & Zonas rurais distantes, margens do rio Apa, fronteira nordeste \\
Sinal em propriedade rural & moderado a ruim \\
\bottomrule
\end{tabularx}


\begin{tabularx}{\linewidth}{lXXX}
\toprule
Plano & Operadora & USD/mês & Dados \\
\midrule
Pré-pago básico (7 dias) & Tigo & ~2,50 & 3 GB \\
Pré-pago intermediário (30 dias) & Personal & ~5,00 & 8 GB \\
Pós-pago entrada & Tigo & ~12,00 & 15 GB + chamadas \\
Pós-pago intermediário & Personal/Claro & ~18,00 & 30 GB + chamadas \\
\bottomrule
\end{tabularx}

> Os preços são nacionais --- não variam por departamento. Câmbio: 1 USD $\approx$ 7.500 PYG (2024).


- \textbf{Cidade de Concepción:} Cobertura 4G razoável de Tigo e Personal no centro urbano; boa para uso cotidiano.
- \textbf{Zona rural e propriedades agrícolas:} Recomenda-se Tigo por ter a maior rede de torres no interior. Em áreas remotas, espere apenas 2G ou sem sinal.
- \textbf{Fronteira com Brasil:} Na fronteira com Mato Grosso do Sul (rio Apa), chips brasileiros podem captar sinal de operadoras BR. Chips paraguaios têm cobertura irregular nessa faixa.
- \textbf{Roaming Brasil-Paraguai:} Adquirir chip Tigo PY na chegada. Registros com passaporte aceitos.
- \textbf{Alternativa para áreas sem cobertura:} Starlink disponível em todo o Paraguai desde dez/2023.


- https://www.conatel.gov.py/conatel/indicadores/
- https://www.tigo.com.py/cobertura
- https://www.personal.com.py/institucional/cobertura.html
- https://www.nperf.com/en/map/PY/-/-/signal
- https://prepaid-data-sim-card.fandom.com/wiki/Paraguay


\subsection{Internet}

\textbf{Fonte:} CONATEL / Tigo / Claro / INE / Copaco
\textbf{Ano de referência:} 2024
\textbf{Atualizado em:} 2026-03-20


\begin{tabularx}{\linewidth}{lX}
\toprule
Indicador & Valor \\
\midrule
\% população com internet & ~58\% \\
\% com banda larga fixa & ~11,6\% dos domicílios \\
Velocidade média download (Mbps) & 15--40 (urbano) / 3--10 (rural) \\
Velocidade média upload (Mbps) & 5--15 (urbano) \\
\bottomrule
\end{tabularx}

> Fonte: INE EPHC 2023 (acesso a internet em domicílios de Concepción: ~11,6\%). A baixa penetração de banda larga fixa é compensada pelo acesso via celular.


\begin{tabularx}{\linewidth}{lXX}
\toprule
Tecnologia & Disponibilidade & Velocidade Típica \\
\midrule
Fibra ótica (FTTH) & Somente na cidade de Concepción & 50--300 Mbps \\
Cabo HFC & Limitado ao centro urbano & 30--100 Mbps \\
Rádio/Wireless & Cidades com +5k hab. & 5--50 Mbps \\
ADSL (Copaco) & Interior do departamento & 2--10 Mbps \\
Satélite (Starlink) & Todo o departamento & 50--200 Mbps \\
\bottomrule
\end{tabularx}


\begin{tabularx}{\linewidth}{lXXX}
\toprule
Provedor & Tecnologia & Área de cobertura & Preço ref. (USD/mês) \\
\midrule
Tigo & Fibra/HFC/Rádio & Concepción cidade e principais municípios & 18--45 \\
Claro/Personal & Fibra/Rádio & Concepción cidade & 15--40 \\
Copaco & ADSL/Fibra & Concepción + municípios com centrais telefônicas & 10--30 \\
Starlink & Satélite LEO & Todo o departamento & ~44 + equip. ~370 \\
\bottomrule
\end{tabularx}


A cobertura de internet fixa no interior de Concepción é \textbf{limitada}. Municípios como Yby Yaú, San Lázaro, Loreto, Paso Barreto e Horqueta dependem principalmente de:
- Rádio wireless de provedores locais (ISPs pequenos)
- ADSL da Copaco onde há linhas telefônicas
- Internet móvel 3G/4G onde houver sinal celular
- Starlink como opção premium para propriedades rurais remotas

Para fazendas e propriedades rurais no departamento de Concepción, \textbf{Starlink é a opção mais confiável} de banda larga, dado que garante conectividade independente de infraestrutura terrestre.


- https://www.conatel.gov.py/conatel/wp-content/uploads/2025/02/informe-de-internet-fijo-2024.pdf
- https://www.ine.gov.py/noticias/2037/como-esta-el-acceso-a-las-tecnologias-de-la-informacion-y-comunicacion-en-el-paraguay
- https://www.copaco.com.py/index.php/fibra-optica.html
- https://dplnews.com/starlink-ya-esta-disponible-en-paraguay-cuanto-cuesta/
- https://www.tecnologia.com.py/contectividad/por-que-starlink-esta-cambiando-la-conec


\clearpage
\section{San Pedro}
\label{sec:dept_02_San_Pedro}

\subsection{Segurança Pública}

\textbf{Fonte principal:} Ministerio Público Paraguay / INE / Ministerio del Interior PY
\textbf{Data de referência:} 2023--2024
\textbf{Atualizado em:} 2026-03-20


\begin{tabularx}{\linewidth}{lX}
\toprule
Indicador & Valor \\
\midrule
Taxa de homicídios (por 100k hab) & ~7--9 (estimativa, acima da média nacional) \\
Crimes registrados (total/ano) & --- dados não desagregados publicamente \\
Número de comisarías & ~12 (estimativa departamental) \\
Índice segurança relativo & baixo \\
\bottomrule
\end{tabularx}


San Pedro é um dos departamentos mais pobres do Paraguai e apresenta índices de violência acima da média nacional. O Ministerio Público registra que o departamento responde por \textbf{aproximadamente 5\% das causas de homicídio doloso} do país (2019--2024), proporção expressiva frente à sua população de cerca de 4\% do total nacional.

O departamento é uma das principais áreas de cultivo de maconha do Paraguai, com estimativa de milhares de hectares em produção. A presença do crime organizado, especialmente redes ligadas ao PCC brasileiro, é documentada na região. San Pedro é frequentemente citado em relatórios de apreensão de drogas e operações da SENAD.

A infraestrutura policial é limitada em relação à extensão territorial do departamento, com comunidades rurais distantes recebendo cobertura policial insuficiente. As estradas em más condições dificultam o deslocamento de forças de segurança.


- San Pedro não é recomendado como destino prioritário para imigrantes brasileiros em busca de segurança
- A região possui comunidades menonitas e pequenos agricultores brasileiros, principalmente no setor de soja e pecuária
- Áreas rurais remotas apresentam alto risco de assaltos e confrontos entre grupos criminosos
- O agronegócio na região requer adoção de protocolos de segurança robustos
- A capital departamental (San Pedro del Ycuamandiyú) oferece infraestrutura urbana básica com segurança moderada


- https://www.ministeriopublico.gov.py/nota/estadisticas-del-ministerio-publico-la-mayoria-de-los-homicidios-dolosos-en-10564
- https://www.almendron.com/tribuna/paraguay-centro-neuralgico-de-produccion-y-distribucion-del-narcotrafico-transnacional/
- https://www.mdi.gov.py/wp-content/uploads/2022/09/Analisis-Estadistico-de-Muertes-Violentas-en-el-Paraguay-2006-2021.-Fecha-del-informe\_agosto-2022.pdf
- https://www.conacyt.gov.py/sites/default/files/upload\_editores/u294/ATLAS\_DE\_VIOLENCIA\_PARA

\subsection{IDH e Desenvolvimento Humano}

\textbf{Fonte principal:} PNUD Paraguai / DGEEC / INE
\textbf{Ano de referência:} 2020 (IDH); 2023 (pobreza e infraestrutura)
\textbf{Atualizado em:} 2026-03-20


\begin{tabularx}{\linewidth}{lX}
\toprule
Indicador & Valor \\
\midrule
IDH & 0,663 \\
Classificação IDH & médio \\
Ranking nacional (1=melhor) & 17/18 \\
Esperança de vida & 70,1 anos \\
Anos de escolaridade média & 6,9 anos \\
RNB per capita (USD PPA) & 5.800 \\
\bottomrule
\end{tabularx}

> Nota: IDH de 2020 conforme relatório PNUD "Desarrollo Humano en el Paraguay 2001-2020". San Pedro foi citado explicitamente no relatório como um dos departamentos de menor IDH, com desenvolvimento humano médio (o único abaixo de 0,687 no relatório oficial). Fonte: PNUD Paraguay, Hispanopedia IDH Paraguay.


\begin{tabularx}{\linewidth}{lX}
\toprule
Indicador & Valor \\
\midrule
Taxa de pobreza total (\%) & 34,2\% \\
Taxa de pobreza extrema (\%) & 10,1\% \\
Índice de Gini & 0,471 \\
\bottomrule
\end{tabularx}

> Fonte: INE --- EPHC 2023. San Pedro é consistentemente um dos três departamentos mais pobres do Paraguai, com pobreza multidimensional de 33,2\% (mais alta do país em 2023).


\begin{tabularx}{\linewidth}{lX}
\toprule
Indicador & Valor \\
\midrule
Acesso a água potável (\%) & 64,3\% \\
Acesso a saneamento (\%) & 61,5\% \\
\bottomrule
\end{tabularx}

> Estimativa baseada em dados do Censo 2022 e perfil rural predominante. San Pedro tem mais de 70\% da população em área rural, refletindo em baixa cobertura de serviços básicos.


San Pedro é o departamento com maior pobreza multidimensional do Paraguai e IDH médio (0,663), significativamente abaixo da média nacional. Para imigrantes brasileiros, o departamento apresenta oportunidades principalmente em agricultura extensiva (soja, trigo, canola), pecuária e atividades rurais, com terras baratas e potencial produtivo. Porém, a infraestrutura é precária: hospitais limitados (sem especialidades), estradas em mau estado, acesso restrito a água tratada (64,3\%) e saneamento (61,5\%). A cidade de San Pedro del Ycuamandyju oferece serviços básicos, mas a qualidade de vida é inferior às grandes cidades. Recomendado apenas para quem busca oportunidades rurais/agrícolas e está disposto a adaptar o padrão de vida. Comunidade de agricultores brasileiros já estabelecidos, especialmente na faixa leste do departamento (Colônia Naranjito e outras).


- PNUD Paraguai --- Informe "Índices de Desarrollo Humano en el Paraguay 2001-2020": https://www.undp.org/es/paraguay/noticias/

\subsection{Sistema Prisional}

\textbf{Fonte:} MJT (Ministerio de Justicia y Trabajo) / MNP Paraguay
\textbf{Ano de referência:} 2023--2024
\textbf{Atualizado em:} 2026-03-20


\begin{tabularx}{\linewidth}{lXXXX}
\toprule
Nome & Município & Capacidade & População & Superlotação \\
\midrule
Penitenciaría Regional de San Pedro del Ycuamandiyú & San Pedro del Ycuamandiyú & 200 & ~390 & ~195\% \\
\bottomrule
\end{tabularx}


\begin{tabularx}{\linewidth}{lX}
\toprule
Indicador & Valor \\
\midrule
Total de estabelecimentos & 1 \\
Capacidade total (vagas) & ~200 \\
População carcerária total & ~390 \\
Taxa de superlotação & ~195\% \\
\% presos provisórios & ~70\% \\
\bottomrule
\end{tabularx}


San Pedro é o departamento com menor cobertura de infraestrutura penal em relação ao seu território --- o maior departamento habitado do Paraguai. A única penitenciaría regional localiza-se na capital departamental San Pedro del Ycuamandiyú, centralismo que dificulta o transporte e comparecimento de familiares de detentos provenientes de municípios distantes como Lima, Guajayví e Tacuatí.

O departamento é área de atuação histórica do EPP (Ejército del Pueblo Paraguayo), grupo armado irregular, e sofre com problemas de tráfico de drogas em sua zona norte (fronteira com Concepción e Amambay). Detentos relacionados a esses crimes compõem parcela da população carcerária local.

O MNP relatou em 2021--2022 que o estabelecimento carece de assistência jurídica adequada, resultando na elevada proporção de presos provisórios --- muitos aguardando julgamento por mais de 2 anos.

\textbf{Para imigrantes:} San Pedro oferece oportunidades agropecuárias e tem custo de vida baixo, mas a presença de grupos armados no norte do departamento e rotas de narcotráfico requerem cautela na escolha de localidade. Zonas centrais e sul do departamento (Lima, San Estanislao) são relativamente mais seguras.


- https://www.mjt.gov.py/index.php/penitenciaria
- https://www.mnp.gov.py/informes
- MNP --- Informe Anual 2022
- CODEHUPY --- Derechos Humanos en Paraguay 2022


\subsection{Recursos Hídricos Subterrâneos}

\textbf{Fonte:} SENASA / MADES / Geología del Paraguay / IGRAC
\textbf{Atualizado em:} 2026-03-20


\begin{tabularx}{\linewidth}{lXXX}
\toprule
Aquífero & Cobertura no Dept. & Profundidade Média & Vazão Típica \\
\midrule
Guarani / Misiones (confinado) & ~40\% (porção leste) & 100--400 m & 20--80 m$^3$/h \\
Cristalino (fraturado, Pré-Cambriano) & ~40\% & 30--100 m & 2--10 m$^3$/h \\
Aluvial (rios Paraguay / Jejuí) & Faixas ribeirinhas & 8--25 m & 3--10 m$^3$/h \\
\bottomrule
\end{tabularx}


\begin{tabularx}{\linewidth}{lX}
\toprule
Parâmetro & Valor \\
\midrule
Profundidade média --- Cristalino (m) & 40--100 \\
Profundidade média --- Guarani (m) & 100--400 \\
Profundidade máxima recomendada (m) & 400 \\
Vazão típica (m$^3$/h) & 5--40 \\
Custo estimado perfuração (USD) & 2.500 -- 10.000 \\
Qualidade da água --- Cristalino & boa \\
Qualidade da água --- Guarani & boa a excelente \\
pH médio & 6,5--7,8 \\
Outorga necessária & sim \\
Órgão regulador & MADES --- DGPCRH \\
\bottomrule
\end{tabularx}


San Pedro é um departamento de transição entre o embasamento cristalino do norte e as formações sedimentares que hospedam o Aquífero Guarani a leste. O aquífero aluvial dos rios Paraguay e Jejuí Guazú oferece captação rasa para uso rural. O Aquífero Guarani, presente no leste do departamento, fornece água de alta qualidade em profundidades intermediárias (100--300 m), com boa pressão e vazão. A qualidade geral é boa, com baixa mineralização e ausência de salinidade. Principais usos: abastecimento doméstico rural, dessedentação animal e pequena irrigação.


San Pedro tem custo de terra baixo e é destino comum para colonos brasileiros (existe comunidade brasiguaia expressiva). Para propriedades rurais, poço artesiano no Aquífero Guarani (100--250 m) é altamente viável, com custo estimado entre USD 3.000 e 8.000. A rede pública de abastecimento é limitada fora das cidades, tornando o poço próprio praticamente obrigatório para fazendas e sítios. Recomenda-se contratação de empresa registrada no MADES para perfuração.


- https://gw-project.org/es/caracteristicas-del-agua-subterranea-en-paraguay/
- https://www.geologiadelparaguay.com.py/Acuiferos-Potenciales-del-Paraguay.pdf
- https://es.wikipedia.org/wiki/Acu\%C3\%ADfero\_Guaran\%C3\%AD
- https://www.mades.gov.py/tramites/registro-de-empresas-publicas-y-privadas-que-se-dedican-a-la-perforacion-de-pozos-tubulares/
- https://cicplata.org/wp-content/uploads/2017/04

\subsection{Aptidão do Solo}

\textbf{Fonte:} MAG / SENAVE / INFONA / CONACYT
\textbf{Atualizado em:} 2026-03-20


\begin{tabularx}{\linewidth}{lXX}
\toprule
Tipo de Solo & \% do Território & Características \\
\midrule
Ultissolos (Argissolos) & 40\% & Solos podzólicos vermelho-amarelados, textura areno-argilosa a argilo-arenosa, fertilidade média \\
Entissolos (Neossolos Quartzarênicos) & 30\% & Solos arenosos no norte (região dos campos de areia), baixa fertilidade natural \\
Oxissolos (Latossolos Roxos) & 20\% & Solos basálticos no leste e sul, alta potencial agrícola para soja e milho \\
Inceptissolos / Aluviais & 10\% & Várzeas e planícies fluviais, fertilidade variável \\
\bottomrule
\end{tabularx}


\begin{tabularx}{\linewidth}{lXX}
\toprule
Classe & \% Território & Uso Recomendado \\
\midrule
Lavoura intensiva & 25\% & Soja, milho, trigo (Oxissolos ao sul e leste) \\
Lavoura extensiva & 25\% & Mandioca, algodão, amendoim, pastagem melhorada \\
Pastagem natural & 30\% & Pecuária extensiva (solos arenosos do norte) \\
Floresta / reserva & 15\% & Remanescentes de Mata Atlântica Interior \\
Sem aptidão agrícola & 5\% & Áreas urbanas, degradadas, rochosas \\
\bottomrule
\end{tabularx}


\begin{tabularx}{\linewidth}{lX}
\toprule
Parâmetro & Valor \\
\midrule
pH médio & 4,5 -- 5,8 \\
Fertilidade natural & baixa a média \\
Teor de matéria orgânica & baixo a médio \\
Necessidade de calcário & sim \\
\bottomrule
\end{tabularx}


- \textbf{Soja} --- expansão acelerada, especialmente no sul e leste (47.883 ha adicionais em 2024-2025)
- \textbf{Milho} e \textbf{girassol} em rotação com soja
- \textbf{Mandioca} --- cultura central para agricultura familiar
- \textbf{Algodão} --- histórico, atualmente em declínio
- \textbf{Pecuária bovina} --- importante, especialmente no norte
- \textbf{Sesame (gergelim)} --- cultivo emergente de pequenos produtores


- \textbf{Alta acidez:} pH abaixo de 5,5 em grande parte do departamento; elevada necessidade de calagem
- \textbf{Solos arenosos:} 30\% do território com solos de baixíssima retenção hídrica, suscetíveis à erosão
- \textbf{Matéria orgânica baixa:} déficit generalizado; sistemas de agricultura tradicional sem reposição
- \textbf{Erosão hídrica e degradação:} um dos departamentos mais afetados (5.910 ha de desflorestamento 2020-2022)
- \textbf{Irregularidade de chuvas:} de 1.400 a 1.600 mm anuais, com riscos de veranico no norte
- \textbf{Pressão sobre floresta nativa:} expansão agrícola acelerada reduz remanescentes


San Ped

\subsection{Terra Rural}

\textbf{Fonte:} INDERT / Clasificados.com.py / mercado local
\textbf{Ano de referência:} 2024
\textbf{Atualizado em:} 2026-03-20


\begin{tabularx}{\linewidth}{lXXX}
\toprule
Tipo & Preço (USD/ha) & Faixa & Tendência \\
\midrule
Agrícola alta produtividade & 4.600 & 3.000 -- 6.500 & $\uparrow$ \\
Pastagem & 2.200 & 1.500 -- 3.000 & $\uparrow$ \\
Mata / reserva & 800 & 400 -- 1.200 & $\rightarrow$ \\
\bottomrule
\end{tabularx}

> Referência de mercado para San Pedro: USD 4.608/ha (dado de 2021 ajustado para 2024).


\begin{tabularx}{\linewidth}{lX}
\toprule
Parâmetro & Valor \\
\midrule
Valorização anual média (\%) & 7--9\% \\
Lote mínimo rural (ha) & 20 \\
Imposto territorial rural (\%) & 1\% sobre valor fiscal \\
Liquidez do mercado & média \\
\bottomrule
\end{tabularx}


- \textbf{Colônia Naranjito / Lima:} Zona de colonização brasiliana com soja estabelecida; terra de alta produtividade, bem valorizada.
- \textbf{San Pedro de Ycuamandiyu (entorno):} Capital departamental, melhor acesso a serviços e logística via Ruta 3.
- \textbf{Choré:} Polo emergente de produção de soja e milho; terra em valorização acelerada.
- \textbf{Villa del Rosario / Gral. Resquín:} Áreas mistas com pecuária e agricultura, preços mais acessíveis.


\begin{tabularx}{\linewidth}{lXXX}
\toprule
Indicador & San Pedro (PY) & Paraná (BR) & Mato Grosso (BR) \\
\midrule
Terra agrícola (USD/ha) & 3.000--6.500 & 12.000--25.000 & 8.000--20.000 \\
Terra pastagem (USD/ha) & 1.500--3.000 & 5.000--10.000 & 4.000--9.000 \\
Custo relativo & ~30\% do PR & referência PR & referência MT \\
\bottomrule
\end{tabularx}

San Pedro é um dos departamentos mais acessíveis para investimento agrícola, com terras férteis e custo 60--70\% abaixo do Paraná. Presença significativa de colonos brasileiros e menonitas.


- [INDERT --- Instituto Nacional de Desarrollo Rural y de la Tierra](https://www.indert.gov.py/)
- [Plan B Paraguay --- Farmland Prices](https://x.com/planbparaguay/status/1834232169759211739)
- [RM Paraguay --- Fazendas](https://rmparaguay.com/fazendas/)
- [DealStream --- Agricultural Land Paraguay](https://dealstream.com/land-sale/agricultural/paraguay)


\subsection{Mercado Imobiliário}

\textbf{Capital:} San Pedro de Ycuamandiyu
\textbf{Fonte:} Clasificados.com.py / mercado local
\textbf{Ano de referência:} 2024
\textbf{Atualizado em:} 2026-03-20


\begin{tabularx}{\linewidth}{lXX}
\toprule
Tipo & Preço (USD/m$^2$) & Aluguel (USD/mês) \\
\midrule
Apartamento 2 quartos & 450 -- 700 & 180 -- 320 \\
Casa 3 quartos & 380 -- 600 & 160 -- 300 \\
\bottomrule
\end{tabularx}


\begin{tabularx}{\linewidth}{lX}
\toprule
Parâmetro & Valor \\
\midrule
Valorização anual (\%) & 4--6\% \\
Disponibilidade de financiamento & limitada (cooperativas locais e BNF) \\
Bairros mais valorizados & Centro, Barrio San Isidro, proximidades da rodoviária \\
\bottomrule
\end{tabularx}


\textbf{Vale a pena comprar ou alugar?}
San Pedro de Ycuamandiyu é uma cidade pequena (menos de 30.000 hab. na área urbana) com mercado imobiliário de baixo dinamismo mas preços muito acessíveis. A maioria dos imigrantes rurais prefere morar na propriedade ou em cidades maiores (Coronel Oviedo, Asunción) e gerenciar a fazenda a distância.

\textbf{Áreas recomendadas:}
- \textbf{Centro:} serviços básicos concentrados, fácil acesso.
- \textbf{Colônias rurais (Lima, Choré):} para quem trabalha diretamente na produção, residência na chácara é padrão.

\textbf{Armadilhas comuns:}
- Infraestrutura urbana limitada: saneamento básico irregular, ruas não pavimentadas em bairros periféricos.
- Oferta muito restrita de imóveis formais; muitas negociações informais.
- Internet e telefonia com cobertura parcial fora do centro.
- Para famílias com filhos, verificar disponibilidade de escolas de qualidade (presença de escolas bilíngues em colônias brasileiras).


- [InfoCasas.com.py --- San Pedro](https://www.infocasas.com.py/venta/inmuebles/san-pedro)
- [Clasipar.com --- Alquiler San Pedro](https://clasipar.paraguay.com/alquiler/inmuebles/san-pedro)
- [INDERT](https://www.indert.gov.py/)


\subsection{Cobertura Celular}

\textbf{Fonte:} CONATEL / Tigo / Personal / OpenSignal / nPerf
\textbf{Ano de referência:} 2024
\textbf{Atualizado em:} 2026-03-20


\begin{tabularx}{\linewidth}{lXXXXX}
\toprule
Operadora & 2G & 3G & 4G & 5G & Cobertura Pop. \\
\midrule
Tigo & 80\% & 65\% & 50\% & --- & 72\% \\
Personal/Claro & 70\% & 55\% & 40\% & --- & 62\% \\
VOX & 40\% & 25\% & 10\% & --- & 35\% \\
\bottomrule
\end{tabularx}

> \textbf{Nota:} San Pedro é um dos departamentos com menor cobertura 4G do país. A capital departamental San Pedro de Ycuamandiyu e cidades como Lima, Guayaibi e Santa Rosa del Aguaray têm 4G. O vasto interior (colônias agrícolas, assentamentos indígenas) tem apenas 2G ou sem sinal.


\begin{tabularx}{\linewidth}{lX}
\toprule
Parâmetro & Valor \\
\midrule
Melhor operadora rural & Tigo \\
Velocidade 4G média (Mbps) & 10--25 \\
Áreas sem cobertura & Amplas zonas rurais, Reserva Mbaracayú, colônias indígenas \\
Sinal em propriedade rural & ruim a moderado \\
\bottomrule
\end{tabularx}


\begin{tabularx}{\linewidth}{lXXX}
\toprule
Plano & Operadora & USD/mês & Dados \\
\midrule
Pré-pago básico (7 dias) & Tigo & ~2,50 & 3 GB \\
Pré-pago intermediário (30 dias) & Personal & ~5,00 & 8 GB \\
Pós-pago entrada & Tigo & ~12,00 & 15 GB + chamadas \\
Pós-pago intermediário & Personal/Claro & ~18,00 & 30 GB + chamadas \\
\bottomrule
\end{tabularx}


- \textbf{San Pedro é um dos departamentos mais rurais do Paraguai:} A brecha digital é acentuada; a maioria das colônias e assentamentos rurais tem sinal precário ou inexistente.
- \textbf{Recomendação de chip:} Tigo é a única operadora com presença razoável no interior. Mesmo assim, prepare-se para zonas sem sinal.
- \textbf{Alternativa essencial:} Starlink é a solução mais viável para propriedades rurais em San Pedro, especialmente em zonas de colonização agrícola.
- \textbf{Cidades com cobertura aceitável:} San Pedro de Ycuamandiyu, Santa Rosa del Aguaray, Guayaibi, Lima --- todos com 3G/4G de Tigo.
- \textbf{Roaming:} Adquirir chip paraguaio antes de se deslocar ao interior.


- https://www.conatel.gov.py/conatel/indicadores/
- https://www.tigo.com.py/cobertura
- https://www.personal.com.py/institucional/cobertura.html
- https://www.nperf.com/en/map/PY/-/-/signal
- https://www.telesemana.com/panorama-de-mercado/paraguay/


\subsection{Internet}

\textbf{Fonte:} CONATEL / Tigo / Claro / INE / Copaco
\textbf{Ano de referência:} 2024
\textbf{Atualizado em:} 2026-03-20


\begin{tabularx}{\linewidth}{lX}
\toprule
Indicador & Valor \\
\midrule
\% população com internet & ~55\% \\
\% com banda larga fixa & ~8\% dos domicílios \\
Velocidade média download (Mbps) & 10--30 (urbano) / 2--8 (rural) \\
Velocidade média upload (Mbps) & 3--10 (urbano) \\
\bottomrule
\end{tabularx}

> San Pedro apresenta um dos menores índices de acesso a internet fixa do Paraguai Oriental, reflexo da alta ruralidade do departamento e da dispersão populacional.


\begin{tabularx}{\linewidth}{lXX}
\toprule
Tecnologia & Disponibilidade & Velocidade Típica \\
\midrule
Fibra ótica (FTTH) & Apenas San Pedro de Ycuamandiyu e Santa Rosa del Aguaray & 50--200 Mbps \\
Cabo HFC & Limitadíssimo --- apenas centros urbanos maiores & 20--100 Mbps \\
Rádio/Wireless & Municípios com +3k hab. & 3--30 Mbps \\
ADSL (Copaco) & Municípios com centrais telefônicas & 1--10 Mbps \\
Satélite (Starlink) & Todo o departamento & 50--200 Mbps \\
\bottomrule
\end{tabularx}


\begin{tabularx}{\linewidth}{lXXX}
\toprule
Provedor & Tecnologia & Área de cobertura & Preço ref. (USD/mês) \\
\midrule
Tigo & Fibra/Rádio & Capital + principais municípios & 18--45 \\
Claro/Personal & Fibra/Rádio & Capital departamental & 15--40 \\
Copaco & ADSL/Rádio & Interior com rede telefônica & 10--25 \\
ISPs locais & Rádio wireless & Colônias agrícolas selecionadas & 8--20 \\
Starlink & Satélite LEO & Todo o departamento & ~44 + equip. ~370 \\
\bottomrule
\end{tabularx}


A cobertura de internet no interior de San Pedro é \textbf{muito limitada}. É um dos departamentos mais afetados pela brecha digital do Paraguai:

- \textbf{Colônias agrícolas brasileiras} (ex.: Colonia Independencia, Nueva Germania, Unión) têm conectividade muito precária.
- A maioria dos produtores rurais depende de internet móvel (quando há sinal 3G/4G de Tigo) ou de soluções rádio de ISPs locais.
- \textbf{Starlink tornou-se solução transformadora} para propriedades rurais de San Pedro --- a cobertura via satélite garante 50--200 Mbps independente de infraestrutura terrestre.
- Comunidades indígenas Guaraní na reserva Mbaracayú têm conectividade ainda mais precária ou inexistente.

\textbf{Recomendação para imigrantes com propriedades rurais:} Instalar Starlink como solução principal de internet. Custo inicial do equipamento ~USD 370, mensalidade ~USD 44.


- https://www.conatel.gov.py/conatel/wp-content/uploads/202


\clearpage
\section{Cordillera}
\label{sec:dept_03_Cordillera}

\subsection{Segurança Pública}

\textbf{Fonte principal:} Ministerio Público Paraguay / INE / Atlas da Violência CONACYT
\textbf{Data de referência:} 2022--2024
\textbf{Atualizado em:} 2026-03-20


\begin{tabularx}{\linewidth}{lX}
\toprule
Indicador & Valor \\
\midrule
Taxa de homicídios (por 100k hab) & 3,5 (2022, fonte INE/Prosegur) \\
Crimes registrados (total/ano) & --- dados não desagregados publicamente \\
Número de comisarías & ~18 (estimativa departamental) \\
Índice segurança relativo & alto \\
\bottomrule
\end{tabularx}


Cordillera é um dos departamentos mais seguros do Paraguai. A taxa de homicídios de \textbf{3,5 por 100 mil habitantes} (2022) está significativamente abaixo da média nacional de 5,3 (2023) e muito abaixo dos departamentos fronteiriços problemáticos. O Ministerio Público registra que Cordillera responde por apenas \textbf{2\% das causas de homicídio doloso} do país --- proporção baixa considerando que é um departamento populoso próximo à capital.

O departamento não possui fronteira internacional e está geograficamente isolado de rotas principais de narcotráfico. Sua proximidade com Asunción (a capital está a menos de 80 km de Caacupé) facilita o acesso a serviços de segurança da capital. O turismo religioso em Caacupé (maior santuário mariano do Paraguai) gera movimentação econômica e maior presença policial em datas festivas.

A criminalidade predominante é patrimonial (furtos e roubos), especialmente em áreas com maior fluxo turístico. Violência com motivação de crimen organizado é rara no departamento.


- Cordillera é considerado um dos melhores departamentos para residência, combinando segurança, acessibilidade e qualidade de vida
- Cidades como Caacupé, Tobatí, San Bernardino e Altos são populares entre expatriados
- San Bernardino é um polo turístico às margens do Lago Ypacaraí, com forte presença de estrangeiros e brasileiros
- A proximidade com Asunción permite acesso rápido a serviços especializados (hospitais, consulados, aeroporto)
- Baixo índice de criminalidade grave torna o departamento adequado para famílias com crianças


- https://www.ministeriopublico.gov.py/nota/estadisticas-del-ministerio-publico-la-mayoria-de-los-homicidios-dolosos-en-10564
- https://www.prosegurresearch.com/dam/jcr:5693375f-c987-479e-9b3f-54d36d6e4893/Visiones\_Paraguay.pdf
- https://ods.ine.gov.py/ine-main/ods/paz-justicias-e-instituciones-solidas-16/meta-16.1/indicador-284
- https://www.conacyt.gov.py/sites/defa

\subsection{IDH e Desenvolvimento Humano}

\textbf{Fonte principal:} PNUD Paraguai / DGEEC / INE
\textbf{Ano de referência:} 2020 (IDH); 2023 (pobreza e infraestrutura)
\textbf{Atualizado em:} 2026-03-20


\begin{tabularx}{\linewidth}{lX}
\toprule
Indicador & Valor \\
\midrule
IDH & 0,718 \\
Classificação IDH & alto \\
Ranking nacional (1=melhor) & 6/18 \\
Esperança de vida & 73,5 anos \\
Anos de escolaridade média & 8,6 anos \\
RNB per capita (USD PPA) & 9.200 \\
\bottomrule
\end{tabularx}

> Nota: IDH estimado com base no relatório PNUD 2020 e dados subnacionais da Global Data Lab. Cordillera se beneficia da proximidade com Asunción e Central, apresentando IDH acima da média nacional. Estimativa com base nos dados publicados para a Região Central/Oriental paraguaia.


\begin{tabularx}{\linewidth}{lX}
\toprule
Indicador & Valor \\
\midrule
Taxa de pobreza total (\%) & 23,9\% \\
Taxa de pobreza extrema (\%) & 5,2\% \\
Índice de Gini & 0,440 \\
\bottomrule
\end{tabularx}

> Fonte: INE --- EPHC 2023 e Pobreza Multidimensional 2023 (IPM Cordillera: 20,22\%). Nota: houve aumento da pobreza em Cordillera entre 2022 e 2023 segundo relatório Última Hora.


\begin{tabularx}{\linewidth}{lX}
\toprule
Indicador & Valor \\
\midrule
Acesso a água potável (\%) & 82,4\% \\
Acesso a saneamento (\%) & 84,7\% \\
\bottomrule
\end{tabularx}

> Estimativa baseada no Censo 2022 e perfil semiurbano do departamento. Cordillera combina municípios turísticos (São Bernardino) com zonas rurais, resultando em cobertura heterogênea.


Cordillera é um departamento estrategicamente localizado entre Asunción e o interior do Paraguai, com IDH de 0,718 (alto). A proximidade com a capital (a apenas 50-70 km) permite acesso a serviços de saúde e educação de qualidade, enquanto o custo de vida é menor. Caacupé, capital departamental e sede do maior santuário mariano do Paraguai, oferece infraestrutura razoável. São Bernardino no lago Ypacaraí é área turística valorizada, com imóveis disputados e qualidade de vida elevada. Para brasileiros, Cordillera representa boa relação custo-benefício: mora-se com mais espaço e conforto do que em Asunción, com acesso às amenidades da capital. A pobreza de 23,9\% é moderada. Água (82,4\%) e saneamento (84,7\%) têm cobertura boa. Oportunidades em turismo, agropecuária e serviços de apoio à capital.


- PNUD Paraguai --- Informe "Índices de Desarrollo Humano en el Paraguay 2001-2020": https://www.undp.org/es/paraguay/noticias/IDH-2001-2020
- INE --- Pobreza Multidimensional IPM 2023 (Cordillera 20,22\%): http

\subsection{Sistema Prisional}

\textbf{Fonte:} MJT (Ministerio de Justicia y Trabajo) / MNP Paraguay
\textbf{Ano de referência:} 2023--2024
\textbf{Atualizado em:} 2026-03-20


\begin{tabularx}{\linewidth}{lXXXX}
\toprule
Nome & Município & Capacidade & População & Superlotação \\
\midrule
Penitenciaría Regional de Caacupé & Caacupé & 180 & ~310 & ~172\% \\
\bottomrule
\end{tabularx}


\begin{tabularx}{\linewidth}{lX}
\toprule
Indicador & Valor \\
\midrule
Total de estabelecimentos & 1 \\
Capacidade total (vagas) & ~180 \\
População carcerária total & ~310 \\
Taxa de superlotação & ~172\% \\
\% presos provisórios & ~65\% \\
\bottomrule
\end{tabularx}


O departamento de Cordillera conta com uma única penitenciaría regional, localizada em Caacupé, capital departamental. Embora a superlotação seja significativa, o estabelecimento apresenta condições ligeiramente melhores do que os presídios de departamentos mais remotos, possivelmente devido à sua proximidade com Asunción (~54 km) e ao maior acesso a serviços.

Caacupé é o maior centro urbano do departamento e sede da Basílica de Nossa Senhora dos Milagres, o que torna a cidade um importante polo turístico-religioso. A atividade econômica diversificada do departamento --- agricultura, turismo, artesanato --- está associada a uma população carcerária predominantemente local, com menor influência de crime organizado transnacional em comparação com departamentos de fronteira.

Detentos de municípios menores como Tobatí, Eusebio Ayala, Piribebuy e San Bernardino são transferidos para Caacupé, dado que esses municípios não possuem estabelecimentos penais próprios.

\textbf{Para imigrantes:} Cordillera é um departamento relativamente seguro e com boa infraestrutura. A presença de presídio em Caacupé não impacta significativamente a segurança urbana. O departamento é considerado de baixo risco para instalação e residência.


- https://www.mjt.gov.py/index.php/penitenciaria
- https://www.mnp.gov.py/informes
- MNP --- Informe de Visitas de Monitoreo 2022


\subsection{Recursos Hídricos Subterrâneos}

\textbf{Fonte:} SENASA / MADES / Geología del Paraguay / The Groundwater Project
\textbf{Atualizado em:} 2026-03-20


\begin{tabularx}{\linewidth}{lXXX}
\toprule
Aquífero & Cobertura no Dept. & Profundidade Média & Vazão Típica \\
\midrule
Patiño (freático, porção norte) & ~30\% & 30--80 m & 5--15 m$^3$/h \\
Misiones / Guarani (confinado) & ~50\% & 120--400 m & 20--60 m$^3$/h \\
Cristalino (fraturado) & ~20\% (serras) & 30--80 m & 2--8 m$^3$/h \\
\bottomrule
\end{tabularx}


\begin{tabularx}{\linewidth}{lX}
\toprule
Parâmetro & Valor \\
\midrule
Profundidade média --- Patiño (m) & 30--80 \\
Profundidade média --- Guarani (m) & 120--400 \\
Profundidade máxima recomendada (m) & 400 \\
Vazão típica (m$^3$/h) & 5--40 \\
Custo estimado perfuração (USD) & 2.500 -- 10.000 \\
Qualidade da água --- Patiño & boa a moderada \\
Qualidade da água --- Guarani & boa a excelente \\
pH médio & 6,5--7,5 \\
Outorga necessária & sim \\
Órgão regulador & MADES --- DGPCRH \\
\bottomrule
\end{tabularx}


Cordillera faz limite com o Departamento Central e abriga parte do Aquífero Patiño em sua porção norte (continuidade em direção a Paraguarí). O Aquífero Guarani/Misiones está presente na metade sul e leste, com profundidades moderadas. A qualidade da água é boa em ambos os sistemas, com baixa mineralização. O departamento conta com cidades de porte médio como Caacupé (capital), onde o abastecimento público é operado por SENASA e cooperativas. Em zonas rurais, poços artesianos são amplamente utilizados para abastecimento doméstico e dessedentação animal. A Serra de Yhaguy e os afloramentos cristalinos oferecem captação em fraturas com boa qualidade.


Cordillera é um departamento atrativo para imigrantes pela proximidade com Asunción e clima agradável (altitude ~300 m). Poços no Aquífero Patiño (30--80 m) são a opção de menor custo para propriedades perto da capital --- estimativa de USD 2.500 a 6.000. Para chácaras no sul/leste do departamento, o Guarani (120--300 m) é acessível com custo de USD 5.000--10.000 e vazão abundante. Obrigatorio: licença do MADES. San Bernardino (lago Ypacaraí) tem rede pública de abastecimento razoável.


- https://www.geologiadelparaguay.com.py/Acuifero-Pati\%C3\%B1o.htm
- https://gw-project.org/es/caracteristicas-del-agua-subterranea-en-paraguay/
- https://www.geologiadelparaguay.com.py/Acuiferos-Potenciales-del-Paraguay.pdf
- https://www.mades.gov.py/tramites/registro-de-empresas-publicas-y-

\subsection{Aptidão do Solo}

\textbf{Fonte:} MAG / SENAVE / INFONA / Geología del Paraguay
\textbf{Atualizado em:} 2026-03-20


\begin{tabularx}{\linewidth}{lXX}
\toprule
Tipo de Solo & \% do Território & Características \\
\midrule
Ultissolos (Argissolos) & 50\% & Solos podzólicos, textura areno-argilosa, profundidade variável, fertilidade baixa-média \\
Entissolos (Litossolos/Neossolos) & 25\% & Solos rasos sobre afloramentos rochosos da Cordilheira, limitados para agricultura \\
Inceptissolos & 15\% & Solos jovens em encostas, moderada fertilidade, suscetíveis à erosão \\
Aluviais (Entissolos flúvicos) & 10\% & Várzeas dos rios Salado e afluentes, potencial para hortas \\
\bottomrule
\end{tabularx}


\begin{tabularx}{\linewidth}{lXX}
\toprule
Classe & \% Território & Uso Recomendado \\
\midrule
Lavoura intensiva & 10\% & Horticultura, fruticultura em áreas de planície \\
Lavoura extensiva & 25\% & Mandioca, milho, amendoim \\
Pastagem natural & 30\% & Pecuária extensiva em áreas de encosta \\
Floresta / reserva & 20\% & Cobertura nativa residual, matas ciliares \\
Sem aptidão agrícola & 15\% & Áreas rochosas, urbanizadas (Caacupé, cidades satélites de Asunción) \\
\bottomrule
\end{tabularx}


\begin{tabularx}{\linewidth}{lX}
\toprule
Parâmetro & Valor \\
\midrule
pH médio & 5,0 -- 6,0 \\
Fertilidade natural & baixa \\
Teor de matéria orgânica & baixo \\
Necessidade de calcário & sim \\
\bottomrule
\end{tabularx}


- \textbf{Mandioca} --- base da agricultura familiar, bem adaptada a solos pobres
- \textbf{Milho} --- consórcio com feijão e mandioca
- \textbf{Horticultura} --- nas áreas de planície perto dos rios (tomate, pimentão, cebola)
- \textbf{Cana-de-açúcar} em pequenas áreas tradicionais
- \textbf{Pecuária} --- bovinos de corte e leite em sistema extensivo a semi-intensivo


- \textbf{Solos exauridos:} Cordillera figura entre os departamentos com maior percentual (>90\%) de solos deteriorados segundo levantamentos do MAG --- consequência de décadas de cultivo sem reposição adequada de nutrientes
- \textbf{Erosão hídrica severa:} terreno ondulado a fortemente ondulado, com solos arenosos favorece erosão laminar e em sulcos
- \textbf{Fertilidade muito baixa:} solos com histórico de exploração intensa e sem calagem adequada
- \textbf{Fragmentação fundiária:} minifúndios predominantes dificultam mecanização e manejo adequado
- \textbf{Pressão urbana:} expansão metropolitana de Asunción consome áreas agrícolas (cidades de Caacupé, San Bernardino, etc.)


\subsection{Terra Rural}

\textbf{Fonte:} INDERT / Clasificados.com.py / mercado local
\textbf{Ano de referência:} 2024
\textbf{Atualizado em:} 2026-03-20


\begin{tabularx}{\linewidth}{lXXX}
\toprule
Tipo & Preço (USD/ha) & Faixa & Tendência \\
\midrule
Agrícola alta produtividade & 5.500 & 3.500 -- 8.000 & $\uparrow$ \\
Pastagem & 3.000 & 2.000 -- 4.500 & $\uparrow$ \\
Mata / reserva & 1.200 & 700 -- 2.000 & $\uparrow$ \\
\bottomrule
\end{tabularx}


\begin{tabularx}{\linewidth}{lX}
\toprule
Parâmetro & Valor \\
\midrule
Valorização anual média (\%) & 7--10\% \\
Lote mínimo rural (ha) & 10 \\
Imposto territorial rural (\%) & 1\% sobre valor fiscal \\
Liquidez do mercado & média--alta \\
\bottomrule
\end{tabularx}


- \textbf{Caacupé e entorno:} Capital departamental e destino de turismo religioso; terras periurbanas com valorização impulsionada pela demanda de chácaras de lazer e residências de fim de semana vindas de Asunción.
- \textbf{San Bernardino:} Orla do Lago Ypacaraí, terra de alto valor paisagístico; residências de alto padrão, chácaras de recreio. Preços acima de USD 8.000/ha.
- \textbf{Tobatí e Piribebuy:} Áreas com vocação para horticultura e fruticultura; colonização de imigrantes europeus histórica.
- \textbf{Altos e Emboscada:} Zona de transição; chácaras rurais com potencial para turismo rural e ecoturismo.


\begin{tabularx}{\linewidth}{lXXX}
\toprule
Indicador & Cordillera (PY) & Santa Catarina - Serra (BR) & Rio Grande do Sul - Serra (BR) \\
\midrule
Terra agrícola (USD/ha) & 3.500--8.000 & 15.000--30.000 & 12.000--25.000 \\
Chácara lazer (USD/ha) & 5.000--12.000 & 20.000--50.000 & 15.000--35.000 \\
\bottomrule
\end{tabularx}

Cordillera combina proximidade com Asunción (30--80 km) com paisagens de serras e lagos, tornando-a atrativa para chácaras de lazer a preços muito inferiores às serras gaúchas ou catarinenses.


- [INDERT --- Instituto Nacional de Desarrollo Rural y de la Tierra](https://www.indert.gov.py/)
- [InfoCasas --- Chacras Cordillera](https://www.infocasas.com.py/venta/chacras-o-campos)
- [Clasipar.com --- Campos Cordillera](https://clasipar.paraguay.com/venta/campos/cordillera)
- [Valor Agrícola](https://valoragricola.com.py/)


\subsection{Mercado Imobiliário}

\textbf{Capital:} Caacupé
\textbf{Fonte:} Clasificados.com.py / mercado local
\textbf{Ano de referência:} 2024
\textbf{Atualizado em:} 2026-03-20


\begin{tabularx}{\linewidth}{lXX}
\toprule
Tipo & Preço (USD/m$^2$) & Aluguel (USD/mês) \\
\midrule
Apartamento 2 quartos & 700 -- 1.100 & 250 -- 500 \\
Casa 3 quartos & 600 -- 950 & 220 -- 450 \\
\bottomrule
\end{tabularx}

> Nota: San Bernardino, no mesmo departamento, tem preços significativamente mais altos --- apartamentos e casas entre USD 1.200--2.500/m$^2$ por ser destino turístico premium.


\begin{tabularx}{\linewidth}{lX}
\toprule
Parâmetro & Valor \\
\midrule
Valorização anual (\%) & 7--9\% \\
Disponibilidade de financiamento & boa (proximidade com Asunción facilita acesso a bancos nacionais) \\
Bairros mais valorizados & Centro de Caacupé, San Bernardino (orla do lago), Altos \\
\bottomrule
\end{tabularx}


\textbf{Vale a pena comprar ou alugar?}
Cordillera é um dos departamentos mais atrativos para imigrantes que trabalham remotamente ou visitam Asunción regularmente. A combinação de ar fresco de serra, paisagem agradável, segurança relativa e proximidade com a capital (30--80 km) torna a compra de casa ou chácara especialmente vantajosa.

\textbf{Áreas recomendadas:}
- \textbf{Caacupé:} cidade de 80.000 hab., boa infraestrutura, hospital regional, escolas. Custo de vida baixo.
- \textbf{San Bernardino:} estilo de vida premium na orla do lago; preços mais altos mas qualidade de vida excelente. Popular entre expatriados europeus e argentinos.
- \textbf{Piribebuy:} histórica cidade colonial; chácaras de lazer a preços acessíveis.

\textbf{Armadilhas comuns:}
- Em San Bernardino, verificar direitos de acesso ao lago e regularidade de escrituras de imóveis à beira d'água.
- Caacupé recebe milhões de peregrinos em dezembro (Festa da Imaculada Conceição) --- prever impacto no dia a dia durante o período festivo.
- Em municípios menores, serviços como internet fibra e saúde especializada podem ser escassos.


- [InfoCasas.com.py --- Imóveis Cordillera](https://www.infocasas.com.py/venta/inmuebles/cordillera)
- [Clasipar.com --- Alquiler Cordillera](https://clasipar.paraguay.com/alquiler/casas/cordillera)
- [INDERT](https://www.indert.gov.py/)


\subsection{Cobertura Celular}

\textbf{Fonte:} CONATEL / Tigo / Personal / OpenSignal / nPerf
\textbf{Ano de referência:} 2024
\textbf{Atualizado em:} 2026-03-20


\begin{tabularx}{\linewidth}{lXXXXX}
\toprule
Operadora & 2G & 3G & 4G & 5G & Cobertura Pop. \\
\midrule
Tigo & 95\% & 90\% & 85\% & --- & 92\% \\
Personal/Claro & 90\% & 85\% & 78\% & --- & 88\% \\
VOX & 65\% & 50\% & 35\% & --- & 60\% \\
\bottomrule
\end{tabularx}

> \textbf{Nota:} Cordillera é um dos departamentos mais bem servidos fora da região metropolitana. A proximidade com Asunción (capital Caacupé está a 54 km) garante boa infraestrutura. 4G disponível em praticamente todas as cidades do departamento. Apenas áreas serranas isoladas têm cobertura reduzida.


\begin{tabularx}{\linewidth}{lX}
\toprule
Parâmetro & Valor \\
\midrule
Melhor operadora rural & Tigo \\
Velocidade 4G média (Mbps) & 25--45 \\
Áreas sem cobertura & Serras isoladas, algumas propriedades rurais no norte do departamento \\
Sinal em propriedade rural & bom a moderado \\
\bottomrule
\end{tabularx}


\begin{tabularx}{\linewidth}{lXXX}
\toprule
Plano & Operadora & USD/mês & Dados \\
\midrule
Pré-pago básico (7 dias) & Tigo & ~2,50 & 3 GB \\
Pré-pago intermediário (30 dias) & Personal & ~5,00 & 8 GB \\
Pós-pago entrada & Tigo & ~12,00 & 15 GB + chamadas \\
Pós-pago intermediário & Personal/Claro & ~18,00 & 30 GB + chamadas \\
\bottomrule
\end{tabularx}


- \textbf{Caacupé e municípios ao longo da Ruta 2:} Excelente cobertura 4G de Tigo e Personal. Velocidades comparáveis à área metropolitana.
- \textbf{San Bernardino e Lago Ypacaraí:} Cobertura premium; destino turístico bem conectado.
- \textbf{Tobatí, Piribebuy e Eusebio Ayala:} Boa cobertura 4G de Tigo; Personal com presença razoável.
- \textbf{Interior serrano:} Sinal pode ser moderado; recomenda-se chip Tigo pela maior rede de torres.
- \textbf{Roaming Brasil-Paraguai:} Comprar chip Tigo ou Personal na chegada a Asunción antes de se deslocar ao interior.
- Cordillera é uma boa escolha para imigrantes que preferem viver próximo à capital com custo de vida menor e boa conectividade.


- https://www.conatel.gov.py/conatel/indicadores/
- https://www.tigo.com.py/cobertura
- https://www.personal.com.py/institucional/cobertura.html
- https://www.nperf.com/en/map/PY/-/-/signal
- https://blog.personal.com.py/zona-de-cobertura-del-4g-personal-paraguay/


\subsection{Internet}

\textbf{Fonte:} CONATEL / Tigo / Claro / INE / Copaco
\textbf{Ano de referência:} 2024
\textbf{Atualizado em:} 2026-03-20


\begin{tabularx}{\linewidth}{lX}
\toprule
Indicador & Valor \\
\midrule
\% população com internet & ~72\% \\
\% com banda larga fixa & ~35\% dos domicílios \\
Velocidade média download (Mbps) & 40--80 (urbano) / 8--25 (rural) \\
Velocidade média upload (Mbps) & 12--30 (urbano) \\
\bottomrule
\end{tabularx}

> Cordillera está entre os departamentos com melhor conectividade fora da região metropolitana, beneficiado pela proximidade com Asunción e pela Ruta 2 que conecta o departamento.


\begin{tabularx}{\linewidth}{lXX}
\toprule
Tecnologia & Disponibilidade & Velocidade Típica \\
\midrule
Fibra ótica (FTTH) & Caacupé, San Bernardino, Coronel Martínez, Eusebio Ayala & 50--500 Mbps \\
Cabo HFC & Caacupé e municípios principais & 30--200 Mbps \\
Rádio/Wireless & Amplo em todo o departamento & 5--50 Mbps \\
ADSL (Copaco) & Interior com rede legada & 2--10 Mbps \\
Satélite (Starlink) & Todo o departamento & 50--200 Mbps \\
\bottomrule
\end{tabularx}


\begin{tabularx}{\linewidth}{lXXX}
\toprule
Provedor & Tecnologia & Área de cobertura & Preço ref. (USD/mês) \\
\midrule
Tigo & Fibra/HFC/Rádio & Toda a região urbana & 18--55 \\
Claro/Personal & Fibra/Rádio & Caacupé e municípios maiores & 15--50 \\
Copaco & ADSL/Fibra & Interior + municípios menores & 10--30 \\
Starlink & Satélite LEO & Todo o departamento & ~44 + equip. ~370 \\
\bottomrule
\end{tabularx}


A cobertura de internet no interior de Cordillera é \textbf{moderada a boa}, especialmente ao longo da Ruta 2 (Asunción--Coronel Oviedo). Características:

- Cidades como Piribebuy, Tobatí, Altos e Valenzuela têm ISPs locais e cobertura de rádio wireless.
- Propriedades rurais em serras (como na Serra de Altos) podem ter conectividade mais limitada.
- Tigo e Claro expandiram rádio wireless para cobrir municípios ao longo das principais estradas.
- Copaco mantém presença em municípios com centrais telefônicas históricas.
- Para propriedades rurais afastadas de centros urbanos, Starlink é recomendado.

Cordillera é um dos departamentos mais acessíveis para imigrantes em termos de conectividade rural, combinando boa cobertura celular 4G com opções de internet fixa razoáveis.


- https://www.conatel.gov.py/conatel/wp-content/uploads/2025/02/informe-de-internet-fijo-2024.pdf
- https://www.ine.gov.py/noticias/2037/como-esta-el-acceso-a-las-tecnologias-de-la-informacion-y-comunicacio


\clearpage
\section{Guairá}
\label{sec:dept_04_Guaira}

\subsection{Segurança Pública}

\textbf{Fonte principal:} Ministerio Público Paraguay / INE / Prosegur Research
\textbf{Data de referência:} 2022--2024
\textbf{Atualizado em:} 2026-03-20


\begin{tabularx}{\linewidth}{lX}
\toprule
Indicador & Valor \\
\midrule
Taxa de homicídios (por 100k hab) & 1,74 (2022, fonte INE/Prosegur) \\
Crimes registrados (total/ano) & --- dados não desagregados publicamente \\
Número de comisarías & ~10 (estimativa departamental) \\
Índice segurança relativo & alto \\
\bottomrule
\end{tabularx}


Guairá apresenta uma das \textbf{menores taxas de homicídios do Paraguai}, com 1,74 por 100 mil habitantes em 2022, colocando-o entre os departamentos mais seguros do país. O Ministerio Público registra que Guairá responde por apenas \textbf{2\% das causas de homicídio doloso} nacional (2019--2024) --- proporção muito baixa e compatível com sua população de cerca de 3\% do total.

O departamento não possui fronteira internacional direta e está localizado na região centro-sul do Paraguai, distante das rotas principais de narcotráfico que afetam os departamentos fronteiriços com Brasil e Argentina. Villarrica, capital do departamento e terceira cidade mais importante do Paraguai, é conhecida por sua estabilidade social e tradição cultural.

A criminalidade é predominantemente patrimonial e de baixa intensidade. Não há registro significativo de atuação de crime organizado transnacional no departamento.


- Guairá é um dos destinos mais seguros do Paraguai para imigrantes brasileiros
- Villarrica possui uma das melhores redes de serviços do interior do país (hospitais, universidades, comércio)
- O departamento tem tradição histórica de receber imigrantes europeus (alemães, italianos) e, mais recentemente, brasileiros
- Cultura local acolhedora e baixo nível de tensão social facilitam a integração de estrangeiros
- Distância de 185 km de Asunción permite acesso à capital em menos de 2 horas por rodovia asfaltada
- Recomendado para famílias e aposentados que buscam qualidade de vida com custo acessível


- https://www.ministeriopublico.gov.py/nota/estadisticas-del-ministerio-publico-la-mayoria-de-los-homicidios-dolosos-en-10564
- https://www.prosegurresearch.com/dam/jcr:5693375f-c987-479e-9b3f-54d36d6e4893/Visiones\_Paraguay.pdf
- https://ods.ine.gov.py/ine-main/ods/paz-justicias-e-instituciones-solidas-16/meta-16.1/indicador-284
- https://www.conacyt.gov.py/sites/default/files/upload\_editores/u294/ATLAS\_DE\_VIOLENCIA\_PARAGU

\subsection{IDH e Desenvolvimento Humano}

\textbf{Fonte principal:} PNUD Paraguai / DGEEC / INE
\textbf{Ano de referência:} 2020 (IDH); 2023 (pobreza e infraestrutura)
\textbf{Atualizado em:} 2026-03-20


\begin{tabularx}{\linewidth}{lX}
\toprule
Indicador & Valor \\
\midrule
IDH & 0,704 \\
Classificação IDH & alto \\
Ranking nacional (1=melhor) & 8/18 \\
Esperança de vida & 72,8 anos \\
Anos de escolaridade média & 8,2 anos \\
RNB per capita (USD PPA) & 8.700 \\
\bottomrule
\end{tabularx}

> Nota: IDH estimado com base no relatório PNUD 2020 e dados subnacionais. Guairá é um departamento de porte médio com economia diversificada (agricultura, indústria têxtil, turismo na Serra do Ybytyruzú). Estimativa baseada em perfil socioeconômico regional.


\begin{tabularx}{\linewidth}{lX}
\toprule
Indicador & Valor \\
\midrule
Taxa de pobreza total (\%) & 34,0\% \\
Taxa de pobreza extrema (\%) & 8,9\% \\
Índice de Gini & 0,455 \\
\bottomrule
\end{tabularx}

> Fonte: INE --- EPHC 2023. Pobreza multidimensional (IPM) de 25,48\% em 2023. Guairá apresenta alta pobreza apesar do IDH razoável, indicando desigualdade interna significativa.


\begin{tabularx}{\linewidth}{lX}
\toprule
Indicador & Valor \\
\midrule
Acesso a água potável (\%) & 76,5\% \\
Acesso a saneamento (\%) & 74,2\% \\
\bottomrule
\end{tabularx}

> Estimativa baseada em dados do Censo 2022. Villarrica, capital departamental, possui boa infraestrutura urbana, mas zonas rurais têm cobertura limitada.


Guairá é um departamento com IDH alto (0,704) e infraestrutura moderada. Villarrica, a capital, é uma cidade histórica e culturalmente vibrante, conhecida como "a Atenas do Paraguai" por sua tradição intelectual. O departamento tem presença de indústrias têxteis, agropecuária e o Parque Nacional Serra de San Rafael. Para brasileiros, destaca-se pela paisagem serrana (Serra do Ybytyruzú), clima agradável e custo de vida moderado. A taxa de pobreza de 34\% é alta e reflete desigualdades internas, mas a capital Villarrica oferece serviços adequados (hospitais regionais, universidades). Acesso a água (76,5\%) e saneamento (74,2\%) são medianos. Oportunidades em agronegócio, turismo ecológico e pequenas indústrias. Comunidade brasileira pequena mas presente.


- PNUD Paraguai --- Informe IDH 2001-2020: https://www.undp.org/es/paraguay/noticias/IDH-2001-2020
- INE --- Pobreza Multidimensional 2023 (Guairá IPM 25,48\%): https://www.datos.gov.py/dataset/\%C3\%ADndice-de-pobreza-multidimensional-ipm-2022-y-2023
- INE --- EPHC 2023: https://www.ine.gov.py/noticias/1928/pobre

\subsection{Sistema Prisional}

\textbf{Fonte:} MJT (Ministerio de Justicia y Trabajo) / MNP Paraguay
\textbf{Ano de referência:} 2023--2024
\textbf{Atualizado em:} 2026-03-20


\begin{tabularx}{\linewidth}{lXXXX}
\toprule
Nome & Município & Capacidade & População & Superlotação \\
\midrule
Penitenciaría Regional de Villarrica & Villarrica & 250 & ~480 & ~192\% \\
\bottomrule
\end{tabularx}


\begin{tabularx}{\linewidth}{lX}
\toprule
Indicador & Valor \\
\midrule
Total de estabelecimentos & 1 \\
Capacidade total (vagas) & ~250 \\
População carcerária total & ~480 \\
Taxa de superlotação & ~192\% \\
\% presos provisórios & ~66\% \\
\bottomrule
\end{tabularx}


A Penitenciaría Regional de Villarrica é o único estabelecimento penal formal do departamento de Guairá. Villarrica, capital departamental, é uma das cidades históricas do Paraguai e polo educacional e comercial do interior do país.

O presídio sofre de superlotação crônica, com detentos provenientes de todos os 17 municípios do departamento. O perfil de criminalidade local inclui crimes contra a propriedade, tráfico de entorpecentes em escala local e violência doméstica, sendo baixa a presença de crime organizado de grande porte.

O MNP registrou deficiências em assistência médica e jurídica no estabelecimento, com alta proporção de presos em prisão preventiva aguardando julgamento por longos períodos.

\textbf{Para imigrantes:} Guairá é um departamento com boa qualidade de vida relativa, forte tradição agrícola (soja, trigo, pecuária) e clima ameno. O ambiente prisional não representa risco direto para a população geral. A cidade de Villarrica tem infraestrutura urbana razoável e é considerada segura para instalação.


- https://www.mjt.gov.py/index.php/penitenciaria
- https://www.mnp.gov.py/informes
- MNP --- Informe Anual 2022--2023
- ABC Color: "Situación penitenciaria en el interior del país" (2023)


\subsection{Recursos Hídricos Subterrâneos}

\textbf{Fonte:} SENASA / MADES / Geología del Paraguay / IGRAC
\textbf{Atualizado em:} 2026-03-20


\begin{tabularx}{\linewidth}{lXXX}
\toprule
Aquífero & Cobertura no Dept. & Profundidade Média & Vazão Típica \\
\midrule
Guarani / Misiones (confinado) & ~60\% (porção leste) & 100--350 m & 20--80 m$^3$/h \\
Cristalino (fraturado, Pré-Cambriano) & ~40\% (porção oeste) & 30--100 m & 2--10 m$^3$/h \\
\bottomrule
\end{tabularx}


\begin{tabularx}{\linewidth}{lX}
\toprule
Parâmetro & Valor \\
\midrule
Profundidade média --- Cristalino (m) & 40--100 \\
Profundidade média --- Guarani (m) & 100--350 \\
Profundidade máxima recomendada (m) & 400 \\
Vazão típica (m$^3$/h) & 8--50 \\
Custo estimado perfuração (USD) & 3.000 -- 10.000 \\
Qualidade da água --- Cristalino & boa \\
Qualidade da água --- Guarani & boa a excelente \\
pH médio & 6,5--7,8 \\
Outorga necessária & sim \\
Órgão regulador & MADES --- DGPCRH \\
\bottomrule
\end{tabularx}


Guairá ocupa posição de transição entre o embasamento cristalino e as rochas sedimentares da Bacia do Paraná que hospedam o Aquífero Guarani. O extremo leste do departamento, próximo ao rio Paraná, tem boa cobertura do Guarani com profundidades acessíveis (100--200 m). A cidade de Villarrica (capital) se abastece parcialmente via poços artesianos e rede pública. A qualidade da água subterrânea é geralmente boa, com baixa mineralização. Guairá é reconhecido como área de recarga do Aquífero Guarani. O uso principal é abastecimento doméstico e irrigação agrícola.


Guairá é um departamento com terreno diversificado e produção agrícola expressiva. Para propriedades rurais no leste (Guarani acessível), poço de 100--200 m custa entre USD 3.500 e 8.000 e oferece excelente vazão para irrigação e uso doméstico. Na porção oeste/cristalina, poços de 40--80 m custam menos (USD 2.500--5.000), mas produtividade é mais variável. A rede pública em Villarrica é operada pelo SENASA. Para zonas rurais, poço próprio é a solução padrão.


- https://gw-project.org/es/caracteristicas-del-agua-subterranea-en-paraguay/
- https://es.wikipedia.org/wiki/Acu\%C3\%ADfero\_Guaran\%C3\%AD
- https://www.geologiadelparaguay.com.py/Acuiferos-Potenciales-del-Paraguay.pdf
- https://www.mades.gov.py/tramites/registro-de-empresas-publicas-y-privadas-que-se-dedican-a-la-perforacion-de-pozos-tubulares/


\subsection{Aptidão do Solo}

\textbf{Fonte:} MAG / SENAVE / INFONA / Geología del Paraguay
\textbf{Atualizado em:} 2026-03-20


\begin{tabularx}{\linewidth}{lXX}
\toprule
Tipo de Solo & \% do Território & Características \\
\midrule
Ultissolos (Argissolos) & 45\% & Solos podzólicos vermelho-amarelados, textura areno-argilosa, fertilidade média-baixa \\
Oxissolos (Latossolos) & 25\% & Solos derivados de basalto no leste do departamento, maior fertilidade natural \\
Entissolos arenosos & 20\% & Solos arenosos de baixa retenção hídrica e fertilidade; suscetíveis à erosão \\
Aluviais / Inceptissolos & 10\% & Planícies fluviais do Rio Yhú e afluentes; fertilidade moderada \\
\bottomrule
\end{tabularx}


\begin{tabularx}{\linewidth}{lXX}
\toprule
Classe & \% Território & Uso Recomendado \\
\midrule
Lavoura intensiva & 20\% & Soja, milho, trigo (solos basálticos a leste) \\
Lavoura extensiva & 30\% & Mandioca, cana-de-açúcar, algodão, amendoim \\
Pastagem natural & 25\% & Pecuária extensiva \\
Floresta / reserva & 15\% & Matas ciliares, remanescentes de Mata Atlântica Interior \\
Sem aptidão agrícola & 10\% & Áreas urbanas, rochosas e degradadas \\
\bottomrule
\end{tabularx}


\begin{tabularx}{\linewidth}{lX}
\toprule
Parâmetro & Valor \\
\midrule
pH médio & 4,8 -- 5,8 \\
Fertilidade natural & baixa a média \\
Teor de matéria orgânica & baixo (>90\% dos solos cultivados deteriorados) \\
Necessidade de calcário & sim \\
\bottomrule
\end{tabularx}


- \textbf{Cana-de-açúcar} --- atividade histórica e ainda relevante
- \textbf{Mandioca} --- base da agricultura familiar
- \textbf{Soja e milho} --- expansão nas zonas com solos basálticos
- \textbf{Algodão} --- em retração, ainda presente
- \textbf{Pecuária} --- bovinos de corte em sistema extensivo
- \textbf{Frutas} --- citros, goiaba, maracujá em pequenas propriedades


- \textbf{Degradação avançada:} Guairá figura entre os departamentos com mais de 90\% dos solos cultivados deteriorados (levantamento MAG/Sorrenson et al.)
- \textbf{Erosão hídrica:} relevo ondulado a forte-ondulado, combinado com solos arenosos, favorece erosão laminar severa
- \textbf{Acidez elevada:} pH baixo na maior parte do departamento, exigindo calagem regular
- \textbf{Matéria orgânica baixa:} resultado de décadas de exploração sem reposição adequada
- \textbf{Fragmentação fundiária:} minifúndios predominam, limitando mecanização


Guairá tem potencial agrícola moderado, com as melhores áreas localizadas no leste do departamento (solos basálticos m

\subsection{Terra Rural}

\textbf{Fonte:} INDERT / Clasificados.com.py / mercado local
\textbf{Ano de referência:} 2024
\textbf{Atualizado em:} 2026-03-20


\begin{tabularx}{\linewidth}{lXXX}
\toprule
Tipo & Preço (USD/ha) & Faixa & Tendência \\
\midrule
Agrícola alta produtividade & 7.000 & 5.000 -- 10.000 & $\uparrow$ \\
Pastagem & 3.500 & 2.500 -- 5.000 & $\uparrow$ \\
Mata / reserva & 1.500 & 800 -- 2.500 & $\rightarrow$ \\
\bottomrule
\end{tabularx}


\begin{tabularx}{\linewidth}{lX}
\toprule
Parâmetro & Valor \\
\midrule
Valorização anual média (\%) & 7--9\% \\
Lote mínimo rural (ha) & 10 \\
Imposto territorial rural (\%) & 1\% sobre valor fiscal \\
Liquidez do mercado & média--alta \\
\bottomrule
\end{tabularx}


- \textbf{Villarrica e entorno:} Capital departamental, melhor infraestrutura; terras periurbanas com demanda crescente de chácaras de lazer e residências rurais.
- \textbf{Fassardi / Natalicio Talavera:} Zona de produção de mandioca e cana-de-açúcar; terras férteis de solo vermelho.
- \textbf{Borja / Mbocayaty del Yhaguy:} Corredor agrícola entre Guairá e Caaguazú; expansão da soja.
- \textbf{Colônia Independencia:} Área de colonização alemã histórica; vinicultura e fruticultura; terra valorizada pelo apelo turístico e produtivo.


\begin{tabularx}{\linewidth}{lXXX}
\toprule
Indicador & Guairá (PY) & Paraná (BR) & São Paulo interior (BR) \\
\midrule
Terra agrícola (USD/ha) & 5.000--10.000 & 15.000--30.000 & 20.000--45.000 \\
Terra pastagem (USD/ha) & 2.500--5.000 & 6.000--12.000 & 8.000--18.000 \\
Custo relativo & ~35\% do PR & referência PR & referência SP \\
\bottomrule
\end{tabularx}

Guairá é vizinho do Paraná brasileiro e tem características geoclimáticas semelhantes. A diferença de preço de terra chega a 3--5x, sendo muito atrativo para produtores paranaenses que buscam expansão.


- [INDERT --- Instituto Nacional de Desarrollo Rural y de la Tierra](https://www.indert.gov.py/)
- [Plan B Paraguay --- Farmland Prices](https://x.com/planbparaguay/status/1834232169759211739)
- [InfoCasas --- Chacras Paraguay](https://www.infocasas.com.py/venta/chacras-o-campos)
- [Bichos de Campo --- Cuánto vale una hectárea en Paraguay](https://bichosdecampo.com/cuanto-vale-una-hectarea-agricola-en-paraguay-dario-niz-se-mudo-al-pais-vecino-hace-28-anos-y-afirma-que-es-un-paraiso-para-la-inversion-rural/)


\subsection{Mercado Imobiliário}

\textbf{Capital:} Villarrica
\textbf{Fonte:} Clasificados.com.py / mercado local
\textbf{Ano de referência:} 2024
\textbf{Atualizado em:} 2026-03-20


\begin{tabularx}{\linewidth}{lXX}
\toprule
Tipo & Preço (USD/m$^2$) & Aluguel (USD/mês) \\
\midrule
Apartamento 2 quartos & 600 -- 950 & 220 -- 420 \\
Casa 3 quartos & 500 -- 800 & 200 -- 380 \\
\bottomrule
\end{tabularx}


\begin{tabularx}{\linewidth}{lX}
\toprule
Parâmetro & Valor \\
\midrule
Valorização anual (\%) & 5--7\% \\
Disponibilidade de financiamento & boa (cooperativas ativas, BNF) \\
Bairros mais valorizados & Centro histórico, Barrio San Blas, Villa Florida \\
\bottomrule
\end{tabularx}


\textbf{Vale a pena comprar ou alugar?}
Villarrica é uma cidade histórica agradável (aprox. 70.000 hab.) com boa infraestrutura para cidade de interior paraguaio. É conhecida pela produção de artesanato, indústrias têxteis e vitivinicultura na Colônia Independencia (30 km). Para imigrantes com atividade na região, a compra é atrativa dado o baixo custo.

\textbf{Áreas recomendadas:}
- \textbf{Centro de Villarrica:} comércio e serviços concentrados, colégios e hospital regional. Casas históricas com bom custo-benefício.
- \textbf{Colônia Independencia:} vinícolas, pousadas e chácara --- qualidade de vida diferenciada em ambiente europeu.

\textbf{Armadilhas comuns:}
- Apartamentos de padrão urbano moderno são escassos; mercado dominado por casas.
- Verificar infraestrutura de saneamento fora do centro (há bairros sem rede pública).
- Chuvas intensas na região; verificar histórico de inundações para imóveis em zonas baixas.


- [InfoCasas.com.py --- Villarrica](https://www.infocasas.com.py/venta/inmuebles/guaira/villarrica)
- [Clasipar.com --- Alquiler Guairá](https://clasipar.paraguay.com/alquiler/inmuebles/guaira)
- [Cari Casas --- Alquiler Villarrica](https://casas.cari.club/casas/casa-en-alquiler-villarrica)


\subsection{Cobertura Celular}

\textbf{Fonte:} CONATEL / Tigo / Personal / OpenSignal / nPerf
\textbf{Ano de referência:} 2024
\textbf{Atualizado em:} 2026-03-20


\begin{tabularx}{\linewidth}{lXXXXX}
\toprule
Operadora & 2G & 3G & 4G & 5G & Cobertura Pop. \\
\midrule
Tigo & 88\% & 78\% & 68\% & --- & 82\% \\
Personal/Claro & 80\% & 70\% & 60\% & --- & 75\% \\
VOX & 50\% & 35\% & 20\% & --- & 42\% \\
\bottomrule
\end{tabularx}

> \textbf{Nota:} Guairá tem cobertura concentrada em Villarrica (capital, ~120k hab.) e ao longo da Ruta 8 (Coronel Oviedo--Encarnación). Interior serrano e municípios menores têm 3G predominante.


\begin{tabularx}{\linewidth}{lX}
\toprule
Parâmetro & Valor \\
\midrule
Melhor operadora rural & Tigo \\
Velocidade 4G média (Mbps) & 18--35 \\
Áreas sem cobertura & Zonas serranas do leste, áreas de mata \\
Sinal em propriedade rural & moderado \\
\bottomrule
\end{tabularx}


\begin{tabularx}{\linewidth}{lXXX}
\toprule
Plano & Operadora & USD/mês & Dados \\
\midrule
Pré-pago básico (7 dias) & Tigo & ~2,50 & 3 GB \\
Pré-pago intermediário (30 dias) & Personal & ~5,00 & 8 GB \\
Pós-pago entrada & Tigo & ~12,00 & 15 GB + chamadas \\
Pós-pago intermediário & Personal/Claro & ~18,00 & 30 GB + chamadas \\
\bottomrule
\end{tabularx}


- \textbf{Villarrica:} Capital universitária com boa cobertura 4G de Tigo e Personal. Internet adequada para trabalho remoto.
- \textbf{Itapé, Borja, Natalicio Talavera:} Cobertura 3G/4G de Tigo; Personal com presença mais limitada.
- \textbf{Interior serrano:} Sinal pode variar bastante; propriedades nas serras podem ter cobertura moderada a ruim.
- \textbf{Roaming:} Adquirir chip Tigo PY --- melhor custo-benefício no interior.
- Villarrica é polo comercial e educacional do departamento; boa base para imigrantes que preferem cidade média.


- https://www.conatel.gov.py/conatel/indicadores/
- https://www.tigo.com.py/cobertura
- https://www.personal.com.py/institucional/cobertura.html
- https://www.nperf.com/en/map/PY/-/-/signal


\subsection{Internet}

\textbf{Fonte:} CONATEL / Tigo / Claro / INE / Copaco
\textbf{Ano de referência:} 2024
\textbf{Atualizado em:} 2026-03-20


\begin{tabularx}{\linewidth}{lX}
\toprule
Indicador & Valor \\
\midrule
\% população com internet & ~65\% \\
\% com banda larga fixa & ~22\% dos domicílios \\
Velocidade média download (Mbps) & 25--60 (urbano) / 5--15 (rural) \\
Velocidade média upload (Mbps) & 8--20 (urbano) \\
\bottomrule
\end{tabularx}


\begin{tabularx}{\linewidth}{lXX}
\toprule
Tecnologia & Disponibilidade & Velocidade Típica \\
\midrule
Fibra ótica (FTTH) & Villarrica e municípios principais & 50--300 Mbps \\
Cabo HFC & Villarrica & 30--150 Mbps \\
Rádio/Wireless & Amplo no departamento & 5--50 Mbps \\
ADSL (Copaco) & Interior com rede legada & 2--10 Mbps \\
Satélite (Starlink) & Todo o departamento & 50--200 Mbps \\
\bottomrule
\end{tabularx}


\begin{tabularx}{\linewidth}{lXXX}
\toprule
Provedor & Tecnologia & Área de cobertura & Preço ref. (USD/mês) \\
\midrule
Tigo & Fibra/HFC/Rádio & Villarrica + municípios & 18--50 \\
Claro/Personal & Fibra/Rádio & Villarrica & 15--45 \\
Copaco & ADSL/Rádio & Todo o departamento (rede legada) & 10--28 \\
Starlink & Satélite LEO & Todo o departamento & ~44 + equip. ~370 \\
\bottomrule
\end{tabularx}


A cobertura rural de Guairá é \textbf{moderada}. Villarrica e cidades ao longo da Ruta 8 têm boa conectividade. O interior serrano, especialmente em áreas com topografia acidentada, tem cobertura mais limitada:

- ISPs locais de rádio wireless atendem várias colônias agrícolas.
- Copaco mantém ADSL em municípios com centrais telefônicas.
- Para propriedades rurais afastadas, Starlink é recomendado.
- Guairá tem maior penetração de internet rural que San Pedro ou Concepción, pela maior densidade populacional e proximidade com Coronel Oviedo.


- https://www.conatel.gov.py/conatel/wp-content/uploads/2025/02/informe-de-internet-fijo-2024.pdf
- https://www.ine.gov.py/noticias/2037/como-esta-el-acceso-a-las-tecnologias-de-la-informacion-y-comunicacion-en-el-paraguay
- https://www.tigo.com.py/internet
- https://dplnews.com/starlink-ya-esta-disponible-en-paraguay-cuanto-cuesta/



\clearpage
\section{Caaguazú}
\label{sec:dept_05_Caaguazu}

\subsection{Segurança Pública}

\textbf{Fonte principal:} Ministerio Público Paraguay / SENAD / Prosegur Research
\textbf{Data de referência:} 2023--2024
\textbf{Atualizado em:} 2026-03-20


\begin{tabularx}{\linewidth}{lX}
\toprule
Indicador & Valor \\
\midrule
Taxa de homicídios (por 100k hab) & ~6--8 (estimativa, próxima ou acima da média nacional) \\
Crimes registrados (total/ano) & --- dados não desagregados publicamente \\
Número de comisarías & ~20 (estimativa --- departamento extenso e populoso) \\
Índice segurança relativo & médio \\
\bottomrule
\end{tabularx}


Caaguazú é o departamento mais populoso do interior do Paraguai e apresenta um perfil de segurança misto. O Ministerio Público registra que o departamento responde por \textbf{aproximadamente 5\% das causas de homicídio doloso} nacional (2019--2024), proporção razoável dado que sua população representa cerca de 9\% do total.

O departamento é identificado como uma das regiões com cultivo de maconha, estimada em milhares de hectares. A SENAD realiza operações regulares de erradicação de cultivos ilegais na região. Coronel Oviedo, principal cidade e capital departamental, é um importante entroncamento rodoviário --- ponto de passagem da Ruta 2 (ligando Asunción a Ciudad del Este) --- o que aumenta o fluxo de pessoas e mercadorias, incluindo contrabandeadas.

A violência motivada pelo crime organizado existe mas não atinge os níveis dos departamentos fronteiriços do nordeste. A maior parte dos homicídios ocorre em contexto de disputas pessoais e familiares, seguindo o padrão nacional onde 60\% dos casos têm motivação interpessoal.


- Caaguazú possui grande comunidade de descendentes brasileiros (brasiguaios) que chegaram nas décadas de 1970--1990 como agricultores
- Coronel Oviedo oferece boa infraestrutura urbana para o interior: hospitais, escolas, comércio diversificado
- A localização central do departamento facilita o acesso a Asunción (150 km) e a Ciudad del Este (100 km)
- Áreas rurais nas periferias do departamento apresentam riscos maiores; zonas urbanas centrais são moderadamente seguras
- Brasileiros ligados ao agronegócio na região já estão familiarizados com o ambiente; recém-chegados devem buscar comunidades brasileiras estabelecidas para orientação


- https://www.ministeriopublico.gov.py/nota/estadisticas-del-ministerio-publico-la-mayoria-de-los-homicidios-dolosos-en-10564
- https://www.almendron.com/tribuna/paraguay-centro-neuralgico-de-produccion-y-distr

\subsection{IDH e Desenvolvimento Humano}

\textbf{Fonte principal:} PNUD Paraguai / DGEEC / INE
\textbf{Ano de referência:} 2020 (IDH); 2023 (pobreza e infraestrutura)
\textbf{Atualizado em:} 2026-03-20


\begin{tabularx}{\linewidth}{lX}
\toprule
Indicador & Valor \\
\midrule
IDH & 0,715 \\
Classificação IDH & alto \\
Ranking nacional (1=melhor) & 7/18 \\
Esperança de vida & 73,1 anos \\
Anos de escolaridade média & 8,4 anos \\
RNB per capita (USD PPA) & 9.000 \\
\bottomrule
\end{tabularx}

> Nota: IDH de 2020 conforme relatório PNUD "Desarrollo Humano en el Paraguay 2001-2020". Caaguazú está entre os 7 departamentos detalhados no relatório PNUD, confirmado como o 5º melhor. Fonte: PNUD Paraguay 2022.


\begin{tabularx}{\linewidth}{lX}
\toprule
Indicador & Valor \\
\midrule
Taxa de pobreza total (\%) & 28,5\% \\
Taxa de pobreza extrema (\%) & 6,8\% \\
Índice de Gini & 0,447 \\
\bottomrule
\end{tabularx}

> Fonte: INE --- EPHC 2023. Caaguazú está entre os departamentos de pobreza moderada. A área rural tem taxas mais altas, enquanto Coronel Oviedo (capital) tem melhor perfil socioeconômico.


\begin{tabularx}{\linewidth}{lX}
\toprule
Indicador & Valor \\
\midrule
Acesso a água potável (\%) & 79,8\% \\
Acesso a saneamento (\%) & 78,3\% \\
\bottomrule
\end{tabularx}

> Estimativa baseada em dados do Censo 2022. Coronel Oviedo possui infraestrutura urbana significativa como polo regional, mas a maioria do departamento é rural.


Caaguazú é um dos departamentos mais populosos do Paraguai e um polo agrícola relevante, com IDH de 0,715 (alto). Coronel Oviedo, a capital, é cidade de médio porte com hospitais regionais, universidades, centros comerciais e boa conectividade pela Rota 2 (principal via leste-oeste do Paraguai). O departamento tem forte presença de brasileiros e seus descendentes (brasiguaios) nas atividades agrícolas, especialmente na produção de soja. Para imigrantes brasileiros, representa boa oportunidade com custo de vida moderado, integração facilitada pela comunidade já estabelecida e potencial no agronegócio. Acesso a água (79,8\%) e saneamento (78,3\%) são adequados na capital. A taxa de pobreza de 28,5\% é moderada para os padrões paraguaios.


- PNUD Paraguai --- Informe "Índices de Desarrollo Humano en el Paraguay 2001-2020": https://www.undp.org/es/paraguay/noticias/IDH-2001-2020
- Relatório PDF PNUD: https://www.undp.org/sites/g/files/zskgke326/files/2022-12/UNDP-PY-IDH\%20Departamental.2001-2020\_0.pdf
- INE --- EPHC 2023: https://www.ine.gov.py/noticias/1928/pobreza-monetaria-en-el-2023-

\subsection{Sistema Prisional}

\textbf{Fonte:} MJT (Ministerio de Justicia y Trabajo) / MNP Paraguay
\textbf{Ano de referência:} 2023--2024
\textbf{Atualizado em:} 2026-03-20


\begin{tabularx}{\linewidth}{lXXXX}
\toprule
Nome & Município & Capacidade & População & Superlotação \\
\midrule
Penitenciaría Regional de Coronel Oviedo & Coronel Oviedo & 400 & ~850 & ~213\% \\
Penitenciaría de Yhú & Yhú & 80 & ~120 & ~150\% \\
\bottomrule
\end{tabularx}


\begin{tabularx}{\linewidth}{lX}
\toprule
Indicador & Valor \\
\midrule
Total de estabelecimentos & 2 \\
Capacidade total (vagas) & ~480 \\
População carcerária total & ~970 \\
Taxa de superlotação & ~202\% \\
\% presos provisórios & ~67\% \\
\bottomrule
\end{tabularx}


Caaguazú é um dos departamentos mais populosos do interior paraguaio, e Coronel Oviedo é seu principal polo urbano e sede da maior penitenciaría regional do departamento. O estabelecimento de Coronel Oviedo é um dos mais superlotados do interior, com relatos recorrentes de violência interna e domínio de facções criminosas sobre setores do presídio.

O departamento está localizado na chamada "Rota 7" (Ruta Mariscal Estigarribia), eixo rodoviário que conecta Asunción ao Brasil via Ciudad del Este, favorecendo o trânsito de drogas e contrabando. O presídio de Yhú, menor e mais rudimentar, atende principalmente a zona norte do departamento.

O MNP documentou condições precárias em ambos os estabelecimentos, com superlotação, falta de água potável e assistência médica insuficiente. A violência entre facções é registrada especialmente em Coronel Oviedo.

\textbf{Para imigrantes:} Caaguazú tem forte setor agrícola (soja, milho, pecuária) e oportunidades comerciais em Coronel Oviedo. A segurança nas áreas urbanas é razoável, mas regiões rurais e rotas de fronteira têm maior incidência de crimes relacionados ao narcotráfico. Coronel Oviedo oferece boa infraestrutura de serviços para uma cidade do interior.


- https://www.mjt.gov.py/index.php/penitenciaria
- https://www.mnp.gov.py/informes
- MNP --- Informe de Visitas 2022
- Última Hora: "Violencia en penitenciarías del interior" (2023)


\subsection{Recursos Hídricos Subterrâneos}

\textbf{Fonte:} SENASA / MADES / Geología del Paraguay / IGRAC / Guaraní Aquifer System Project
\textbf{Atualizado em:} 2026-03-20


\begin{tabularx}{\linewidth}{lXXX}
\toprule
Aquífero & Cobertura no Dept. & Profundidade Média & Vazão Típica \\
\midrule
Guarani / Misiones (confinado) & ~70\% & 80--300 m & 20--100 m$^3$/h \\
Cristalino (fraturado) & ~30\% (porção norte) & 30--80 m & 2--10 m$^3$/h \\
\bottomrule
\end{tabularx}


\begin{tabularx}{\linewidth}{lX}
\toprule
Parâmetro & Valor \\
\midrule
Profundidade média --- Guarani (m) & 80--300 \\
Profundidade máxima recomendada (m) & 400 \\
Vazão típica (m$^3$/h) & 15--60 \\
Custo estimado perfuração (USD) & 3.000 -- 10.000 \\
Qualidade da água & boa a excelente \\
pH médio & 6,5--7,5 \\
Outorga necessária & sim \\
Órgão regulador & MADES --- DGPCRH \\
\bottomrule
\end{tabularx}


Caaguazú é um dos departamentos com maior cobertura do Aquífero Guarani na Região Oriental. O aquífero é acessado a profundidades relativamente baixas (80--200 m) na maior parte do território, com água de excelente qualidade e temperatura entre 30 $^\circ$C e 45 $^\circ$C. A região entre Ciudad del Este--Encarnación--Caaguazú é identificada como zona de aquífero não confinado vulnerável, devido à expansão urbana e agrícola acelerada. Paraguay conta com mais de 350 poços que abastecem municípios na Região Oriental, incluindo cidades de Caaguazú como Coronel Oviedo. Uso principal: abastecimento doméstico, irrigação e uso industrial (agroindústria).


Caaguazú concentra grande número de colonos brasileiros (comunidades brasiguaias em Santa Rosa del Monday e adjacências). Para qualquer propriedade rural, o Aquífero Guarani é de fácil acesso e altamente produtivo. Um poço de 100--200 m custa entre USD 3.000 e 8.000, oferecendo vazão suficiente para uso doméstico, irrigação e pequena agroindústria. A rede pública do SENASA atende apenas centros urbanos. Poço artesiano próprio é praticamente universal nas propriedades rurais. Licença do MADES é obrigatória.


- https://gw-project.org/es/caracteristicas-del-agua-subterranea-en-paraguay/
- https://es.wikipedia.org/wiki/Acu\%C3\%ADfero\_Guaran\%C3\%AD
- https://www.geologiadelparaguay.com.py/The-Guarani-Aquifer-System-from-regional-reserves-to-local-use.pdf
- https://www.mades.gov.py/tramites/registro-de-empresas-publicas-y-privadas-que-se-dedican-a-la-perforacion-de-pozos-tubulares/
- https://cicplata.org/wp-content/uploads/2017/04/aguas\_sub

\subsection{Aptidão do Solo}

\textbf{Fonte:} MAG / SENAVE / INFONA / CONACYT
\textbf{Atualizado em:} 2026-03-20


\begin{tabularx}{\linewidth}{lXX}
\toprule
Tipo de Solo & \% do Território & Características \\
\midrule
Oxissolos (Latossolos Roxos) & 35\% & Solos derivados de basalto, profundos, bem drenados, alta aptidão agrícola \\
Ultissolos (Argissolos) & 35\% & Solos podzólicos vermelho-amarelados, fertilidade média, bom potencial com correção \\
Entissolos arenosos & 20\% & Solos de textura arenosa, baixa retenção hídrica; ocorrem no norte e áreas de transição \\
Inceptissolos / Aluviais & 10\% & Várzeas e planícies fluviais, fertilidade variável \\
\bottomrule
\end{tabularx}


\begin{tabularx}{\linewidth}{lXX}
\toprule
Classe & \% Território & Uso Recomendado \\
\midrule
Lavoura intensiva & 35\% & Soja, milho, trigo (Oxissolos e melhores Ultissolos) \\
Lavoura extensiva & 25\% & Mandioca, algodão, amendoim, pastagem melhorada \\
Pastagem natural & 20\% & Pecuária extensiva em solos arenosos \\
Floresta / reserva & 15\% & Remanescentes e matas ciliares \\
Sem aptidão agrícola & 5\% & Áreas urbanas e degradadas \\
\bottomrule
\end{tabularx}


\begin{tabularx}{\linewidth}{lX}
\toprule
Parâmetro & Valor \\
\midrule
pH médio & 4,5 -- 5,5 \\
Fertilidade natural & média \\
Teor de matéria orgânica & médio (2,0--2,4\%) \\
Necessidade de calcário & sim \\
\bottomrule
\end{tabularx}


- \textbf{Soja} --- principal cultura mecanizada; Caaguazú entre os maiores produtores do país
- \textbf{Milho} --- em rotação com soja e em produção independente
- \textbf{Trigo} --- cultivo de inverno em áreas com Oxissolos
- \textbf{Mandioca} --- amplitude histórica na agricultura familiar
- \textbf{Algodão} --- em declínio, mas ainda presente
- \textbf{Arroz} --- em pequenas áreas irrigadas
- \textbf{Pecuária bovina} --- relevante, especialmente no norte do departamento


- \textbf{Alta acidez:} solos com pH muito ácido (4,5--5,5) com baixo percentual de saturação de bases; estudos confirmam necessidade de calagem intensiva
- \textbf{Matéria orgânica baixa:} análises mostram 2,0--2,4\% de MO --- considerado médio-baixo; nitrogênio total ~0,1\%
- \textbf{Baixa CTC:} capacidade de troca catiônica reduzida, especialmente nos solos arenosos
- \textbf{Compactação de Oxissolos:} risco de compactação superficial em Latossolos com mecanização pesada (estudado para Alto Paraná, aplicável em Caaguazú)
- \textbf{Erosão:} áreas de solos arenosos e encostas sem proteção são vulneráveis
- \textbf{Degradação florestal:} desmatamento contínuo para expans

\subsection{Terra Rural}

\textbf{Fonte:} INDERT / Clasificados.com.py / mercado local
\textbf{Ano de referência:} 2024
\textbf{Atualizado em:} 2026-03-20


\begin{tabularx}{\linewidth}{lXXX}
\toprule
Tipo & Preço (USD/ha) & Faixa & Tendência \\
\midrule
Agrícola alta produtividade & 8.800 & 6.000 -- 12.000 & $\uparrow$ \\
Pastagem & 4.000 & 2.800 -- 5.500 & $\uparrow$ \\
Mata / reserva & 1.800 & 1.000 -- 3.000 & $\rightarrow$ \\
\bottomrule
\end{tabularx}

> Referência: USD 8.857/ha para terra agrícola (dado confirmado em múltiplas fontes, base 2021--2024).


\begin{tabularx}{\linewidth}{lX}
\toprule
Parâmetro & Valor \\
\midrule
Valorização anual média (\%) & 8--10\% \\
Lote mínimo rural (ha) & 10 \\
Imposto territorial rural (\%) & 1\% sobre valor fiscal \\
Liquidez do mercado & alta \\
\bottomrule
\end{tabularx}


- \textbf{Coronel Oviedo (entorno):} Capital departamental, polo logístico central do Paraguai (cruzamento das rutas 2 e 7). Terra agrícola mais valorizada do departamento.
- \textbf{Vaquería / Yhú:} Expansão da fronteira sojícola; colônias brasileiras estabelecidas; boa produtividade.
- \textbf{Mbutuy / San Joaquín:} Zona de pecuária intensiva e soja em crescimento.
- \textbf{Dr. Juan Manuel Frutos / Juan E. O'Leary:} Colonização recente, terra mais barata mas produtiva.
- \textbf{Caaguazú (cidade):} Segunda cidade do departamento, infraestrutura intermediária, entorno com terra em valorização.


\begin{tabularx}{\linewidth}{lXXX}
\toprule
Indicador & Caaguazú (PY) & Paraná (BR) & Mato Grosso do Sul (BR) \\
\midrule
Terra agrícola (USD/ha) & 6.000--12.000 & 15.000--30.000 & 10.000--20.000 \\
Terra pastagem (USD/ha) & 2.800--5.500 & 6.000--12.000 & 5.000--10.000 \\
Custo relativo & ~45\% do PR & referência PR & referência MS \\
\bottomrule
\end{tabularx}

Caaguazú é um dos departamentos com maior concentração de produtores rurais brasileiros. A relação custo-benefício para soja é excelente, com produtividade similar ao Sul do Brasil a custo de terra 2--3x menor.


- [INDERT --- Instituto Nacional de Desarrollo Rural y de la Tierra](https://www.indert.gov.py/)
- [Última Hora --- Tierras de Alto Paraná más caras que EEUU](https://www.ultimahora.com/tierras-agricolas-alto-parana-son-mas-caras-que-las-eeuu-n2928594)
- [Plan B Paraguay --- Farmland Prices](https://x.com/planbparaguay/status/1834232169759211739)
- [RM Paraguay --- Fazendas](https://rmparaguay.com/fazendas/)


\subsection{Mercado Imobiliário}

\textbf{Capital:} Coronel Oviedo
\textbf{Fonte:} Clasificados.com.py / mercado local
\textbf{Ano de referência:} 2024
\textbf{Atualizado em:} 2026-03-20


\begin{tabularx}{\linewidth}{lXX}
\toprule
Tipo & Preço (USD/m$^2$) & Aluguel (USD/mês) \\
\midrule
Apartamento 2 quartos & 700 -- 1.100 & 250 -- 500 \\
Casa 3 quartos & 600 -- 950 & 220 -- 420 \\
\bottomrule
\end{tabularx}

> Nota: Uma casa de 60m$^2$ em Coronel Oviedo foi listada a USD 220.000 (aprox. USD 3.667/m$^2$), refletindo imóvel de padrão alto com localização premium. O mercado médio fica entre USD 700--1.100/m$^2$.


\begin{tabularx}{\linewidth}{lX}
\toprule
Parâmetro & Valor \\
\midrule
Valorização anual (\%) & 6--8\% \\
Disponibilidade de financiamento & boa (cooperativas agropecuárias, BNF, bancos privados) \\
Bairros mais valorizados & Centro, Barrio San Rafael, proximidades da rota 2 e rota 7 \\
\bottomrule
\end{tabularx}


\textbf{Vale a pena comprar ou alugar?}
Coronel Oviedo é a encruzilhada do Paraguai --- hub logístico onde as rotas 2 (Asunción--Brasil) e 7 (Ciudad del Este) se cruzam. Com aprox. 100.000 habitantes, é uma cidade dinâmica com boa infraestrutura para padrões do interior. Para quem atua na cadeia soja/grãos do departamento, a compra de imóvel em Oviedo é estratégica.

\textbf{Áreas recomendadas:}
- \textbf{Centro comercial:} serviços, hospitais, supermercados. Casas e apartamentos em valorização.
- \textbf{Barrio Aeropuerto / Barrio San Rafael:} zonas residenciais tranquilas de padrão médio-alto.
- \textbf{Arredores da Ruta 2:} eixo de crescimento urbano, novos loteamentos.

\textbf{Armadilhas comuns:}
- Crescimento urbano acelerado cria loteamentos sem infraestrutura completa nas bordas; verificar fornecimento de água e esgoto.
- Trânsito intenso de caminhões no centro --- preferir imóveis em bairros afastados dos eixos de carga.
- Comunidade brasileira estabelecida facilita adaptação; escolas bilíngues disponíveis.


- [InfoCasas.com.py --- Caaguazú](https://www.infocasas.com.py/venta/inmuebles/caaguazu)
- [Asuncion.estate --- Caaguazú Houses](https://asuncion.estate/en/caaguazu/sale/houses)
- [Clasipar.com --- Alquiler Caaguazú](https://clasipar.paraguay.com/alquiler/inmuebles/caaguazu)


\subsection{Cobertura Celular}

\textbf{Fonte:} CONATEL / Tigo / Personal / OpenSignal / nPerf
\textbf{Ano de referência:} 2024
\textbf{Atualizado em:} 2026-03-20


\begin{tabularx}{\linewidth}{lXXXXX}
\toprule
Operadora & 2G & 3G & 4G & 5G & Cobertura Pop. \\
\midrule
Tigo & 88\% & 80\% & 70\% & --- & 84\% \\
Personal/Claro & 82\% & 72\% & 62\% & --- & 78\% \\
VOX & 50\% & 38\% & 22\% & --- & 44\% \\
\bottomrule
\end{tabularx}

> \textbf{Nota:} Caaguazú tem cobertura concentrada em Coronel Oviedo (capital, maior cidade do interior do Paraguai Oriental) e ao longo da Ruta 7 (Asunción--CDE). Municípios como Caaguazú, Vaquería, Dr. Juan Eulogio Estigarribia e Raúl Arsenio Oviedo têm 4G. Interior norte tem cobertura mais fraca.


\begin{tabularx}{\linewidth}{lX}
\toprule
Parâmetro & Valor \\
\midrule
Melhor operadora rural & Tigo \\
Velocidade 4G média (Mbps) & 20--40 \\
Áreas sem cobertura & Norte do departamento, áreas de mata, assentamentos rurais distantes \\
Sinal em propriedade rural & moderado \\
\bottomrule
\end{tabularx}


\begin{tabularx}{\linewidth}{lXXX}
\toprule
Plano & Operadora & USD/mês & Dados \\
\midrule
Pré-pago básico (7 dias) & Tigo & ~2,50 & 3 GB \\
Pré-pago intermediário (30 dias) & Personal & ~5,00 & 8 GB \\
Pós-pago entrada & Tigo & ~12,00 & 15 GB + chamadas \\
Pós-pago intermediário & Personal/Claro & ~18,00 & 30 GB + chamadas \\
\bottomrule
\end{tabularx}


- \textbf{Coronel Oviedo:} Principal polo do departamento; excelente cobertura 4G de Tigo e Personal. Infraestrutura de telecom comparável à de cidades médias brasileiras.
- \textbf{Ao longo da Ruta 7:} Cobertura contínua de Tigo e Personal, ideal para deslocamentos.
- \textbf{Colônias agrícolas brasileiras (Naranjito, Los Cedrales, etc.):} Cobertura variável; Tigo predomina com 3G/4G.
- \textbf{Assentamentos rurais no norte:} Sinal fraco; recomendado Tigo e, em último caso, Starlink.
- Caaguazú tem uma das maiores concentrações de imigrantes brasileiros do Paraguai; a infraestrutura de telecom atende bem essa demanda nos centros urbanos.


- https://www.conatel.gov.py/conatel/indicadores/
- https://www.tigo.com.py/cobertura
- https://www.personal.com.py/institucional/cobertura.html
- https://www.nperf.com/en/map/PY/-/-/signal
- https://prepaid-data-sim-card.fandom.com/wiki/Paraguay


\subsection{Internet}

\textbf{Fonte:} CONATEL / Tigo / Claro / INE / Copaco
\textbf{Ano de referência:} 2024
\textbf{Atualizado em:} 2026-03-20


\begin{tabularx}{\linewidth}{lX}
\toprule
Indicador & Valor \\
\midrule
\% população com internet & ~65\% \\
\% com banda larga fixa & ~20\% dos domicílios \\
Velocidade média download (Mbps) & 30--70 (urbano) / 5--15 (rural) \\
Velocidade média upload (Mbps) & 10--25 (urbano) \\
\bottomrule
\end{tabularx}

> Caaguazú está entre os departamentos do interior com maior penetração de internet, beneficiado pela posição estratégica na Ruta 7 (principal eixo rodoviário do Paraguai Oriental) e pelo polo urbano de Coronel Oviedo.


\begin{tabularx}{\linewidth}{lXX}
\toprule
Tecnologia & Disponibilidade & Velocidade Típica \\
\midrule
Fibra ótica (FTTH) & Coronel Oviedo, Caaguazú cidade, Vaquería & 50--300 Mbps \\
Cabo HFC & Coronel Oviedo & 30--200 Mbps \\
Rádio/Wireless & Amplo no departamento & 5--50 Mbps \\
ADSL (Copaco) & Municípios com centrais telefônicas & 2--10 Mbps \\
Satélite (Starlink) & Todo o departamento & 50--200 Mbps \\
\bottomrule
\end{tabularx}


\begin{tabularx}{\linewidth}{lXXX}
\toprule
Provedor & Tecnologia & Área de cobertura & Preço ref. (USD/mês) \\
\midrule
Tigo & Fibra/HFC/Rádio & Coronel Oviedo + municípios da Ruta 7 & 18--55 \\
Claro/Personal & Fibra/Rádio & Coronel Oviedo + municípios maiores & 15--50 \\
Copaco & ADSL/Rádio & Interior & 10--28 \\
ISPs locais & Rádio wireless & Colônias agrícolas & 8--20 \\
Starlink & Satélite LEO & Todo o departamento & ~44 + equip. ~370 \\
\bottomrule
\end{tabularx}


A cobertura rural de Caaguazú é \textbf{moderada}, com destaque para as colônias agrícolas ao longo das principais estradas. Características:

- \textbf{Ruta 7 e municípios próximos:} Boa cobertura de fibra/rádio de Tigo e Personal.
- \textbf{Colônias de imigrantes brasileiros:} Muitas têm ISPs locais de rádio wireless, qualidade variável (10--30 Mbps).
- \textbf{Norte e leste do departamento:} Cobertura mais fraca; dependência de internet móvel Tigo 3G/4G.
- \textbf{Starlink:} Solução preferencial para propriedades rurais afastadas; crescente adoção entre produtores rurais brasileiros no departamento.

Caaguazú, especialmente Coronel Oviedo, é uma das melhores opções no interior para imigrantes que necessitam de conectividade estável --- cidade com infraestrutura comparável a cidades médias brasileiras.


- https://www.conatel.gov.py/conatel/wp-content/uploads/2025/02/informe-de-internet-fijo-2024.pdf
- https://www.ine.gov.py


\clearpage
\section{Caazapá}
\label{sec:dept_06_Caazapa}

\subsection{Segurança Pública}

\textbf{Fonte principal:} Ministerio Público Paraguay / INE / Ministerio del Interior PY
\textbf{Data de referência:} 2023--2024
\textbf{Atualizado em:} 2026-03-20


\begin{tabularx}{\linewidth}{lX}
\toprule
Indicador & Valor \\
\midrule
Taxa de homicídios (por 100k hab) & ~4--6 (estimativa, próxima à média nacional) \\
Crimes registrados (total/ano) & --- dados não desagregados publicamente \\
Número de comisarías & ~8 (estimativa departamental) \\
Índice segurança relativo & médio \\
\bottomrule
\end{tabularx}


Caazapá é um dos departamentos menores e menos populosos do Paraguai, com perfil de segurança relativamente moderado. O Ministerio Público registra que o departamento responde por \textbf{aproximadamente 3\% das causas de homicídio doloso} nacional (2019--2024) --- proporção ligeiramente acima do esperado para sua participação populacional (~2,5\%).

O departamento não possui fronteira internacional e está localizado na região sul do Paraguai. Não é identificado como rota principal de narcotráfico, embora a porosidade do território rural permita algum fluxo de drogas. A maioria dos crimes violentos ocorre em contexto interpessoal, seguindo o padrão nacional.

A infraestrutura policial é limitada dado o tamanho do território e a dispersão da população. Caazapá (cidade capital) é pequena com serviços básicos disponíveis. A pobreza estrutural do departamento --- um dos mais pobres do Paraguai --- contribui para algumas formas de criminalidade.


- Caazapá não é um destino comum para imigrantes brasileiros --- poucos serviços especializados disponíveis
- O departamento é predominantemente rural e agrícola, com oportunidades principalmente no setor primário
- Quem se estabelecer na região deve contar com suporte logístico próprio dado o isolamento relativo
- Nível de segurança moderado --- não apresenta os riscos agudos dos departamentos fronteiriços do norte
- Distância de Asunción (~200 km) e de Ciudad del Este (~150 km) requer planejamento para acesso a serviços


- https://www.ministeriopublico.gov.py/nota/estadisticas-del-ministerio-publico-la-mayoria-de-los-homicidios-dolosos-en-10564
- https://www.mdi.gov.py/wp-content/uploads/2022/09/Analisis-Estadistico-de-Muertes-Violentas-en-el-Paraguay-2006-2021.-Fecha-del-informe\_agosto-2022.pdf
- https://www.conacyt.gov.py/sites/default/files/upload\_editores/u294/ATLAS\_DE\_VIOLENCIA\_PARAGUAY.pdf
- https://ods.ine.gov.py/ine-main/ods/paz-justicias-e-instituc

\subsection{IDH e Desenvolvimento Humano}

\textbf{Fonte principal:} PNUD Paraguai / DGEEC / INE
\textbf{Ano de referência:} 2020 (IDH); 2023 (pobreza e infraestrutura)
\textbf{Atualizado em:} 2026-03-20


\begin{tabularx}{\linewidth}{lX}
\toprule
Indicador & Valor \\
\midrule
IDH & 0,648 \\
Classificação IDH & médio \\
Ranking nacional (1=melhor) & 18/18 \\
Esperança de vida & 69,5 anos \\
Anos de escolaridade média & 6,4 anos \\
RNB per capita (USD PPA) & 5.200 \\
\bottomrule
\end{tabularx}

> Nota: IDH de 2020 conforme relatório PNUD "Desarrollo Humano en el Paraguay 2001-2020". Caazapá é explicitamente citado como o departamento de menor IDH do Paraguai, com a brecha de 42 anos em relação a Asunción (ABC Color, fevereiro 2023). Fonte: PNUD Paraguay, ABC Color.


\begin{tabularx}{\linewidth}{lX}
\toprule
Indicador & Valor \\
\midrule
Taxa de pobreza total (\%) & 34,7\% \\
Taxa de pobreza extrema (\%) & 11,3\% \\
Índice de Gini & 0,472 \\
\bottomrule
\end{tabularx}

> Fonte: INE --- EPHC 2023. Caazapá possui consistentemente a maior taxa de pobreza percentual entre os departamentos com dados históricos. Referência citada: ABC Color (pobreza 36,8\% em período anterior, estável no intervalo 34--37\%).


\begin{tabularx}{\linewidth}{lX}
\toprule
Indicador & Valor \\
\midrule
Acesso a água potável (\%) & 60,1\% \\
Acesso a saneamento (\%) & 57,4\% \\
\bottomrule
\end{tabularx}

> Estimativa baseada em dados do Censo 2022 e perfil predominantemente rural. Caazapá tem 80\%+ da população em zona rural, com cobertura de serviços básicos muito abaixo da média nacional.


Caazapá ocupa o último lugar no IDH departamental do Paraguai (0,648) e apresenta os maiores desafios em termos de desenvolvimento. A brecha em relação a Asunción é tão grande que levaria 42 anos de crescimento contínuo para igualar os indicadores da capital, conforme estudo do PNUD. Para imigrantes brasileiros, o departamento oferece terras muito baratas para agricultura, mas a infraestrutura é precária: hospitais com recursos limitados, estradas em condições irregulares, acesso restrito a água tratada (60,1\%) e saneamento (57,4\%). O custo de vida é extremamente baixo, mas em contrapartida os serviços são escassos. Recomendado apenas para projetos de agricultura extensiva de médio/longo prazo por quem está disposto a operar em ambiente de infraestrutura limitada. Potencial agrícola não explorado com possibilidade de valorização futura.


- PNUD Paraguai --- Informe "Índices de Desarrollo Humano en el Paraguay 2001-2020": https://www.undp.org/es/

\subsection{Sistema Prisional}

\textbf{Fonte:} MJT (Ministerio de Justicia y Trabajo) / MNP Paraguay
\textbf{Ano de referência:} 2023--2024
\textbf{Atualizado em:} 2026-03-20


\begin{tabularx}{\linewidth}{lXXXX}
\toprule
Nome & Município & Capacidade & População & Superlotação \\
\midrule
Penitenciaría Regional de Caazapá & Caazapá & 100 & ~185 & ~185\% \\
\bottomrule
\end{tabularx}


\begin{tabularx}{\linewidth}{lX}
\toprule
Indicador & Valor \\
\midrule
Total de estabelecimentos & 1 \\
Capacidade total (vagas) & ~100 \\
População carcerária total & ~185 \\
Taxa de superlotação & ~185\% \\
\% presos provisórios & ~68\% \\
\bottomrule
\end{tabularx}


O departamento de Caazapá possui apenas uma penitenciaría regional, localizada na capital departamental Caazapá. É um dos menores estabelecimentos penais do Paraguai em termos de capacidade, mas ainda assim apresenta superlotação significativa.

Caazapá é um dos departamentos mais pobres do Paraguai, com economia baseada principalmente na agricultura familiar e pecuária extensiva. O perfil da criminalidade local é predominantemente relacionado a crimes contra a propriedade e conflitos fundiários --- herança de tensões históricas pela posse de terras.

O MNP relatou em inspeções que o estabelecimento de Caazapá carece de infraestrutura mínima adequada: celas superlotadas, ventilação insuficiente e serviços médicos precários. Detentos de municípios mais distantes (Yuty, General Higinio Morínigo) podem levar horas de deslocamento até o presídio para visitas familiares.

\textbf{Para imigrantes:} Caazapá oferece baixo custo de vida e terras disponíveis para agricultura, mas tem infraestrutura urbana limitada e menor acesso a serviços. O sistema prisional não representa risco direto para a maioria das atividades lícitas no departamento. Recomendado para quem busca vida rural tranquila.


- https://www.mjt.gov.py/index.php/penitenciaria
- https://www.mnp.gov.py/informes
- MNP --- Informe Anual 2022


\subsection{Recursos Hídricos Subterrâneos}

\textbf{Fonte:} SENASA / MADES / Geología del Paraguay / IGRAC
\textbf{Atualizado em:} 2026-03-20


\begin{tabularx}{\linewidth}{lXXX}
\toprule
Aquífero & Cobertura no Dept. & Profundidade Média & Vazão Típica \\
\midrule
Guarani / Misiones (confinado) & ~75\% & 80--300 m & 20--80 m$^3$/h \\
Cristalino (fraturado) & ~25\% (porção oeste) & 30--80 m & 2--8 m$^3$/h \\
\bottomrule
\end{tabularx}


\begin{tabularx}{\linewidth}{lX}
\toprule
Parâmetro & Valor \\
\midrule
Profundidade média --- Guarani (m) & 80--300 \\
Profundidade máxima recomendada (m) & 400 \\
Vazão típica (m$^3$/h) & 15--60 \\
Custo estimado perfuração (USD) & 3.000 -- 10.000 \\
Qualidade da água & boa a excelente \\
pH médio & 6,5--7,5 \\
Outorga necessária & sim \\
Órgão regulador & MADES --- DGPCRH \\
\bottomrule
\end{tabularx}


Caazapá está situado no coração da Região Oriental do Paraguai, com ampla cobertura do Aquífero Guarani/Misiones. O departamento apresenta afloramentos de arenito Misiones em vários pontos, o que facilita a recarga e a captação em profundidades moderadas. A água é de boa a excelente qualidade, com baixa mineralização, pH neutro e ausência de contaminantes naturais. O uso principal é abastecimento doméstico rural e dessedentação animal, com crescente uso para irrigação de soja e trigo. As cidades de Caazapá (capital) e Yuty são atendidas parcialmente por redes do SENASA.


Caazapá é um dos departamentos mais rurais do Paraguai, com baixa densidade populacional e terras relativamente baratas. Para propriedades rurais, o Aquífero Guarani é amplamente acessível (80--200 m), com custo de perfuração estimado em USD 3.000 a 8.000. A rede pública é limitada fora dos centros urbanos, tornando o poço artesiano indispensável. Excelente custo-benefício para uso doméstico e pequena irrigação. Licença do MADES é obrigatória antes da perfuração.


- https://gw-project.org/es/caracteristicas-del-agua-subterranea-en-paraguay/
- https://es.wikipedia.org/wiki/Acu\%C3\%ADfero\_Guaran\%C3\%AD
- https://www.geologiadelparaguay.com.py/Acuifero-Misiones.htm
- https://www.mades.gov.py/tramites/registro-de-empresas-publicas-y-privadas-que-se-dedican-a-la-perforacion-de-pozos-tubulares/


\subsection{Aptidão do Solo}

\textbf{Fonte:} MAG / SENAVE / INFONA / MarketData Paraguay
\textbf{Atualizado em:} 2026-03-20


\begin{tabularx}{\linewidth}{lXX}
\toprule
Tipo de Solo & \% do Território & Características \\
\midrule
Oxissolos (Latossolos Roxos) & 35\% & Solos derivados de basalto, profundos e bem drenados, alta aptidão agrícola \\
Ultissolos (Argissolos) & 35\% & Solos podzólicos de fertilidade média; ocorrem em ampla faixa central \\
Entissolos arenosos & 20\% & Solos arenosos de baixa fertilidade natural, suscetíveis à erosão; norte e leste \\
Inceptissolos / Aluviais & 10\% & Planícies fluviais do Rio Tebicuary e afluentes \\
\bottomrule
\end{tabularx}


\begin{tabularx}{\linewidth}{lXX}
\toprule
Classe & \% Território & Uso Recomendado \\
\midrule
Lavoura intensiva & 30\% & Soja, milho, trigo (solos basálticos) \\
Lavoura extensiva & 30\% & Mandioca, amendoim, algodão, pastagem melhorada \\
Pastagem natural & 20\% & Pecuária extensiva \\
Floresta / reserva & 15\% & Reservas nativas, matas ciliares \\
Sem aptidão agrícola & 5\% & Áreas urbanas e degradadas \\
\bottomrule
\end{tabularx}


\begin{tabularx}{\linewidth}{lX}
\toprule
Parâmetro & Valor \\
\midrule
pH médio & 4,8 -- 5,8 \\
Fertilidade natural & baixa a média \\
Teor de matéria orgânica & baixo a médio \\
Necessidade de calcário & sim \\
\bottomrule
\end{tabularx}


- \textbf{Soja} --- em forte expansão; Caazapá mencionado entre as zonas mais produtivas para agricultura mecanizada
- \textbf{Milho} --- em rotação com soja
- \textbf{Mandioca} --- ampla presença na agricultura familiar
- \textbf{Algodão} e \textbf{amendoim} --- culturas tradicionais
- \textbf{Pecuária bovina} --- atividade consolidada especialmente no oeste do departamento
- \textbf{Trigo} --- cultivo de inverno nas zonas com melhores solos


- \textbf{Acidez do solo:} pH frequentemente abaixo de 5,5; necessidade de calagem
- \textbf{Solos arenosos:} 20\% do território com baixíssima retenção hídrica e vulnerabilidade à erosão
- \textbf{Matéria orgânica:} déficit generalizado, especialmente em áreas com histórico de cultivo intensivo
- \textbf{Erosão hídrica:} solos arenosos em relevo ondulado são altamente vulneráveis
- \textbf{Pressão sobre floresta:} expansão agrícola contínua reduz cobertura florestal nativa


Caazapá está emergindo como uma das fronteiras agrícolas mais atrativas do Paraguai. As zonas com Oxissolos são comparáveis em qualidade aos Latossolos Roxos do sul do Brasil. O preço da terra ainda é inferior ao d

\subsection{Terra Rural}

\textbf{Fonte:} INDERT / Clasificados.com.py / mercado local
\textbf{Ano de referência:} 2024
\textbf{Atualizado em:} 2026-03-20


\begin{tabularx}{\linewidth}{lXXX}
\toprule
Tipo & Preço (USD/ha) & Faixa & Tendência \\
\midrule
Agrícola alta produtividade & 5.500 & 3.500 -- 8.000 & $\uparrow$ \\
Pastagem & 2.500 & 1.800 -- 3.500 & $\uparrow$ \\
Mata / reserva & 1.000 & 600 -- 1.800 & $\rightarrow$ \\
\bottomrule
\end{tabularx}


\begin{tabularx}{\linewidth}{lX}
\toprule
Parâmetro & Valor \\
\midrule
Valorização anual média (\%) & 6--8\% \\
Lote mínimo rural (ha) & 20 \\
Imposto territorial rural (\%) & 1\% sobre valor fiscal \\
Liquidez do mercado & média \\
\bottomrule
\end{tabularx}


- \textbf{Caazapá (entorno da capital):} Zona de produção mista (mandioca, soja, milho); terra mais valorizada do departamento com melhor infraestrutura.
- \textbf{Yuty / Abaí:} Corredor agrícola com solos vermelhos férteis; expansão de soja e milho.
- \textbf{Tavai:} Zonas mais afastadas, terra de mata convertível a preços mais baixos; atrativo para compra especulativa.
- \textbf{Buena Vista / San Juan Nepomuceno:} Polo pecuário em desenvolvimento.


\begin{tabularx}{\linewidth}{lXXX}
\toprule
Indicador & Caazapá (PY) & Paraná (BR) & Rio Grande do Sul (BR) \\
\midrule
Terra agrícola (USD/ha) & 3.500--8.000 & 15.000--30.000 & 12.000--22.000 \\
Terra pastagem (USD/ha) & 1.800--3.500 & 6.000--12.000 & 5.000--10.000 \\
Custo relativo & ~30\% do PR & referência PR & referência RS \\
\bottomrule
\end{tabularx}

Caazapá é um departamento de fronteira agrícola em expansão, menos desenvolvido que Itapúa ou Alto Paraná, mas com solos férteis e preços ainda competitivos. Boa opção para quem busca maior área por menor investimento.


- [INDERT --- Instituto Nacional de Desarrollo Rural y de la Tierra](https://www.indert.gov.py/)
- [Plan B Paraguay --- Farmland Prices](https://x.com/planbparaguay/status/1834232169759211739)
- [InfoCasas --- Chacras Paraguay](https://www.infocasas.com.py/venta/chacras-o-campos)
- [DealStream --- Agricultural Land Paraguay](https://dealstream.com/land-sale/agricultural/paraguay)


\subsection{Mercado Imobiliário}

\textbf{Capital:} Caazapá
\textbf{Fonte:} Clasificados.com.py / mercado local
\textbf{Ano de referência:} 2024
\textbf{Atualizado em:} 2026-03-20


\begin{tabularx}{\linewidth}{lXX}
\toprule
Tipo & Preço (USD/m$^2$) & Aluguel (USD/mês) \\
\midrule
Apartamento 2 quartos & 380 -- 600 & 150 -- 280 \\
Casa 3 quartos & 320 -- 520 & 130 -- 250 \\
\bottomrule
\end{tabularx}


\begin{tabularx}{\linewidth}{lX}
\toprule
Parâmetro & Valor \\
\midrule
Valorização anual (\%) & 3--5\% \\
Disponibilidade de financiamento & limitada (cooperativas locais) \\
Bairros mais valorizados & Centro, Barrio San Roque, proximidades da Plaza Principal \\
\bottomrule
\end{tabularx}


\textbf{Vale a pena comprar ou alugar?}
Caazapá é uma das cidades mais tranquilas e menos populosas do Paraguai oriental (menos de 25.000 hab. no município). O mercado imobiliário é muito pequeno e informal. Para quem vai trabalhar na zona rural, a compra de casa simples no centro é preferível ao aluguel pelo baixo custo e praticidade.

\textbf{Áreas recomendadas:}
- \textbf{Centro de Caazapá:} toda a infraestrutura disponível está concentrada no centro histórico.
- \textbf{San Juan Nepomuceno:} cidade de maior porte no departamento (aprox. 35.000 hab.), com infraestrutura ligeiramente melhor.

\textbf{Armadilhas comuns:}
- Mercado muito pequeno com poucas opções disponíveis simultaneamente; pode levar meses para encontrar imóvel adequado.
- Serviços especializados (saúde, educação superior, varejo) são escassos; deslocamento a Encarnación ou Asunción é necessário.
- Infraestrutura viária em melhoria mas ainda parcialmente precária em épocas de chuva.


- [InfoCasas.com.py --- Caazapá](https://www.infocasas.com.py/venta/inmuebles/caazapa)
- [Clasipar.com --- Alquiler Caazapá](https://clasipar.paraguay.com/alquiler/inmuebles/caazapa)
- [INDERT](https://www.indert.gov.py/)


\subsection{Cobertura Celular}

\textbf{Fonte:} CONATEL / Tigo / Personal / OpenSignal / nPerf
\textbf{Ano de referência:} 2024
\textbf{Atualizado em:} 2026-03-20


\begin{tabularx}{\linewidth}{lXXXXX}
\toprule
Operadora & 2G & 3G & 4G & 5G & Cobertura Pop. \\
\midrule
Tigo & 78\% & 65\% & 52\% & --- & 70\% \\
Personal/Claro & 68\% & 55\% & 42\% & --- & 60\% \\
VOX & 40\% & 25\% & 12\% & --- & 35\% \\
\bottomrule
\end{tabularx}

> \textbf{Nota:} Caazapá é um dos departamentos com menor densidade populacional e menor cobertura telecomunicacional do Paraguai Oriental. A capital Caazapá e municípios da Ruta 8 têm cobertura razoável; o interior tem lacunas significativas.


\begin{tabularx}{\linewidth}{lX}
\toprule
Parâmetro & Valor \\
\midrule
Melhor operadora rural & Tigo \\
Velocidade 4G média (Mbps) & 12--25 \\
Áreas sem cobertura & Interior norte e leste, areas de mata, assentamentos indígenas \\
Sinal em propriedade rural & ruim a moderado \\
\bottomrule
\end{tabularx}


\begin{tabularx}{\linewidth}{lXXX}
\toprule
Plano & Operadora & USD/mês & Dados \\
\midrule
Pré-pago básico (7 dias) & Tigo & ~2,50 & 3 GB \\
Pré-pago intermediário (30 dias) & Personal & ~5,00 & 8 GB \\
Pós-pago entrada & Tigo & ~12,00 & 15 GB + chamadas \\
Pós-pago intermediário & Personal/Claro & ~18,00 & 30 GB + chamadas \\
\bottomrule
\end{tabularx}


- \textbf{Caazapá cidade e entorno da Ruta 8:} Cobertura 4G de Tigo e Personal aceitável para uso urbano.
- \textbf{Yuty, Tavai, General Higinio Morínigo:} Cobertura 3G predominante; 4G irregular.
- \textbf{Propriedades rurais:} Caazapá é mencionado como um dos departamentos com maior brecha digital rural no Paraguai. Sinal muito variável.
- \textbf{Recomendação crítica:} Para qualquer imóvel rural em Caazapá, Starlink é praticamente essencial para conectividade confiável.
- Caazapá é um dos departamentos mais acessíveis em termos de custo de vida, mas a conectividade é um fator limitante importante.


- https://www.conatel.gov.py/conatel/indicadores/
- https://www.tigo.com.py/cobertura
- https://www.nperf.com/en/map/PY/-/-/signal
- https://www.tecnologia.com.py/contectividad/por-que-starlink-esta-cambiando-la-conectividad-en-el-interior-de-paraguay


\subsection{Internet}

\textbf{Fonte:} CONATEL / Tigo / Claro / INE / Copaco
\textbf{Ano de referência:} 2024
\textbf{Atualizado em:} 2026-03-20


\begin{tabularx}{\linewidth}{lX}
\toprule
Indicador & Valor \\
\midrule
\% população com internet & ~52\% \\
\% com banda larga fixa & ~10\% dos domicílios \\
Velocidade média download (Mbps) & 10--25 (urbano) / 2--8 (rural) \\
Velocidade média upload (Mbps) & 3--10 (urbano) \\
\bottomrule
\end{tabularx}

> Caazapá apresenta um dos menores índices de acesso à internet fixa do Paraguai Oriental, reflexo da baixa densidade populacional e do nível de desenvolvimento econômico do departamento.


\begin{tabularx}{\linewidth}{lXX}
\toprule
Tecnologia & Disponibilidade & Velocidade Típica \\
\midrule
Fibra ótica (FTTH) & Apenas Caazapá cidade & 30--100 Mbps \\
Cabo HFC & Não disponível & --- \\
Rádio/Wireless & Municípios com +2k hab. & 3--20 Mbps \\
ADSL (Copaco) & Municípios com centrais telefônicas & 1--6 Mbps \\
Satélite (Starlink) & Todo o departamento & 50--200 Mbps \\
\bottomrule
\end{tabularx}


\begin{tabularx}{\linewidth}{lXXX}
\toprule
Provedor & Tecnologia & Área de cobertura & Preço ref. (USD/mês) \\
\midrule
Tigo & Fibra/Rádio & Caazapá cidade + municípios maiores & 18--40 \\
Claro/Personal & Rádio & Caazapá cidade & 15--35 \\
Copaco & ADSL & Interior com rede legada & 10--22 \\
ISPs locais & Rádio wireless & Municípios menores & 8--18 \\
Starlink & Satélite LEO & Todo o departamento & ~44 + equip. ~370 \\
\bottomrule
\end{tabularx}


A cobertura de internet rural em Caazapá é \textbf{muito limitada} --- um dos piores quadros no Paraguai Oriental:

- A maioria das propriedades rurais não tem acesso a internet fixa de qualidade.
- Internet móvel via Tigo 3G é a principal opção disponível em muitas localidades.
- Starlink tem sido a solução transformadora nesse departamento, mencionado explicitamente como caso de uso prioritário pela mídia tecnológica paraguaia.
- A chegada de Starlink representou uma mudança real para produtores rurais e moradores do interior de Caazapá.

\textbf{Recomendação para imigrantes com propriedades em Caazapá:} Orçar Starlink como item obrigatório na instalação. Equip. ~USD 370 (pagamento único) + ~USD 44/mês garante conectividade de 50--200 Mbps em qualquer ponto do departamento.


- https://www.conatel.gov.py/conatel/wp-content/uploads/2025/02/informe-de-internet-fijo-2024.pdf
- https://www.tecnologia.com.py/contectividad/por-que-starlink-esta-cambiando-la-conectividad-en-el-interior-de-paraguay
- h


\clearpage
\section{Itapúa}
\label{sec:dept_07_Itapua}

\subsection{Segurança Pública}

\textbf{Fonte principal:} Ministerio Público Paraguay / InSight Crime / Municipalidade de Encarnación
\textbf{Data de referência:} 2023--2024
\textbf{Atualizado em:} 2026-03-20


\begin{tabularx}{\linewidth}{lX}
\toprule
Indicador & Valor \\
\midrule
Taxa de homicídios (por 100k hab) & ~4--6 (estimativa, próxima à média nacional) \\
Crimes registrados (total/ano) & --- dados não desagregados publicamente \\
Número de comisarías & ~25 (estimativa --- departamento populoso com fronteira) \\
Índice segurança relativo & médio \\
\bottomrule
\end{tabularx}


Itapúa é um dos departamentos mais desenvolvidos do sul do Paraguai, com Encarnación como sua capital e principal cidade turística. O Ministerio Público registra que o departamento responde por \textbf{aproximadamente 5\% das causas de homicídio doloso} nacional (2019--2024), compatível com sua participação populacional (~8\% do total).

O departamento faz fronteira com a Argentina (Provincia de Misiones, via Ponte Internacional San Roque González de Santa Cruz/Posadas). Esta fronteira é identificada como rota de contrabando e tráfico de drogas, com cocaína e maconha sendo transportadas entre os dois países. O relatório do InSight Crime (2024) documenta presença do PCC em Encarnación, com cerca de 100 membros encarcerados localmente e atividades de lavagem de dinheiro no setor imobiliário, turístico e de construção.

Encarnación experimenta crescimento turístico acelerado (reconhecida internacionalmente por seu carnaval), o que aumenta o risco de tráfico de pessoas, com casos documentados de tráfico sexual e trabalho forçado de menores especialmente durante eventos sazonais.


- Encarnación é uma das cidades paraguaias mais atraentes para imigrantes --- turismo, gastronomia e qualidade de vida elevada para os padrões paraguaios
- A presença de crime organizado é real mas não afeta diretamente o cotidiano de residentes nas áreas centrais
- Bairros turísticos e residenciais de classe média são considerados seguros
- A conexão com Posadas (Argentina) facilita acesso a serviços e bens que não estão disponíveis localmente
- Brasileiros devem ter atenção ao ambiente de negócios: lavagem de dinheiro é fenômeno estrutural que pode afetar transações imobiliárias e comerciais
- Recomenda-se verificar a idoneidade de parceiros comerciais locais antes de fechar negócios


- https://insightcrime.org/paraguay-organized-crime-news/itapua-paraguay/
- https://www.ministe

\subsection{IDH e Desenvolvimento Humano}

\textbf{Fonte principal:} PNUD Paraguai / DGEEC / INE
\textbf{Ano de referência:} 2020 (IDH); 2023 (pobreza e infraestrutura)
\textbf{Atualizado em:} 2026-03-20


\begin{tabularx}{\linewidth}{lX}
\toprule
Indicador & Valor \\
\midrule
IDH & 0,728 \\
Classificação IDH & alto \\
Ranking nacional (1=melhor) & 4/18 \\
Esperança de vida & 74,0 anos \\
Anos de escolaridade média & 8,9 anos \\
RNB per capita (USD PPA) & 10.100 \\
\bottomrule
\end{tabularx}

> Nota: IDH de 2020 conforme relatório PNUD "Desarrollo Humano en el Paraguay 2001-2020". Itapúa está entre os 7 departamentos detalhados no relatório PNUD como 4º lugar, atrás apenas de Asunción, Central e Alto Paraná. Fonte: PNUD Paraguay 2022.


\begin{tabularx}{\linewidth}{lX}
\toprule
Indicador & Valor \\
\midrule
Taxa de pobreza total (\%) & 22,1\% \\
Taxa de pobreza extrema (\%) & 4,8\% \\
Índice de Gini & 0,432 \\
\bottomrule
\end{tabularx}

> Fonte: INE --- EPHC 2023. Itapúa tem uma das menores taxas de pobreza entre os departamentos do interior, reflexo de sua economia diversificada (agronegócio, turismo, indústria).


\begin{tabularx}{\linewidth}{lX}
\toprule
Indicador & Valor \\
\midrule
Acesso a água potável (\%) & 84,6\% \\
Acesso a saneamento (\%) & 83,1\% \\
\bottomrule
\end{tabularx}

> Estimativa baseada em dados do Censo 2022. Encarnación possui infraestrutura urbana sólida, com a reformulação da orla do rio Paraná (Costanera de Encarnación) elevando a qualidade de vida.


Itapúa é um dos departamentos mais desenvolvidos do Paraguai (IDH 0,728) e Encarnación, sua capital, é frequentemente classificada como a segunda ou terceira melhor cidade para se viver no país. A relação com o Brasil é intensa --- a cidade faz fronteira com Posadas (Argentina), mas tem forte fluxo comercial com brasileiros. O departamento concentra grande produção de soja, trigo e canola, com forte presença de cooperativas (muitas de origem alemã, como as Menonitas). A pobreza de 22,1\% é moderada e o acesso a serviços é bom. Encarnación tem universidades, hospitais modernos, shopping centers e a famosa Costanera revitalizada. Para brasileiros, representa excelente opção de radicação com qualidade de vida alta, custo acessível, e proximidade cultural (muitos moradores têm ascendência alemã, italiana ou brasiguaia). Oportunidades no agronegócio, turismo e serviços.


- PNUD Paraguai --- Informe "Índices de Desarrollo Humano en el Paraguay 2001-2020": https://www.undp.org/es/paraguay/noticias/IDH-2001-2020
- Relatório PDF PNUD: htt

\subsection{Sistema Prisional}

\textbf{Fonte:} MJT (Ministerio de Justicia y Trabajo) / MNP Paraguay
\textbf{Ano de referência:} 2023--2024
\textbf{Atualizado em:} 2026-03-20


\begin{tabularx}{\linewidth}{lXXXX}
\toprule
Nome & Município & Capacidade & População & Superlotação \\
\midrule
Penitenciaría Regional de Encarnación & Encarnación & 350 & ~720 & ~206\% \\
Penitenciaría de Fram & Fram & 60 & ~90 & ~150\% \\
\bottomrule
\end{tabularx}


\begin{tabularx}{\linewidth}{lX}
\toprule
Indicador & Valor \\
\midrule
Total de estabelecimentos & 2 \\
Capacidade total (vagas) & ~410 \\
População carcerária total & ~810 \\
Taxa de superlotação & ~198\% \\
\% presos provisórios & ~64\% \\
\bottomrule
\end{tabularx}


Encarnación, capital de Itapúa e segunda maior cidade do Paraguai, sedia a principal penitenciaría regional do departamento. O estabelecimento atende também detentos oriundos de municípios como Hohenau, Bella Vista, Edelira e Coronel Bogado.

Itapúa faz fronteira com a Argentina (Posadas, através da Ponte Internacional San Roque González de Santa Cruz) e com o Brasil, o que cria dinâmicas de crime transfronteiriço: tráfico de drogas, armas e mercadorias contrabandeadas. O perfil carcerário inclui casos ligados a esses crimes, além de crimes comuns.

A Penitenciaría de Fram é um estabelecimento menor e mais rudimentar que atende a região central do departamento. O MNP registrou superlotação e condições estruturais precárias em ambos os estabelecimentos.

\textbf{Para imigrantes:} Itapúa é um dos departamentos mais atrativos do Paraguai para imigrantes, especialmente brasileiros e argentinos. Encarnación tem excelente infraestrutura urbana, forte turismo (Carnaval, praias do rio Paraná), e o departamento possui forte setor agrícola e agroindustrial. O sistema prisional não impacta negativamente a maioria das áreas residenciais e comerciais de Encarnación.


- https://www.mjt.gov.py/index.php/penitenciaria
- https://www.mnp.gov.py/informes
- MNP --- Informe Anual 2022--2023
- La Nación Paraguay: "Situación penitenciaria en Encarnación" (2023)


\subsection{Recursos Hídricos Subterrâneos}

\textbf{Fonte:} SENASA / MADES / Geología del Paraguay / IGRAC / Projeto Piloto SAG Itapúa
\textbf{Atualizado em:} 2026-03-20


\begin{tabularx}{\linewidth}{lXXX}
\toprule
Aquífero & Cobertura no Dept. & Profundidade Média & Vazão Típica \\
\midrule
Guarani / Misiones (afloramento e confinado) & ~85\% & 70--300 m & 20--100 m$^3$/h \\
Basáltico (fraturado, Serra San Rafael) & ~15\% & 40--120 m & 5--20 m$^3$/h \\
\bottomrule
\end{tabularx}


\begin{tabularx}{\linewidth}{lX}
\toprule
Parâmetro & Valor \\
\midrule
Profundidade média (m) & 70--120 \\
Profundidade máxima recomendada (m) & 350 \\
Vazão típica (m$^3$/h) & 20--80 \\
Custo estimado perfuração (USD) & 3.000 -- 12.000 \\
Qualidade da água & boa a excelente \\
pH médio & 6,5--7,5 \\
Outorga necessária & sim \\
Órgão regulador & MADES --- DGPCRH \\
\bottomrule
\end{tabularx}


Itapúa é um dos departamentos mais importantes para o Aquífero Guarani no Paraguai. O Projeto Piloto do Departamento de Itapúa (financiado pelo GEF/Banco Mundial) identificou cerca de 60 poços registrados, com profundidades principais entre 70--120 m, embora alguns poços alcancem mais de 300 m em áreas com grande espessura de basalto cobrindo o aquífero. O afloramento do Aquífero Guarani representa cerca de 40\% da área departamental. A água tem qualidade excelente: baixa mineralização, pH neutro, temperatura moderada (25--40 $^\circ$C nas profundidades típicas). A cidade de Encarnación (capital) é abastecida por combinação de poços artesianos do Guarani e captação superficial do rio Paraná. Uso principal: abastecimento doméstico, irrigação (soja, trigo, arroz) e uso urbano.


Itapúa é destino tradicional de imigrantes brasileiros e europeus (colonização alemã, japonesa, italiana histórica). O Aquífero Guarani é de fácil acesso em 70--120 m na maior parte do território, com custo de perfuração entre USD 3.500 e 10.000. Para propriedades agrícolas com necessidade de irrigação, a vazão típica de 20--80 m$^3$/h é mais que suficiente. A relação custo-benefício é excelente. Recomenda-se poço artesiano próprio para qualquer propriedade rural. Licença obrigatória no MADES.


- https://www.oas.org/DSD/WaterResources/Pastprojects/Guarani\_esp.asp
- https://gw-project.org/es/caracteristicas-del-agua-subterranea-en-paraguay/
- https://es.wikipedia.org/wiki/Acu\%C3\%ADfero\_Guaran\%C3\%AD
- https://www.gwp.org/globalassets/global/toolbox/case-studies/americas-and-caribbean/transboundary.-grou

\subsection{Aptidão do Solo}

\textbf{Fonte:} MAG / SENAVE / INFONA / CONACYT / CAPECO
\textbf{Atualizado em:} 2026-03-20


\begin{tabularx}{\linewidth}{lXX}
\toprule
Tipo de Solo & \% do Território & Características \\
\midrule
Oxissolos (Latossolos Roxos / Terra Roxa) & 55\% & Solos derivados de basalto, profundos (>180 cm), bem drenados, alta fertilidade, cor vermelha escura \\
Ultissolos (Argissolos Vermelhos) & 25\% & Solos podzólicos com horizonte B textural, fertilidade média-alta \\
Entissolos / Aluviais & 12\% & Planícies do Rio Paraná e afluentes, alta fertilidade natural \\
Inceptissolos & 8\% & Solos jovens em áreas de relevo suave-ondulado, boa aptidão \\
\bottomrule
\end{tabularx}


\begin{tabularx}{\linewidth}{lXX}
\toprule
Classe & \% Território & Uso Recomendado \\
\midrule
Lavoura intensiva & 55\% & Soja, milho, trigo, canola (solos basálticos) \\
Lavoura extensiva & 20\% & Mandioca, girassol, sorgo \\
Pastagem natural/melhorada & 10\% & Pecuária bovina de corte e leite \\
Floresta / reserva & 10\% & Remanescentes de Mata Atlântica, matas ciliares \\
Sem aptidão agrícola & 5\% & Áreas urbanas (Encarnación), degradadas \\
\bottomrule
\end{tabularx}


\begin{tabularx}{\linewidth}{lX}
\toprule
Parâmetro & Valor \\
\midrule
pH médio & 5,0 -- 6,2 \\
Fertilidade natural & alta \\
Teor de matéria orgânica & médio-alto (2,5--4,0\%) \\
Necessidade de calcário & parcial (solos basálticos têm menor acidez) \\
\bottomrule
\end{tabularx}


- \textbf{Soja} --- principal cultura; Itapúa é o 2º maior produtor nacional (692.000 ha, 19\% da área total)
- \textbf{Trigo} --- cultivo de inverno extremamente expressivo; Itapúa é o principal produtor de trigo do Paraguai
- \textbf{Milho} --- em rotação com soja e trigo
- \textbf{Canola} --- em expansão como cultura de inverno
- \textbf{Girassol} --- tradicional, ainda relevante
- \textbf{Erva-mate} --- zona sul do departamento (Trinidad, Pilar do Sul)
- \textbf{Mandioca} e \textbf{amendoim} --- agricultura familiar


- \textbf{Compactação:} intenso uso de maquinário pesado pode causar compactação nos Oxissolos
- \textbf{Erosão laminar:} em áreas de relevo ondulado; mitigada por plantio direto consolidado na região
- \textbf{Acidez residual:} mesmo em solos basálticos, 35--40\% das amostras mostram MO < 25 g/kg; calagem parcial necessária
- \textbf{Pressão sobre remanescentes:} floresta nativa reduzida; área já altamente convertida em agricultura
- \textbf{Variabilidade climática:} chuvas irregulares em alguns anos causam perdas


It

\subsection{Terra Rural}

\textbf{Fonte:} INDERT / Clasificados.com.py / mercado local
\textbf{Ano de referência:} 2024
\textbf{Atualizado em:} 2026-03-20


\begin{tabularx}{\linewidth}{lXXX}
\toprule
Tipo & Preço (USD/ha) & Faixa & Tendência \\
\midrule
Agrícola alta produtividade & 6.900 & 5.000 -- 10.000 & $\uparrow$ \\
Pastagem & 3.200 & 2.200 -- 4.500 & $\uparrow$ \\
Mata / reserva & 1.400 & 800 -- 2.200 & $\rightarrow$ \\
\bottomrule
\end{tabularx}

> Referência: USD 6.945/ha para terra agrícola confirmado em múltiplas fontes (base 2021--2024).


\begin{tabularx}{\linewidth}{lX}
\toprule
Parâmetro & Valor \\
\midrule
Valorização anual média (\%) & 7--9\% \\
Lote mínimo rural (ha) & 10 \\
Imposto territorial rural (\%) & 1\% sobre valor fiscal \\
Liquidez do mercado & alta \\
\bottomrule
\end{tabularx}


- \textbf{Encarnación (entorno):} Capital departamental e cidade mais próspera do Sul do Paraguai; terra periurbana valorizada; hub de serviços e turismo de compras (turistas argentinos).
- \textbf{Coronel Bogado / Hohenau / Bella Vista:} Polo de colonização alemã e japonesa; produção de soja, erva-mate, soja; terras premium.
- \textbf{Trinidad / Jesús:} Zona histórica Jesuítica (Patrimônio da UNESCO); turismo e agricultura mista.
- \textbf{Natalio / Edelira:} Expansão da fronteira sojícola; colônias brasileiras ativas.
- \textbf{Fram / Mayor Otaño:} Corredor sul, terras férteis com boa logística para exportação via Argentina.


\begin{tabularx}{\linewidth}{lXXX}
\toprule
Indicador & Itapúa (PY) & Rio Grande do Sul (BR) & Paraná Sul (BR) \\
\midrule
Terra agrícola (USD/ha) & 5.000--10.000 & 12.000--22.000 & 15.000--28.000 \\
Terra pastagem (USD/ha) & 2.200--4.500 & 5.000--10.000 & 6.000--12.000 \\
Custo relativo & ~40\% do RS & referência RS & referência PR \\
\bottomrule
\end{tabularx}

Itapúa é o departamento com maior similaridade cultural ao Sul do Brasil (colonização alemã, brasileira e japonesa), combinando produtividade agrícola elevada com custo de terra significativamente menor.


- [INDERT --- Instituto Nacional de Desarrollo Rural y de la Tierra](https://www.indert.gov.py/)
- [Última Hora --- Tierras de Alto Paraná más caras que EEUU](https://www.ultimahora.com/tierras-agricolas-alto-parana-son-mas-caras-que-las-eeuu-n2928594)
- [Plan B Paraguay --- Farmland Prices](https://x.com/planbparaguay/status/1834232169759211739)
- [RM Paraguay --- Fazendas](https://rmparaguay.com/fazendas/)


\subsection{Mercado Imobiliário}

\textbf{Capital:} Encarnación
\textbf{Fonte:} Clasificados.com.py / InfoCasas / Zonaprop / mercado local
\textbf{Ano de referência:} 2024
\textbf{Atualizado em:} 2026-03-20


\begin{tabularx}{\linewidth}{lXX}
\toprule
Tipo & Preço (USD/m$^2$) & Aluguel (USD/mês) \\
\midrule
Apartamento 1 quarto & 900 -- 1.400 & 350 -- 600 \\
Apartamento 2 quartos & 1.000 -- 1.600 & 500 -- 900 \\
Casa 3 quartos & 800 -- 1.200 & 400 -- 750 \\
\bottomrule
\end{tabularx}

> Exemplos reais 2024: apartamento 52m$^2$ por USD 95.000 (USD 1.827/m$^2$); aluguel 2 quartos area Poti'y: USD 800/mês; apartamento 123m$^2$ (3q/3b): aluguel USD 800/mês.


\begin{tabularx}{\linewidth}{lX}
\toprule
Parâmetro & Valor \\
\midrule
Valorização anual (\%) & 8--12\% \\
Disponibilidade de financiamento & boa (bancos paraguaios e cooperativas alemãs/japonesas locais) \\
Bairros mais valorizados & Areguá (beira do Rio Paraná), Av. Costanera, Centro, Barrio San Pedro \\
\bottomrule
\end{tabularx}


\textbf{Vale a pena comprar ou alugar?}
Encarnación é a "Ibiza do Paraguai" --- cidade com praias fluviais no Rio Paraná, intensa vida noturna no verão e fortíssimo turismo argentino e brasileiro. O mercado imobiliário teve valorização expressiva após obras de reurbanização da costanera (2010--2024). A compra de apartamento próximo à orla é um investimento sólido com potencial de renda por aluguel de temporada.

\textbf{Áreas recomendadas:}
- \textbf{Costanera e Centro:} orla reformada, alta demanda turística, apartamentos modernos. Mais caros mas melhor liquidez.
- \textbf{Barrio San Pedro / Villa Alta:} residencial consolidado, bom custo-benefício para moradia permanente.
- \textbf{Capitán Miranda / Hohenau:} cidades próximas com qualidade de vida elevada e comunidade alemã estabelecida.

\textbf{Armadilhas comuns:}
- Imóveis na orla costanera têm custo mais elevado de manutenção pela proximidade com o rio (umidade, cheias históricas).
- Verificar regularidade de imóveis em zonas que passaram por reurbanização --- alguns títulos foram afetados pelas desapropriações da Itaipu/Yacyretá.
- Encarnación recebe turistas argentinos em massa no verão (dez--fev) --- aluguel de temporada pode ser muito rentável mas requer gestão.
- Apartamentos "en pozo" muito comuns; verificar histórico e solidez das construtoras.


- [InfoCasas.com.py --- Aluguel Encarnación](https://www.infocasas.com.py/alquiler/departamentos/itapua/encarnacion)
- [InfoCasas.com.py --- Venda Encarnación](https://www.infocasas.com.py/venta/departamentos/itapua/encarnacion)


\subsection{Cobertura Celular}

\textbf{Fonte:} CONATEL / Tigo / Personal / OpenSignal / nPerf
\textbf{Ano de referência:} 2024
\textbf{Atualizado em:} 2026-03-20


\begin{tabularx}{\linewidth}{lXXXXX}
\toprule
Operadora & 2G & 3G & 4G & 5G & Cobertura Pop. \\
\midrule
Tigo & 92\% & 85\% & 78\% & --- & 88\% \\
Personal/Claro & 88\% & 80\% & 72\% & --- & 84\% \\
VOX & 55\% & 42\% & 28\% & --- & 48\% \\
\bottomrule
\end{tabularx}

> \textbf{Nota:} Itapúa é um dos departamentos mais desenvolvidos do interior, com Encarnación sendo a terceira maior cidade do Paraguai (~120k hab.). A cobertura 4G é excelente em Encarnación e municípios ao longo da Ruta 1 e da fronteira com Argentina. Interior norte tem cobertura mais fraca.


\begin{tabularx}{\linewidth}{lX}
\toprule
Parâmetro & Valor \\
\midrule
Melhor operadora rural & Tigo \\
Velocidade 4G média (Mbps) & 25--50 \\
Áreas sem cobertura & Interior norte e noroeste do departamento, zonas de mata \\
Sinal em propriedade rural & bom a moderado \\
\bottomrule
\end{tabularx}


\begin{tabularx}{\linewidth}{lXXX}
\toprule
Plano & Operadora & USD/mês & Dados \\
\midrule
Pré-pago básico (7 dias) & Tigo & ~2,50 & 3 GB \\
Pré-pago intermediário (30 dias) & Personal & ~5,00 & 8 GB \\
Pós-pago entrada & Tigo & ~12,00 & 15 GB + chamadas \\
Pós-pago intermediário & Personal/Claro & ~18,00 & 30 GB + chamadas \\
\bottomrule
\end{tabularx}


- \textbf{Encarnación:} Excelente cobertura 4G de Tigo e Personal; cidade com infraestrutura de telecom comparável a centros médios brasileiros. Proximidade com Posadas (Argentina) oferece alternativa de roaming.
- \textbf{Municípios da faixa de fronteira (Capitán Miranda, Coronel Bogado, Carmen del Paraná):} Boa cobertura 4G de Tigo.
- \textbf{Colônias agrícolas (Jesus, Trinidad, Obligado):} Cobertura 3G/4G razoável de Tigo; Personal com presença crescente.
- \textbf{Interior norte (Natalio, San Juan del Paraná):} Cobertura mais variável; predomina 3G.
- Itapúa é uma das melhores opções no interior do Paraguai para imigrantes, combinando boa conectividade com qualidade de vida.


- https://www.conatel.gov.py/conatel/indicadores/
- https://www.tigo.com.py/cobertura
- https://www.personal.com.py/institucional/cobertura.html
- https://www.nperf.com/en/map/PY/-/-/signal


\subsection{Internet}

\textbf{Fonte:} CONATEL / Tigo / Claro / INE / Copaco
\textbf{Ano de referência:} 2024
\textbf{Atualizado em:} 2026-03-20


\begin{tabularx}{\linewidth}{lX}
\toprule
Indicador & Valor \\
\midrule
\% população com internet & 78,2\% \\
\% com banda larga fixa & ~30\% dos domicílios \\
Velocidade média download (Mbps) & 40--80 (urbano) / 8--20 (rural) \\
Velocidade média upload (Mbps) & 12--30 (urbano) \\
\bottomrule
\end{tabularx}

> Fonte: INE EPHC 2024 --- 78,2\% de penetração de internet em Itapúa, o mais alto entre os departamentos do interior do Paraguai.


\begin{tabularx}{\linewidth}{lXX}
\toprule
Tecnologia & Disponibilidade & Velocidade Típica \\
\midrule
Fibra ótica (FTTH) & Encarnación, Coronel Bogado, Capitán Miranda & 50--500 Mbps \\
Cabo HFC & Encarnación & 30--200 Mbps \\
Rádio/Wireless & Amplo no departamento & 5--50 Mbps \\
ADSL (Copaco) & Interior com rede legada & 2--10 Mbps \\
Satélite (Starlink) & Todo o departamento & 50--200 Mbps \\
\bottomrule
\end{tabularx}


\begin{tabularx}{\linewidth}{lXXX}
\toprule
Provedor & Tecnologia & Área de cobertura & Preço ref. (USD/mês) \\
\midrule
Tigo & Fibra/HFC/Rádio & Encarnación + municípios & 18--55 \\
Claro/Personal & Fibra/Rádio & Encarnación + municípios maiores & 15--50 \\
Copaco & ADSL/Fibra & Interior + municípios menores & 10--30 \\
ISPs locais & Rádio wireless & Colônias agrícolas & 8--20 \\
Starlink & Satélite LEO & Todo o departamento & ~44 + equip. ~370 \\
\bottomrule
\end{tabularx}


A cobertura rural de Itapúa é \textbf{moderada a boa}, a melhor entre os departamentos do interior distante de Asunción:

- \textbf{Faixa de fronteira com Argentina} (ao longo do Rio Paraná): Excelente cobertura, beneficiada por infraestrutura de Encarnación.
- \textbf{Municípios colonizados por imigrantes alemães e japoneses} (Fram, Obligado, Pirapó): ISPs locais bem estabelecidos, rádio wireless com qualidade razoável.
- \textbf{Interior norte} (Natalio, General Delgado): Cobertura mais limitada; rádio wireless e internet móvel.
- Itapúa tem tradição de colonização europeia e japonesa, com comunidades relativamente bem organizadas em termos de infraestrutura.

\textbf{Para imigrantes brasileiros:} Encarnación e municípios da faixa sul oferecem excelente conectividade. Interior mais ao norte requer avaliação prévia da cobertura na propriedade específica.


- https://www.conatel.gov.py/conatel/wp-content/uploads/2025/02/informe-de-internet-fijo-2024.pdf
- https://marketdata.com.py/laboratorio/analisis/acceso-a-internet-


\clearpage
\section{Misiones}
\label{sec:dept_08_Misiones}

\subsection{Segurança Pública}

\textbf{Fonte principal:} Ministerio Público Paraguay / INE / Atlas da Violência CONACYT
\textbf{Data de referência:} 2023--2024
\textbf{Atualizado em:} 2026-03-20


\begin{tabularx}{\linewidth}{lX}
\toprule
Indicador & Valor \\
\midrule
Taxa de homicídios (por 100k hab) & ~2--4 (estimativa, abaixo da média nacional) \\
Crimes registrados (total/ano) & --- dados não desagregados publicamente \\
Número de comisarías & ~8 (estimativa departamental) \\
Índice segurança relativo & alto \\
\bottomrule
\end{tabularx}


Misiones é um dos departamentos mais tranquilos do Paraguai em termos de segurança pública. O Ministerio Público registra que o departamento responde por apenas \textbf{1\% das causas de homicídio doloso} nacional (2019--2024) --- proporção muito baixa e compatível com sua pequena população (cerca de 1,7\% do total nacional).

O departamento faz fronteira com Itapúa ao sul e com Ñeembucú a sudoeste, e está situado entre os rios Tebicuary e Paraná. Embora a fronteira sul com Argentina represente potencialmente uma rota de contrabando, o InSight Crime identifica o trecho Itapúa-Misiones como área de cruzamento, mas não Misiones como epicentro de violência. San Juan Bautista, capital departamental, é uma cidade pequena e tranquila.

A criminalidade é predominantemente de baixa intensidade, com crimes patrimoniais esporádicos. A baixa densidade populacional e o caráter rural do departamento contribuem para a baixa taxa de violência.


- Misiones é adequado para quem busca tranquilidade e ambiente rural com boa segurança
- O departamento tem presença histórica de agricultores de diversas origens, incluindo brasileiros
- San Juan Bautista oferece serviços básicos, mas para serviços mais especializados é necessário viajar a Asunción (~170 km) ou Encarnación (~100 km)
- Baixíssima taxa de criminalidade grave o torna atrativo para famílias e aposentados
- Oportunidades principais no setor agropecuário e turismo rural


- https://www.ministeriopublico.gov.py/nota/estadisticas-del-ministerio-publico-la-mayoria-de-los-homicidios-dolosos-en-10564
- https://insightcrime.org/paraguay-organized-crime-news/itapua-paraguay/
- https://www.conacyt.gov.py/sites/default/files/upload\_editores/u294/ATLAS\_DE\_VIOLENCIA\_PARAGUAY.pdf
- https://ods.ine.gov.py/ine-main/ods/paz-justicias-e-instituciones-solidas-16/meta-16.1/indicador-284


\subsection{IDH e Desenvolvimento Humano}

\textbf{Fonte principal:} PNUD Paraguai / DGEEC / INE
\textbf{Ano de referência:} 2020 (IDH); 2023 (pobreza e infraestrutura)
\textbf{Atualizado em:} 2026-03-20


\begin{tabularx}{\linewidth}{lX}
\toprule
Indicador & Valor \\
\midrule
IDH & 0,712 \\
Classificação IDH & alto \\
Ranking nacional (1=melhor) & 7/18 (empate estimado com Caaguazú) \\
Esperança de vida & 73,2 anos \\
Anos de escolaridade média & 8,3 anos \\
RNB per capita (USD PPA) & 8.900 \\
\bottomrule
\end{tabularx}

> Nota: IDH estimado com base no relatório PNUD 2020 e dados subnacionais da Global Data Lab para a região sul do Paraguai. Misiones tem perfil agropecuário consolidado e é o menor departamento em população da Região Oriental.


\begin{tabularx}{\linewidth}{lX}
\toprule
Indicador & Valor \\
\midrule
Taxa de pobreza total (\%) & 26,4\% \\
Taxa de pobreza extrema (\%) & 5,9\% \\
Índice de Gini & 0,442 \\
\bottomrule
\end{tabularx}

> Fonte: INE --- EPHC 2023 (dado de pobreza 26,4\% para Misiones conforme dados regionais Sul). Departamento com pobreza moderada para os padrões paraguaios.


\begin{tabularx}{\linewidth}{lX}
\toprule
Indicador & Valor \\
\midrule
Acesso a água potável (\%) & 80,3\% \\
Acesso a saneamento (\%) & 78,9\% \\
\bottomrule
\end{tabularx}

> Estimativa baseada em dados do Censo 2022 e perfil semiurbano/rural do departamento. San Juan Bautista, capital, tem boa infraestrutura para cidade de porte pequeno.


Misiones é o menor departamento em população da Região Oriental com IDH alto (0,712) e economia baseada em pecuária, agricultura e turismo histórico (Ruínas Jesuíticas). San Juan Bautista, a capital, é uma cidade tranquila com bons serviços básicos para seu tamanho. Para brasileiros, o departamento representa uma opção de vida rural com qualidade, distante dos centros urbanos. O turismo em torno das Missões Jesuíticas (San Cosme y Damián, Santa Rosa) oferece oportunidades de empreendimento. A pobreza de 26,4\% é moderada e o acesso a água (80,3\%) e saneamento (78,9\%) são satisfatórios. O departamento não tem grandes cidades, o que limita oportunidades de emprego formal mas reduz custos de vida. Indicado para quem busca tranquilidade rural, agronegócio ou turismo histórico-cultural.


- PNUD Paraguai --- Informe IDH 2001-2020: https://www.undp.org/es/paraguay/noticias/IDH-2001-2020
- INE --- EPHC 2023: https://www.ine.gov.py/noticias/1928/pobreza-monetaria-en-el-2023-se-ubico-en-el-227
- INE --- Censo 2022: https://www.ine.gov.py/noticias/2101/principales-resultados-fi

\subsection{Sistema Prisional}

\textbf{Fonte:} MJT (Ministerio de Justicia y Trabajo) / MNP Paraguay
\textbf{Ano de referência:} 2023--2024
\textbf{Atualizado em:} 2026-03-20


\begin{tabularx}{\linewidth}{lXXXX}
\toprule
Nome & Município & Capacidade & População & Superlotação \\
\midrule
Penitenciaría Regional de San Juan Bautista & San Juan Bautista & 120 & ~210 & ~175\% \\
\bottomrule
\end{tabularx}


\begin{tabularx}{\linewidth}{lX}
\toprule
Indicador & Valor \\
\midrule
Total de estabelecimentos & 1 \\
Capacidade total (vagas) & ~120 \\
População carcerária total & ~210 \\
Taxa de superlotação & ~175\% \\
\% presos provisórios & ~65\% \\
\bottomrule
\end{tabularx}


Misiones é um dos departamentos menos populosos do Paraguai, com economia baseada em pecuária e agricultura. O único estabelecimento penal localiza-se em San Juan Bautista, capital departamental, e atende detentos de todos os 10 municípios do departamento.

O perfil de criminalidade em Misiones é relativamente baixo em comparação com departamentos de fronteira. O tráfico de gado e crimes contra a propriedade rural são as infrações mais comuns que alimentam a população carcerária local.

O departamento faz fronteira com a Argentina ao sul (rio Paraná), mas o volume de comércio e tráfego nessa fronteira é menor do que em Itapúa ou Alto Paraná, reduzindo a pressão do crime organizado transnacional.

\textbf{Para imigrantes:} Misiones é um departamento tranquilo e com baixo custo de vida. San Juan Bautista tem infraestrutura modesta mas funcional. O sistema prisional tem impacto mínimo na vida cotidiana da maioria dos municípios. Indicado para quem busca ambiente rural e tranquilo.


- https://www.mjt.gov.py/index.php/penitenciaria
- https://www.mnp.gov.py/informes
- MNP --- Informe de Visitas 2022


\subsection{Recursos Hídricos Subterrâneos}

\textbf{Fonte:} SENASA / MADES / Geología del Paraguay / IGRAC
\textbf{Atualizado em:} 2026-03-20


\begin{tabularx}{\linewidth}{lXXX}
\toprule
Aquífero & Cobertura no Dept. & Profundidade Média & Vazão Típica \\
\midrule
Guarani / Misiones (afloramento e confinado) & ~80\% & 60--250 m & 20--80 m$^3$/h \\
Aluvial (rio Tebicuary e afluentes) & Faixas ribeirinhas & 8--25 m & 5--15 m$^3$/h \\
\bottomrule
\end{tabularx}


\begin{tabularx}{\linewidth}{lX}
\toprule
Parâmetro & Valor \\
\midrule
Profundidade média --- Guarani (m) & 60--200 \\
Profundidade máxima recomendada (m) & 300 \\
Vazão típica (m$^3$/h) & 15--60 \\
Custo estimado perfuração (USD) & 2.500 -- 9.000 \\
Qualidade da água & boa a excelente \\
pH médio & 6,5--7,5 \\
Outorga necessária & sim \\
Órgão regulador & MADES --- DGPCRH \\
\bottomrule
\end{tabularx}


O departamento de Misiones dá nome à própria Formação Misiones, unidade geológica que compõe o Aquífero Guarani nesta região. Os afloramentos dessa formação são frequentes, com arenitos de alta permeabilidade que facilitam a recarga do sistema. As profundidades de captação são relativamente baixas (60--150 m), tornando a perfuração mais econômica que em outros departamentos. A água é de boa qualidade, com baixa mineralização, pH neutro e ausência de contaminantes naturais. Uso principal: abastecimento doméstico, pecuária e irrigação. San Juan Bautista (capital) tem rede pública do SENASA.


Misiones é um departamento com economia agropecuária forte e terrenos acessíveis. O Aquífero Guarani é amplamente disponível em profundidades relativamente rasas (60--150 m), o que reduz o custo de perfuração para USD 2.500--7.000. Excelente para propriedades rurais com uso doméstico e irrigação. A rede pública cobre apenas centros urbanos. Poço artesiano é a solução padrão e economicamente vantajosa para qualquer imóvel rural. Licença do MADES obrigatória.


- https://www.geologiadelparaguay.com.py/Acuifero-Misiones.htm
- https://gw-project.org/es/caracteristicas-del-agua-subterranea-en-paraguay/
- https://es.wikipedia.org/wiki/Acu\%C3\%ADfero\_Guaran\%C3\%AD
- https://www.mades.gov.py/tramites/registro-de-empresas-publicas-y-privadas-que-se-dedican-a-la-perforacion-de-pozos-tubulares/


\subsection{Aptidão do Solo}

\textbf{Fonte:} MAG / SENAVE / CONACYT / Geología del Paraguay
\textbf{Atualizado em:} 2026-03-20


\begin{tabularx}{\linewidth}{lXX}
\toprule
Tipo de Solo & \% do Território & Características \\
\midrule
Ultissolos / Planossolos (Alfissolos) & 40\% & Solos com horizonte B argiloso-imperfeito; limitações físicas para agricultura intensiva \\
Entissolos arenosos (Quartzarênicos) & 30\% & Solos derivados de arenito, baixa fertilidade natural, suscetíveis à erosão e seca \\
Oxissolos (Latossolos) & 15\% & Solos basálticos no nordeste do departamento, melhor aptidão agrícola \\
Inceptissolos / Aluviais & 15\% & Planícies do Rio Paraguai e Rio Pilcomayo; solos jovens de boa fertilidade \\
\bottomrule
\end{tabularx}


\begin{tabularx}{\linewidth}{lXX}
\toprule
Classe & \% Território & Uso Recomendado \\
\midrule
Lavoura intensiva & 15\% & Soja e milho nas áreas basálticas do nordeste \\
Lavoura extensiva & 25\% & Mandioca, sorgo, girassol, amendoim \\
Pastagem natural/melhorada & 40\% & Pecuária extensiva e semi-intensiva (atividade principal) \\
Floresta / reserva & 15\% & Remanescentes nativos, matas ciliares \\
Sem aptidão agrícola & 5\% & Áreas alagáveis, urbanas, degradadas \\
\bottomrule
\end{tabularx}


\begin{tabularx}{\linewidth}{lX}
\toprule
Parâmetro & Valor \\
\midrule
pH médio & 5,0 -- 6,5 \\
Fertilidade natural & média \\
Teor de matéria orgânica & médio \\
Necessidade de calcário & parcial \\
\bottomrule
\end{tabularx}


- \textbf{Pecuária bovina} --- atividade dominante; Misiones é região de excelência para gado de corte
- \textbf{Soja} --- nas zonas nordeste com solos basálticos
- \textbf{Mandioca} --- base da agricultura familiar
- \textbf{Sorgo} e \textbf{girassol} --- como forrageiros e oleaginosas
- \textbf{Amendoim} --- tradicional nas áreas arenosas
- \textbf{Cana-de-açúcar} --- pequenas áreas


- \textbf{Solos com limitações físicas:} planossolos com B textural argiloso criam camada impermeável; drenagem interna deficiente
- \textbf{Solos arenosos:} baixa retenção hídrica; agricultura intensiva sem irrigação é arriscada
- \textbf{Alagamentos:} zonas baixas próximas aos rios Tebicuary e Paraná sujeitas a inundações sazonais
- \textbf{Apenas 20\% de área cultivável:} levantamentos indicam que somente ~20\% do território tem aptidão real para lavoura
- \textbf{Erosão em solos arenosos:} alto risco em relevo suave-ondulado sem proteção vegetal


Misiones não é o destino prioritário para imigrantes brasileiros com foco em a

\subsection{Terra Rural}

\textbf{Fonte:} INDERT / Clasificados.com.py / mercado local
\textbf{Ano de referência:} 2024
\textbf{Atualizado em:} 2026-03-20


\begin{tabularx}{\linewidth}{lXXX}
\toprule
Tipo & Preço (USD/ha) & Faixa & Tendência \\
\midrule
Agrícola alta produtividade & 4.500 & 3.000 -- 6.500 & $\uparrow$ \\
Pastagem & 2.200 & 1.500 -- 3.200 & $\uparrow$ \\
Mata / reserva & 900 & 500 -- 1.500 & $\rightarrow$ \\
\bottomrule
\end{tabularx}


\begin{tabularx}{\linewidth}{lX}
\toprule
Parâmetro & Valor \\
\midrule
Valorização anual média (\%) & 5--7\% \\
Lote mínimo rural (ha) & 20 \\
Imposto territorial rural (\%) & 1\% sobre valor fiscal \\
Liquidez do mercado & média \\
\bottomrule
\end{tabularx}


- \textbf{San Juan Bautista (entorno):} Capital departamental com melhor infraestrutura; terras mistas de agricultura e pecuária.
- \textbf{San Ignacio:} Segunda cidade do departamento; zona próxima às Ruínas Jesuíticas (turismo + agricultura).
- \textbf{Santa Rosa / Santiago:} Corredor de pecuária extensiva; terras de boa fertilidade.
- \textbf{Ayolas:} Cidade da Hidrelétrica Yacyretá; economia impulsionada pela energia; terra valorizada.
- \textbf{Villa Florida:} Destino turístico às margens do Rio Tebicuary; chácaras de lazer em valorização.


\begin{tabularx}{\linewidth}{lXXX}
\toprule
Indicador & Misiones (PY) & Rio Grande do Sul (BR) & Mato Grosso do Sul (BR) \\
\midrule
Terra agrícola (USD/ha) & 3.000--6.500 & 12.000--22.000 & 10.000--20.000 \\
Terra pastagem (USD/ha) & 1.500--3.200 & 5.000--9.000 & 4.000--8.000 \\
Custo relativo & ~28\% do RS & referência RS & referência MS \\
\bottomrule
\end{tabularx}

Misiones é um departamento com vocação pecuária e turística. Os preços de terra são competitivos, especialmente em zonas de pastagem para gado de corte.


- [INDERT --- Instituto Nacional de Desarrollo Rural y de la Tierra](https://www.indert.gov.py/)
- [Plan B Paraguay --- Farmland Prices](https://x.com/planbparaguay/status/1834232169759211739)
- [InfoCasas --- Chacras Paraguay](https://www.infocasas.com.py/venta/chacras-o-campos)
- [Hendyla.com --- Campos Paraguay](https://www.hendyla.com/inmuebles/campos/)


\subsection{Mercado Imobiliário}

\textbf{Capital:} San Juan Bautista
\textbf{Fonte:} Clasificados.com.py / mercado local
\textbf{Ano de referência:} 2024
\textbf{Atualizado em:} 2026-03-20


\begin{tabularx}{\linewidth}{lXX}
\toprule
Tipo & Preço (USD/m$^2$) & Aluguel (USD/mês) \\
\midrule
Apartamento 2 quartos & 500 -- 800 & 180 -- 350 \\
Casa 3 quartos & 420 -- 700 & 160 -- 320 \\
\bottomrule
\end{tabularx}


\begin{tabularx}{\linewidth}{lX}
\toprule
Parâmetro & Valor \\
\midrule
Valorização anual (\%) & 4--6\% \\
Disponibilidade de financiamento & limitada (cooperativas e BNF) \\
Bairros mais valorizados & Centro, proximidades do Parque Industrial, bairros residenciais com infraestrutura de Yacyretá \\
\bottomrule
\end{tabularx}


\textbf{Vale a pena comprar ou alugar?}
San Juan Bautista é uma cidade de porte moderado (aprox. 50.000 hab.) com boa qualidade de vida para padrões paraguaios. A presença da Usina Yacyretá cria uma infraestrutura de serviços acima da média regional. Para quem atua na agroindústria ou pecuária regional, compra de casa é vantajosa.

\textbf{Áreas recomendadas:}
- \textbf{San Juan Bautista centro:} concentração dos serviços básicos.
- \textbf{San Ignacio:} cidade turística com ruínas jesuíticas, boa qualidade de vida; preços similares.
- \textbf{Ayolas:} cidade planejada para funcionários de Yacyretá, infraestrutura de qualidade superior; imóveis escassos mas bem mantidos.
- \textbf{Villa Florida:} opção de chácara/rancho turístico para moradia e renda adicional.

\textbf{Armadilhas comuns:}
- Mercado imobiliário muito pequeno; pouca oferta formal.
- Ayolas tem imóveis controlados pela EBY (Entidade Binacional Yacyretá); verificar titulação antes de qualquer negociação.
- Conectividade limitada com outras regiões; acesso principal pela Ruta 1.


- [InfoCasas.com.py --- Misiones](https://www.infocasas.com.py/venta/inmuebles/misiones)
- [Clasipar.com --- Alquiler Misiones](https://clasipar.paraguay.com/alquiler/inmuebles/misiones)
- [INDERT](https://www.indert.gov.py/)


\subsection{Cobertura Celular}

\textbf{Fonte:} CONATEL / Tigo / Personal / OpenSignal / nPerf
\textbf{Ano de referência:} 2024
\textbf{Atualizado em:} 2026-03-20


\begin{tabularx}{\linewidth}{lXXXXX}
\toprule
Operadora & 2G & 3G & 4G & 5G & Cobertura Pop. \\
\midrule
Tigo & 82\% & 72\% & 60\% & --- & 76\% \\
Personal/Claro & 75\% & 65\% & 52\% & --- & 70\% \\
VOX & 45\% & 30\% & 15\% & --- & 38\% \\
\bottomrule
\end{tabularx}

> \textbf{Nota:} Misiones é um departamento pequeno com população relativamente concentrada na capital San Juan Bautista e ao longo da Ruta 1 (Asunción--Encarnación). A Ruta 1 garante boa cobertura linear; interior mais afastado tem sinal reduzido.


\begin{tabularx}{\linewidth}{lX}
\toprule
Parâmetro & Valor \\
\midrule
Melhor operadora rural & Tigo \\
Velocidade 4G média (Mbps) & 15--35 \\
Áreas sem cobertura & Interior oeste, margens do rio Paraguai, zonas de mata \\
Sinal em propriedade rural & moderado \\
\bottomrule
\end{tabularx}


\begin{tabularx}{\linewidth}{lXXX}
\toprule
Plano & Operadora & USD/mês & Dados \\
\midrule
Pré-pago básico (7 dias) & Tigo & ~2,50 & 3 GB \\
Pré-pago intermediário (30 dias) & Personal & ~5,00 & 8 GB \\
Pós-pago entrada & Tigo & ~12,00 & 15 GB + chamadas \\
Pós-pago intermediário & Personal/Claro & ~18,00 & 30 GB + chamadas \\
\bottomrule
\end{tabularx}


- \textbf{San Juan Bautista e Ruta 1:} Boa cobertura 4G de Tigo e Personal. San Juan Bautista é uma cidade tranquila (~25k hab.) com conectividade urbana razoável.
- \textbf{San Ignacio, Santa Rosa, Villa Florida:} Cobertura 3G/4G de Tigo; Personal com presença razoável.
- \textbf{Área do Pantanal sul e margens do rio Paraguai:} Cobertura muito fraca a inexistente.
- \textbf{Para turismo no Pantanal:} Sinal muito irregular; recomendado não contar com cobertura celular confiável.
- Misiones é destino turístico (Ruínas Jesuíticas); a infraestrutura urbana de San Ignacio e Santa María tem boa cobertura celular para padrão de cidades pequenas.


- https://www.conatel.gov.py/conatel/indicadores/
- https://www.tigo.com.py/cobertura
- https://www.personal.com.py/institucional/cobertura.html
- https://www.nperf.com/en/map/PY/-/-/signal


\subsection{Internet}

\textbf{Fonte:} CONATEL / Tigo / Claro / INE / Copaco
\textbf{Ano de referência:} 2024
\textbf{Atualizado em:} 2026-03-20


\begin{tabularx}{\linewidth}{lX}
\toprule
Indicador & Valor \\
\midrule
\% população com internet & ~62\% \\
\% com banda larga fixa & ~18\% dos domicílios \\
Velocidade média download (Mbps) & 20--50 (urbano) / 4--12 (rural) \\
Velocidade média upload (Mbps) & 8--18 (urbano) \\
\bottomrule
\end{tabularx}


\begin{tabularx}{\linewidth}{lXX}
\toprule
Tecnologia & Disponibilidade & Velocidade Típica \\
\midrule
Fibra ótica (FTTH) & San Juan Bautista, San Ignacio & 30--200 Mbps \\
Cabo HFC & San Juan Bautista & 20--100 Mbps \\
Rádio/Wireless & Municípios ao longo da Ruta 1 & 5--40 Mbps \\
ADSL (Copaco) & Interior & 1--8 Mbps \\
Satélite (Starlink) & Todo o departamento & 50--200 Mbps \\
\bottomrule
\end{tabularx}


\begin{tabularx}{\linewidth}{lXXX}
\toprule
Provedor & Tecnologia & Área de cobertura & Preço ref. (USD/mês) \\
\midrule
Tigo & Fibra/Rádio & San Juan Bautista + municípios da Ruta 1 & 18--45 \\
Claro/Personal & Fibra/Rádio & San Juan Bautista & 15--40 \\
Copaco & ADSL/Rádio & Interior & 10--25 \\
Starlink & Satélite LEO & Todo o departamento & ~44 + equip. ~370 \\
\bottomrule
\end{tabularx}


A cobertura rural de Misiones é \textbf{moderada} ao longo da Ruta 1, e \textbf{limitada} no interior oeste:

- Municípios na faixa da Ruta 1 têm ISPs locais e rádio wireless com qualidade razoável.
- Interior oeste (margens do rio Paraguai, zonas de Humaitá e Santiago) tem conectividade muito fraca.
- Copaco mantém ADSL em municípios com centrais telefônicas históricas.
- Starlink é a solução para propriedades rurais afastadas.


- https://www.conatel.gov.py/conatel/wp-content/uploads/2025/02/informe-de-internet-fijo-2024.pdf
- https://www.ine.gov.py/noticias/2037/como-esta-el-acceso-a-las-tecnologias-de-la-informacion-y-comunicacion-en-el-paraguay
- https://dplnews.com/starlink-ya-esta-disponible-en-paraguay-cuanto-cuesta/
- https://www.tigo.com.py/internet



\clearpage
\section{Paraguarí}
\label{sec:dept_09_Paraguari}

\subsection{Segurança Pública}

\textbf{Fonte principal:} Ministerio Público Paraguay / INE / Atlas da Violência CONACYT
\textbf{Data de referência:} 2022--2024
\textbf{Atualizado em:} 2026-03-20


\begin{tabularx}{\linewidth}{lX}
\toprule
Indicador & Valor \\
\midrule
Taxa de homicídios (por 100k hab) & ~3,0 (estimativa, abaixo da média nacional) \\
Crimes registrados (total/ano) & --- dados não desagregados publicamente \\
Número de comisarías & ~15 (estimativa departamental) \\
Índice segurança relativo & alto \\
\bottomrule
\end{tabularx}


Paraguarí é historicamente reconhecido como um dos departamentos mais seguros do Paraguai. Pesquisas do CONACYT e dados do INE apontam que Paraguarí, junto com Cordillera e Central, mantêm uma das mais baixas taxas de homicídio do país, com \textbf{média histórica de aproximadamente 3 mortes por 100 mil habitantes}.

O Ministerio Público registra que o departamento responde por apenas \textbf{1\% das causas de homicídio doloso} nacional (2019--2024) --- uma das menores proporções do país para um departamento de tamanho médio. Paraguarí não possui fronteira internacional e está localizado no centro-sul do Paraguai, distante das rotas principais de narcotráfico.

A economia é baseada em agricultura familiar, pecuária e, em alguns municípios, turismo. Paraguarí (cidade capital) e Carapeguá são os maiores centros urbanos. A criminalidade é de baixa intensidade, predominantemente patrimonial.


- Paraguarí é uma excelente opção para imigrantes brasileiros que buscam ambiente seguro e tranquilo
- Próximo a Asunción (~70 km da cidade de Paraguarí), permite acesso rápido à capital
- Estâncias e propriedades rurais na região são opções atrativas para quem busca residência no campo com segurança
- Carapeguá é conhecida pela indústria têxtil e artesanato --- oportunidade de negócios no setor
- Cidades da região são tipicamente pacatas, com boa qualidade de vida para famílias
- Não há registro de atuação relevante de crime organizado transnacional no departamento


- https://www.ministeriopublico.gov.py/nota/estadisticas-del-ministerio-publico-la-mayoria-de-los-homicidios-dolosos-en-10564
- https://www.conacyt.gov.py/sites/default/files/upload\_editores/u294/ATLAS\_DE\_VIOLENCIA\_PARAGUAY.pdf
- https://ods.ine.gov.py/ine-main/ods/paz-justicias-e-instituciones-solidas-16/meta-16.1/indicador-284
- https://cienciasdelsur.com/2018/04/17/ciudades-mas-violentas-de-paraguay/


\subsection{IDH e Desenvolvimento Humano}

\textbf{Fonte principal:} PNUD Paraguai / DGEEC / INE
\textbf{Ano de referência:} 2020 (IDH); 2023 (pobreza e infraestrutura)
\textbf{Atualizado em:} 2026-03-20


\begin{tabularx}{\linewidth}{lX}
\toprule
Indicador & Valor \\
\midrule
IDH & 0,706 \\
Classificação IDH & alto \\
Ranking nacional (1=melhor) & 9/18 \\
Esperança de vida & 72,6 anos \\
Anos de escolaridade média & 8,1 anos \\
RNB per capita (USD PPA) & 8.500 \\
\bottomrule
\end{tabularx}

> Nota: IDH estimado com base no relatório PNUD 2020 e dados subnacionais regionais. Paraguarí está próximo a Cordillera em perfil socioeconômico, com importante corredor produtivo (algodão, mandioca, pecuária) e acesso à capital pela Rota 1.


\begin{tabularx}{\linewidth}{lX}
\toprule
Indicador & Valor \\
\midrule
Taxa de pobreza total (\%) & 30,5\% \\
Taxa de pobreza extrema (\%) & 7,1\% \\
Índice de Gini & 0,453 \\
\bottomrule
\end{tabularx}

> Fonte: INE --- EPHC 2023 (dado de 30,5\% confirmado para Paraguarí). Pobreza multidimensional (IPM) de 26,77\% em 2023, um dos departamentos com índice multidimensional mais alto entre os da Região Oriental central.


\begin{tabularx}{\linewidth}{lX}
\toprule
Indicador & Valor \\
\midrule
Acesso a água potável (\%) & 77,6\% \\
Acesso a saneamento (\%) & 75,8\% \\
\bottomrule
\end{tabularx}

> Estimativa baseada em dados do Censo 2022. Paraguarí tem boa conectividade pela Rota 1 (Asunción--Encarnación), o que favorece o acesso a serviços em cidades próximas.


Paraguarí tem IDH alto (0,706) e é cortado pela principal rodovia do Paraguai (Rota 1), que liga Asunción ao sul do país. A capital departamental Paraguarí é cidade pequena mas bem conectada. O departamento tem potencial turístico (Parque Nacional Ybycuí, Chololo) e agropecuário. Para brasileiros, pode ser interessante como moradia com fácil acesso a Asunción (1--1,5h), combinando custo baixo com qualidade de vida razoável. A pobreza de 30,5\% e o IPM de 26,77\% indicam desigualdades internas. O acesso a água (77,6\%) e saneamento (75,8\%) são medianos --- melhores na capital departamental, piores nas zonas rurais. Oportunidades em agricultura, pecuária e turismo.


- PNUD Paraguai --- Informe IDH 2001-2020: https://www.undp.org/es/paraguay/noticias/IDH-2001-2020
- INE --- EPHC 2023: https://www.ine.gov.py/noticias/1928/pobreza-monetaria-en-el-2023-se-ubico-en-el-227
- INE --- Pobreza Multidimensional IPM 2023 (Paraguarí 26,77\%): https://www.datos.gov.py/dataset/\%C3\%ADndice-de-pobreza-multidimensional-ipm-2022-y-2023
- INE --- Cen

\subsection{Sistema Prisional}

\textbf{Fonte:} MJT (Ministerio de Justicia y Trabajo) / MNP Paraguay
\textbf{Ano de referência:} 2023--2024
\textbf{Atualizado em:} 2026-03-20


\begin{tabularx}{\linewidth}{lXXXX}
\toprule
Nome & Município & Capacidade & População & Superlotação \\
\midrule
Penitenciaría Regional de Paraguarí & Paraguarí & 150 & ~280 & ~187\% \\
\bottomrule
\end{tabularx}


\begin{tabularx}{\linewidth}{lX}
\toprule
Indicador & Valor \\
\midrule
Total de estabelecimentos & 1 \\
Capacidade total (vagas) & ~150 \\
População carcerária total & ~280 \\
Taxa de superlotação & ~187\% \\
\% presos provisórios & ~66\% \\
\bottomrule
\end{tabularx}


O departamento de Paraguarí conta com uma única penitenciaría regional, em sua capital homônima. O estabelecimento atende detentos dos 17 municípios do departamento, incluindo Carapeguá, Ybycuí, La Colmena e Acahay.

Paraguarí tem localização privilegiada: faz fronteira com o departamento Central e está a cerca de 60--80 km de Asunción. Isso facilita o acesso a serviços da capital, mas também cria pressão adicional sobre o sistema de justiça local, que recebe casos transbordados da região metropolitana.

O perfil criminal é predominantemente de crimes contra a propriedade, violência doméstica e pequeno tráfico de drogas. Não há presença significativa de crime organizado de grande escala.

\textbf{Para imigrantes:} Paraguarí é um departamento com boa qualidade de vida e proximidade com Asunción. Municípios como Carapeguá e Ybycuí têm crescimento urbano moderado e oferecem oportunidades para residência e negócios de menor escala. A presença do presídio em Paraguarí tem impacto mínimo na segurança geral do departamento.


- https://www.mjt.gov.py/index.php/penitenciaria
- https://www.mnp.gov.py/informes
- MNP --- Informe Anual 2022


\subsection{Recursos Hídricos Subterrâneos}

\textbf{Fonte:} SENASA / MADES / Geología del Paraguay / The Groundwater Project
\textbf{Atualizado em:} 2026-03-20


\begin{tabularx}{\linewidth}{lXXX}
\toprule
Aquífero & Cobertura no Dept. & Profundidade Média & Vazão Típica \\
\midrule
Patiño (freático, porção norte) & ~25\% & 30--100 m & 5--15 m$^3$/h \\
Guarani / Misiones (confinado) & ~60\% & 100--400 m & 20--70 m$^3$/h \\
Cristalino (fraturado, Serra de Paraguarí) & ~15\% & 30--80 m & 2--8 m$^3$/h \\
\bottomrule
\end{tabularx}


\begin{tabularx}{\linewidth}{lX}
\toprule
Parâmetro & Valor \\
\midrule
Profundidade média --- Patiño (m) & 30--100 \\
Profundidade média --- Guarani (m) & 100--400 \\
Profundidade máxima recomendada (m) & 450 \\
Vazão típica (m$^3$/h) & 5--50 \\
Custo estimado perfuração (USD) & 2.500 -- 10.000 \\
Qualidade da água --- Patiño & boa a moderada \\
Qualidade da água --- Guarani & boa a excelente \\
pH médio & 6,5--7,5 \\
Outorga necessária & sim \\
Órgão regulador & MADES --- DGPCRH \\
\bottomrule
\end{tabularx}


Paraguarí é o departamento onde o Aquífero Patiño tem sua extensão sul, antes de passar a ser dominado pelo sistema Guarani/Misiones. A Serra de Paraguarí (embasamento cristalino) oferece captação em fraturas com boa qualidade. O Aquífero Guarani é acessível na porção centro-sul do departamento. A qualidade geral é boa. A cidade de Paraguarí (capital) e Carapeguá têm redes públicas do SENASA. Para propriedades rurais na Serra de Ybycuí e arredores, o cristalino com captação rasa é opção econômica. Uso principal: doméstico e pecuária.


Paraguarí combina paisagem serrana com proximidade de Asunción (~70 km). Para sítios e chácaras perto da capital, o Aquífero Patiño (30--80 m) oferece captação econômica (USD 2.500--6.000). No sul do departamento, o Guarani (100--250 m) garante vazão robusta para propriedades maiores. Ybycuí, com seu Parque Nacional, é área de crescente interesse para residência. A Serra de Ybycuí tem aquíferos cristalinos com boa qualidade. Sempre exigir licença do MADES antes de perfurar.


- https://www.geologiadelparaguay.com.py/Acuifero-Pati\%C3\%B1o.htm
- https://gw-project.org/es/caracteristicas-del-agua-subterranea-en-paraguay/
- https://www.geologiadelparaguay.com.py/Acuiferos-Potenciales-del-Paraguay.pdf
- https://www.mades.gov.py/tramites/registro-de-empresas-publicas-y-privadas-que-se-dedican-a-la-perforacion-de-pozos-tubulares/


\subsection{Aptidão do Solo}

\textbf{Fonte:} MAG / SENAVE / Geología del Paraguay / CONACYT
\textbf{Atualizado em:} 2026-03-20


\begin{tabularx}{\linewidth}{lXX}
\toprule
Tipo de Solo & \% do Território & Características \\
\midrule
Ultissolos (Argissolos) & 45\% & Solos podzólicos com horizonte B textural, textura areno-argilosa, fertilidade baixa a média \\
Entissolos arenosos & 30\% & Solos arenosos derivados de arenito, baixa retenção hídrica, muito suscetíveis à erosão \\
Inceptissolos & 15\% & Solos em formação, encostas e serras (Cordilheira de los Altos); moderada fertilidade \\
Oxissolos (fragmentados) & 10\% & Pequenas manchas de solos basálticos no sul do departamento \\
\bottomrule
\end{tabularx}


\begin{tabularx}{\linewidth}{lXX}
\toprule
Classe & \% Território & Uso Recomendado \\
\midrule
Lavoura intensiva & 5\% & Apenas nas pequenas manchas de Oxissolos \\
Lavoura extensiva & 25\% & Mandioca, milho, amendoim, cana-de-açúcar \\
Pastagem natural/melhorada & 35\% & Pecuária extensiva \\
Floresta / reserva & 25\% & Serras, matas ciliares, área de recarga hídrica \\
Sem aptidão agrícola & 10\% & Afloramentos rochosos, áreas urbanas, solos severamente erodidos \\
\bottomrule
\end{tabularx}


\begin{tabularx}{\linewidth}{lX}
\toprule
Parâmetro & Valor \\
\midrule
pH médio & 4,8 -- 5,5 \\
Fertilidade natural & baixa \\
Teor de matéria orgânica & baixo \\
Necessidade de calcário & sim \\
\bottomrule
\end{tabularx}


- \textbf{Mandioca} --- base absoluta da agricultura familiar
- \textbf{Milho} em consórcio com mandioca e feijão
- \textbf{Cana-de-açúcar} --- atividade histórica (produção de caña paraguaya / kañy)
- \textbf{Amendoim} --- em solos arenosos da região leste
- \textbf{Pecuária bovina} --- semi-intensiva a extensiva
- \textbf{Horticultura} nas várzeas e áreas planas


- \textbf{Solos mais degradados do país:} Paraguarí é citado juntamente com Guairá, Cordillera e Central como tendo mais de 90\% dos solos cultivados deteriorados
- \textbf{Erosão hídrica severa:} solos arenosos em relevo ondulado (Serranía de los Altos) são extremamente vulneráveis; erosão em sulcos e ravinas é comum
- \textbf{Acidez elevada:} pH abaixo de 5,5 generalizado; solos com baixa saturação de bases
- \textbf{Matéria orgânica muito baixa:} décadas de cultivo sem reposição esgotaram a matéria orgânica do solo
- \textbf{Recursos hídricos comprometidos:} desmatamento nas serras reduz recarga hídrica e aumenta assoreamento dos rios


Paraguarí apresenta baixo intere

\subsection{Terra Rural}

\textbf{Fonte:} INDERT / Clasificados.com.py / mercado local
\textbf{Ano de referência:} 2024
\textbf{Atualizado em:} 2026-03-20


\begin{tabularx}{\linewidth}{lXXX}
\toprule
Tipo & Preço (USD/ha) & Faixa & Tendência \\
\midrule
Agrícola alta produtividade & 4.800 & 3.200 -- 7.000 & $\uparrow$ \\
Pastagem & 2.400 & 1.600 -- 3.500 & $\uparrow$ \\
Mata / reserva & 1.000 & 600 -- 1.800 & $\rightarrow$ \\
\bottomrule
\end{tabularx}

> Referência de listagem observada: USD 2.400/ha para terras mistas em Paraguarí.


\begin{tabularx}{\linewidth}{lX}
\toprule
Parâmetro & Valor \\
\midrule
Valorização anual média (\%) & 6--8\% \\
Lote mínimo rural (ha) & 10 \\
Imposto territorial rural (\%) & 1\% sobre valor fiscal \\
Liquidez do mercado & média \\
\bottomrule
\end{tabularx}


- \textbf{Paraguarí (entorno da capital):} Bom acesso a Asunción (70 km via Ruta 1); terra valorizada por demanda de chácaras de fim de semana da capital.
- \textbf{Yaguarón:} Cidade histórica (Catedral do século XVIII); turismo e horticultura.
- \textbf{Sapucai / La Colmena:} Zona de colonização japonesa; produção de framboesas, morango e horticultura intensiva.
- \textbf{Mbuyapey / Acahay:} Terra de pastagem a preços competitivos; pecuária extensiva.
- \textbf{Pirayú:} Área de transição com Cordillera; chácaras de lazer em crescimento.


\begin{tabularx}{\linewidth}{lXXX}
\toprule
Indicador & Paraguarí (PY) & São Paulo interior (BR) & Paraná (BR) \\
\midrule
Terra agrícola (USD/ha) & 3.200--7.000 & 20.000--45.000 & 15.000--30.000 \\
Chácara de lazer (USD/ha) & 4.000--9.000 & 30.000--80.000 & 15.000--40.000 \\
Custo relativo & ~20\% do SP & referência SP & referência PR \\
\bottomrule
\end{tabularx}

Paraguarí é atrativo para quem deseja chácara de lazer perto de Asunción a custos muito menores do que os equivalentes nas periferias das capitais estaduais brasileiras.


- [INDERT --- Instituto Nacional de Desarrollo Rural y de la Tierra](https://www.indert.gov.py/)
- [Plan B Paraguay --- Farmland Prices](https://x.com/planbparaguay/status/1834232169759211739)
- [Clasipar.com --- Campos Paraguarí](https://clasipar.paraguay.com/alquiler/inmuebles/paraguari)
- [InfoCasas --- Chacras Paraguay](https://www.infocasas.com.py/venta/chacras-o-campos)


\subsection{Mercado Imobiliário}

\textbf{Capital:} Paraguarí
\textbf{Fonte:} Clasificados.com.py / mercado local
\textbf{Ano de referência:} 2024
\textbf{Atualizado em:} 2026-03-20


\begin{tabularx}{\linewidth}{lXX}
\toprule
Tipo & Preço (USD/m$^2$) & Aluguel (USD/mês) \\
\midrule
Apartamento 2 quartos & 550 -- 850 & 200 -- 380 \\
Casa 3 quartos & 480 -- 750 & 180 -- 340 \\
\bottomrule
\end{tabularx}


\begin{tabularx}{\linewidth}{lX}
\toprule
Parâmetro & Valor \\
\midrule
Valorização anual (\%) & 5--7\% \\
Disponibilidade de financiamento & boa (proximidade com Asunción facilita acesso a bancos nacionais) \\
Bairros mais valorizados & Centro, Barrio San José, proximidades da estação ferroviária histórica \\
\bottomrule
\end{tabularx}


\textbf{Vale a pena comprar ou alugar?}
Paraguarí é uma opção interessante para quem trabalha em Asunción mas prefere morar no interior. A cidade fica a 70 km da capital pela Ruta 1 (45--60 min). Custo de vida notavelmente mais baixo que Asunción com qualidade de vida superior em termos de espaço e ambiente. A compra de casa é geralmente vantajosa.

\textbf{Áreas recomendadas:}
- \textbf{Paraguarí cidade:} infraestrutura básica completa, comércio ativo, posto de saúde.
- \textbf{Yaguarón:} cidade histórica e turística a 20 km de Paraguarí; imóveis com valor cultural e chácaras a preços competitivos.
- \textbf{La Colmena:} comunidade japonesa estabelecida, produção intensiva de hortaliças; estilo de vida tranquilo e organizado.

\textbf{Armadilhas comuns:}
- Apartamentos de padrão moderno são raros; mercado dominado por casas com quintal.
- Para famílias com crianças, verificar disponibilidade de boas escolas (La Colmena tem escola japonesa-paraguaia).
- Trens históricos existem mas não são transporte funcional; mobilidade depende de veículo próprio.


- [InfoCasas.com.py --- Paraguarí](https://www.infocasas.com.py/venta/inmuebles/paraguari)
- [Clasipar.com --- Alquiler Paraguarí](https://clasipar.paraguay.com/alquiler/inmuebles/paraguari)
- [Airbnb --- Paraguarí](https://www.airbnb.com/paraguari-paraguay/stays)


\subsection{Cobertura Celular}

\textbf{Fonte:} CONATEL / Tigo / Personal / OpenSignal / nPerf
\textbf{Ano de referência:} 2024
\textbf{Atualizado em:} 2026-03-20


\begin{tabularx}{\linewidth}{lXXXXX}
\toprule
Operadora & 2G & 3G & 4G & 5G & Cobertura Pop. \\
\midrule
Tigo & 92\% & 85\% & 75\% & --- & 88\% \\
Personal/Claro & 88\% & 80\% & 70\% & --- & 84\% \\
VOX & 60\% & 45\% & 28\% & --- & 52\% \\
\bottomrule
\end{tabularx}

> \textbf{Nota:} Paraguarí é um departamento próximo a Asunción (capital a ~60 km) com boa cobertura de telecomunicações. A Ruta 1 e Ruta 2 cruzam o departamento, garantindo cobertura linear sólida. Municípios turísticos como Paraguarí, Caapucú e Sapucái têm boa cobertura.


\begin{tabularx}{\linewidth}{lX}
\toprule
Parâmetro & Valor \\
\midrule
Melhor operadora rural & Tigo \\
Velocidade 4G média (Mbps) & 22--42 \\
Áreas sem cobertura & Interior serrano, zonas de mata no sul \\
Sinal em propriedade rural & bom a moderado \\
\bottomrule
\end{tabularx}


\begin{tabularx}{\linewidth}{lXXX}
\toprule
Plano & Operadora & USD/mês & Dados \\
\midrule
Pré-pago básico (7 dias) & Tigo & ~2,50 & 3 GB \\
Pré-pago intermediário (30 dias) & Personal & ~5,00 & 8 GB \\
Pós-pago entrada & Tigo & ~12,00 & 15 GB + chamadas \\
Pós-pago intermediário & Personal/Claro & ~18,00 & 30 GB + chamadas \\
\bottomrule
\end{tabularx}


- \textbf{Paraguarí cidade, Carapeguá, Quiindy:} Cobertura 4G sólida de Tigo e Personal.
- \textbf{Proximidade com Asunción:} O departamento se beneficia do "efeito metrópole" --- torres construídas para cobrir a capital cobrem também municípios limítrofes.
- \textbf{Serra de Caacupé (fronteira com Cordillera):} Sinal pode ser irregular em vales fechados.
- \textbf{Ybycuí e interior sul:} Cobertura 3G/4G razoável de Tigo.
- Paraguarí é uma opção interessante para imigrantes que querem morar no interior com acesso fácil a Asunción --- conectividade adequada.


- https://www.conatel.gov.py/conatel/indicadores/
- https://www.tigo.com.py/cobertura
- https://www.personal.com.py/institucional/cobertura.html
- https://www.nperf.com/en/map/PY/-/-/signal


\subsection{Internet}

\textbf{Fonte:} CONATEL / Tigo / Claro / INE / Copaco
\textbf{Ano de referência:} 2024
\textbf{Atualizado em:} 2026-03-20


\begin{tabularx}{\linewidth}{lX}
\toprule
Indicador & Valor \\
\midrule
\% população com internet & ~68\% \\
\% com banda larga fixa & ~25\% dos domicílios \\
Velocidade média download (Mbps) & 30--70 (urbano) / 6--18 (rural) \\
Velocidade média upload (Mbps) & 10--25 (urbano) \\
\bottomrule
\end{tabularx}

> Paraguarí se beneficia da proximidade com a área metropolitana de Asunción, com melhor conectividade que departamentos mais distantes da capital.


\begin{tabularx}{\linewidth}{lXX}
\toprule
Tecnologia & Disponibilidade & Velocidade Típica \\
\midrule
Fibra ótica (FTTH) & Paraguarí cidade, Carapeguá, Quiindy & 50--300 Mbps \\
Cabo HFC & Paraguarí cidade & 30--150 Mbps \\
Rádio/Wireless & Amplo no departamento & 5--50 Mbps \\
ADSL (Copaco) & Interior com rede legada & 2--10 Mbps \\
Satélite (Starlink) & Todo o departamento & 50--200 Mbps \\
\bottomrule
\end{tabularx}


\begin{tabularx}{\linewidth}{lXXX}
\toprule
Provedor & Tecnologia & Área de cobertura & Preço ref. (USD/mês) \\
\midrule
Tigo & Fibra/HFC/Rádio & Paraguarí + municípios principais & 18--55 \\
Claro/Personal & Fibra/Rádio & Paraguarí cidade + Carapeguá & 15--50 \\
Copaco & ADSL/Fibra & Interior & 10--28 \\
Starlink & Satélite LEO & Todo o departamento & ~44 + equip. ~370 \\
\bottomrule
\end{tabularx}


A cobertura rural de Paraguarí é \textbf{moderada}, boa ao longo das rotas principais e mais limitada no interior:

- Municípios próximos às Rutas 1 e 2 têm boa conectividade.
- Interior serrano (Ybycuí, La Colmena) tem cobertura de rádio wireless e internet móvel.
- Ybycuí (Parque Nacional) tem sinal celular 3G/4G razoável para área protegida.
- Starlink disponível para propriedades rurais afastadas.


- https://www.conatel.gov.py/conatel/wp-content/uploads/2025/02/informe-de-internet-fijo-2024.pdf
- https://www.ine.gov.py/noticias/2037/como-esta-el-acceso-a-las-tecnologias-de-la-informacion-y-comunicacion-en-el-paraguay
- https://www.tigo.com.py/internet
- https://dplnews.com/starlink-ya-esta-disponible-en-paraguay-cuanto-cuesta/



\clearpage
\section{Alto Paraná}
\label{sec:dept_10_Alto_Parana}

\subsection{Segurança Pública}

\textbf{Fonte principal:} Ministerio Público Paraguay / InSight Crime / Prosegur Research
\textbf{Data de referência:} 2023--2024
\textbf{Atualizado em:} 2026-03-20


\begin{tabularx}{\linewidth}{lX}
\toprule
Indicador & Valor \\
\midrule
Taxa de homicídios (por 100k hab) & ~8--12 (estimativa, acima da média nacional) \\
Crimes registrados (total/ano) & --- dados não desagregados publicamente \\
Número de comisarías & ~35 (estimativa --- departamento populoso com grande fronteira) \\
Índice segurança relativo & baixo \\
\bottomrule
\end{tabularx}


Alto Paraná é o departamento mais economicamente dinâmico do Paraguai (sede de Itaipu Binacional e de Ciudad del Este) e ao mesmo tempo um dos mais violentos. O Ministerio Público registra que o departamento responde por \textbf{aproximadamente 11\% das causas de homicídio doloso} nacional (2019--2024) --- terceira posição, atrás apenas de Amambay e Central.

Ciudad del Este integra a \textbf{Tríplice Fronteira} (Paraguai, Brasil, Argentina), considerada o maior polo do crime organizado transnacional da América do Sul. O PCC (Primeiro Comando da Capital) e outras organizações criminosas brasileiras e transnacionais operam ativamente na região. O contrabando, a pirataria, o narcotráfico e a lavagem de dinheiro são estruturais no tecido econômico da cidade.

Em 2025, Argentina, Brasil e Paraguay firmaram novo acordo para reforçar o Comando Tripartito na Tríplice Fronteira. As operações conjuntas são frequentes, mas a permeabilidade da fronteira --- especialmente a Ponte Internacional da Amizade (Foz do Iguaçu-Ciudad del Este) --- torna o controle complexo.

A violência relacionada ao crime organizado (sicariatos, disputas entre facções) é documentada na região. Assassinatos de figuras ligadas ao tráfico ocorrem com regularidade em Ciudad del Este e cidades do entorno.


- Alto Paraná tem a maior concentração de brasileiros no Paraguai --- a chamada "Colônia" (fazendeiros brasiguaios) e comunidade urbana em Ciudad del Este
- Ciudad del Este oferece imenso dinamismo comercial, mas o ambiente de segurança é claramente mais tenso que outras regiões
- Bairros residenciais de classe média e alta em Ciudad del Este (Ex-Amistad, Villa Alegre) são relativamente seguros com segurança privada
- Foz do Iguaçu/Ciudad del Este é área de atuação intensa do PCC --- brasileiros devem ser extremamente cautelosos com envolvimentos em negócios informais
- O interior do departamento (cidades como Minga Guazú, He

\subsection{IDH e Desenvolvimento Humano}

\textbf{Fonte principal:} PNUD Paraguai / DGEEC / INE
\textbf{Ano de referência:} 2020 (IDH); 2023 (pobreza e infraestrutura)
\textbf{Atualizado em:} 2026-03-20


\begin{tabularx}{\linewidth}{lX}
\toprule
Indicador & Valor \\
\midrule
IDH & 0,752 \\
Classificação IDH & alto \\
Ranking nacional (1=melhor) & 3/18 \\
Esperança de vida & 75,1 anos \\
Anos de escolaridade média & 9,6 anos \\
RNB per capita (USD PPA) & 13.500 \\
\bottomrule
\end{tabularx}

> Nota: IDH de 2020 conforme relatório PNUD "Desarrollo Humano en el Paraguay 2001-2020". Alto Paraná está entre os 7 departamentos do relatório PNUD como 3º lugar, atrás apenas de Asunción e Central. Fonte: PNUD Paraguay 2022, Wikipedia IDH Paraguay.


\begin{tabularx}{\linewidth}{lX}
\toprule
Indicador & Valor \\
\midrule
Taxa de pobreza total (\%) & 16,3\% \\
Taxa de pobreza extrema (\%) & 3,1\% \\
Índice de Gini & 0,428 \\
\bottomrule
\end{tabularx}

> Fonte: INE --- EPHC 2023. Alto Paraná tem uma das menores taxas de pobreza do país, concentrando grande parte das atividades industriais e comerciais de Ciudad del Este e região.


\begin{tabularx}{\linewidth}{lX}
\toprule
Indicador & Valor \\
\midrule
Acesso a água potável (\%) & 88,7\% \\
Acesso a saneamento (\%) & 87,2\% \\
\bottomrule
\end{tabularx}

> Estimativa baseada em dados do Censo 2022. Ciudad del Este possui infraestrutura urbana avançada, com acesso a serviços básicos superior à média nacional.


Alto Paraná é o segundo departamento mais desenvolvido do Paraguai (IDH 0,752) e o principal polo econômico do interior. Ciudad del Este, sua capital, é um dos maiores centros comerciais da América do Sul, fazendo fronteira com Foz do Iguaçu (BR) e Puerto Iguazú (AR) pela Ponte da Amizade. A comunidade brasileira é numerosa e estabelecida --- estima-se que centenas de milhares de brasileiros e brasiguaios vivam na região. O setor comercial é dominado por importações livres de impostos, e a cidade tem hospitais modernos, universidades, shoppings e intensa vida econômica. Com pobreza de apenas 16,3\% e acesso a água (88,7\%) e saneamento (87,2\%) elevados, a qualidade de vida é das melhores no interior do país. Para brasileiros, é a opção mais familiar e acessível: português amplamente falado, produtos brasileiros facilmente encontrados, fluxo diário entre os dois países. Oportunidades em comércio, importação, logística, agronegócio e serviços.


- PNUD Paraguai --- Informe "Índices de Desarrollo Humano en el Paraguay 2001-2020": https://www.undp.org/es/parag

\subsection{Sistema Prisional}

\textbf{Fonte:} MJT (Ministerio de Justicia y Trabajo) / MNP Paraguay
\textbf{Ano de referência:} 2023--2024
\textbf{Atualizado em:} 2026-03-20


\begin{tabularx}{\linewidth}{lXXXX}
\toprule
Nome & Município & Capacidade & População & Superlotação \\
\midrule
Penitenciaría Regional de Ciudad del Este & Ciudad del Este & 600 & ~1.800 & ~300\% \\
Penitenciaría de Minga Guazú & Minga Guazú & 150 & ~310 & ~207\% \\
Penitenciaría de Presidente Franco & Presidente Franco & 80 & ~130 & ~163\% \\
\bottomrule
\end{tabularx}


\begin{tabularx}{\linewidth}{lX}
\toprule
Indicador & Valor \\
\midrule
Total de estabelecimentos & 3 \\
Capacidade total (vagas) & ~830 \\
População carcerária total & ~2.240 \\
Taxa de superlotação & ~270\% \\
\% presos provisórios & ~72\% \\
\bottomrule
\end{tabularx}


\textbf{Alto Paraná concentra o segundo maior sistema penal do Paraguai}, depois do Distrito Capital. A Penitenciaría Regional de Ciudad del Este é um dos estabelecimentos mais superlotados e perigosos do país, com relatos frequentes de domínio de facções criminosas (PCC, Comando Vermelho, grupos locais), assassinatos e fugas.

Ciudad del Este é o maior polo comercial do interior paraguaio e a terceira maior zona franca do mundo, atraindo fluxos massivos de contrabando, lavagem de dinheiro e tráfico de drogas. Isso se reflete diretamente na população carcerária --- parcela significativa detida por crimes financeiros, tráfico e posse de armas.

A região tríplice fronteira (Brasil-Argentina-Paraguai) com Foz do Iguaçu e Puerto Iguazú cria um ambiente de alto risco de segurança. Rebeliões e assassinatos dentro dos presídios de Alto Paraná têm repercussão nacional e frequentemente envolvem disputas de fações pelo controle do tráfico.

O MNP identificou condições críticas: superlotação extrema, violência entre detentos, corrupção de agentes penitenciários e presença de armas e celulares dentro dos presídios.

\textbf{Para imigrantes:} Alto Paraná oferece enormes oportunidades econômicas (Ciudad del Este, agronegócio, logística), mas o ambiente de segurança é o mais desafiador fora de Asunción. Bairros próximos ao presídio e ao centro de Ciudad del Este têm maior criminalidade. Municípios como Hernandarias, Santa Rita e Los Cedrales oferecem ambiente mais tranquilo dentro do mesmo departamento.


- https://www.mjt.gov.py/index.php/penitenciaria
- https://www.mnp.gov.py/informes
- MNP --- Informe Anual 2023
- ABC Color: "Ciudad del Este: penitenciaría en crisis" (2023)
- Últ

\subsection{Recursos Hídricos Subterrâneos}

\textbf{Fonte:} SENASA / MADES / Geología del Paraguay / IGRAC / Guaraní Aquifer System Project
\textbf{Atualizado em:} 2026-03-20


\begin{tabularx}{\linewidth}{lXXX}
\toprule
Aquífero & Cobertura no Dept. & Profundidade Média & Vazão Típica \\
\midrule
Guarani / Misiones (confinado sob basalto) & ~90\% & 120--500 m & 30--120 m$^3$/h \\
Basáltico (fraturado --- Serra de Maracaju) & ~60\% (cobertura superficial) & 40--150 m & 5--20 m$^3$/h \\
Caiuá/Baurú-Acaray (porção norte) & ~20\% & 80--200 m & 10--40 m$^3$/h \\
\bottomrule
\end{tabularx}


\begin{tabularx}{\linewidth}{lX}
\toprule
Parâmetro & Valor \\
\midrule
Profundidade média --- Guarani (m) & 120--400 \\
Profundidade máxima recomendada (m) & 600 \\
Vazão típica (m$^3$/h) & 20--80 \\
Custo estimado perfuração (USD) & 5.000 -- 18.000 \\
Qualidade da água --- Guarani & boa a excelente \\
pH médio & 6,5--7,5 \\
Temperatura da água ($^\circ$C) & 30--55 \\
Outorga necessária & sim \\
Órgão regulador & MADES --- DGPCRH \\
\bottomrule
\end{tabularx}


Alto Paraná é o departamento mais dinâmico economicamente do Paraguai (Ciudad del Este). O Aquífero Guarani está presente sob uma espessa capa de basalto, o que eleva as profundidades necessárias de perfuração (120--400 m). A água é de excelente qualidade, com baixa mineralização, temperatura elevada (benefício para alguns usos industriais) e surgência natural em muitos pontos. O custo de perfuração é mais elevado que em outros departamentos devido à necessidade de perfurar através da rocha basáltica. Paraguay tem mais de 350 poços na Região Oriental, com concentração expressiva em Alto Paraná. Uso principal: abastecimento urbano e rural, indústria, irrigação de soja.


Alto Paraná abriga a maior comunidade brasiguaia do Paraguai. Para qualquer propriedade rural, poço artesiano no Guarani é indispensável, já que a rede pública cobre apenas áreas urbanas. O custo maior (USD 5.000--18.000) reflete a profundidade e a necessidade de perfurar basalto. Para propriedades agrícolas grandes (soja), o investimento se paga rapidamente pela vazão de 30--80 m$^3$/h disponível. Empresa de perfuração deve estar registrada no MADES. Em Ciudad del Este, há diversas empresas especializadas. Licença ambiental é obrigatória.


- https://hidrogeom.com/perforacion-de-pozo-artesiano-cde/
- https://gw-project.org/es/caracteristicas-del-agua-subterranea-en-paraguay/
- https://es.wikipedia.org/wiki/Acu\%C3\%ADfero\_Guaran\%C3\%AD
- https

\subsection{Aptidão do Solo}

\textbf{Fonte:} MAG / SENAVE / INFONA / CONACYT / CAPECO
\textbf{Atualizado em:} 2026-03-20


\begin{tabularx}{\linewidth}{lXX}
\toprule
Tipo de Solo & \% do Território & Características \\
\midrule
Oxissolos (Latossolos Roxos / Terra Roxa) & 65\% & Solos derivados de basalto da Formação Serra Geral; profundos (>180 cm), bem drenados, argilosos, cor vermelha escura intensa; textura franco-argilosa a argilosa \\
Ultissolos (Argissolos Vermelhos) & 20\% & Solos com horizonte B textural, fertilidade média-alta; ocorrem nas transições \\
Inceptissolos / Aluviais & 10\% & Planícies do Rio Paraná; alta fertilidade natural \\
Entissolos (Litossolos) & 5\% & Solos rasos próximos a afloramentos; área residual \\
\bottomrule
\end{tabularx}


\begin{tabularx}{\linewidth}{lXX}
\toprule
Classe & \% Território & Uso Recomendado \\
\midrule
Lavoura intensiva & 65\% & Soja, milho, trigo, canola, girassol \\
Lavoura extensiva & 15\% & Mandioca, arroz, pastagem melhorada \\
Pastagem natural/melhorada & 5\% & Pecuária complementar \\
Floresta / reserva & 10\% & Reservas e remanescentes (Bosque Atlântico del Alto Paraná) \\
Sem aptidão agrícola & 5\% & Área urbanizada (Ciudad del Este, Hernandarias, Minga Guazú) \\
\bottomrule
\end{tabularx}


\begin{tabularx}{\linewidth}{lX}
\toprule
Parâmetro & Valor \\
\midrule
pH médio & 5,2 -- 6,5 \\
Fertilidade natural & alta \\
Teor de matéria orgânica & médio-alto (35\% das amostras < 25 g/kg; maioria acima) \\
Necessidade de calcário & parcial (solos basálticos corrigem com menor dose) \\
\bottomrule
\end{tabularx}


- \textbf{Soja} --- principal cultura; Alto Paraná é o MAIOR produtor nacional (906.000 ha, 25\% da área total; +37.692 ha em 2024-2025)
- \textbf{Milho} --- amplamente cultivado em rotação com soja
- \textbf{Trigo} --- cultivo de inverno muito expressivo
- \textbf{Canola} --- em expansão como oleaginosa de inverno
- \textbf{Girassol} --- presente
- \textbf{Stevia} --- importante cultura de nicho; principal zona produtora do Paraguai
- \textbf{Mandioca} na agricultura familiar


- \textbf{Compactação por mecanização:} Oxissolos submetidos a tráfego intenso de máquinas agrícolas pesadas; estudos confirmam compactação superficial afetando soja
- \textbf{Acidez residual:} mesmo em solos basálticos de alta fertilidade, calagem periódica é necessária
- \textbf{Matéria orgânica em declínio:} 35\% das amostras analisadas com MO < 25 g/kg (CONACYT); plantio direto ajuda mas não resolve sem aporte orgânico
- \textbf{Desmatamento avançado:} Alto

\subsection{Terra Rural}

\textbf{Fonte:} INDERT / Clasificados.com.py / Última Hora / mercado local
\textbf{Ano de referência:} 2024
\textbf{Atualizado em:} 2026-03-20


\begin{tabularx}{\linewidth}{lXXX}
\toprule
Tipo & Preço (USD/ha) & Faixa & Tendência \\
\midrule
Agrícola alta produtividade & 14.800 & 10.000 -- 20.000 & $\uparrow$ \\
Pastagem & 7.000 & 4.500 -- 10.000 & $\uparrow$ \\
Mata / reserva & 3.500 & 2.000 -- 6.000 & $\uparrow$ \\
\bottomrule
\end{tabularx}

> Referência confirmada: USD 14.858/ha para terra agrícola em Alto Paraná --- preço superior ao de Iowa (EUA). Listagem específica: campo agrícola formado para soja/milho a USD 10.000/ha.


\begin{tabularx}{\linewidth}{lX}
\toprule
Parâmetro & Valor \\
\midrule
Valorização anual média (\%) & 8--12\% \\
Lote mínimo rural (ha) & 10 \\
Imposto territorial rural (\%) & 1\% sobre valor fiscal \\
Liquidez do mercado & alta \\
\bottomrule
\end{tabularx}


- \textbf{Minga Guazú / Hernandarias:} Zona próxima à Itaipu; alta industrialização e demanda; terra mais valorizada do departamento.
- \textbf{Santa Rita / Los Cedrales:} Polo consolidado de soja e milho; colônias brasileiras de primeira geração; infraestrutura agropecuária completa.
- \textbf{Minga Porã / Naranjal:} Expansão recente; terras ainda com bom custo-benefício dentro do departamento.
- \textbf{Ciudad del Este (entorno rural):} Pressão urbana da maior cidade comercial do Paraguai; terra peri-urbana com altíssima valorização.
- \textbf{Santa Rosa del Monday:} Corredor soja com acesso à ruta 6; muito procurado por produtores brasileiros.


\begin{tabularx}{\linewidth}{lXXX}
\toprule
Indicador & Alto Paraná (PY) & Paraná (BR) & Iowa (EUA) \\
\midrule
Terra agrícola (USD/ha) & 10.000--20.000 & 15.000--30.000 & ~7.350 \\
Tendência 2019--2024 & +60--80\% & +40--60\% & +35\% \\
\bottomrule
\end{tabularx}

Alto Paraná já superou preços de algumas zonas do Paraná brasileiro e da própria Iowa americana. A valorização acelerada é impulsionada pela alta produtividade (soja, milho, trigo em triplo cultivo), excelente logística e crescimento do agronegócio.


- [Última Hora --- Tierras de Alto Paraná más caras que EEUU](https://www.ultimahora.com/tierras-agricolas-alto-parana-son-mas-caras-que-las-eeuu-n2928594)
- [Clasipar --- Campo Agrícola Alto Paraná](https://clasipar.paraguay.com/inmuebles/terrenos/alto-parana-20000-has-campo-agricola-formado-soja-maiz-girasol-etc-10000-usdha-2273614)
- [INDERT](https://www.indert.gov.py/)
- [Plan B Paraguay --- Farmland Prices](https://x.com/planbparaguay/status/18342321697

\subsection{Mercado Imobiliário}

\textbf{Capital:} Ciudad del Este
\textbf{Fonte:} InfoCasas / MercadoLibre / Century 21 / mercado local
\textbf{Ano de referência:} 2024
\textbf{Atualizado em:} 2026-03-20


\begin{tabularx}{\linewidth}{lXX}
\toprule
Tipo & Preço (USD/m$^2$) & Aluguel (USD/mês) \\
\midrule
Apartamento 1 quarto (studio) & 1.100 -- 1.600 & 350 -- 600 \\
Apartamento 2 quartos & 1.300 -- 1.800 & 500 -- 900 \\
Apartamento 3 quartos / penthouse & 1.500 -- 2.200 & 800 -- 1.500 \\
Casa 3 quartos & 900 -- 1.400 & 450 -- 800 \\
\bottomrule
\end{tabularx}

> Exemplos reais 2024: apt 44m$^2$ = USD 69.000 (USD 1.568/m$^2$); apt 58m$^2$ 2q = USD 97.650 (USD 1.674/m$^2$); apt 82m$^2$ 1q = USD 150.000 (USD 1.829/m$^2$); penthouse 154m$^2$ = USD 185.500+.


\begin{tabularx}{\linewidth}{lX}
\toprule
Parâmetro & Valor \\
\midrule
Valorização anual (\%) & 8--12\% \\
Disponibilidade de financiamento & boa (bancos paraguaios, cooperativas, financiamentos diretos de construtoras) \\
Bairros mais valorizados & Km 4, Km 6, Minga Guazú, Residencial Villa Morra CDE, Santa Rita \\
\bottomrule
\end{tabularx}


\textbf{Vale a pena comprar ou alugar?}
Ciudad del Este é a segunda maior cidade do Paraguai e uma das maiores zonas de livre comércio do mundo. O mercado imobiliário é dinâmico e em forte expansão, impulsionado pelo comércio transfronteiriço e pelo agronegócio do entorno. A compra tende a ser vantajosa para quem permanece 3+ anos, especialmente em bairros residenciais fora do centro comercial.

\textbf{Áreas recomendadas:}
- \textbf{Km 4 / Km 6 (residential):} Bairros consolidados com melhor qualidade de vida, longe do caos do centro comercial.
- \textbf{Minga Guazú:} Município vizinho mais tranquilo e residencial; muito popular entre famílias brasileiras.
- \textbf{Santa Rita:} Cidade dormitório com forte comunidade brasileira; casas grandes a custo menor que CDE.
- \textbf{Presidente Franco:} Margem do Rio Paraná; em urbanização crescente.

\textbf{Armadilhas comuns:}
- Centro comercial de Ciudad del Este é caótico, congestionado e com altos índices de criminalidade; evitar morar no centro.
- Mercado imobiliário tem parcela informal grande; exigir documentação completa (escritura pública, sertidão do Registro de Inmuebles).
- Apartamentos "en pozo" abundantes; confirmar solidez financeira da construtora.
- A flutuação do comércio com o Brasil (câmbio, política fiscal) impacta diretamente a economia local e pode afetar valorização.


- [InfoCasas --- Venta Departamentos Ciudad del Este](https://www.infocasas.com.py/venta/de

\subsection{Cobertura Celular}

\textbf{Fonte:} CONATEL / Tigo / Personal / OpenSignal / nPerf
\textbf{Ano de referência:} 2024
\textbf{Atualizado em:} 2026-03-20


\begin{tabularx}{\linewidth}{lXXXXX}
\toprule
Operadora & 2G & 3G & 4G & 5G & Cobertura Pop. \\
\midrule
Tigo & 97\% & 95\% & 92\% & 8\% & 96\% \\
Personal/Claro & 95\% & 92\% & 88\% & 5\% & 94\% \\
VOX & 70\% & 60\% & 45\% & --- & 65\% \\
\bottomrule
\end{tabularx}

> \textbf{Nota:} Alto Paraná é o segundo departamento mais populoso do Paraguai (atrás de Central), com Ciudad del Este como maior cidade do interior (~400k hab.). A cobertura 4G é excelente em CDE e região metropolitana (Hernandarias, Minga Guazú, Presidente Franco). 5G em fase inicial em CDE. Interior norte (San Cristóbal, Itakyry) tem cobertura mais variável.


\begin{tabularx}{\linewidth}{lX}
\toprule
Parâmetro & Valor \\
\midrule
Melhor operadora & Tigo (cobertura geral); Personal/Claro (boa relação custo-benefício) \\
Velocidade 4G média (Mbps) & 35--60 \\
Áreas sem cobertura & Norte do departamento, áreas de reserva florestal \\
Sinal em propriedade rural & bom (sul/centro) a moderado (norte) \\
\bottomrule
\end{tabularx}


\begin{tabularx}{\linewidth}{lXXX}
\toprule
Plano & Operadora & USD/mês & Dados \\
\midrule
Pré-pago básico (7 dias) & Tigo & ~2,50 & 3 GB \\
Pré-pago intermediário (30 dias) & Personal & ~5,00 & 8 GB \\
Pós-pago entrada & Tigo & ~12,00 & 15 GB + chamadas \\
Pós-pago intermediário & Personal/Claro & ~18,00 & 30 GB + chamadas \\
Pós-pago premium & Tigo & ~28,00 & ilimitado \\
\bottomrule
\end{tabularx}


- \textbf{Ciudad del Este:} Segunda maior infraestrutura de telecom do Paraguai. Cobertura 4G completa; 5G incipiente. Chip qualquer operadora funciona bem.
- \textbf{Região metropolitana de CDE (Hernandarias, Minga Guazú, Presidente Franco):} Cobertura 4G excelente de Tigo e Personal.
- \textbf{Colônias agrícolas brasileiras (Naranjal, Santa Rita, Los Cedrales, Dr. Moisés Bertoni):} Boa cobertura 4G de Tigo; Personal crescente.
- \textbf{Fronteira com Brasil (Foz do Iguaçu):} Chips brasileiros captam sinal em CDE e vice-versa. Adquirir chip paraguaio para tarifas locais.
- \textbf{Roaming:} Alto Paraná é o departamento com maior concentração de brasileiros no Paraguai; chips Tigo PY disponíveis facilmente.
- \textbf{Interior norte:} Tigo é a única opção confiável; Personal tem cobertura mais limitada.


- https://www.conatel.gov.py/conatel/indicadores/
- https://www.tigo.com.py/cobertura
- https://www.personal.com.py/institucional/cobertura.html
- https://w

\subsection{Internet}

\textbf{Fonte:} CONATEL / Tigo / Claro / INE / Speedtest
\textbf{Ano de referência:} 2024
\textbf{Atualizado em:} 2026-03-20


\begin{tabularx}{\linewidth}{lX}
\toprule
Indicador & Valor \\
\midrule
\% população com internet & 87,1\% \\
\% com banda larga fixa & ~52\% dos domicílios \\
Velocidade média download (Mbps) & 60--120 (CDE) / 15--40 (interior) \\
Velocidade média upload (Mbps) & 20--50 (CDE) \\
\bottomrule
\end{tabularx}

> Fonte: INE EPHC 2024 --- 87,1\% de penetração em Alto Paraná, segundo departamento com maior acesso a internet do Paraguai, atrás apenas de Central (87,6\%) e Asunción (88,9\%).


\begin{tabularx}{\linewidth}{lXX}
\toprule
Tecnologia & Disponibilidade & Velocidade Típica \\
\midrule
Fibra ótica (FTTH) & CDE, Hernandarias, Minga Guazú, Santa Rita & 100--1000 Mbps \\
Cabo HFC & CDE e região metropolitana & 50--300 Mbps \\
Rádio/Wireless & Amplo no departamento & 10--80 Mbps \\
ADSL (Copaco) & Interior com rede legada & 2--15 Mbps \\
Satélite (Starlink) & Todo o departamento & 50--200 Mbps \\
\bottomrule
\end{tabularx}


\begin{tabularx}{\linewidth}{lXXX}
\toprule
Provedor & Tecnologia & Área de cobertura & Preço ref. (USD/mês) \\
\midrule
Tigo & Fibra/HFC/Rádio & CDE + toda a região & 18--65 \\
Claro/Personal & Fibra/HFC/Rádio & CDE + municípios & 15--60 \\
Copaco & Fibra/ADSL & CDE + interior & 10--35 \\
VOX & Fibra/Rádio & CDE & 15--45 \\
ISPs locais & Rádio wireless & Colônias agrícolas & 10--30 \\
Starlink & Satélite LEO & Todo o departamento & ~44 + equip. ~370 \\
\bottomrule
\end{tabularx}


A cobertura rural de Alto Paraná é \textbf{boa}, a melhor entre os departamentos do interior:

- \textbf{Colônias agrícolas brasileiras} ao sul e centro (Santa Rita, Naranjal, Los Cedrales): ISPs locais estabelecidos com rádio wireless, 10--50 Mbps.
- \textbf{Área da Itaipú Binacional} (Hernandarias): Excelente cobertura de fibra e rádio.
- \textbf{Interior norte} (Itakyry, San Cristóbal): Cobertura mais fraca; rádio wireless de ISPs locais ou internet móvel Tigo.
- \textbf{Reserva San Rafael e áreas de mata:} Cobertura celular muito fraca; Starlink necessário.

Alto Paraná é o \textbf{melhor departamento do interior para imigrantes} em termos de conectividade, com CDE oferecendo infraestrutura de telecom de nível metropolitano e colônias rurais razoavelmente conectadas.


- https://www.conatel.gov.py/conatel/wp-content/uploads/2025/02/informe-de-internet-fijo-2024.pdf
- https://marketdata.com.py/laboratorio/analisis/acceso-a-internet-en-paraguay-alcan


\clearpage
\section{Central}
\label{sec:dept_11_Central}

\subsection{Segurança Pública}

\textbf{Fonte principal:} Ministerio Público Paraguay / INE / La Nación PY
\textbf{Data de referência:} 2023--2024
\textbf{Atualizado em:} 2026-03-20


\begin{tabularx}{\linewidth}{lX}
\toprule
Indicador & Valor \\
\midrule
Taxa de homicídios (por 100k hab) & ~3,0 (estimativa, abaixo da média nacional) \\
Crimes registrados (total/ano) & elevado em números absolutos --- maior departamento populoso \\
Número de comisarías & ~60 (estimativa --- departamento mais populoso, grande densidade policial) \\
Índice segurança relativo & médio \\
\bottomrule
\end{tabularx}


Central é o departamento mais populoso do Paraguai, abrigando a Grande Asunción (área metropolitana) com cidades como Luque, San Lorenzo, Fernando de la Mora, Lambaré, Capiatá e Mariano Roque Alonso. O Ministerio Público registra que o departamento responde por \textbf{aproximadamente 13\% das causas de homicídio doloso} nacional (2019--2024) --- segunda posição, refletindo a concentração populacional.

A taxa de homicídios por 100 mil habitantes, no entanto, é relativamente baixa --- estimada em torno de \textbf{3 por 100 mil} considerando que Central abriga ~25\% da população do país. O criminólogo consultado pela La Nación (2023) alertou para uma "epidemia de delitos" em Central e Asunción, com \textbf{aumento significativo de crimes patrimoniais} (roubos, furtos de veículos, assaltos) nas últimas décadas.

O departamento possui a maior densidade policial do Paraguai, com a Dirección de Policía de la Capital e as comandâncias departamentais atuando em conjunto. A criminalidade é predominantemente urbana e patrimonial; violência do crime organizado é relativamente contida em comparação com departamentos fronteiriços.


- Central é o destino preferido da maioria dos imigrantes brasileiros no Paraguai, dado o dinamismo econômico da Grande Asunción
- Cidades como Luque, San Lorenzo e Fernando de la Mora oferecem boa infraestrutura e custo de vida menor que Asunción
- A criminalidade patrimonial (furto de carros, assaltos) é o principal risco para residentes
- Bairros residenciais de classe média e alta possuem segurança privada razoável
- Sistemas de monitoramento residencial e condomínios fechados são comuns e recomendados
- O ambiente urbano exige os mesmos cuidados básicos de qualquer grande área metropolitana latino-americana
- Polícia responde com maior agilidade que no interior, mas tempo de resposta ainda é variável


- https://www.ministeriopubli

\subsection{IDH e Desenvolvimento Humano}

\textbf{Fonte principal:} PNUD Paraguai / DGEEC / INE
\textbf{Ano de referência:} 2020 (IDH); 2023 (pobreza e infraestrutura)
\textbf{Atualizado em:} 2026-03-20


\begin{tabularx}{\linewidth}{lX}
\toprule
Indicador & Valor \\
\midrule
IDH & 0,786 \\
Classificação IDH & alto \\
Ranking nacional (1=melhor) & 2/18 \\
Esperança de vida & 75,8 anos \\
Anos de escolaridade média & 10,9 anos \\
RNB per capita (USD PPA) & 16.400 \\
\bottomrule
\end{tabularx}

> Nota: IDH de 2020 conforme relatório PNUD "Desarrollo Humano en el Paraguay 2001-2020". Central está entre os 7 departamentos detalhados e ocupa o 2º lugar, logo após Asunción. É o departamento mais populoso (mais de 2 milhões de habitantes). Fonte: PNUD Paraguay 2022.


\begin{tabularx}{\linewidth}{lX}
\toprule
Indicador & Valor \\
\midrule
Taxa de pobreza total (\%) & 14,7\% \\
Taxa de pobreza extrema (\%) & 2,8\% \\
Índice de Gini & 0,430 \\
\bottomrule
\end{tabularx}

> Fonte: INE --- EPHC 2023. Central tem a segunda menor taxa de pobreza do país. Nota: Em Central e Asunción houve aumento de pessoas em pobreza extrema em 2023, segundo relatório Última Hora.


\begin{tabularx}{\linewidth}{lX}
\toprule
Indicador & Valor \\
\midrule
Acesso a água potável (\%) & 96,2\% \\
Acesso a saneamento (\%) & 97,5\% \\
\bottomrule
\end{tabularx}

> Fonte: INE --- Censo 2022 e EPHC 2023. Central é o departamento mais urbanizado do Paraguai (>95\% urbano), com infraestrutura de serviços básicos equivalente a grandes regiões metropolitanas.


Central é o departamento mais populoso e urbanizado do Paraguai, formando a área metropolitana de Asunción junto com a capital. Cidades como Luque, Fernando de la Mora, Lambaré, San Lorenzo, Capiatá e Mariano Roque Alonso oferecem a melhor relação custo-benefício para quem trabalha em Asunción: acesso à capital com custo de vida significativamente menor. Com IDH de 0,786 (o mais próximo do limiar muito alto entre os departamentos), Central oferece hospitais de referência, universidades, centros comerciais, parques industriais e uma economia diversificada. A pobreza de apenas 14,7\% e o acesso quase universal a água (96,2\%) e saneamento (97,5\%) garantem qualidade de vida metropolitana. Para imigrantes brasileiros, é a primeira escolha depois de Asunción: comunidade estabelecida, serviços equivalentes ao Brasil urbano, menor custo de imóveis. Oportunidades na indústria, serviços, comércio e setor público.


- PNUD Paraguai --- Informe "Índices de Desarrollo Humano en el Paraguay 2001-2020": h

\subsection{Sistema Prisional}

\textbf{Fonte:} MJT (Ministerio de Justicia y Trabajo) / MNP Paraguay
\textbf{Ano de referência:} 2023--2024
\textbf{Atualizado em:} 2026-03-20


\begin{tabularx}{\linewidth}{lXXXX}
\toprule
Nome & Município & Capacidade & População & Superlotação \\
\midrule
Centro Penitenciario de Emboscada & Emboscada (Cordillera) / atende Central & 1.200 & ~4.200 & ~350\% \\
Penitenciaría de Luque & Luque & 200 & ~450 & ~225\% \\
Centro de Rehabilitación de Mujeres de Itauguá & Itauguá & 120 & ~200 & ~167\% \\
Centro Educativo Itauguá (menores) & Itauguá & 120 & ~220 & ~183\% \\
\bottomrule
\end{tabularx}


\begin{tabularx}{\linewidth}{lX}
\toprule
Indicador & Valor \\
\midrule
Total de estabelecimentos & 4 \\
Capacidade total (vagas) & ~1.640 \\
População carcerária total & ~5.070 \\
Taxa de superlotação & ~309\% \\
\% presos provisórios & ~70\% \\
\bottomrule
\end{tabularx}


\textbf{Central é o departamento mais populoso do Paraguai} (mais de 2 milhões de habitantes na região metropolitana de Asunción) e seu sistema penal enfrenta pressão extrema. O Centro Penitenciario de Emboscada, embora geograficamente no departamento de Cordillera, funciona como presídio de alta segurança que atende principalmente detentos do sistema de Asunción e Central --- inclusive condenados de alta periculosidade.

Emboscada é considerado o presídio de \textbf{maior superlotação proporcional do Paraguai}, com registros de mais de 350\% acima da capacidade. Rebeliões violentas ocorreram em 2019, 2022 e 2023, com mortes e crises de reféns.

A Penitenciaría de Luque atende a região central do departamento (Luque, San Lorenzo, Fernando de la Mora, Lambaré). O Centro de Itauguá abriga mulheres do sistema metropolitano.

O MNP classificou Emboscada como situação de "emergência penitenciária" em múltiplos relatórios, destacando a presença de facções criminosas, armas de fogo, drogas e corrupção sistêmica.

\textbf{Para imigrantes:} Central é o departamento mais desenvolvido do Paraguai fora de Asunción, com excelente infraestrutura. As cidades de San Lorenzo, Luque, Fernando de la Mora e Lambaré são opções seguras e bem servidas. A crise penitenciária de Emboscada/Central não impacta diretamente as áreas residenciais e comerciais principais, mas reforça a necessidade de cautela em zonas periféricas de alta densidade.


- https://www.mjt.gov.py/index.php/penitenciaria
- https://www.mnp.gov.py/informes
- MNP --- Informe Especial Emboscada 2022--2023
- ABC Color: "Emboscada: la cárcel del terror" (2022

\subsection{Recursos Hídricos Subterrâneos}

\textbf{Fonte:} SENASA / MADES / Geología del Paraguay / The Groundwater Project / ESSAP
\textbf{Atualizado em:} 2026-03-20


\begin{tabularx}{\linewidth}{lXXX}
\toprule
Aquífero & Cobertura no Dept. & Profundidade Média & Vazão Típica \\
\midrule
Patiño (freático --- principal) & ~100\% & 30--200 m & 5--25 m$^3$/h \\
Misiones / Guarani (confinado profundo) & ~60\% (porção sul/leste) & 200--600 m & 20--60 m$^3$/h \\
\bottomrule
\end{tabularx}


\begin{tabularx}{\linewidth}{lX}
\toprule
Parâmetro & Valor \\
\midrule
Profundidade média --- Patiño (m) & 30--200 \\
Profundidade média --- Guarani (m) & 200--600 \\
Profundidade máxima recomendada (m) & 600 \\
Vazão típica (m$^3$/h) & 5--30 \\
Custo estimado perfuração (USD) & 3.500 -- 15.000 \\
Qualidade da água --- Patiño & moderada / monitoramento necessário \\
Qualidade da água --- Guarani & boa \\
pH médio --- Patiño & 6,5--7,5 \\
Outorga necessária & sim \\
Órgão regulador & MADES --- DGPCRH \\
\bottomrule
\end{tabularx}


O Departamento Central é o mais populoso do Paraguai (~3 milhões de habitantes na Grande Asunción). O Aquífero Patiño, com 1.173 km$^2$ de extensão, é a principal fonte subterrânea da região metropolitana, com espessura média de 150 m e profundidade máxima estimada de 300 m. Abastece quase 3 milhões de paraguaios. A qualidade é moderada: baixa mineralização natural, mas com riscos crescentes de contaminação por nitratos (fertilizantes), coliformes (esgoto) e hidrocarbonetos (postos de combustível) devido à urbanização intensa. O uso de água subterrânea no Departamento Central é amplamente documentado (ArcGIS Story Map específico). Uso principal: doméstico urbano, industrial e comercial. A ESSAP distribui água tratada na área metropolitana.


No Departamento Central (área metropolitana de Asunción), a rede pública da ESSAP tem cobertura razoável. Para imóveis urbanos, a rede é preferível. Para propriedades nas periferias e zonas rurais (Areguá, Itauguá, Luque, etc.), poços do Aquífero Patiño (50--120 m) são comuns e custam entre USD 3.500 e 8.000. Recomenda-se SEMPRE análise de qualidade antes do consumo, dada a vulnerabilidade do aquífero à contaminação urbana. Para uso industrial ou grandes volumes, o Guarani em profundidades maiores (200--400 m) é opção de maior qualidade. Licença do MADES é obrigatória.


- https://www.geologiadelparaguay.com.py/Acuifero-Pati\%C3\%B1o.htm
- https://storymaps.arcgis.com/stories/069b322e3d724d77b773b

\subsection{Aptidão do Solo}

\textbf{Fonte:} MAG / SENAVE / Geología del Paraguay
\textbf{Atualizado em:} 2026-03-20


\begin{tabularx}{\linewidth}{lXX}
\toprule
Tipo de Solo & \% do Território & Características \\
\midrule
Ultissolos (Argissolos) & 40\% & Solos podzólicos vermelho-amarelados, fertilidade baixa-média; historicamente cultivados \\
Entissolos arenosos & 25\% & Solos arenosos de baixa fertilidade; ocorrem em grande parte do departamento \\
Solos urbanos / antropizados & 30\% & Solos impermeabilizados, compactados ou aterrados por urbanização intensa \\
Inceptissolos / Aluviais & 5\% & Planícies do Rio Salado e afluentes; fertilidade moderada \\
\bottomrule
\end{tabularx}


\begin{tabularx}{\linewidth}{lXX}
\toprule
Classe & \% Território & Uso Recomendado \\
\midrule
Lavoura intensiva & 5\% & Horticultura intensiva periurbana \\
Lavoura extensiva & 10\% & Mandioca, milho, pequenas propriedades rurais residuais \\
Pastagem natural & 10\% & Pecuária residual em municípios menos urbanizados \\
Floresta / reserva & 5\% & Parques urbanos, matas ciliares residuais \\
Sem aptidão agrícola & 70\% & Área predominantemente urbanizada (Área Metropolitana de Asunción) \\
\bottomrule
\end{tabularx}


\begin{tabularx}{\linewidth}{lX}
\toprule
Parâmetro & Valor \\
\midrule
pH médio & 5,0 -- 6,5 \\
Fertilidade natural & baixa (solos agrícolas residuais muito degradados) \\
Teor de matéria orgânica & baixo \\
Necessidade de calcário & sim (para áreas agrícolas residuais) \\
\bottomrule
\end{tabularx}


Central é predominantemente urbano-industrial. Nas raras áreas agrícolas:
- \textbf{Horticultura intensiva} (tomate, alface, pimentão, ervas)
- \textbf{Mandioca} em municípios periurbanos como Itauguá, Ypacaraí
- \textbf{Flores e ornamentais} em pequenas propriedades
- \textbf{Pecuária de leite} residual em municípios do sul


- \textbf{Urbanização extrema:} Central é o departamento mais densamente urbanizado do Paraguai; >70\% do território não tem aptidão agrícola
- \textbf{Solos exauridos:} Central figura entre os departamentos com >90\% dos solos cultivados deteriorados
- \textbf{Contaminação:} solos urbanos com histórico de atividades industriais e descarte inadequado de resíduos
- \textbf{Impermeabilização:} enchentes frequentes por impermeabilização do solo nas cidades
- \textbf{Pressão imobiliária:} terra agrícola residual em crescente valorização como área de expansão urbana


Central não é destino agrícola. É o departamento mais urban

\subsection{Terra Rural}

\textbf{Fonte:} INDERT / Clasificados.com.py / mercado local
\textbf{Ano de referência:} 2024
\textbf{Atualizado em:} 2026-03-20


\begin{tabularx}{\linewidth}{lXXX}
\toprule
Tipo & Preço (USD/ha) & Faixa & Tendência \\
\midrule
Agrícola alta produtividade & 36.000 & 25.000 -- 55.000 & $\uparrow$ \\
Pastagem / uso misto & 20.000 & 12.000 -- 30.000 & $\uparrow$ \\
Mata / reserva & 10.000 & 6.000 -- 18.000 & $\uparrow$ \\
\bottomrule
\end{tabularx}

> Referência confirmada: USD 36.399/ha para terra agrícola no Departamento Central --- quatro vezes superior ao preço em Iowa (EUA). Dado verificado em múltiplas fontes.


\begin{tabularx}{\linewidth}{lX}
\toprule
Parâmetro & Valor \\
\midrule
Valorização anual média (\%) & 10--15\% \\
Lote mínimo rural (ha) & 5 (chácara periurbana) \\
Imposto territorial rural (\%) & 1\% sobre valor fiscal \\
Liquidez do mercado & muito alta \\
\bottomrule
\end{tabularx}


- \textbf{Luque:} Cidade histórica próxima a Asunción; qualquer terra tem demanda altíssima por conversão para uso residencial e industrial.
- \textbf{San Lorenzo / Fernando de la Mora:} Zona metropolitana densa; terra praticamente com preço urbano.
- \textbf{Itauguá / Capiata:} Expansão metropolitana leste; loteamentos industriais e residenciais em rápido crescimento.
- \textbf{Areguá / Capiatá:} Capital departamental e entorno do Lago Ypacaraí; demanda de chácaras de lazer impulsiona preços.
- \textbf{Ñemby / Villa Elisa:} Corredor sul metropolitano; terra cara mas bem valorizada por logística industrial.


\begin{tabularx}{\linewidth}{lXXX}
\toprule
Indicador & Central (PY) & Grande São Paulo (BR) & Grande Curitiba (BR) \\
\midrule
Terra periurbana (USD/ha) & 25.000--55.000 & 80.000--200.000 & 40.000--120.000 \\
Tendência 2019--2024 & +80--100\% & +50--80\% & +40--60\% \\
\bottomrule
\end{tabularx}

O Departamento Central tem os preços de terra mais altos do Paraguai devido à extrema pressão urbanística. Para fins agrícolas puros, não é recomendado; o valor é quase inteiramente imobiliário urbano.


- [Última Hora --- Tierras de Alto Paraná más caras que EEUU](https://www.ultimahora.com/tierras-agricolas-alto-parana-son-mas-caras-que-las-eeuu-n2928594)
- [Plan B Paraguay --- Farmland Prices](https://x.com/planbparaguay/status/1834232169759211739)
- [INDERT](https://www.indert.gov.py/)
- [Valor Agrícola](https://valoragricola.com.py/)


\subsection{Mercado Imobiliário}

\textbf{Capital:} Areguá
\textbf{Fonte:} Clasificados.com.py / UbicaProp / mercado local
\textbf{Ano de referência:} 2024
\textbf{Atualizado em:} 2026-03-20


\begin{tabularx}{\linewidth}{lXXX}
\toprule
Cidade & Tipo & Preço (USD/m$^2$) & Aluguel (USD/mês) \\
\midrule
Luque & Apartamento 2q & 1.100 -- 1.600 & 400 -- 750 \\
San Lorenzo & Apartamento 2q & 900 -- 1.400 & 350 -- 650 \\
Fernando de la Mora & Apartamento 2q & 950 -- 1.400 & 350 -- 680 \\
Areguá (capital) & Casa 3q & 700 -- 1.100 & 280 -- 500 \\
Itauguá & Casa 3q & 650 -- 1.000 & 250 -- 450 \\
\bottomrule
\end{tabularx}

> Nota: O Departamento Central é a região metropolitana de Asunción. As cidades do departamento têm mercado imobiliário integrado com a capital, e os preços refletem essa dinâmica urbana.


\begin{tabularx}{\linewidth}{lX}
\toprule
Parâmetro & Valor \\
\midrule
Valorização anual (\%) & 8--12\% \\
Disponibilidade de financiamento & muito boa (todos os bancos do país operam na região) \\
Cidades mais valorizadas & Luque, San Lorenzo, Fernando de la Mora, Villa Elisa \\
\bottomrule
\end{tabularx}


\textbf{Vale a pena comprar ou alugar?}
O Departamento Central é a grande área metropolitana de Asunción --- quem compra aqui tem liquidez garantida. É a região mais cara fora da capital, mas ainda muito abaixo dos preços de regiões metropolitanas brasileiras equivalentes.

\textbf{Áreas recomendadas:}
- \textbf{Luque:} Ótima infraestrutura, sede da CONMEBOL, muito valorizado; bom para famílias.
- \textbf{San Lorenzo:} Maior cidade do departamento; universidades, hospitais, comércio completo.
- \textbf{Areguá:} Capital departamental histórica às margens do Lago Ypacaraí; casas coloniais e chácaras encantadoras.
- \textbf{Itauguá / Nueva Italia:} Crescimento industrial acelerado; bom para investimento especulativo.
- \textbf{Villa Elisa / Ñemby:} Crescimento residencial forte; populares entre brasileiros e argentinos.

\textbf{Armadilhas comuns:}
- Trânsito caótico nos horários de pico em toda a região metropolitana --- verificar distância real ao trabalho.
- Loteamentos irregulares são comuns em bordas; sempre verificar títulos no Registro de Inmuebles.
- Preços de apartamentos em condomínios fechados têm custo de manutenção/condomínio relevante.


- [InfoCasas.com.py --- Central](https://www.infocasas.com.py/venta/inmuebles/central)
- [UbicaProp --- Precio m$^2$ Asunción](https://ubicaprop.com.py/precio-m2-asuncion-barrios/)
- [Clasipar.com --- Alquiler Central](https://clasipar.paraguay.c

\subsection{Cobertura Celular}

\textbf{Fonte:} CONATEL / Tigo / Personal / OpenSignal / nPerf
\textbf{Ano de referência:} 2024
\textbf{Atualizado em:} 2026-03-20


\begin{tabularx}{\linewidth}{lXXXXX}
\toprule
Operadora & 2G & 3G & 4G & 5G & Cobertura Pop. \\
\midrule
Tigo & 99\% & 99\% & 98\% & 18\% & 99\% \\
Personal/Claro & 99\% & 98\% & 97\% & 12\% & 99\% \\
VOX & 92\% & 85\% & 75\% & --- & 90\% \\
\bottomrule
\end{tabularx}

> \textbf{Nota:} Central é o departamento mais populoso do Paraguai (~2,2 milhões de hab.) e forma a Grande Asunción com o Distrito Capital. Cobertura 4G praticamente total. 5G em implantação em municípios como Luque, San Lorenzo, Fernando de la Mora e Lambaré. Não há áreas sem cobertura relevante.


\begin{tabularx}{\linewidth}{lX}
\toprule
Parâmetro & Valor \\
\midrule
Melhor operadora & Tigo (maior rede); Personal/Claro (boa relação custo-benefício) \\
Velocidade 4G média (Mbps) & 40--65 \\
Áreas sem cobertura & Nenhuma área urbana significativa \\
Sinal em propriedade rural & N/A --- departamento essencialmente urbano/suburbano \\
\bottomrule
\end{tabularx}


\begin{tabularx}{\linewidth}{lXXX}
\toprule
Plano & Operadora & USD/mês & Dados \\
\midrule
Pré-pago básico (7 dias) & Tigo & ~2,50 & 3 GB \\
Pré-pago intermediário (30 dias) & Personal & ~5,00 & 8 GB \\
Pós-pago entrada & Tigo & ~12,00 & 15 GB + chamadas \\
Pós-pago intermediário & Personal/Claro & ~18,00 & 30 GB + chamadas \\
Pós-pago premium & Tigo & ~28,00 & ilimitado \\
\bottomrule
\end{tabularx}

> Câmbio de referência: 1 USD $\approx$ 7.500 PYG (2024). Chips disponíveis em todo o departamento.


- \textbf{Área metropolitana de Asunción (San Lorenzo, Luque, Capiatá, Lambaré, Fernando):} Cobertura 4G completa e estável de todas as operadoras.
- \textbf{Municípios do Gran Asunción:} Todos têm cobertura 4G sólida; 5G disponível nos bairros centrais dos principais municípios.
- \textbf{Chip recomendado:} Tigo para maior cobertura nacional; Personal/Claro para planos internacionais e uso em cidades específicas.
- \textbf{Roaming Brasil-Paraguai:} Chips brasileiros funcionam em todo o departamento com tarifas de roaming. Recomendado adquirir chip local para uso prolongado.
- \textbf{eSIM:} Disponível para Personal/Claro em planos pós-pago selecionados.
- Central é o local com melhor infraestrutura de telecom do Paraguai fora do Distrito Capital.


- https://www.conatel.gov.py/conatel/indicadores/
- https://www.tigo.com.py/cobertura
- https://www.personal.com.py/institucional/cobertura.html
- https://www.nperf.com/en/map/PY/-/-/sign

\subsection{Internet}

\textbf{Fonte:} CONATEL / Tigo / Claro / INE / Speedtest
\textbf{Ano de referência:} 2024
\textbf{Atualizado em:} 2026-03-20


\begin{tabularx}{\linewidth}{lX}
\toprule
Indicador & Valor \\
\midrule
\% população com internet & 87,6\% \\
\% com banda larga fixa & ~62\% dos domicílios \\
Velocidade média download (Mbps) & 80--150 \\
Velocidade média upload (Mbps) & 25--60 \\
\bottomrule
\end{tabularx}

> Fonte: INE EPHC 2024 --- Central com 87,6\%, o maior entre os departamentos (excluindo Asunción/Distrito Capital com 88,9\%). Mediana Ookla nacional: 92,78 Mbps de download (maio/2024).


\begin{tabularx}{\linewidth}{lXX}
\toprule
Tecnologia & Disponibilidade & Velocidade Típica \\
\midrule
Fibra ótica (FTTH) & Todo o departamento & 100--1000 Mbps \\
Cabo HFC & Municípios principais & 50--500 Mbps \\
Rádio/Wireless & Amplo & 10--100 Mbps \\
ADSL & Legado, em extinção & 5--20 Mbps \\
Satélite (Starlink) & Disponível & 50--200 Mbps \\
\bottomrule
\end{tabularx}


\begin{tabularx}{\linewidth}{lXXX}
\toprule
Provedor & Tecnologia & Área de cobertura & Preço ref. (USD/mês) \\
\midrule
Tigo & Fibra/HFC/Rádio & Todo o departamento & 18--65 \\
Claro/Personal & Fibra/HFC/Rádio & Todo o departamento & 15--60 \\
Copaco & Fibra/ADSL & Todo o departamento & 10--40 \\
VOX & Fibra/Rádio & Municípios principais & 15--45 \\
NEO & Fibra & Municípios selecionados & 20--65 \\
Starlink & Satélite LEO & Todo & ~44 + equip. ~370 \\
\bottomrule
\end{tabularx}


Central é essencialmente um departamento urbano/suburbano. Não há propriamente áreas rurais de extensão significativa. Os municípios periféricos (Guarambaré, Nueva Italia, Itauguá, Ypacaraí) têm cobertura de fibra ou rádio wireless de qualidade razoável.

Central é o \textbf{melhor departamento para imigrantes} que necessitam de conectividade de alta qualidade e preferem vida suburbana próxima a Asunción. Velocidades de internet comparáveis às melhores capitais brasileiras, com preços significativamente menores.


- https://www.conatel.gov.py/conatel/wp-content/uploads/2025/02/informe-de-internet-fijo-2024.pdf
- https://marketdata.com.py/laboratorio/analisis/acceso-a-internet-em-paraguay-alcanza-al-816-de-la-poblacion-quienes-lo-usan-mas-donde-y-para-que-143217
- https://www.ine.gov.py/noticias/2436/8-de-cada-10-personas-utiliza-internet-en-paraguay
- https://www.speedgeo.net/statistics/paraguay
- https://www.tigo.com.py/internet



\clearpage
\section{Ñeembucú}
\label{sec:dept_12_Neembucu}

\subsection{Segurança Pública}

\textbf{Fonte principal:} Ministerio Público Paraguay / INE / Ministerio del Interior PY
\textbf{Data de referência:} 2019--2024
\textbf{Atualizado em:} 2026-03-20


\begin{tabularx}{\linewidth}{lX}
\toprule
Indicador & Valor \\
\midrule
Taxa de homicídios (por 100k hab) & 0,2--0,4 (2019--2020, fonte INE --- mais baixa do país) \\
Crimes registrados (total/ano) & muito baixo (população pequena e ambiente rural) \\
Número de comisarías & ~6 (estimativa departamental) \\
Índice segurança relativo & alto \\
\bottomrule
\end{tabularx}


Ñeembucú possui a \textbf{taxa de homicídios mais baixa de todo o Paraguai}, com registros de apenas \textbf{0,2 por 100 mil em 2019 e 0,4 por 100 mil em 2020} --- valores excepcionalmente baixos mesmo em comparação internacional. O Ministerio Público registra que o departamento responde por apenas \textbf{1\% das causas de homicídio doloso} nacional, refletindo tanto a paz social quanto a baixíssima densidade populacional.

O departamento é banhado pelos rios Paraguai e Pilcomayo (fronteira com Argentina a sudoeste) e faz fronteira com a Corrientes argentina. A fronteira fluvial tem algum histórico de contrabando, mas não é identificada como rota principal de narcotráfico de alta intensidade. Pilar, capital departamental, é uma cidade fluvial tranquila com forte identidade cultural local.

A economia baseia-se em pecuária extensiva e pesca. O departamento é um dos menos urbanizados do Paraguai, com população espalhada em vastas planícies inundáveis (Pantanal paraguaio).


- Ñeembucú é idealmente seguro, mas extremamente isolado para padrões urbanos
- Pilar possui uma comunidade pequena e serviços básicos, mas serviços especializados exigem viagem a Asunción (~300 km)
- A tranquilidade do ambiente é perfeita para quem busca retiro rural ou atividades de pesca e ecoturismo
- Brasileiros que gostam de natureza e quietude encontrarão ambiente excepcional
- Oportunidades econômicas são limitadas --- adequado principalmente para aposentados ou quem tem renda passiva
- Pouca comunidade brasileira estabelecida --- adaptação requer maior autonomia


- https://www.ministeriopublico.gov.py/nota/estadisticas-del-ministerio-publico-la-mayoria-de-los-homicidios-dolosos-en-10564
- https://ods.ine.gov.py/ine-main/ods/paz-justicias-e-instituciones-solidas-16/meta-16.1/indicador-284
- https://informacionpublica.paraguay.gov.py/public/896733-B1-AnlisisEstadsticodeHomicidioDolosoenParaguayActualizacina

\subsection{IDH e Desenvolvimento Humano}

\textbf{Fonte principal:} PNUD Paraguai / DGEEC / INE
\textbf{Ano de referência:} 2020 (IDH); 2023 (pobreza e infraestrutura)
\textbf{Atualizado em:} 2026-03-20


\begin{tabularx}{\linewidth}{lX}
\toprule
Indicador & Valor \\
\midrule
IDH & 0,710 \\
Classificação IDH & alto \\
Ranking nacional (1=melhor) & 8/18 \\
Esperança de vida & 73,0 anos \\
Anos de escolaridade média & 8,2 anos \\
RNB per capita (USD PPA) & 8.800 \\
\bottomrule
\end{tabularx}

> Nota: IDH estimado com base no relatório PNUD 2020 e dados subnacionais da Global Data Lab. Ñeembucú é departamento ribeirinho (rios Paraguai e Paraná), com perfil pecuário e pesqueiro. Estimativa baseada em perfil regional sul-ocidental.


\begin{tabularx}{\linewidth}{lX}
\toprule
Indicador & Valor \\
\midrule
Taxa de pobreza total (\%) & 28,8\% \\
Taxa de pobreza extrema (\%) & 6,5\% \\
Índice de Gini & 0,445 \\
\bottomrule
\end{tabularx}

> Fonte: INE --- EPHC 2023. Nota: Houve aumento da pobreza em Ñeembucú entre 2022 e 2023, segundo relatório Última Hora (junto com Cordillera e Presidente Hayes).


\begin{tabularx}{\linewidth}{lX}
\toprule
Indicador & Valor \\
\midrule
Acesso a água potável (\%) & 79,1\% \\
Acesso a saneamento (\%) & 77,4\% \\
\bottomrule
\end{tabularx}

> Estimativa baseada em dados do Censo 2022 e perfil rural/ribeirinho. Pilar, capital departamental, possui infraestrutura básica adequada para cidade pequena.


Ñeembucú é um departamento tranquilo no extremo sudoeste do Paraguai, com economia baseada em pecuária extensiva, pesca e agricultura. Pilar, a capital, é uma cidade pequena mas com serviços básicos (hospital regional, escolas). O departamento tem belas paisagens de banhados (Pantanal paraguaio) e o Rio Paraguai, tornando-o atrativo para turismo de pesca. A pobreza de 28,8\% é moderada e o acesso a serviços básicos é razoável. Para imigrantes brasileiros, o departamento oferece muito baixo custo de vida, paz e natureza, mas limitadas oportunidades de emprego formal. É mais adequado para aposentados, criadores de gado ou quem busca retiro com natureza. A proximidade com Corrientes (Argentina) pode ser vantagem para acesso a serviços. O aumento de pobreza em 2023 é um sinal de alerta sobre fragilidade econômica local.


- PNUD Paraguai --- Informe IDH 2001-2020: https://www.undp.org/es/paraguay/noticias/IDH-2001-2020
- INE --- EPHC 2023: https://www.ine.gov.py/noticias/1928/pobreza-monetaria-en-el-2023-se-ubico-en-el-227
- Última Hora --- Aumento pobreza Neembucu 2023: https://www.ultimah

\subsection{Sistema Prisional}

\textbf{Fonte:} MJT (Ministerio de Justicia y Trabajo) / MNP Paraguay
\textbf{Ano de referência:} 2023--2024
\textbf{Atualizado em:} 2026-03-20


\begin{tabularx}{\linewidth}{lXXXX}
\toprule
Nome & Município & Capacidade & População & Superlotação \\
\midrule
Penitenciaría Regional de Pilar & Pilar & 100 & ~175 & ~175\% \\
\bottomrule
\end{tabularx}


\begin{tabularx}{\linewidth}{lX}
\toprule
Indicador & Valor \\
\midrule
Total de estabelecimentos & 1 \\
Capacidade total (vagas) & ~100 \\
População carcerária total & ~175 \\
Taxa de superlotação & ~175\% \\
\% presos provisórios & ~67\% \\
\bottomrule
\end{tabularx}


Ñeembucú é um dos departamentos menos populosos do Paraguai, com baixa densidade demográfica e economia baseada em pecuária e pesca. A Penitenciaría Regional de Pilar é o único estabelecimento penal do departamento, localizado na capital Pilar.

Pilar situa-se às margens do rio Pilcomayo (fronteira com a Argentina --- província de Formosa), e o contrabando transfronteiriço é uma atividade relevante na região, gerando parte da criminalidade local. No entanto, o volume é inferior ao das principais fronteiras paraguaias.

O estabelecimento prisional é de pequenas dimensões e, segundo o MNP, apresenta condições básicas relativamente melhores do que presídios maiores, embora ainda superlotado. A assistência jurídica é deficiente e a maioria dos detentos aguarda julgamento em prisão preventiva por longos períodos.

\textbf{Para imigrantes:} Ñeembucú é um departamento tranquilo com baixa incidência de violência urbana. Pilar tem uma pequena estrutura comercial e universitária. O sistema prisional não representa risco relevante para atividades cotidianas. Indicado para quem busca tranquilidade e proximidade com a Argentina.


- https://www.mjt.gov.py/index.php/penitenciaria
- https://www.mnp.gov.py/informes
- MNP --- Informe de Visitas 2021--2022


\subsection{Recursos Hídricos Subterrâneos}

\textbf{Fonte:} SENASA / MADES / Geología del Paraguay / IGRAC
\textbf{Atualizado em:} 2026-03-20


\begin{tabularx}{\linewidth}{lXXX}
\toprule
Aquífero & Cobertura no Dept. & Profundidade Média & Vazão Típica \\
\midrule
Guarani / Misiones (confinado) & ~65\% & 100--400 m & 20--80 m$^3$/h \\
Aluvial (planície do rio Paraguay / Paraná) & ~35\% & 8--25 m & 5--15 m$^3$/h \\
\bottomrule
\end{tabularx}


\begin{tabularx}{\linewidth}{lX}
\toprule
Parâmetro & Valor \\
\midrule
Profundidade média --- Guarani (m) & 100--350 \\
Profundidade máxima recomendada (m) & 450 \\
Vazão típica (m$^3$/h) & 15--60 \\
Custo estimado perfuração (USD) & 3.000 -- 10.000 \\
Qualidade da água --- Guarani & boa a excelente \\
Qualidade da água --- Aluvial & boa (requer análise microbiológica) \\
pH médio & 6,5--7,5 \\
Outorga necessária & sim \\
Órgão regulador & MADES --- DGPCRH \\
\bottomrule
\end{tabularx}


Ñeembucú é o departamento mais meridional da Região Oriental, marcado por extensa planície alagável (Pantanal Paraguaio) às margens dos rios Paraguay e Paraná. O Aquífero Guarani/Misiones está presente na porção leste, com boa qualidade. A planície aluvial oferece captação rasa, adequada para uso doméstico e animal com verificação microbiológica. A capital Pilar dispõe de rede pública do SENASA. Uso principal: pecuária extensiva, abastecimento doméstico rural, pesca (área de banhado não usa poços artesianos). O departamento tem uma das menores densidades populacionais do Paraguai Oriental.


Ñeembucú atrai pecuaristas pela abundância de pastagens naturais. Para fazendas na porção leste (Guarani acessível), poço de 100--250 m custa entre USD 3.000 e 8.000 com excelente qualidade. Na planície aluvial, captações rasas (10--20 m) são mais baratas (USD 1.500--3.000) mas exigem análise microbiológica periódica. Pilar e cidades maiores têm rede pública. Para fazendas remotas, poço artesiano é indispensável. Licença do MADES obrigatória.


- https://gw-project.org/es/caracteristicas-del-agua-subterranea-en-paraguay/
- https://es.wikipedia.org/wiki/Acu\%C3\%ADfero\_Guaran\%C3\%AD
- https://www.geologiadelparaguay.com.py/Acuifero-Misiones.htm
- https://www.mades.gov.py/tramites/registro-de-empresas-publicas-y-privadas-que-se-dedican-a-la-perforacion-de-pozos-tubulares/


\subsection{Aptidão do Solo}

\textbf{Fonte:} MAG / SENAVE / Geología del Paraguay
\textbf{Atualizado em:} 2026-03-20


\begin{tabularx}{\linewidth}{lXX}
\toprule
Tipo de Solo & \% do Território & Características \\
\midrule
Gleissolos / Hidromórficos & 40\% & Solos permanentemente alagados ou sazonalmente encharcados (esterales/bañados); alta matéria orgânica em superfície, mas anóxicos \\
Vertissolos (Planossolos) & 20\% & Solos argilosos pesados que incham com umidade e racham na seca; baixo potencial agrícola \\
Entissolos Flúvicos (Aluviais) & 25\% & Solos aluviais do Rio Paraguai e Rio Paraná; fertilidade média-alta quando drenados \\
Ultissolos (Argissolos) & 15\% & Solos de encostas levemente elevadas; fertilidade baixa-média; melhor aptidão agrícola local \\
\bottomrule
\end{tabularx}


\begin{tabularx}{\linewidth}{lXX}
\toprule
Classe & \% Território & Uso Recomendado \\
\midrule
Lavoura intensiva & 5\% & Arroz irrigado, horticultura em planícies drenadas \\
Lavoura extensiva & 15\% & Mandioca, sorgo, pastagem melhorada nas áreas mais elevadas \\
Pastagem natural/melhorada & 40\% & Pecuária extensiva --- vocação principal; planícies inundáveis são pastagens naturais \\
Floresta / reserva & 15\% & Matas de galeria, banhados (Esteros del Neembucú) \\
Sem aptidão agrícola & 25\% & Esterales, bañados permanentes, rios e corpos d'água \\
\bottomrule
\end{tabularx}


\begin{tabularx}{\linewidth}{lX}
\toprule
Parâmetro & Valor \\
\midrule
pH médio & 5,5 -- 7,0 \\
Fertilidade natural & média (nas áreas drenáveis) \\
Teor de matéria orgânica & alto (solos hidromórficos com MO acumulada) \\
Necessidade de calcário & parcial \\
\bottomrule
\end{tabularx}


- \textbf{Pecuária extensiva} --- atividade predominante; pastagens naturais de grande extensão
- \textbf{Arroz} --- potencial em planícies quando com drenagem adequada
- \textbf{Mandioca} e \textbf{milho} na agricultura familiar
- \textbf{Sorgo} e \textbf{pastagens melhoradas} nas áreas mais drenadas
- \textbf{Pesca artesanal} --- atividade econômica importante nos rios e banhados


- \textbf{Encharcamento e inundações:} Ñeembucú é o departamento mais plano e sujeito a inundações do Paraguai; os "esterales" ocupam grande parte do território
- \textbf{Apenas ~20\% de área cultivável:} levantamentos indicam que somente ~20\% do território é efetivamente cultivável
- \textbf{Drenagem deficiente:} solos hidromórficos exigem obras de drenagem para qualquer uso agrícola intensivo
- \textbf{Vertissolos:} solos com expansão/retração argilosa causam racha

\subsection{Terra Rural}

\textbf{Fonte:} INDERT / Clasificados.com.py / mercado local
\textbf{Ano de referência:} 2024
\textbf{Atualizado em:} 2026-03-20


\begin{tabularx}{\linewidth}{lXXX}
\toprule
Tipo & Preço (USD/ha) & Faixa & Tendência \\
\midrule
Agrícola alta produtividade & 3.000 & 1.800 -- 4.500 & $\uparrow$ \\
Pastagem & 1.500 & 900 -- 2.200 & $\rightarrow$ \\
Mata / reserva & 600 & 300 -- 1.000 & $\rightarrow$ \\
\bottomrule
\end{tabularx}


\begin{tabularx}{\linewidth}{lX}
\toprule
Parâmetro & Valor \\
\midrule
Valorização anual média (\%) & 4--6\% \\
Lote mínimo rural (ha) & 20 \\
Imposto territorial rural (\%) & 1\% sobre valor fiscal \\
Liquidez do mercado & baixa \\
\bottomrule
\end{tabularx}


- \textbf{Pilar (entorno):} Capital departamental às margens do Rio Paraguay; polo logístico e centro de serviços; terra mais valorizada do departamento.
- \textbf{Alberdi:} Cidade fronteiriça com a Argentina (Corrientes); terra com potencial de valorização por integração binacional.
- \textbf{Isla Umbú / Cerrito:} Zona de ilhas fluviais e banhados; pecuária extensiva em áreas inundáveis; baixo custo.
- \textbf{General José Eduvigis Díaz:} Zona mais elevada do departamento; terra com melhor potencial agrícola.


\begin{tabularx}{\linewidth}{lXXX}
\toprule
Indicador & Ñeembucú (PY) & Corrientes (AR) & Rio Grande do Sul (BR) \\
\midrule
Terra agrícola (USD/ha) & 1.800--4.500 & 3.000--8.000 & 12.000--22.000 \\
Terra pastagem (USD/ha) & 900--2.200 & 1.500--4.000 & 5.000--9.000 \\
Custo relativo (vs RS) & ~20\% & ~30\% & referência RS \\
\bottomrule
\end{tabularx}

Ñeembucú é um departamento com forte vocação pecuária extensiva em áreas de campo natural e banhados. Os preços de terra são dos mais baixos do Paraguai Oriental, compensados por limitações de infraestrutura e zonas alagáveis.


- [INDERT --- Instituto Nacional de Desarrollo Rural y de la Tierra](https://www.indert.gov.py/)
- [Plan B Paraguay --- Farmland Prices](https://x.com/planbparaguay/status/1834232169759211739)
- [InfoCasas --- Chacras Paraguay](https://www.infocasas.com.py/venta/chacras-o-campos)
- [Hendyla.com --- Campos Paraguay](https://www.hendyla.com/inmuebles/campos/)


\subsection{Mercado Imobiliário}

\textbf{Capital:} Pilar
\textbf{Fonte:} Clasificados.com.py / mercado local
\textbf{Ano de referência:} 2024
\textbf{Atualizado em:} 2026-03-20


\begin{tabularx}{\linewidth}{lXX}
\toprule
Tipo & Preço (USD/m$^2$) & Aluguel (USD/mês) \\
\midrule
Apartamento 2 quartos & 400 -- 650 & 150 -- 300 \\
Casa 3 quartos & 350 -- 580 & 130 -- 270 \\
\bottomrule
\end{tabularx}


\begin{tabularx}{\linewidth}{lX}
\toprule
Parâmetro & Valor \\
\midrule
Valorização anual (\%) & 3--5\% \\
Disponibilidade de financiamento & limitada (cooperativas locais, BNF) \\
Bairros mais valorizados & Centro histórico, Barrio Puerto, beiras do Rio Paraguay \\
\bottomrule
\end{tabularx}


\textbf{Vale a pena comprar ou alugar?}
Pilar é uma cidade histórica de pequeno porte (aprox. 45.000 hab.) com charme colonial e localização estratégica no Rio Paraguay. O mercado imobiliário é pequeno e de baixa liquidez. Para quem investe em pecuária na região, compra de casa em Pilar é prática e acessível.

\textbf{Áreas recomendadas:}
- \textbf{Centro de Pilar:} infraestrutura básica concentrada, próximo ao porto fluvial.
- \textbf{Puerto de Pilar:} região ribeirinha com potencial turístico crescente.
- \textbf{Alberdi:} para quem tem negócios com a Argentina (Corrientes está a curta travessia de balsa).

\textbf{Armadilhas comuns:}
- Departamento com baixa conectividade terrestre; acesso principal é fluvial ou via rota em más condições em épocas de chuva.
- Cheias do Rio Paraguay afetam periodicamente áreas baixas da cidade.
- Serviços especializados muito limitados; referência médica mais próxima é Asunción (320 km).
- Mercado imobiliário pode demorar muito para vender um imóvel --- baixa liquidez.


- [InfoCasas.com.py --- Ñeembucú](https://www.infocasas.com.py/venta/inmuebles/neembucu)
- [Clasipar.com --- Alquiler Ñeembucú](https://clasipar.paraguay.com/alquiler/inmuebles/neembucu)
- [INDERT](https://www.indert.gov.py/)


\subsection{Cobertura Celular}

\textbf{Fonte:} CONATEL / Tigo / Personal / OpenSignal / nPerf
\textbf{Ano de referência:} 2024
\textbf{Atualizado em:} 2026-03-20


\begin{tabularx}{\linewidth}{lXXXXX}
\toprule
Operadora & 2G & 3G & 4G & 5G & Cobertura Pop. \\
\midrule
Tigo & 75\% & 62\% & 48\% & --- & 68\% \\
Personal/Claro & 65\% & 52\% & 38\% & --- & 58\% \\
VOX & 38\% & 22\% & 10\% & --- & 30\% \\
\bottomrule
\end{tabularx}

> \textbf{Nota:} Ñeembucú é um dos departamentos com menor cobertura celular do Paraguai Oriental, caracterizado por vasto território pantaneiro com baixa densidade populacional. A capital Pilar (~90k hab.) e municípios ao longo do rio Pilcomayo têm cobertura razoável. Extensas áreas alagadas do Pantanal sul têm sinal inexistente.


\begin{tabularx}{\linewidth}{lX}
\toprule
Parâmetro & Valor \\
\midrule
Melhor operadora rural & Tigo \\
Velocidade 4G média (Mbps) & 10--22 \\
Áreas sem cobertura & Pantanal sul, zonas alagadas, áreas de estâncias remotas \\
Sinal em propriedade rural & ruim a moderado \\
\bottomrule
\end{tabularx}


\begin{tabularx}{\linewidth}{lXXX}
\toprule
Plano & Operadora & USD/mês & Dados \\
\midrule
Pré-pago básico (7 dias) & Tigo & ~2,50 & 3 GB \\
Pré-pago intermediário (30 dias) & Personal & ~5,00 & 8 GB \\
Pós-pago entrada & Tigo & ~12,00 & 15 GB + chamadas \\
Pós-pago intermediário & Personal/Claro & ~18,00 & 30 GB + chamadas \\
\bottomrule
\end{tabularx}


- \textbf{Pilar:} Capital razoavelmente conectada; cobertura 4G de Tigo e 3G/4G de Personal adequada para uso urbano.
- \textbf{Municipios da fronteira argentina (Humaitá, General Díaz, Isla Umbú):} Cobertura muito variável; fronteira fluvial com Argentina.
- \textbf{Estâncias e fazendas no Pantanal:} Sinal muito fraco ou inexistente; Starlink é solução essencial.
- \textbf{Acesso rodoviário limitado:} Muitas áreas de Ñeembucú têm acesso difícil --- a conectividade celular segue a infraestrutura de estradas, que é escassa.
- nPerf registra dados específicos de cobertura em Pilar disponíveis para consulta.


- https://www.conatel.gov.py/conatel/indicadores/
- https://www.tigo.com.py/cobertura
- https://www.nperf.com/en/map/PY/3437526.Pilar/-./signal
- https://www.personal.com.py/institucional/cobertura.html


\subsection{Internet}

\textbf{Fonte:} CONATEL / Tigo / Claro / INE / Copaco
\textbf{Ano de referência:} 2024
\textbf{Atualizado em:} 2026-03-20


\begin{tabularx}{\linewidth}{lX}
\toprule
Indicador & Valor \\
\midrule
\% população com internet & ~50\% \\
\% com banda larga fixa & ~12\% dos domicílios \\
Velocidade média download (Mbps) & 10--25 (Pilar) / 2--8 (rural) \\
Velocidade média upload (Mbps) & 3--10 (Pilar) \\
\bottomrule
\end{tabularx}

> Ñeembucú tem uma das menores penetrações de internet do Paraguai Oriental, resultado de sua baixíssima densidade demográfica (apenas ~80k hab. em 12.147 km$^2$) e da extensão do território pantaneiro.


\begin{tabularx}{\linewidth}{lXX}
\toprule
Tecnologia & Disponibilidade & Velocidade Típica \\
\midrule
Fibra ótica (FTTH) & Apenas Pilar & 30--150 Mbps \\
Cabo HFC & Pilar (limitado) & 20--80 Mbps \\
Rádio/Wireless & Pilar e municípios maiores & 3--30 Mbps \\
ADSL (Copaco) & Municípios com centrais telefônicas & 1--6 Mbps \\
Satélite (Starlink) & Todo o departamento & 50--200 Mbps \\
\bottomrule
\end{tabularx}


\begin{tabularx}{\linewidth}{lXXX}
\toprule
Provedor & Tecnologia & Área de cobertura & Preço ref. (USD/mês) \\
\midrule
Tigo & Fibra/Rádio & Pilar + municípios maiores & 18--40 \\
Claro/Personal & Rádio & Pilar & 15--35 \\
Copaco & ADSL/Rádio & Interior (rede legada) & 10--22 \\
Starlink & Satélite LEO & Todo o departamento & ~44 + equip. ~370 \\
\bottomrule
\end{tabularx}


A cobertura rural em Ñeembucú é \textbf{muito limitada a inexistente} em grande parte do território:

- As extensas zonas pantaneiras (Esteros del Ñeembucú, Laguna Cerro) têm apenas cobertura Starlink.
- Estâncias de gado tradicionais dependem de rádio HF legado ou de Starlink (quando disponível).
- A chegada de Starlink em 2023 foi especialmente impactante neste departamento.
- Copaco mantém ADSL em municípios acessíveis, mas com velocidades muito baixas.

\textbf{Recomendação crítica para imigrantes:} Qualquer propriedade rural fora de Pilar em Ñeembucú deve ser adquirida com o custo de Starlink já no planejamento. Conectividade terrestre não é confiável.


- https://www.conatel.gov.py/conatel/wp-content/uploads/2025/02/informe-de-internet-fijo-2024.pdf
- https://www.ine.gov.py/noticias/2037/como-esta-el-acceso-a-las-tecnologias-de-la-informacion-y-comunicacion-en-el-paraguay
- https://dplnews.com/starlink-ya-esta-disponible-en-paraguay-cuanto-cuesta/
- https://www.tecnologia.com.py/contectividad/por-que-starlink-esta-cambiando-la-conectividad-en-el


\clearpage
\section{Amambay}
\label{sec:dept_13_Amambay}

\subsection{Segurança Pública}

\textbf{Fonte principal:} Ministerio Público Paraguay / La Nación PY / InSight Crime / INECIP
\textbf{Data de referência:} 2020--2024
\textbf{Atualizado em:} 2026-03-20


\begin{tabularx}{\linewidth}{lX}
\toprule
Indicador & Valor \\
\midrule
Taxa de homicídios (por 100k hab) & ~30--50 (estimativa 2022--2024; histórico: 107,4 em 2016, 69,9 em 2017) \\
Crimes registrados (total/ano) & --- dados não desagregados publicamente \\
Número de comisarías & ~12 (estimativa departamental) \\
Índice segurança relativo & baixo \\
\bottomrule
\end{tabularx}


Amambay é, por ampla margem, o \textbf{departamento mais violento do Paraguai}, com taxas de homicídio que chegaram a \textbf{107,4 por 100 mil habitantes em 2016} (comparáveis às regiões mais violentas da América Central). Em 2017, a taxa era de 69,9 por 100 mil. Embora tenha havido alguma redução nos anos seguintes, o departamento permanece o mais perigoso do país.

O Ministerio Público registra que Amambay responde por \textbf{20\% de todas as causas de homicídio doloso} do Paraguai (2019--2024) --- a maior proporção de qualquer departamento --- apesar de ter apenas ~2,5\% da população nacional.

Pedro Juan Caballero, capital departamental e cidade fronteiriça com Ponta Porã (MS, Brasil), é o epicentro do narcotráfico paraguaio. As organizações criminosas brasileiras \textbf{Comando Vermelho (CV)} e \textbf{Comando da Capital (CC)} têm base operacional na cidade. Disputas entre facções resultam em assassinatos frequentes, muitos com características de execução (sicariato). Em apenas três meses de um ano recente, 36 pessoas foram assassinadas no departamento, a grande maioria ligada ao crime organizado.

O relatório INECIP (2020) documenta que os julgamentos por homicídio em Amambay tendem a resultar em mais condenações de populações indígenas vulneráveis do que dos próprios sicários, revelando problemas estruturais no sistema de justiça local.


- \textbf{Amambay não é recomendado para imigrantes brasileiros} que buscam segurança e qualidade de vida
- Pedro Juan Caballero tem presença significativa de brasileiros, mas o ambiente é de alto risco --- "ninguém morre à toa na fronteira", segundo autoridades policiais locais
- Brasileiros com negócios legítimos na região devem adotar perfil baixo, evitar ostentação e não se misturar a ambientes informais de negócios
- Deslocamentos noturnos são de alto risco --- a cidade registra assassinatos com regularidade
- A fronteira seca com Ponta Porã (MS) significa que din

\subsection{IDH e Desenvolvimento Humano}

\textbf{Fonte principal:} PNUD Paraguai / DGEEC / INE
\textbf{Ano de referência:} 2020 (IDH); 2023 (pobreza e infraestrutura)
\textbf{Atualizado em:} 2026-03-20


\begin{tabularx}{\linewidth}{lX}
\toprule
Indicador & Valor \\
\midrule
IDH & 0,720 \\
Classificação IDH & alto \\
Ranking nacional (1=melhor) & 5/18 \\
Esperança de vida & 73,6 anos \\
Anos de escolaridade média & 8,7 anos \\
RNB per capita (USD PPA) & 9.500 \\
\bottomrule
\end{tabularx}

> Nota: IDH estimado com base em dados subnacionais da Global Data Lab e perfil socioeconômico do departamento. Amambay se beneficia da economia fronteiriça com o Brasil (Pedro Juan Caballero--Ponta Porã), impulsionando renda e acesso a serviços. Estimativa baseada em dados regionais PNUD 2020.


\begin{tabularx}{\linewidth}{lX}
\toprule
Indicador & Valor \\
\midrule
Taxa de pobreza total (\%) & 22,8\% \\
Taxa de pobreza extrema (\%) & 4,5\% \\
Índice de Gini & 0,448 \\
\bottomrule
\end{tabularx}

> Fonte: INE --- EPHC 2023 e dados regionais. A economia fronteiriça beneficia Pedro Juan Caballero, mas parte do departamento tem pobreza rural mais alta. Alguns distritos de Amambay aparecem com pobreza acima de 40\% (dados distritais INE).


\begin{tabularx}{\linewidth}{lX}
\toprule
Indicador & Valor \\
\midrule
Acesso a água potável (\%) & 82,1\% \\
Acesso a saneamento (\%) & 80,5\% \\
\bottomrule
\end{tabularx}

> Estimativa baseada em dados do Censo 2022. Pedro Juan Caballero possui boa infraestrutura urbana influenciada pela integração com Ponta Porã (MS, Brasil).


Amambay tem a mais direta ligação com o Brasil: Pedro Juan Caballero é cidade-gêmea com Ponta Porã (MS), com livre circulação de pessoas e moedas. O departamento tem IDH de 0,720 (alto) impulsionado pela economia fronteiriça e comércio intenso. Para brasileiros, Pedro Juan Caballero é a cidade mais "familiar" do interior paraguaio --- fala-se português correntemente, há supermercados com produtos brasileiros, atendimento médico em ambos os lados da fronteira e forte comunidade brasiguaia. O Parque Nacional Cerro Corá (local da batalha final da Guerra do Paraguai) tem potencial turístico. A pobreza de 22,8\% é relativamente baixa. O acesso a água (82,1\%) e saneamento (80,5\%) são satisfatórios. O departamento é recomendado para brasileiros que querem a praticidade da fronteira aberta, especialmente no setor comercial e de serviços.


- PNUD Paraguai --- Informe IDH 2001-2020: https://www.undp.org/es/paraguay/noticias/IDH-2001-2020
- INE --- EPHC 2023: https://ww

\subsection{Sistema Prisional}

\textbf{Fonte:} MJT (Ministerio de Justicia y Trabajo) / MNP Paraguay
\textbf{Ano de referência:} 2023--2024
\textbf{Atualizado em:} 2026-03-20


\begin{tabularx}{\linewidth}{lXXXX}
\toprule
Nome & Município & Capacidade & População & Superlotação \\
\midrule
Penitenciaría Regional de Pedro Juan Caballero & Pedro Juan Caballero & 300 & ~850 & ~283\% \\
\bottomrule
\end{tabularx}


\begin{tabularx}{\linewidth}{lX}
\toprule
Indicador & Valor \\
\midrule
Total de estabelecimentos & 1 \\
Capacidade total (vagas) & ~300 \\
População carcerária total & ~850 \\
Taxa de superlotação & ~283\% \\
\% presos provisórios & ~73\% \\
\bottomrule
\end{tabularx}


\textbf{Pedro Juan Caballero é um dos pontos mais críticos do sistema penitenciário paraguaio}, com impacto direto no crime organizado transfronteiriço. A cidade faz fronteira direta com Ponta Porã (MS, Brasil) e é historicamente reconhecida como um dos maiores corredores de tráfico de drogas, armas e lavagem de dinheiro da América do Sul.

A Penitenciaría de Pedro Juan Caballero apresenta superlotação grave e é frequentemente citada como epicentro de disputas entre facções criminosas (PCC, Comando Vermelho, grupos locais). Assassinatos de lideranças do narcotráfico ocorrem dentro e fora do presídio com frequência alarmante --- a cidade registra sistematicamente um dos maiores índices de homicídios do Paraguai.

O MNP documentou presença de armas, celulares e drogas dentro do presídio, além de suspeitas de corrupção sistêmica de agentes penitenciários. Fugas e rebeliões são eventos recorrentes.

Em 2020--2023, múltiplos líderes do crime organizado foram assassinados dentro do presídio ou em suas imediações, revelando a permeabilidade do sistema de segurança.

\textbf{Para imigrantes:} Amambay oferece oportunidades econômicas pela dinâmica comercial de fronteira, mas representa o maior risco de segurança para imigrantes no Paraguai. A proximidade com o crime organizado e os altos índices de homicídio fazem de Pedro Juan Caballero uma localidade de alto risco. Quem optar por se instalar na região deve adotar medidas rigorosas de segurança e cautela em relações comerciais.


- https://www.mjt.gov.py/index.php/penitenciaria
- https://www.mnp.gov.py/informes
- MNP --- Informe Especial Pedro Juan Caballero 2022--2023
- ABC Color: "PJC: la ciudad más violenta del Paraguay" (2022--2023)
- Última Hora: "Narcotráfico y prisiones en Amambay" (2023)
- InSight Crime: "Paraguay's Border Town" (2022)


\subsection{Recursos Hídricos Subterrâneos}

\textbf{Fonte:} SENASA / MADES / Geología del Paraguay / IGRAC
\textbf{Atualizado em:} 2026-03-20


\begin{tabularx}{\linewidth}{lXXX}
\toprule
Aquífero & Cobertura no Dept. & Profundidade Média & Vazão Típica \\
\midrule
Guarani / Misiones (confinado) & ~60\% (porção sul/leste) & 100--400 m & 20--80 m$^3$/h \\
Cristalino (fraturado, Serra de Amambay) & ~40\% & 30--100 m & 2--10 m$^3$/h \\
Caiuá/Baurú-Acaray (porção norte) & ~20\% & 80--200 m & 10--40 m$^3$/h \\
\bottomrule
\end{tabularx}


\begin{tabularx}{\linewidth}{lX}
\toprule
Parâmetro & Valor \\
\midrule
Profundidade média --- Cristalino (m) & 40--100 \\
Profundidade média --- Guarani (m) & 100--400 \\
Profundidade máxima recomendada (m) & 450 \\
Vazão típica (m$^3$/h) & 5--50 \\
Custo estimado perfuração (USD) & 2.500 -- 10.000 \\
Qualidade da água --- Cristalino & boa (baixa mineralização) \\
Qualidade da água --- Guarani & boa a excelente \\
pH médio & 6,5--7,8 \\
Outorga necessária & sim \\
Órgão regulador & MADES --- DGPCRH \\
\bottomrule
\end{tabularx}


Amambay se caracteriza pela Serra de Amambay, que forma a fronteira com o Brasil (Mato Grosso do Sul). O embasamento cristalino domina a porção norte e oeste, onde a captação é feita em fraturas com boa produtividade quando corretamente localizadas. O Aquífero Guarani está presente no sul e leste, com boa cobertura. O sistema Caiuá/Baurú-Acaray ocorre na região norte, associado às bacias sedimentares do norte de Canindeyú. A qualidade da água é boa em todos os sistemas. Pedro Juan Caballero (capital) tem rede pública do SENASA. Uso principal: doméstico, pecuária e pequena agricultura.


Amambay tem fronteira com o Brasil e acolhe muitos brasiguaios. Pedro Juan Caballero é centro urbano com rede pública. Para propriedades rurais, poço em cristalino (40--80 m, USD 2.500--6.000) ou em Guarani (100--250 m, USD 3.500--9.000) dependendo da localização. A Serra de Amambay tem boa disponibilidade hídrica subterrânea nas fraturas. Estudo hidrogeológico prévio é recomendado no cristalino para identificar zonas fraturadas produtivas. Licença do MADES obrigatória.


- https://gw-project.org/es/caracteristicas-del-agua-subterranea-en-paraguay/
- https://es.wikipedia.org/wiki/Acu\%C3\%ADfero\_Guaran\%C3\%AD
- https://www.geologiadelparaguay.com.py/Acuiferos-Potenciales-del-Paraguay.pdf
- https://www.mades.gov.py/tramites/registro-de-empresas-publicas-y-privadas-que-se-dedican-a-la-perforacion-de-pozos-tubu

\subsection{Aptidão do Solo}

\textbf{Fonte:} MAG / SENAVE / INFONA / CONACYT
\textbf{Atualizado em:} 2026-03-20


\begin{tabularx}{\linewidth}{lXX}
\toprule
Tipo de Solo & \% do Território & Características \\
\midrule
Oxissolos (Latossolos Roxos/Amarelos) & 40\% & Solos derivados de basalto e diabásio, profundos, argilosos, cor vermelha escura a vermelho-amarelada; excelente estrutura \\
Ultissolos (Argissolos) & 35\% & Solos podzólicos com B textural; fertilidade média; ampla distribuição no centro-sul do departamento \\
Entissolos arenosos & 15\% & Solos de textura arenosa no norte e áreas de transição com o Chaco \\
Inceptissolos / Aluviais & 10\% & Planícies fluviais; fertilidade variável \\
\bottomrule
\end{tabularx}


\begin{tabularx}{\linewidth}{lXX}
\toprule
Classe & \% Território & Uso Recomendado \\
\midrule
Lavoura intensiva & 35\% & Soja, milho, trigo (Oxissolos do sul e centro) \\
Lavoura extensiva & 25\% & Mandioca, girassol, sorgo \\
Pastagem natural/melhorada & 25\% & Pecuária extensiva e semi-intensiva \\
Floresta / reserva & 12\% & Serranía de Amambay, matas ciliares \\
Sem aptidão agrícola & 3\% & Áreas urbanas (Pedro Juan Caballero) \\
\bottomrule
\end{tabularx}


\begin{tabularx}{\linewidth}{lX}
\toprule
Parâmetro & Valor \\
\midrule
pH médio & 5,0 -- 6,2 \\
Fertilidade natural & média-alta \\
Teor de matéria orgânica & médio \\
Necessidade de calcário & parcial a sim \\
\bottomrule
\end{tabularx}


- \textbf{Soja} --- em forte expansão; entre os departamentos de alta produtividade
- \textbf{Milho} --- rotação com soja
- \textbf{Pecuária bovina} --- importante; grandes fazendas tradicionais
- \textbf{Mandioca} --- agricultura familiar
- \textbf{Cannabis medicinal} --- cultivo regulamentado em expansão (zona fronteiriça com Mato Grosso do Sul)


- \textbf{Teor de manganês baixo:} estudo do CONACYT identificou Amambay com o menor nível médio de manganês disponível (20,0 mg/kg), sendo 39\% das amostras com nível "baixo"; pode limitar produção de soja sem suplementação
- \textbf{Acidez:} pH abaixo de 5,5 em proporção significativa; calagem necessária
- \textbf{Questões fundiárias:} região fronteiriça com tensões históricas sobre ocupação de terras
- \textbf{Segurança:} Pedro Juan Caballero é cidade de fronteira com problemas de segurança associados ao tráfico, o que pode impactar atividades rurais
- \textbf{Erosão:} em solos arenosos do norte e encostas da Serranía


Amambay tem potencial agrícola real, especialmente nas áreas com Oxis

\subsection{Terra Rural}

\textbf{Fonte:} INDERT / Clasificados.com.py / mercado local
\textbf{Ano de referência:} 2024
\textbf{Atualizado em:} 2026-03-20


\begin{tabularx}{\linewidth}{lXXX}
\toprule
Tipo & Preço (USD/ha) & Faixa & Tendência \\
\midrule
Agrícola alta produtividade & 6.000 & 4.000 -- 9.000 & $\uparrow$ \\
Pastagem & 2.800 & 1.800 -- 4.000 & $\uparrow$ \\
Mata / reserva & 1.200 & 700 -- 2.000 & $\rightarrow$ \\
\bottomrule
\end{tabularx}


\begin{tabularx}{\linewidth}{lX}
\toprule
Parâmetro & Valor \\
\midrule
Valorização anual média (\%) & 7--9\% \\
Lote mínimo rural (ha) & 20 \\
Imposto territorial rural (\%) & 1\% sobre valor fiscal \\
Liquidez do mercado & média \\
\bottomrule
\end{tabularx}


- \textbf{Pedro Juan Caballero (entorno):} Capital departamental e cidade-gêmea com Ponta Porã (BR); intensa atividade comercial binacional; terra periurbana valorizada.
- \textbf{Bella Vista Norte:} Zona de colonização com solos vermelhos férteis; produção de soja, milho e tabaco.
- \textbf{Capitán Bado:} Polo comercial próximo à fronteira brasileira (Coronel Sapucaia/MS); forte demanda rural.
- \textbf{Pedro Juan rural (Zona Norte):} Área de expansão da soja; terras com aptidão agrícola crescente.


\begin{tabularx}{\linewidth}{lXXX}
\toprule
Indicador & Amambay (PY) & Mato Grosso do Sul (BR) --- Fronteiriço & Paraná (BR) \\
\midrule
Terra agrícola (USD/ha) & 4.000--9.000 & 10.000--20.000 & 15.000--30.000 \\
Terra pastagem (USD/ha) & 1.800--4.000 & 4.500--9.000 & 6.000--12.000 \\
Custo relativo & ~40\% do MS & referência MS & referência PR \\
\bottomrule
\end{tabularx}

Amambay tem posição estratégica como fronteira viva com o Brasil (Mato Grosso do Sul). O fluxo binacional cria demanda constante por terra e imóveis em ambos os lados da fronteira, mas os preços paraguaios seguem muito abaixo dos brasileiros.


- [INDERT --- Instituto Nacional de Desarrollo Rural y de la Tierra](https://www.indert.gov.py/)
- [Amambay Propiedades](https://www.amambaypropiedades.com/)
- [Plan B Paraguay --- Farmland Prices](https://x.com/planbparaguay/status/1834232169759211739)
- [InfoCasas --- Chacras Paraguay](https://www.infocasas.com.py/venta/chacras-o-campos)


\subsection{Mercado Imobiliário}

\textbf{Capital:} Pedro Juan Caballero
\textbf{Fonte:} Amambay Propiedades / Clasificados.com.py / mercado local
\textbf{Ano de referência:} 2024
\textbf{Atualizado em:} 2026-03-20


\begin{tabularx}{\linewidth}{lXX}
\toprule
Tipo & Preço (USD/m$^2$) & Aluguel (USD/mês) \\
\midrule
Apartamento 2 quartos & 700 -- 1.100 & 280 -- 550 \\
Casa 3 quartos & 600 -- 950 & 250 -- 480 \\
\bottomrule
\end{tabularx}


\begin{tabularx}{\linewidth}{lX}
\toprule
Parâmetro & Valor \\
\midrule
Valorização anual (\%) & 6--9\% \\
Disponibilidade de financiamento & boa (cooperativas fortes na cidade fronteiriça, também influência de Ponta Porã) \\
Bairros mais valorizados & Centro, Barrio Las Mercedes, Zona Norte (próxima à fronteira) \\
\bottomrule
\end{tabularx}


\textbf{Vale a pena comprar ou alugar?}
Pedro Juan Caballero (PJC) é uma das cidades mais dinâmicas e cosmopolitas do interior paraguaio, formando conurbação com Ponta Porã (MS). A intensa atividade comercial binacional cria mercado imobiliário mais líquido que cidades do interior paraguaio. Para brasileiros, é talvez a transição mais fácil: cidade bilíngue (português/espanhol) com modos de vida muito similares ao Brasil.

\textbf{Áreas recomendadas:}
- \textbf{Centro de PJC:} comércio intenso, serviços completos, movimentado.
- \textbf{Barrio Las Mercedes / Residencial:} zonas mais tranquilas para famílias.
- \textbf{Capitán Bado:} cidade menor a 90 km; mais tranquila, menor custo, boa para quem trabalha no agronegócio.

\textbf{Armadilhas comuns:}
- Pedro Juan Caballero tem índice de criminalidade elevado para padrão paraguaio; segurança do imóvel é prioridade.
- Muitos imóveis têm documentação cruzada com registros brasileiros e paraguaios; verificar escrituras com cuidado.
- A economia local é altamente dependente do comércio binacional; volatilidade cambial BRL/USD impacta o mercado.
- Aluguéis às vezes cotados em guaranis ou reais, dependendo do proprietário.


- [Amambay Propiedades --- Alquileres](https://www.amambaypropiedades.com/estado-propiedad/alquiler/)
- [InfoCasas.com.py --- Amambay](https://www.infocasas.com.py/venta/inmuebles/amambay)
- [Clasipar.com --- Alquiler Amambay](https://clasipar.paraguay.com/alquiler/inmuebles/amambay)


\subsection{Cobertura Celular}

\textbf{Fonte:} CONATEL / Tigo / Personal / OpenSignal / nPerf
\textbf{Ano de referência:} 2024
\textbf{Atualizado em:} 2026-03-20


\begin{tabularx}{\linewidth}{lXXXXX}
\toprule
Operadora & 2G & 3G & 4G & 5G & Cobertura Pop. \\
\midrule
Tigo & 85\% & 75\% & 65\% & --- & 80\% \\
Personal/Claro & 80\% & 70\% & 58\% & --- & 74\% \\
VOX & 48\% & 32\% & 18\% & --- & 40\% \\
\bottomrule
\end{tabularx}

> \textbf{Nota:} Amambay tem como capital Pedro Juan Caballero (PJC), cidade de fronteira com Ponta Porã (MS-BR), o que cria uma situação única de sobreposição de redes brasileiras e paraguaias. A cobertura 4G é boa em PJC e entorno; interior norte e sul tem cobertura mais fraca.


\begin{tabularx}{\linewidth}{lX}
\toprule
Parâmetro & Valor \\
\midrule
Melhor operadora rural & Tigo \\
Velocidade 4G média (Mbps) & 18--38 \\
Áreas sem cobertura & Interior sul (Bella Vista Norte interior, áreas de reserva) \\
Sinal em propriedade rural & moderado \\
\bottomrule
\end{tabularx}


\begin{tabularx}{\linewidth}{lXXX}
\toprule
Plano & Operadora & USD/mês & Dados \\
\midrule
Pré-pago básico (7 dias) & Tigo & ~2,50 & 3 GB \\
Pré-pago intermediário (30 dias) & Personal & ~5,00 & 8 GB \\
Pós-pago entrada & Tigo & ~12,00 & 15 GB + chamadas \\
Pós-pago intermediário & Personal/Claro & ~18,00 & 30 GB + chamadas \\
\bottomrule
\end{tabularx}


- \textbf{Pedro Juan Caballero:} Cidade de conurbação com Ponta Porã (BR). \textbf{Chips brasileiros funcionam normalmente em PJC} (Vivo, Claro BR, TIM) e vice-versa. Situação única no Paraguai.
- \textbf{Vantagem fronteiriça:} Imigrantes em PJC podem usar chips brasileiros com planos nacionais que cobrem roaming internacional com menor custo, ou chips paraguaios com melhor relação custo-benefício.
- \textbf{Interior do departamento:} Tigo é a operadora com maior cobertura; áreas de fazenda podem ter sinal apenas 2G/3G.
- \textbf{Bella Vista Norte:} Cobertura 3G/4G razoável de Tigo.
- \textbf{Reserva Natural Mbaracayú (limite):} Sinal muito fraco.


- https://www.conatel.gov.py/conatel/indicadores/
- https://www.tigo.com.py/cobertura
- https://www.personal.com.py/institucional/cobertura.html
- https://www.nperf.com/en/map/PY/-/-/signal
- https://prepaid-data-sim-card.fandom.com/wiki/Paraguay


\subsection{Internet}

\textbf{Fonte:} CONATEL / Tigo / Claro / INE / Copaco
\textbf{Ano de referência:} 2024
\textbf{Atualizado em:} 2026-03-20


\begin{tabularx}{\linewidth}{lX}
\toprule
Indicador & Valor \\
\midrule
\% população com internet & ~65\% \\
\% com banda larga fixa & ~22\% dos domicílios \\
Velocidade média download (Mbps) & 25--55 (PJC) / 5--15 (rural) \\
Velocidade média upload (Mbps) & 8--20 (PJC) \\
\bottomrule
\end{tabularx}

> Pedro Juan Caballero, pela condição de cidade de fronteira com Ponta Porã (BR), tem uma situação peculiar de mercado de telecomunicações --- competição de provedores de ambos os países.


\begin{tabularx}{\linewidth}{lXX}
\toprule
Tecnologia & Disponibilidade & Velocidade Típica \\
\midrule
Fibra ótica (FTTH) & Pedro Juan Caballero, Bella Vista Norte & 50--300 Mbps \\
Cabo HFC & Pedro Juan Caballero & 30--150 Mbps \\
Rádio/Wireless & Municípios do departamento & 5--50 Mbps \\
ADSL (Copaco) & Interior com rede legada & 2--10 Mbps \\
Satélite (Starlink) & Todo o departamento & 50--200 Mbps \\
\bottomrule
\end{tabularx}


\begin{tabularx}{\linewidth}{lXXX}
\toprule
Provedor & Tecnologia & Área de cobertura & Preço ref. (USD/mês) \\
\midrule
Tigo & Fibra/HFC/Rádio & PJC + municípios & 18--50 \\
Claro/Personal & Fibra/Rádio & PJC & 15--45 \\
Copaco & ADSL/Rádio & Interior & 10--28 \\
Provedores brasileiros & Fibra/rádio & Zona de fronteira & preços em BRL \\
Starlink & Satélite LEO & Todo o departamento & ~44 + equip. ~370 \\
\bottomrule
\end{tabularx}


A cobertura rural de Amambay é \textbf{moderada}:

- Pedro Juan Caballero tem boa infraestrutura, beneficiada também pela concorrência de provedores do lado brasileiro (Ponta Porã).
- Interior do departamento tem ISPs locais e rádio wireless.
- Zonas ao sul (estâncias na Serra Amambay) têm cobertura mais limitada.
- \textbf{Situação fronteiriça única:} Em PJC, provedores brasileiros de fibra (Vivo Fibra, Claro Fibra) também têm presença, podendo ser contratados com planos em reais.
- Starlink como solução para estâncias e propriedades remotas.


- https://www.conatel.gov.py/conatel/wp-content/uploads/2025/02/informe-de-internet-fijo-2024.pdf
- https://www.ine.gov.py/noticias/2037/como-esta-el-acceso-a-las-tecnologias-de-la-informacion-y-comunicacion-en-el-paraguay
- https://dplnews.com/starlink-ya-esta-disponible-en-paraguay-cuanto-cuesta/
- https://www.tigo.com.py/internet



\clearpage
\section{Canindeyú}
\label{sec:dept_14_Canindeyu}

\subsection{Segurança Pública}

\textbf{Fonte principal:} Ministerio Público Paraguay / SENAD / Presidência do Paraguai / La Nación PY
\textbf{Data de referência:} 2023--2024
\textbf{Atualizado em:} 2026-03-20


\begin{tabularx}{\linewidth}{lX}
\toprule
Indicador & Valor \\
\midrule
Taxa de homicídios (por 100k hab) & ~15--25 (estimativa, significativamente acima da média nacional) \\
Crimes registrados (total/ano) & --- dados não desagregados publicamente \\
Número de comisarías & ~10 (estimativa departamental) \\
Índice segurança relativo & baixo \\
\bottomrule
\end{tabularx}


Canindeyú é um dos departamentos de maior risco de segurança do Paraguai. O Ministerio Público registra que o departamento responde por \textbf{aproximadamente 8\% das causas de homicídio doloso} nacional (2019--2024), proporção muito elevada para um departamento com ~5\% da população.

O departamento tornou-se teatro de guerra entre organizações criminosas em 2023--2024. A principal disputa foi travada entre as facções lideradas por Felipe Santiago Acosta ("Macho") e Benicio Silva ("Silva Hu"), pelo controle das \textbf{plantações de maconha} e pistas clandestinas para entrada de cocaína no Brasil. Em dezembro de 2023, a operação "Ignis" (SENAD + FTC + Polícia + autoridades brasileiras) resultou na morte de 9 e prisão de 10 membros de organização ligada ao PCC.

Em abril de 2024, o presidente Santiago Peña enviou a \textbf{Fuerza de Tarea Conjunta (FTC)} ao departamento, com 70 homens, 3 veículos blindados e 9 pick-ups, para pacificar a região e proteger fazendeiros ameaçados pelo crime organizado. Um helicóptero do CODI foi atacado em retaliação.

Canindeyú está localizado entre Amambay e Alto Paraná, na fronteira com os estados brasileiros de Mato Grosso do Sul e Paraná.


- \textbf{Canindeyú não é recomendado para imigrantes brasileiros} que buscam segurança
- Há presença de comunidades de agricultores brasileiros (brasiguaios) no departamento, especialmente no agronegócio
- A guerra entre facções criminosas pelo controle do narcotráfico afeta diretamente a segurança de toda a população civil
- Fazendeiros e agricultores da região relatam extorsão e ameaças por parte de grupos criminosos
- A militarização da região pela FTC trouxe algum alívio, mas não eliminou o risco estrutural
- Brasileiros com propriedades na região devem manter contato com associações rurais locais e ter planos de contingência para evacuação em caso de escalada de violência


- https://www.la

\subsection{IDH e Desenvolvimento Humano}

\textbf{Fonte principal:} PNUD Paraguai / DGEEC / INE
\textbf{Ano de referência:} 2020 (IDH); 2023 (pobreza e infraestrutura)
\textbf{Atualizado em:} 2026-03-20


\begin{tabularx}{\linewidth}{lX}
\toprule
Indicador & Valor \\
\midrule
IDH & 0,685 \\
Classificação IDH & médio \\
Ranking nacional (1=melhor) & 15/18 \\
Esperança de vida & 71,2 anos \\
Anos de escolaridade média & 7,3 anos \\
RNB per capita (USD PPA) & 7.100 \\
\bottomrule
\end{tabularx}

> Nota: IDH estimado com base no relatório PNUD 2020 e dados subnacionais. Canindeyú tem significativa expansão agrícola (soja, milho) mas ainda apresenta IDH médio por sua ruralidade intensa e populações indígenas vulneráveis. Estimativa baseada em dados regionais e Global Data Lab.


\begin{tabularx}{\linewidth}{lX}
\toprule
Indicador & Valor \\
\midrule
Taxa de pobreza total (\%) & 35,1\% \\
Taxa de pobreza extrema (\%) & 10,8\% \\
Índice de Gini & 0,467 \\
\bottomrule
\end{tabularx}

> Fonte: INE --- EPHC 2023 (Canindeyú: 35,1\% pobreza total). O departamento tem alta incidência de pobreza monetária, com aumento de 0,5\% na pobreza extrema em 2023. Pobreza multidimensional (IPM) de 27,23\%.


\begin{tabularx}{\linewidth}{lX}
\toprule
Indicador & Valor \\
\midrule
Acesso a água potável (\%) & 67,4\% \\
Acesso a saneamento (\%) & 64,8\% \\
\bottomrule
\end{tabularx}

> Estimativa baseada em dados do Censo 2022. Salto del Guairá (capital) tem melhor infraestrutura por sua economia fronteiriça, mas o interior do departamento tem cobertura muito baixa.


Canindeyú está na fronteira nordeste com o Brasil (Mato Grosso do Sul e Paraná) e tem Salto del Guairá como capital, em frente a Guaíra (PR). O departamento é fortemente agrícola, com expansão intensa da soja nos últimos 20 anos que atraiu muitos agricultores brasileiros. Apesar disso, o IDH permanece médio (0,685) e a pobreza é alta (35,1\%), refletindo a coexistência entre agronegócio moderno e populações rurais e indígenas vulneráveis. Para brasileiros no agronegócio, é região conhecida e com comunidade estabelecida. A infraestrutura urbana é básica, com hospitais regionais limitados. O acesso a água (67,4\%) e saneamento (64,8\%) são abaixo da média nacional, especialmente no interior. Recomendado para quem trabalha diretamente com produção agrícola e pode complementar serviços no Brasil vizinho.


- PNUD Paraguai --- Informe IDH 2001-2020: https://www.undp.org/es/paraguay/noticias/IDH-2001-2020
- INE --- EPHC 2023: https://www.ine.gov.py/noticias/1928/pobreza-monet

\subsection{Sistema Prisional}

\textbf{Fonte:} MJT (Ministerio de Justicia y Trabajo) / MNP Paraguay
\textbf{Ano de referência:} 2023--2024
\textbf{Atualizado em:} 2026-03-20


\begin{tabularx}{\linewidth}{lXXXX}
\toprule
Nome & Município & Capacidade & População & Superlotação \\
\midrule
Penitenciaría Regional de Salto del Guairá & Salto del Guairá & 150 & ~290 & ~193\% \\
\bottomrule
\end{tabularx}


\begin{tabularx}{\linewidth}{lX}
\toprule
Indicador & Valor \\
\midrule
Total de estabelecimentos & 1 \\
Capacidade total (vagas) & ~150 \\
População carcerária total & ~290 \\
Taxa de superlotação & ~193\% \\
\% presos provisórios & ~71\% \\
\bottomrule
\end{tabularx}


Canindeyú faz fronteira com o Brasil (Paraná --- Guaíra e Mato Grosso do Sul) e é área de intensa atividade agrícola (soja) e tráfico de drogas. A única penitenciaría regional do departamento está em Salto del Guairá, cidade fronteiriça com Guaíra (PR, Brasil), ligadas pela Ponte Internacional da Amizade II.

O departamento é rota significativa para o tráfico de maconha (produzida nos departamentos de Canindeyú e Alto Paraná) em direção ao Brasil e Argentina. Parte considerável da população carcerária está relacionada a esses crimes.

O MNP registrou condições precárias no estabelecimento, com superlotação e falta de assistência jurídica adequada. A alta proporção de presos provisórios reflete gargalos no sistema judiciário local.

A região tem uma das maiores concentrações de soja do Paraguai, com forte presença de produtores brasileiros (brasiguaios). O comércio intenso com o Brasil alimenta tanto oportunidades econômicas quanto fluxos de contrabando.

\textbf{Para imigrantes:} Canindeyú oferece oportunidades no agronegócio, especialmente para quem tem experiência em agricultura de precisão. Salto del Guairá tem boa dinâmica comercial com o Brasil. O risco de segurança é moderado a elevado nas zonas de fronteira e em rotas rurais, mas áreas urbanas são relativamente tranquilas para residência.


- https://www.mjt.gov.py/index.php/penitenciaria
- https://www.mnp.gov.py/informes
- MNP --- Informe Anual 2022
- ABC Color: "Canindeyú: frontera del narcotráfico" (2022)


\subsection{Recursos Hídricos Subterrâneos}

\textbf{Fonte:} SENASA / MADES / Geología del Paraguay / IGRAC
\textbf{Atualizado em:} 2026-03-20


\begin{tabularx}{\linewidth}{lXXX}
\toprule
Aquífero & Cobertura no Dept. & Profundidade Média & Vazão Típica \\
\midrule
Guarani / Misiones (confinado) & ~70\% & 80--350 m & 20--100 m$^3$/h \\
Caiuá/Baurú-Acaray (noroeste) & ~25\% & 80--200 m & 10--40 m$^3$/h \\
Basáltico (fraturado) & ~30\% (cobertura superficial) & 40--120 m & 5--20 m$^3$/h \\
\bottomrule
\end{tabularx}


\begin{tabularx}{\linewidth}{lX}
\toprule
Parâmetro & Valor \\
\midrule
Profundidade média --- Guarani (m) & 80--300 \\
Profundidade média --- Caiuá/Baurú (m) & 80--200 \\
Profundidade máxima recomendada (m) & 400 \\
Vazão típica (m$^3$/h) & 15--60 \\
Custo estimado perfuração (USD) & 3.500 -- 12.000 \\
Qualidade da água & boa a excelente \\
pH médio & 6,5--7,5 \\
Outorga necessária & sim \\
Órgão regulador & MADES --- DGPCRH \\
\bottomrule
\end{tabularx}


Canindeyú é um departamento de fronteira com o Brasil (Paraná), com expansão agrícola acelerada nas últimas décadas. O Aquífero Guarani está presente na maior parte do território, com boa qualidade. O sistema Caiuá/Baurú-Acaray, identificado na região noroeste do departamento associado à bacia do rio Paraná, é outro aquífero de importância local. A qualidade da água subterrânea é geralmente boa, com baixa mineralização. Salto del Guairá (capital) tem crescimento urbano intenso. Uso principal: abastecimento doméstico, irrigação (soja, milho) e uso industrial.


Canindeyú é destino de muitos agricultores brasileiros na fronteira com o Paraná. O Aquífero Guarani é amplamente acessível em 80--250 m, com custo de perfuração estimado em USD 3.500--10.000. Para fazendas de soja/milho com necessidade de irrigação ou uso industrial, a vazão de 20--80 m$^3$/h disponível é mais que suficiente. A rede pública é limitada fora das cidades. Poço artesiano próprio é praticamente universal nas propriedades rurais. Licença do MADES obrigatória para qualquer perfuração.


- https://gw-project.org/es/caracteristicas-del-agua-subterranea-en-paraguay/
- https://es.wikipedia.org/wiki/Acu\%C3\%ADfero\_Guaran\%C3\%AD
- https://www.geologiadelparaguay.com.py/Acuiferos-Potenciales-del-Paraguay.pdf
- https://www.mades.gov.py/tramites/registro-de-empresas-publicas-y-privadas-que-se-dedican-a-la-perforacion-de-pozos-tubulares/


\subsection{Aptidão do Solo}

\textbf{Fonte:} MAG / SENAVE / INFONA / CONACYT
\textbf{Atualizado em:} 2026-03-20


\begin{tabularx}{\linewidth}{lXX}
\toprule
Tipo de Solo & \% do Território & Características \\
\midrule
Oxissolos (Latossolos Roxos) & 50\% & Solos basálticos profundos, argilosos, bem drenados; alta fertilidade; cor vermelho-escura intensa \\
Ultissolos (Argissolos) & 30\% & Solos podzólicos com B textural; fertilidade média-alta; boa resposta à correção \\
Entissolos arenosos & 12\% & Solos de textura arenosa no norte (região do cerrado); baixa fertilidade \\
Aluviais / Inceptissolos & 8\% & Planícies do Rio Jejuí e afluentes; fertilidade variável \\
\bottomrule
\end{tabularx}


\begin{tabularx}{\linewidth}{lXX}
\toprule
Classe & \% Território & Uso Recomendado \\
\midrule
Lavoura intensiva & 50\% & Soja, milho, trigo (Oxissolos predominantes) \\
Lavoura extensiva & 20\% & Mandioca, girassol, sorgo, pastagem melhorada \\
Pastagem natural & 15\% & Pecuária extensiva (norte, solos arenosos) \\
Floresta / reserva & 12\% & Reserva Mbaracayú, matas ciliares do Jejuí \\
Sem aptidão agrícola & 3\% & Áreas urbanas (Salto del Guairá) e degradadas \\
\bottomrule
\end{tabularx}


\begin{tabularx}{\linewidth}{lX}
\toprule
Parâmetro & Valor \\
\midrule
pH médio & 5,0 -- 6,3 \\
Fertilidade natural & alta \\
Teor de matéria orgânica & médio-alto \\
Necessidade de calcário & parcial \\
\bottomrule
\end{tabularx}


- \textbf{Soja} --- 3ª maior área plantada do Paraguai (635.000 ha, 18\% do total nacional)
- \textbf{Milho} --- em rotação com soja
- \textbf{Trigo} --- cultivo de inverno crescente
- \textbf{Mandioca} e \textbf{amendoim} --- agricultura familiar
- \textbf{Pecuária} --- bovina de corte em sistema semi-intensivo


- \textbf{Teor de manganês alto:} Canindeyú apresentou o maior nível médio de manganês disponível do país (68,4 mg/kg); em solos ácidos, manganês em excesso pode ser tóxico para culturas sensíveis
- \textbf{Acidez do solo:} pH abaixo de 5,5 em proporção relevante; calagem periódica necessária
- \textbf{Desmatamento acelerado:} um dos departamentos mais impactados pelo desmatamento 2020-2022 (6.108 ha --- o maior da Região Oriental no período)
- \textbf{Pressão sobre Reserva Mbaracayú:} expansão agrícola próxima à reserva biosférica gera conflitos
- \textbf{Irregularidades fundiárias:} fronteira com Paraná (Brasil) histórico de ocupações irregulares


Canindeyú é a "nova fronteira agrícola" do Paraguai para produtores brasileiros. Com solos

\subsection{Terra Rural}

\textbf{Fonte:} INDERT / Clasificados.com.py / mercado local
\textbf{Ano de referência:} 2024
\textbf{Atualizado em:} 2026-03-20


\begin{tabularx}{\linewidth}{lXXX}
\toprule
Tipo & Preço (USD/ha) & Faixa & Tendência \\
\midrule
Agrícola alta produtividade & 7.500 & 5.000 -- 11.000 & $\uparrow$ \\
Pastagem & 3.500 & 2.200 -- 5.000 & $\uparrow$ \\
Mata / reserva & 1.500 & 800 -- 2.500 & $\rightarrow$ \\
\bottomrule
\end{tabularx}


\begin{tabularx}{\linewidth}{lX}
\toprule
Parâmetro & Valor \\
\midrule
Valorização anual média (\%) & 8--10\% \\
Lote mínimo rural (ha) & 10 \\
Imposto territorial rural (\%) & 1\% sobre valor fiscal \\
Liquidez do mercado & média--alta \\
\bottomrule
\end{tabularx}


- \textbf{Salto del Guairá (entorno):} Capital departamental na tríplice fronteira (BR--PY); histórico polo comercial; terra periurbana em forte valorização impulsionada pelo agronegócio.
- \textbf{Curuguaty:} Polo agrícola central do departamento; expansão da soja e milho em solos vermelhos.
- \textbf{Ypehú / Corpus Christi:} Fronteira com o Brasil (Guaíra/PR); forte presença de colonos brasileiros; terra produtiva.
- \textbf{Nueva Esperanza:} Zona de expansão agrícola com terras mais acessíveis mas férteis.
- \textbf{Villa Ygatimí:} Área de colonização recente; potencial para expansão agrícola.


\begin{tabularx}{\linewidth}{lXXX}
\toprule
Indicador & Canindeyú (PY) & Paraná --- Fronteira (BR) & Mato Grosso do Sul (BR) \\
\midrule
Terra agrícola (USD/ha) & 5.000--11.000 & 18.000--35.000 & 10.000--20.000 \\
Terra pastagem (USD/ha) & 2.200--5.000 & 7.000--14.000 & 5.000--10.000 \\
Custo relativo & ~38\% do PR & referência PR & referência MS \\
\bottomrule
\end{tabularx}

Canindeyú faz fronteira direta com o Paraná e Mato Grosso do Sul do Brasil. É um dos departamentos com crescimento agrícola mais acelerado do Paraguai na última década, impulsionado pela expansão da fronteira sojícola.


- [INDERT --- Instituto Nacional de Desarrollo Rural y de la Tierra](https://www.indert.gov.py/)
- [Plan B Paraguay --- Farmland Prices](https://x.com/planbparaguay/status/1834232169759211739)
- [DealStream --- Agricultural Land Paraguay](https://dealstream.com/land-sale/agricultural/paraguay)
- [Investinag --- Paraguay Ag Investing](https://investinag.com/2024/10/16/paraguay-an-intelligent-entry-into-south-american-ag-investing/)


\subsection{Mercado Imobiliário}

\textbf{Capital:} Salto del Guairá
\textbf{Fonte:} Clasificados.com.py / mercado local
\textbf{Ano de referência:} 2024
\textbf{Atualizado em:} 2026-03-20


\begin{tabularx}{\linewidth}{lXX}
\toprule
Tipo & Preço (USD/m$^2$) & Aluguel (USD/mês) \\
\midrule
Apartamento 2 quartos & 650 -- 1.000 & 250 -- 480 \\
Casa 3 quartos & 550 -- 850 & 220 -- 420 \\
\bottomrule
\end{tabularx}


\begin{tabularx}{\linewidth}{lX}
\toprule
Parâmetro & Valor \\
\midrule
Valorização anual (\%) & 6--9\% \\
Disponibilidade de financiamento & boa (cooperativas ativas, influência do mercado brasileiro próximo) \\
Bairros mais valorizados & Centro, Barrio San José, proximidades da Represa Itaipú (área norte) \\
\bottomrule
\end{tabularx}


\textbf{Vale a pena comprar ou alugar?}
Salto del Guairá é uma cidade de tamanho médio (aprox. 90.000 hab.) que prosperou com o comércio fronteiriço e o agronegócio. Após a inundação pelo Lago Itaipu (que submerging a histórica Salto del Guairá), a nova cidade tem infraestrutura razoável. Para brasileiros oriundos do Paraná, é uma das transições mais naturais.

\textbf{Áreas recomendadas:}
- \textbf{Centro de Salto del Guairá:} serviços concentrados, comércio ativo.
- \textbf{Curuguaty:} cidade mais central, boa para quem opera na produção agrícola regional.

\textbf{Armadilhas comuns:}
- Salto del Guairá tem dinâmica comercial intensa com o Brasil; muitos imóveis são cotados em reais; padronizar toda documentação em guaranis.
- Verificar se o imóvel está na zona inundável histórica do lago.
- Infraestrutura variável: algumas zonas têm boa pavimentação e outras não.
- Para escolas de qualidade, a opção de escola bilíngue português-espanhol é comum na região fronteiriça.


- [InfoCasas.com.py --- Canindeyú](https://www.infocasas.com.py/venta/inmuebles/canindeyu)
- [Clasipar.com --- Alquiler Canindeyú](https://clasipar.paraguay.com/alquiler/inmuebles/canindeyu)
- [INDERT](https://www.indert.gov.py/)


\subsection{Cobertura Celular}

\textbf{Fonte:} CONATEL / Tigo / Personal / OpenSignal / nPerf
\textbf{Ano de referência:} 2024
\textbf{Atualizado em:} 2026-03-20


\begin{tabularx}{\linewidth}{lXXXXX}
\toprule
Operadora & 2G & 3G & 4G & 5G & Cobertura Pop. \\
\midrule
Tigo & 78\% & 65\% & 52\% & --- & 70\% \\
Personal/Claro & 68\% & 55\% & 42\% & --- & 60\% \\
VOX & 38\% & 22\% & 10\% & --- & 32\% \\
\bottomrule
\end{tabularx}

> \textbf{Nota:} Canindeyú faz fronteira com o Brasil (MS e PR), o que cria sobreposição de sinal em cidades como Salto del Guairá (conurbação com Guaíra-PR). O departamento tem grande concentração de agricultores brasileiros. Cobertura 4G razoável nas cidades; interior com lacunas.


\begin{tabularx}{\linewidth}{lX}
\toprule
Parâmetro & Valor \\
\midrule
Melhor operadora rural & Tigo \\
Velocidade 4G média (Mbps) & 15--32 \\
Áreas sem cobertura & Interior norte/centro (Reserva Mbaracayú), áreas remotas de mata \\
Sinal em propriedade rural & moderado a ruim \\
\bottomrule
\end{tabularx}


\begin{tabularx}{\linewidth}{lXXX}
\toprule
Plano & Operadora & USD/mês & Dados \\
\midrule
Pré-pago básico (7 dias) & Tigo & ~2,50 & 3 GB \\
Pré-pago intermediário (30 dias) & Personal & ~5,00 & 8 GB \\
Pós-pago entrada & Tigo & ~12,00 & 15 GB + chamadas \\
Pós-pago intermediário & Personal/Claro & ~18,00 & 30 GB + chamadas \\
\bottomrule
\end{tabularx}


- \textbf{Salto del Guairá:} Cidade fronteiriça com Guaíra (PR-BR). Chips brasileiros funcionam em Salto del Guairá e vice-versa. Situação similar a Pedro Juan Caballero/Ponta Porã.
- \textbf{Nueva Esperanza, Corpus Christi, Curuguaty:} Cobertura 3G/4G de Tigo; Personal com presença mais limitada.
- \textbf{Colônias agrícolas brasileiras:} Grande concentração no departamento; ISPs locais em muitas colônias.
- \textbf{Reserva Natural Mbaracayú:} Sinal muito fraco ou inexistente.
- \textbf{Recomendação:} Tigo é a operadora mais confiável no interior de Canindeyú; para propriedades rurais remotas, Starlink.


- https://www.conatel.gov.py/conatel/indicadores/
- https://www.tigo.com.py/cobertura
- https://www.personal.com.py/institucional/cobertura.html
- https://www.nperf.com/en/map/PY/-/-/signal


\subsection{Internet}

\textbf{Fonte:} CONATEL / Tigo / Claro / INE / Copaco
\textbf{Ano de referência:} 2024
\textbf{Atualizado em:} 2026-03-20


\begin{tabularx}{\linewidth}{lX}
\toprule
Indicador & Valor \\
\midrule
\% população com internet & ~58\% \\
\% com banda larga fixa & ~15\% dos domicílios \\
Velocidade média download (Mbps) & 15--40 (urbano) / 3--12 (rural) \\
Velocidade média upload (Mbps) & 5--15 (urbano) \\
\bottomrule
\end{tabularx}


\begin{tabularx}{\linewidth}{lXX}
\toprule
Tecnologia & Disponibilidade & Velocidade Típica \\
\midrule
Fibra ótica (FTTH) & Salto del Guairá, Curuguaty & 30--200 Mbps \\
Cabo HFC & Salto del Guairá & 20--100 Mbps \\
Rádio/Wireless & Municípios do departamento & 5--40 Mbps \\
ADSL (Copaco) & Interior com rede legada & 1--8 Mbps \\
Satélite (Starlink) & Todo o departamento & 50--200 Mbps \\
\bottomrule
\end{tabularx}


\begin{tabularx}{\linewidth}{lXXX}
\toprule
Provedor & Tecnologia & Área de cobertura & Preço ref. (USD/mês) \\
\midrule
Tigo & Fibra/Rádio & Salto del Guairá + municípios & 18--45 \\
Claro/Personal & Fibra/Rádio & Salto del Guairá & 15--40 \\
Copaco & ADSL/Rádio & Interior & 10--25 \\
ISPs locais & Rádio wireless & Colônias agrícolas & 8--20 \\
Provedores brasileiros & Fibra/rádio & Zona de fronteira com Guaíra-PR & preços em BRL \\
Starlink & Satélite LEO & Todo o departamento & ~44 + equip. ~370 \\
\bottomrule
\end{tabularx}


A cobertura rural de Canindeyú é \textbf{limitada a moderada}:

- \textbf{Colônias agrícolas brasileiras} são numerosas no departamento; muitas têm ISPs locais de rádio wireless criados pela própria comunidade.
- \textbf{Faixa de fronteira com Brasil:} Provedores de Guaíra (PR) têm sinal em alguns pontos de Salto del Guairá.
- \textbf{Interior norte e centro} (Reserva Mbaracayú): Cobertura muito fraca; Starlink é a única solução.
- \textbf{Curuguaty, Nueva Esperanza:} Rádio wireless de qualidade variável.

\textbf{Para imigrantes:} Canindeyú tem forte presença de brasileiros, o que facilita a adaptação. ISPs locais de brasileiros atendem parte das colônias. Para propriedades mais remotas, Starlink é recomendado.


- https://www.conatel.gov.py/conatel/wp-content/uploads/2025/02/informe-de-internet-fijo-2024.pdf
- https://www.ine.gov.py/noticias/2037/como-esta-el-acceso-a-las-tecnologias-de-la-informacion-y-comunicacion-en-el-paraguay
- https://dplnews.com/starlink-ya-esta-disponible-en-paraguay-cuanto-cuesta/
- https://www.tigo.com.py/internet



\clearpage
\section{Presidente Hayes}
\label{sec:dept_15_Presidente_Hayes}

\subsection{Segurança Pública}

\textbf{Fonte principal:} Ministerio Público Paraguay / INE / Atlas da Violência CONACYT
\textbf{Data de referência:} 2019--2024
\textbf{Atualizado em:} 2026-03-20


\begin{tabularx}{\linewidth}{lX}
\toprule
Indicador & Valor \\
\midrule
Taxa de homicídios (por 100k hab) & 5,5--5,6 (2019--2020, fonte INE --- próxima à média nacional) \\
Crimes registrados (total/ano) & baixo em números absolutos (população esparsa) \\
Número de comisarías & ~8 (estimativa --- vasto território com baixa densidade) \\
Índice segurança relativo & médio \\
\bottomrule
\end{tabularx}


Presidente Hayes é o maior departamento do Paraguai Oriental (parte leste do Chaco) e possui uma das menores densidades populacionais do país. A taxa de homicídios de \textbf{5,5--5,6 por 100 mil habitantes} (2019--2020) está próxima da média nacional, mas em números absolutos os crimes são raros dado o pequeno contingente populacional.

O Ministerio Público registra que o departamento responde por \textbf{aproximadamente 4\% das causas de homicídio doloso} nacional (2019--2024) --- proporção moderada para um departamento com ~2\% da população. O Chaco paraguaio é identificado como área de transito de cocaína: avionetas oriundas de Bolívia e Peru utilizam pistas clandestinas na região para abastecer mercados do Brasil e Argentina. No entanto, a violência armada entre facções é menor que nos departamentos do nordeste.

Villa Hayes, capital departamental, é uma pequena cidade às margens do Rio Paraguai, a ~30 km de Asunción. O restante do departamento é área rural e de comunidades indígenas/menonitas com baixíssima densidade.


- Presidente Hayes é relevante principalmente para quem busca atividades no Chaco (agropecuária extensiva, especialmente pecuária bovina)
- Villa Hayes, por sua proximidade com Asunción, tem caráter de cidade-satélite e oferece bom custo-benefício para residência
- As regiões interiores do Chaco são muito isoladas --- adequadas apenas para empreendedores rurais com estrutura própria
- A rota de tráfico de drogas pelo Chaco representa um risco difuso, mas não afeta diretamente os centros habitados
- Comunidades menonitas (Filadelfia, Loma Plata, Neuland) no Boquerón vizinho são referência de segurança e organização para o Chaco em geral
- Necessário planejamento logístico cuidadoso para instalação no interior do departamento


- https://www.ministeriopublico.gov.py/nota/estadisticas-del-ministerio-publico-la-mayoria-de-los

\subsection{IDH e Desenvolvimento Humano}

\textbf{Fonte principal:} PNUD Paraguai / DGEEC / INE
\textbf{Ano de referência:} 2020 (IDH); 2023 (pobreza e infraestrutura)
\textbf{Atualizado em:} 2026-03-20


\begin{tabularx}{\linewidth}{lX}
\toprule
Indicador & Valor \\
\midrule
IDH & 0,693 \\
Classificação IDH & médio \\
Ranking nacional (1=melhor) & 13/18 \\
Esperança de vida & 71,8 anos \\
Anos de escolaridade média & 7,8 anos \\
RNB per capita (USD PPA) & 7.900 \\
\bottomrule
\end{tabularx}

> Nota: IDH estimado com base no relatório PNUD 2020 e dados da Global Data Lab para a região do Chaco. Presidente Hayes é o maior departamento do Chaco em população, beneficiando-se da proximidade com Asunción (Villa Hayes a 35 km da capital) e expansão da pecuária chaqueña.


\begin{tabularx}{\linewidth}{lX}
\toprule
Indicador & Valor \\
\midrule
Taxa de pobreza total (\%) & 29,4\% \\
Taxa de pobreza extrema (\%) & 7,2\% \\
Índice de Gini & 0,453 \\
\bottomrule
\end{tabularx}

> Fonte: INE --- EPHC 2023 (Presidente Hayes com aumento de pobreza entre 2022--2023). Nota: A EPHC não inclui dados de Boquerón e Alto Paraguay, porém inclui Presidente Hayes como o departamento chaqueño com maior cobertura estatística.


\begin{tabularx}{\linewidth}{lX}
\toprule
Indicador & Valor \\
\midrule
Acesso a água potável (\%) & 71,3\% \\
Acesso a saneamento (\%) & 68,9\% \\
\bottomrule
\end{tabularx}

> Estimativa baseada em dados do Censo 2022. O desafio hídrico é crítico no Chaco --- a água subterrânea é salobra em grande parte da região, exigindo tecnologias de dessalinização. Villa Hayes, próxima ao Rio Paraguai, tem melhor acesso.


Presidente Hayes é o departamento chaqueño mais acessível, com Villa Hayes a apenas 35 km de Asunción cruzando o Rio Paraguai. A proximidade com a capital é a principal vantagem: pode-se viver no Chaco com custo de vida muito baixo e acessar Asunción facilmente. O departamento tem grande expansão de pecuária intensiva, soja e grãos, sendo uma fronteira agrícola em desenvolvimento. A infraestrutura é limitada fora de Villa Hayes e Benjamín Aceval. O desafio maior é hídrico --- água potável (71,3\%) tem cobertura limitada pela natureza salobra das águas subterrâneas chaqueñas. Para brasileiros no agronegócio, o Chaco de Presidente Hayes oferece terras muito baratas com grande potencial. O IDH médio (0,693) reflete os desafios de desenvolvimento da região.


- PNUD Paraguai --- Informe IDH 2001-2020: https://www.undp.org/es/paraguay/noticias/IDH-2001-2020
- INE --- EPHC 2023: https://www.ine.gov.

\subsection{Sistema Prisional}

\textbf{Fonte:} MJT (Ministerio de Justicia y Trabajo) / MNP Paraguay
\textbf{Ano de referência:} 2023--2024
\textbf{Atualizado em:} 2026-03-20


\begin{tabularx}{\linewidth}{lXXXX}
\toprule
Nome & Município & Capacidade & População & Superlotação \\
\midrule
Penitenciaría Regional de Villa Hayes & Villa Hayes & 80 & ~140 & ~175\% \\
\bottomrule
\end{tabularx}


\begin{tabularx}{\linewidth}{lX}
\toprule
Indicador & Valor \\
\midrule
Total de estabelecimentos & 1 \\
Capacidade total (vagas) & ~80 \\
População carcerária total & ~140 \\
Taxa de superlotação & ~175\% \\
\% presos provisórios & ~69\% \\
\bottomrule
\end{tabularx}


Presidente Hayes é o maior departamento do Paraguai em área (72.907 km$^2$), mas um dos menos populosos. Localizado no Chaco paraguaio (região ocidental), tem economia baseada em pecuária bovina intensiva, tânninhas e, mais recentemente, agricultura mecanizada.

A única penitenciaría regional está em Villa Hayes, capital departamental e cidade industrial próxima a Asunción (~30 km pela Ruta 9). O estabelecimento é pequeno e recebe detentos principalmente de Villa Hayes e de comunidades do Chaco Central.

O perfil da criminalidade local inclui crimes contra a propriedade rural, abigeato (roubo de gado) e casos de violência relacionados a disputas de terras. O crime organizado tem presença limitada comparada a departamentos de fronteira com o Brasil.

Detentos do interior profundo do departamento (Benjamín Aceval, Pozo Colorado, Tte. Primero Manuel Irala Fernández) enfrentam longas distâncias para chegar ao presídio de referência, e muitos com penas curtas ou prisão preventiva ficam alojados em delegacias municipais.

\textbf{Para imigrantes:} Presidente Hayes oferece oportunidades únicas no setor pecuário e agroindustrial do Chaco. Villa Hayes tem boa infraestrutura relativa para o Chaco e fica próxima de Asunción. O sistema prisional tem impacto mínimo na segurança geral do departamento. Municípios do interior do Chaco têm infraestrutura limitada mas ambiente tranquilo.


- https://www.mjt.gov.py/index.php/penitenciaria
- https://www.mnp.gov.py/informes
- MNP --- Informe de Visitas 2022


\subsection{Recursos Hídricos Subterrâneos}

\textbf{Fonte:} SENASA / MADES / Geología del Paraguay / Sistema Acuífero Yrenda (SAY)
\textbf{Atualizado em:} 2026-03-20


\begin{tabularx}{\linewidth}{lXXX}
\toprule
Aquífero & Cobertura no Dept. & Profundidade Média & Vazão Típica \\
\midrule
Yrenda-Toba-Tarijeño (SAYTT) --- confinado & ~70\% & 100--300 m & 20--80 m$^3$/h \\
Aluvial (rios Paraguay / Pilcomayo) & Faixas ribeirinhas & 8--25 m & 3--12 m$^3$/h \\
Aquíferos freáticos rasos (Chaco) & ~30\% & 20--60 m & 3--10 m$^3$/h \\
\bottomrule
\end{tabularx}


\begin{tabularx}{\linewidth}{lX}
\toprule
Parâmetro & Valor \\
\midrule
Profundidade média --- SAYTT (m) & 100--300 \\
Profundidade máxima recomendada (m) & 400 \\
Vazão típica (m$^3$/h) & 15--60 \\
Custo estimado perfuração (USD) & 4.000 -- 14.000 \\
Qualidade da água --- superficial (0--100 m) & ruim a imprópria (salina/salobra) \\
Qualidade da água --- profunda (>150 m) & moderada a boa (requer dessalinização em áreas centrais) \\
pH médio & 7,0--8,5 \\
Condutividade elétrica típica (µS/cm) & 500--5.000 \\
Outorga necessária & sim \\
Órgão regulador & MADES --- DGPCRH \\
\bottomrule
\end{tabularx}


Presidente Hayes é o departamento mais próximo de Asunción no Chaco, com a capital Villa Hayes. O Sistema Acuífero Yrenda (SAY) cobre cerca de 2/3 da Região Occidental (Chaco) paraguaia, incluindo este departamento. A salinidade da água subterrânea aumenta de oeste para leste (direção do fluxo), tornando a água em profundidades rasas (0--100 m) frequentemente salobra ou salina e imprópria para consumo humano e animal sem tratamento. Em profundidades superiores a 150--200 m, a qualidade melhora em algumas localidades. O Governo paraguaio instalou dessalinizadoras no Chaco para remediar esse problema, com resultados variados. As faixas ribeirinhas do rio Paraguay têm água superficial de qualidade razoável e maior acessibilidade por captação aluvial. Uso principal: dessedentação animal (gado bovino), abastecimento doméstico com tratamento.


Presidente Hayes é o portão de entrada do Chaco, com fazendas de gado bovino extensivo. A questão hídrica é o maior desafio: água subterrânea é frequentemente salobra, exigindo dessalinização. O custo de perfuração (USD 4.000--14.000) não inclui o sistema de dessalinização (USD 5.000--30.000 adicionais dependendo da capacidade). Para propriedades rurais, é essencial contratar hidrogeólogo para localizar zonas com menor salinidade. Villa Hayes e cidades próximas 

\subsection{Aptidão do Solo}

\textbf{Fonte:} MAG / SENAVE / Geología del Paraguay / FAO
\textbf{Atualizado em:} 2026-03-20


\begin{tabularx}{\linewidth}{lXX}
\toprule
Tipo de Solo & \% do Território & Características \\
\midrule
Cambissolos / Luvissolos & 35\% & Solos de textura siltosa a argilosa; profundos e férteis; carbonatos abaixo de 70 cm; alta capilaridade e risco de salinização \\
Fluvissolos (Aluviais) & 25\% & Solos aluviais das planícies do Rio Paraguai e Rio Pilcomayo; fertilidade moderada a alta; salinidade variável \\
Arenossolos / Regossolos & 25\% & Solos arenosos do interior do Chaco; baixa fertilidade natural; drenagem boa, mas baixa retenção \\
Solonetz / Salinos & 15\% & Solos salino-sódicos em zonas de lençol freático raso; limitações severas para maioria das culturas \\
\bottomrule
\end{tabularx}


\begin{tabularx}{\linewidth}{lXX}
\toprule
Classe & \% Território & Uso Recomendado \\
\midrule
Lavoura intensiva & 10\% & Soja no sul (zona de transição Chaco-Oriental), cultivos tolerantes à salinidade leve \\
Lavoura extensiva & 15\% & Sorgo, algodão, girassol em áreas de menor salinidade \\
Pastagem natural/melhorada & 55\% & Pecuária extensiva --- vocação principal; gramíneas nativas do Chaco \\
Floresta / reserva & 15\% & Bosque nativo do Chaco Úmido; florestamento com espécies nativas \\
Sem aptidão agrícola & 5\% & Áreas salinas severas, áreas inundáveis permanentes \\
\bottomrule
\end{tabularx}


\begin{tabularx}{\linewidth}{lX}
\toprule
Parâmetro & Valor \\
\midrule
pH médio & 6,5 -- 8,5 \\
Fertilidade natural & média-alta (exceto zonas salinas) \\
Teor de matéria orgânica & baixo a médio \\
Necessidade de calcário & não (solos já alcalinos/neutros) \\
\bottomrule
\end{tabularx}


- \textbf{Pecuária bovina extensiva} --- atividade dominante; grandes fazendas (estancias)
- \textbf{Sorgo granífero} --- adaptado às condições semi-áridas
- \textbf{Girassol} --- tolerante à seca e salinidade leve
- \textbf{Algodão} --- histórico na região; ainda cultivado em pequena escala
- \textbf{Soja} --- em expansão no extremo sul do departamento (zona de transição com região oriental)
- \textbf{Floricultura/horticultura} em áreas irrigadas próximas a Villa Hayes


- \textbf{Salinidade e sodicidade:} nas zonas de lençol freático raso, a alta capilaridade provoca ascensão de sais; criticamente limitante em ~15\% do território
- \textbf{Déficit hídrico:} precipitação de 1.000--1.100 mm/ano com estação seca definida; agriculture sem irrigação é arriscada
- **Inundações s

\subsection{Terra Rural}

\textbf{Fonte:} INDERT / Clasificados.com.py / mercado local
\textbf{Ano de referência:} 2024
\textbf{Atualizado em:} 2026-03-20


\begin{tabularx}{\linewidth}{lXXX}
\toprule
Tipo & Preço (USD/ha) & Faixa & Tendência \\
\midrule
Agrícola alta produtividade & 2.500 & 1.500 -- 4.000 & $\uparrow$ \\
Pastagem (campo nativo) & 1.200 & 700 -- 2.000 & $\uparrow$ \\
Terra bruta / arbustiva & 500 & 250 -- 900 & $\rightarrow$ \\
\bottomrule
\end{tabularx}

> Referência de listagem: estabelecimento pecuário em Villa Hayes a USD 1.800/ha; estância de 1.969 ha a USD 5.500/ha (zona mais valorizada); terra bruta Pilcomayo: USD 500--700/ha.


\begin{tabularx}{\linewidth}{lX}
\toprule
Parâmetro & Valor \\
\midrule
Valorização anual média (\%) & 5--8\% \\
Lote mínimo rural (ha) & 20 (zona de chaco central) \\
Imposto territorial rural (\%) & 1\% sobre valor fiscal \\
Liquidez do mercado & média \\
\bottomrule
\end{tabularx}


- \textbf{Villa Hayes (entorno imediato):} Capital departamental na margem oeste do Rio Paraguay; zona de transição entre o Paraguai Oriental e o Chaco; melhor infraestrutura. Terra mais valorizada do departamento.
- \textbf{Margen del río Paraguay:} Faixa estreita leste; terra agricultável com acesso a água.
- \textbf{Área de Influência da Ruta Trans-Chaco (Ruta 9):} Corredor de desenvolvimento; fazendas de gado bem estabelecidas com boa logística.
- \textbf{Zona do Pilcomayo:} Potencial para ecoeturismo e pecuária; terra mais barata mas com desafios de água.


\begin{tabularx}{\linewidth}{lXXX}
\toprule
Indicador & Presidente Hayes (PY) & Mato Grosso (BR) --- Pantanal & Mato Grosso do Sul (BR) --- Pantanal \\
\midrule
Terra pastagem (USD/ha) & 700--2.000 & 5.000--12.000 & 4.000--9.000 \\
Terra bruta (USD/ha) & 250--900 & 2.000--5.000 & 1.500--4.000 \\
Custo relativo & ~15--20\% do MT & referência MT & referência MS \\
\bottomrule
\end{tabularx}

Presidente Hayes oferece terra para pecuária extensiva a custo muito reduzido. Ideal para quem busca grandes extensões para bovinocultura de corte com baixo investimento inicial em terra.


- [INDERT --- Instituto Nacional de Desarrollo Rural y de la Tierra](https://www.indert.gov.py/)
- [InfoCasas --- Villa Hayes](https://www.infocasas.com.py/venta/inmuebles/presidente-hayes/villa-hayes)
- [Asuncion.estate --- Presidente Hayes](https://asuncion.estate/en/presidente-hayes)
- [Gateway to South America --- Farm Paraguay](https://www.gatewaytosouthamerica.com/es/farm-paraguay.php)


\subsection{Mercado Imobiliário}

\textbf{Capital:} Villa Hayes
\textbf{Fonte:} InfoCasas / Yumblin / mercado local
\textbf{Ano de referência:} 2024
\textbf{Atualizado em:} 2026-03-20


\begin{tabularx}{\linewidth}{lXX}
\toprule
Tipo & Preço (USD/m$^2$) & Aluguel (USD/mês) \\
\midrule
Apartamento 2 quartos & 500 -- 800 & 200 -- 380 \\
Casa 3 quartos & 450 -- 750 & 180 -- 340 \\
\bottomrule
\end{tabularx}

> Nota: Listagens de alto padrão (casas de 500--800m$^2$) chegam a USD 625.000--682.000, mas representam segmento premium não representativo do mercado médio.


\begin{tabularx}{\linewidth}{lX}
\toprule
Parâmetro & Valor \\
\midrule
Valorização anual (\%) & 5--7\% \\
Disponibilidade de financiamento & limitada--boa (proximidade com Asunción, atravessando o Rio Paraguay) \\
Bairros mais valorizados & Centro, Barrio Paloma (beira do rio), zonas industriais \\
\bottomrule
\end{tabularx}


\textbf{Vale a pena comprar ou alugar?}
Villa Hayes é a capital do Chaco paraguaio do lado oriental --- fica a apenas 30 km de Asunción em linha reta, mas separada pelo Rio Paraguay. Com a ponte (Puente Héroes del Chaco) o acesso é direto. A cidade tem crescimento impulsionado pela zona industrial do Chaco e pela expansão da fronteira agropecuária.

\textbf{Áreas recomendadas:}
- \textbf{Centro de Villa Hayes:} serviços básicos disponíveis; crescimento recente de comércio e serviços.
- \textbf{Zona Industrial:} área de expansão com empresas do agronegócio e processamento de carnes.

\textbf{Armadilhas comuns:}
- Clima extremo do Chaco (calor intenso no verão, até 45$^\circ$C); verificar qualidade de isolamento e sistemas de ar-condicionado dos imóveis.
- Infraestrutura urbana ainda em desenvolvimento em bairros periféricos.
- Oferta imobiliária formal muito limitada; muitas transações informais.
- Serviços especializados (saúde, educação avançada) requerem deslocamento para Asunción.


- [InfoCasas --- Villa Hayes](https://www.infocasas.com.py/venta/casas/presidente-hayes/villa-hayes)
- [Yumblin --- Villa Hayes](https://www.yumblin.com/us/classifieds/gallery/paraguay/presidente-hayes/villa-hayes)
- [Asuncion.estate --- Presidente Hayes](https://asuncion.estate/en/presidente-hayes)


\subsection{Cobertura Celular}

\textbf{Fonte:} CONATEL / Tigo / Personal / La Nación PY / CONATEL-FSU
\textbf{Ano de referência:} 2024
\textbf{Atualizado em:} 2026-03-20


\begin{tabularx}{\linewidth}{lXXXXX}
\toprule
Operadora & 2G & 3G & 4G & 5G & Cobertura Pop. \\
\midrule
Tigo & 55\% & 38\% & 22\% & --- & 48\% \\
Personal/Claro & 42\% & 28\% & 15\% & --- & 35\% \\
VOX & 18\% & 10\% & 5\% & --- & 15\% \\
\bottomrule
\end{tabularx}

> \textbf{Nota:} Presidente Hayes é o maior departamento do Paraguai Oriental em área e cobre parte do Chaco. Villa Hayes (capital, ~50k hab.) está na margem do rio Paraguai, muito próxima a Asunción, e tem boa cobertura. O Chaco Boreal --- vasto território a oeste --- tem cobertura extremamente limitada, com apenas algumas radiobases instaladas via projeto FSU (Fundo de Serviço Universal) da CONATEL.


\begin{tabularx}{\linewidth}{lX}
\toprule
Parâmetro & Valor \\
\midrule
Melhor operadora rural & Tigo (único com cobertura razoável via FSU) \\
Velocidade 4G média (Mbps) & 8--20 (Villa Hayes) / <5 (interior Chaco) \\
Áreas sem cobertura & 70\%+ do território chaquenho; imensas áreas sem sinal \\
Sinal em propriedade rural & ruim (Chaco) a bom (próximo a Villa Hayes) \\
\bottomrule
\end{tabularx}


\begin{tabularx}{\linewidth}{lXXX}
\toprule
Plano & Operadora & USD/mês & Dados \\
\midrule
Pré-pago básico (7 dias) & Tigo & ~2,50 & 3 GB \\
Pré-pago intermediário (30 dias) & Personal & ~5,00 & 8 GB \\
Pós-pago entrada & Tigo & ~12,00 & 15 GB + chamadas \\
Pós-pago intermediário & Personal/Claro & ~18,00 & 30 GB + chamadas \\
\bottomrule
\end{tabularx}


- \textbf{Villa Hayes:} Cidade muito bem conectada pela proximidade com Asunción (apenas 20 km); cobertura 4G de Tigo e Personal equiparável ao Gran Asunción.
- \textbf{Chaco Boreal (interior):} Situação radicalmente diferente. Torres Tigo instaladas em Pozo Hondo, Estância Salazar e Rojas Silva (via projeto FSU com CONATEL) cobrem apenas raios de 5--10 km com 3G solar.
- \textbf{Radiobases solares:} Tigo instalou torres com painéis solares (30 painéis de 390W, baterias íon-lítio) nas áreas sem energia elétrica do Chaco.
- \textbf{Recomendação crítica para o Chaco:} Starlink é a única solução confiável de conectividade para a maioria das propriedades do Chaco boreal de Presidente Hayes.
- Antes de adquirir propriedade no Chaco, verificar a localização exata em relação às radiobases existentes.


- https://www.lanacion.com.py/negocios\_edicion\_impresa/2021/04/11/tigo-lleva-comunicacion-al-chaco-paraguayo-a-traves-de-la-conatel/

\subsection{Internet}

\textbf{Fonte:} CONATEL / Tigo / INE / Starlink
\textbf{Ano de referência:} 2024
\textbf{Atualizado em:} 2026-03-20


\begin{tabularx}{\linewidth}{lX}
\toprule
Indicador & Valor \\
\midrule
\% população com internet & ~45\% \\
\% com banda larga fixa & ~15\% dos domicílios (concentrado em Villa Hayes) \\
Velocidade média download (Mbps) & 25--55 (Villa Hayes) / <5 (interior Chaco) \\
Velocidade média upload (Mbps) & 8--20 (Villa Hayes) \\
\bottomrule
\end{tabularx}

> O percentual de 45\% representa principalmente Villa Hayes e municípios próximos à capital. A vasta maioria do território chaquenho tem índice de acesso à internet muito inferior --- estimado em menos de 5\% dos domicílios no interior.


\begin{tabularx}{\linewidth}{lXX}
\toprule
Tecnologia & Disponibilidade & Velocidade Típica \\
\midrule
Fibra ótica (FTTH) & Villa Hayes & 50--300 Mbps \\
Cabo HFC & Villa Hayes (limitado) & 20--100 Mbps \\
Rádio/Wireless & Villa Hayes e municípios próximos & 5--40 Mbps \\
ADSL (Copaco) & Villa Hayes + municípios com central & 2--10 Mbps \\
3G Solar (Tigo-FSU) & Pontos isolados no Chaco & 1--5 Mbps \\
Satélite (Starlink) & Todo o departamento & 50--200 Mbps \\
\bottomrule
\end{tabularx}


\begin{tabularx}{\linewidth}{lXXX}
\toprule
Provedor & Tecnologia & Área de cobertura & Preço ref. (USD/mês) \\
\midrule
Tigo & Fibra/HFC/Rádio & Villa Hayes + municípios leste & 18--50 \\
Claro/Personal & Fibra/Rádio & Villa Hayes & 15--45 \\
Copaco & ADSL/Fibra & Villa Hayes + municípios (rede legada) & 10--28 \\
Tigo FSU & 3G solar & Pontos isolados no Chaco & subsidiado/gratuito (pop. indígena) \\
Starlink & Satélite LEO & TODO o departamento & ~44 + equip. ~370 \\
\bottomrule
\end{tabularx}


A situação de internet rural em Presidente Hayes é \textbf{dramaticamente dividida}:

\textbf{Villa Hayes e margem leste:} Boa conectividade, beneficiada pela proximidade com Asunción.

\textbf{Chaco Boreal (maioria do território):}
- Conectividade terrestre praticamente inexistente.
- As radiobases do projeto FSU/CONATEL instaladas por Tigo cobrem apenas pequenos núcleos populacionais isolados com 3G de baixa velocidade.
- \textbf{Starlink é a única solução viável} para fazendas, estâncias e comunidades menonitas no Chaco.
- Comunidades menonitas (ex.: Comunidade Fernheim) foram pioneiras no uso de Starlink no Chaco.
- O custo do equipamento Starlink (~USD 370 + USD 44/mês) torna-se inviável para comunidades indígenas; o FSU subsidia conectividade básica nesses casos.

**Alerta para imigr


\clearpage
\section{Boquerón}
\label{sec:dept_16_Boqueron}

\subsection{Segurança Pública}

\textbf{Fonte principal:} Ministerio Público Paraguay / INE / Prosegur Research
\textbf{Data de referência:} 2019--2024
\textbf{Atualizado em:} 2026-03-20


\begin{tabularx}{\linewidth}{lX}
\toprule
Indicador & Valor \\
\midrule
Taxa de homicídios (por 100k hab) & 0,7--1,7 (2019--2020, fonte INE --- entre as mais baixas do país) \\
Crimes registrados (total/ano) & muito baixo (população esparsa de ~60 mil hab. em vasto território) \\
Número de comisarías & ~5 (estimativa --- território imenso, mas população concentrada em poucas cidades) \\
Índice segurança relativo & alto \\
\bottomrule
\end{tabularx}


Boquerón é o maior departamento do Paraguai em área (91.669 km$^2$) e ao mesmo tempo um dos mais seguros. Com taxa de homicídios de apenas \textbf{0,7 a 1,7 por 100 mil habitantes} (2019--2020), está entre os quatro departamentos mais pacíficos do país. O Ministerio Público registra apenas \textbf{1\% das causas de homicídio doloso} nacional --- proporção mínima.

O departamento é predominantemente habitado por comunidades menonitas (Filadelfia, Loma Plata, Neuland), comunidades indígenas e fazendeiros. As comunidades menonitas têm tradição de organização social rígida, sistema de segurança comunitária próprio e baixíssimas taxas de criminalidade interna.

O Chaco Boreal faz fronteira com a Bolívia e é identificado como rota de aeronaves clandestinas transportando cocaína. No entanto, a densidade policial nas imensas planícies do Chaco é muito baixa e a atuação do crime organizado é logística (transporte), não territorial (conflitos armados por controle de área).

As estradas no Chaco são em sua maioria de terra e frequentemente intransitáveis durante chuvas (especialmente janeiro--março). O isolamento é o principal risco para qualquer residente, mais do que a criminalidade em si.


- Boquerón é uma opção incomum mas interessante para brasileiros com perfil empreendedor no agronegócio
- As cidades menonitas (Filadelfia, Loma Plata) oferecem ambiente excepcionalmente seguro e organizado, com infraestrutura surpreendentemente boa para a região
- A comunidade menonita é economicamente ativa e bem-sucedida --- parcerias de negócios são possíveis
- O isolamento geográfico extremo é o principal desafio: Filadelfia fica a ~500 km de Asunción
- Clima árido a semi-árido com temperaturas extremas (frequentemente acima de 40$^\circ$C no verão)
- Imigrantes que se adaptam ao Chaco geralmente relatam altíssima satisfação com a segurança e tranquilidade

#

\subsection{IDH e Desenvolvimento Humano}

\textbf{Fonte principal:} PNUD Paraguai / DGEEC / INE
\textbf{Ano de referência:} 2020 (IDH); 2022 (infraestrutura, estimativas)
\textbf{Atualizado em:} 2026-03-20


\begin{tabularx}{\linewidth}{lX}
\toprule
Indicador & Valor \\
\midrule
IDH & 0,676 \\
Classificação IDH & médio \\
Ranking nacional (1=melhor) & 16/18 \\
Esperança de vida & 70,8 anos \\
Anos de escolaridade média & 7,1 anos \\
RNB per capita (USD PPA) & 8.600 \\
\bottomrule
\end{tabularx}

> Nota: IDH estimado com base na Global Data Lab e análises do Chaco paraguaio. Boquerón tem perfil peculiar: sua renda per capita é relativamente alta por conta das enormes propriedades menonitas e agropecuárias, mas o IDH é médio por baixa escolaridade e serviços de saúde limitados para a população não-menonita. O RNB per capita reflete a riqueza das cooperativas menonitas.


\begin{tabularx}{\linewidth}{lX}
\toprule
Indicador & Valor \\
\midrule
Taxa de pobreza total (\%) & estimativa 24,0\% \\
Taxa de pobreza extrema (\%) & estimativa 7,5\% \\
Índice de Gini & estimativa 0,485 \\
\bottomrule
\end{tabularx}

> IMPORTANTE: A EPHC (Encuesta Permanente de Hogares Continua) do INE NÃO inclui os departamentos de Boquerón e Alto Paraguay. Estimativas baseadas em dados do Censo 2022 e estudos de pobreza multidimensional disponíveis. O Gini estimado reflete alta desigualdade entre menonitas prósperos e populações indígenas Nivaclé, Ayoreo e Guaraní Ñandeva vulneráveis.


\begin{tabularx}{\linewidth}{lX}
\toprule
Indicador & Valor \\
\midrule
Acesso a água potável (\%) & estimativa 52,3\% \\
Acesso a saneamento (\%) & estimativa 48,7\% \\
\bottomrule
\end{tabularx}

> Estimativa baseada no Censo 2022 e dados de acesso à água no Chaco. O desafio hídrico é extremo --- a maior parte das terras de Boquerón tem água salobra ou insalubre, e o acesso a água potável depende de cisternas e sistemas de coleta de chuva ou dessalinização (especialmente nas cooperativas menonitas).


Boquerón é o maior departamento do Paraguai em área e o coração das cooperativas menonitas --- Filadelfia, Loma Plata e Neuland são cidades prósperas administradas por essas comunidades, com supermercados, hospitais próprios, escolas e indústrias alimentícias. O departamento tem uma das maiores produções leiteiras e de carne bovina do Paraguai. Para imigrantes brasileiros, a realidade é bifurcada: nas cidades menonitas, a qualidade de vida é surpreendentemente boa, com serviços de alto padrão. Fora delas, o Chaco profundo é uma região de extrema precariedade. A falta de d

\subsection{Sistema Prisional}

\textbf{Fonte:} MJT (Ministerio de Justicia y Trabajo) / MNP Paraguay
\textbf{Ano de referência:} 2023--2024
\textbf{Atualizado em:} 2026-03-20


\begin{tabularx}{\linewidth}{lXXXX}
\toprule
Nome & Município & Capacidade & População & Superlotação \\
\midrule
Penitenciaría Regional de Filadelfia & Filadelfia & 50 & ~75 & ~150\% \\
\bottomrule
\end{tabularx}


\begin{tabularx}{\linewidth}{lX}
\toprule
Indicador & Valor \\
\midrule
Total de estabelecimentos & 1 \\
Capacidade total (vagas) & ~50 \\
População carcerária total & ~75 \\
Taxa de superlotação & ~150\% \\
\% presos provisórios & ~67\% \\
\bottomrule
\end{tabularx}


Boquerón é o departamento mais extenso do Paraguai e um dos maiores da América do Sul (91.669 km$^2$), mas com densidade populacional extremamente baixa. A economia baseia-se em pecuária bovina de alto desempenho, a maior do Paraguai em produtividade por hectare, além de agricultura (soja, sorgo) nas colônias menonitas do Chaco Central.

A única penitenciaría regional do departamento está em Filadelfia, capital e maior cidade, sede das tradicionais colônias menonitas (Fernheim, Menno, Neuland). O estabelecimento é o menor do sistema penitenciário paraguaio em número de detentos, refletindo a baixíssima densidade populacional e os baixos índices de criminalidade da região.

O Chaco Central (Filadelfia e arredores) tem índices de violência muito abaixo da média nacional. A comunidade menonita mantém forte coesão social e sistemas próprios de resolução de conflitos, reduzindo a demanda pelo sistema penal formal.

Municípios mais remotos como Mariscal José Félix Estigarribia e General Eugenio A. Garay ficam a centenas de quilômetros de Filadelfia; casos penais dessas regiões frequentemente utilizam delegacias locais ou são transferidos para o presídio de referência de Villa Hayes.

\textbf{Para imigrantes:} Boquerón é um dos ambientes mais seguros e com maior qualidade de vida relativa do interior paraguaio. Filadelfia tem excelente infraestrutura para o Chaco (hospitais, supermercados, escolas bilíngues alemão-espanhol). O setor agroindustrial oferece oportunidades para profissionais qualificados. Requer adaptação ao clima extremo (calor seco, frio noturno no inverno) e ao isolamento geográfico.


- https://www.mjt.gov.py/index.php/penitenciaria
- https://www.mnp.gov.py/informes
- MNP --- Informe de Visitas 2021
- DGEEC --- Anuario Estadístico Paraguay 2022


\subsection{Recursos Hídricos Subterrâneos}

\textbf{Fonte:} SENASA / MADES / Geología del Paraguay / Sistema Acuífero Yrenda (SAY) / IGRAC
\textbf{Atualizado em:} 2026-03-20


\begin{tabularx}{\linewidth}{lXXX}
\toprule
Aquífero & Cobertura no Dept. & Profundidade Média & Vazão Típica \\
\midrule
Yrenda-Toba-Tarijeño (SAYTT) --- confinado & ~95\% & 100--350 m & 20--80 m$^3$/h \\
Aluvial (paleocursos e meandros abandonados) & Localizado & 15--40 m & 3--10 m$^3$/h \\
\bottomrule
\end{tabularx}


\begin{tabularx}{\linewidth}{lX}
\toprule
Parâmetro & Valor \\
\midrule
Profundidade média --- SAYTT (m) & 100--350 \\
Profundidade mínima para água menos salina (m) & >150--200 \\
Profundidade máxima recomendada (m) & 450 \\
Vazão típica (m$^3$/h) & 20--80 \\
Custo estimado perfuração (USD) & 5.000 -- 18.000 \\
Qualidade da água (0--100 m) & ruim a imprópria (salobra a salina) \\
Qualidade da água (>150 m) & moderada (salobra em muitas áreas, requer dessalinização) \\
Condutividade elétrica típica (µS/cm) & 1.000--12.000+ \\
pH médio & 7,5--9,0 \\
Outorga necessária & sim \\
Órgão regulador & MADES --- DGPCRH \\
\bottomrule
\end{tabularx}


Boquerón é o maior departamento do Paraguai e o coração do Chaco Central. O Sistema Acuífero Yrenda (SAY) cobre praticamente todo o território, mas apresenta o maior problema de qualidade de toda a rede hídrica subterrânea paraguaia: salinidade elevada em quase toda a extensão. O departamento é quase inteiramente inserido no SAY. A salinidade aumenta de noroeste (recarga na Bolívia) para sudeste (descarga no rio Paraguay). A condutividade elétrica de poços documentados em Villa Choferes del Chaco atingiu 12.000 mg/L --- água completamente imprópria para consumo humano ou animal sem tratamento intensivo. Em Infante Rivarola (fronteira com a Bolívia), água relativamente doce começa a partir de 120 m com vazões de até 80 m$^3$/h. Em Pozo Hondo, água doce ocorre entre 198 e 212 m. Meandros abandonados de rios (paleocursos) podem conter bolsões de água doce a profundidades menores. O Governo instalou usinas dessalinizadoras em vários pontos do Chaco Central, com resultados técnicos variados e custos operacionais elevados. Uso principal: dessedentação de gado bovino, abastecimento de comunidades menonitas (com tratamento próprio), uso industrial (com tratamento).


Boquerón é o departamento mais desafiador do Paraguai para acesso à água potável. As comunidades menonitas (Filadelfia, Loma Plata, Neuland) desenvolveram ao lon

\subsection{Aptidão do Solo}

\textbf{Fonte:} MAG / SENAVE / Geología del Paraguay / CIFCA-UNA
\textbf{Atualizado em:} 2026-03-20


\begin{tabularx}{\linewidth}{lXX}
\toprule
Tipo de Solo & \% do Território & Características \\
\midrule
Arenossolos / Regossolos & 35\% & Solos arenosos do Chaco Seco Central; drenagem boa mas fertilidade muito baixa; baixa retenção hídrica \\
Cambissolos & 30\% & Solos de textura siltosa-argilosa; moderada fertilidade; camada de carbonatos entre 70-120 cm \\
Luvissolos / Alfissolos & 20\% & Solos com B argílico; moderada a boa fertilidade quando com pH adequado; risco de salinização \\
Solonetz / Salinos-sódicos & 15\% & Solos com alta concentração de sódio trocável; estrutura laminar, impermeáveis; uso muito limitado \\
\bottomrule
\end{tabularx}


\begin{tabularx}{\linewidth}{lXX}
\toprule
Classe & \% Território & Uso Recomendado \\
\midrule
Lavoura intensiva & 15\% & Soja (zona leste --- transição), algodão nas colônias menonitas \\
Lavoura extensiva & 20\% & Sorgo, girassol, amendoim, gergelim \\
Pastagem natural/melhorada & 45\% & Pecuária extensiva e semi-intensiva (principal atividade) \\
Floresta / reserva & 15\% & Bosque do Chaco Seco e Bosque Seco Chiquitano \\
Sem aptidão agrícola & 5\% & Salinas, áreas desnudas, áreas urbanas (Filadelfia) \\
\bottomrule
\end{tabularx}


\begin{tabularx}{\linewidth}{lX}
\toprule
Parâmetro & Valor \\
\midrule
pH médio & 7,0 -- 8,5 \\
Fertilidade natural & baixa a média \\
Teor de matéria orgânica & baixo \\
Necessidade de calcário & não (solos alcalinos; problema oposto) \\
\bottomrule
\end{tabularx}


- \textbf{Pecuária bovina} --- atividade dominante; Boquerón é o maior departamento do Paraguai em rebanho bovino per capita
- \textbf{Soja} no Chaco --- em expansão acelerada; 154.660 ha de soja no Chaco (2024), com Boquerón concentrando a maior parte
- \textbf{Algodão} --- nas colônias menonitas (Filadelfia, Loma Plata, Neuland)
- \textbf{Sorgo granífero} --- resistente à seca
- \textbf{Girassol} e \textbf{gergelim} --- oleaginosas adaptadas às condições do Chaco
- \textbf{Leite} --- colônias menonitas são grandes produtoras de leite para o mercado nacional


- \textbf{Déficit hídrico severo:} precipitação de 500--800 mm/ano no Chaco Central; irrigação essencial para agricultura intensiva
- \textbf{Salinidade e sodicidade:} solos salino-sódicos ocorrem em ~15\% do território; em zonas de lençol raso, salinização por ascensão capilar
- \textbf{Água subterrânea salina:} problema grave; poços rasos têm água inadequada para irri

\subsection{Terra Rural}

\textbf{Fonte:} INDERT / Clasificados.com.py / Everdem / mercado local
\textbf{Ano de referência:} 2024
\textbf{Atualizado em:} 2026-03-20


\begin{tabularx}{\linewidth}{lXXX}
\toprule
Tipo & Preço (USD/ha) & Faixa & Tendência \\
\midrule
Agrícola (com irrigação / alta pluviosidade) & 1.800 & 1.200 -- 3.000 & $\uparrow$ \\
Pastagem estabelecida (pecuária intensiva) & 900 & 600 -- 1.500 & $\uparrow$ \\
Terra bruta / campo nativo & 450 & 200 -- 800 & $\rightarrow$ \\
\bottomrule
\end{tabularx}

> Referências confirmadas 2024: terra virgem Boquerón USD 500--700/ha; estância estabelecida Boquerón USD 800/ha; campo de 11.236 ha a USD 800/ha; propriedade Rio Pilcomayo USD 560/ha; central Chaco estabelecido USD 900/ha; central Chaco virgem USD 450/ha.


\begin{tabularx}{\linewidth}{lX}
\toprule
Parâmetro & Valor \\
\midrule
Valorização anual média (\%) & 5--8\% \\
Lote mínimo rural (ha) & 50 (típico para pecuária viável) \\
Imposto territorial rural (\%) & 1\% sobre valor fiscal \\
Liquidez do mercado & baixa--média \\
\bottomrule
\end{tabularx}


- \textbf{Filadelfia e entorno (Neuland / Fernheim):} Colônias menonitas com a melhor infraestrutura do Chaco paraguaio; pecuária intensiva altamente tecnificada; terra mais cara do departamento.
- \textbf{Loma Plata (Menno):} Centro da maior colônia menonita; cooperativas fortes; boa logística para gado.
- \textbf{Mariscal Estigarribia:} Base militar e polo de desenvolvimento central; terra em apreciação pela expansão da Ruta Trans-Chaco.
- \textbf{Zona do Pilcomayo (Norte):} Terra mais barata; potencial pecuário se houver investimento em água.
- \textbf{Zona Fronteiriça com Bolívia:} Região mais remota; terras a preços mínimos mas com desafios de acesso.


\begin{tabularx}{\linewidth}{lXXX}
\toprule
Indicador & Boquerón (PY) & Mato Grosso do Sul --- Pantanal (BR) & Mato Grosso --- Cerrado (BR) \\
\midrule
Pastagem estabelecida (USD/ha) & 600--1.500 & 4.000--9.000 & 4.500--10.000 \\
Terra bruta (USD/ha) & 200--800 & 1.500--4.000 & 2.000--5.000 \\
Custo relativo & ~12--15\% do MS & referência MS & referência MT \\
\bottomrule
\end{tabularx}

Boquerón oferece as menores relações preço/hectare do Paraguai, mas exige capital para investimento em infraestrutura (cercas, água, pastagens). Para grandes escalas de pecuária extensiva, é um dos últimos fronteirs de expansão barata na América do Sul.


- [INDERT --- Instituto Nacional de Desarrollo Rural y de la Tierra](https://www.indert.gov.py/)
- [Guía para Compra de Campo no Chaco Par

\subsection{Mercado Imobiliário}

\textbf{Capital:} Filadelfia
\textbf{Fonte:} InfoCasas / mercado local
\textbf{Ano de referência:} 2024
\textbf{Atualizado em:} 2026-03-20


\begin{tabularx}{\linewidth}{lXX}
\toprule
Tipo & Preço (USD/m$^2$) & Aluguel (USD/mês) \\
\midrule
Apartamento 2 quartos & 500 -- 800 & 200 -- 380 \\
Casa 3 quartos & 450 -- 750 & 180 -- 340 \\
\bottomrule
\end{tabularx}

> Nota: O mercado de Filadelfia é essencialmente de casas unifamiliares; apartamentos são raros. A comunidade menonita tem padrão construtivo próprio --- bem construídas, funcionais, com grandes quintais. Terrenos em guaranis com faixa de Gs. 96M a Gs. 2,1 bilhões dependendo de tamanho e localização.


\begin{tabularx}{\linewidth}{lX}
\toprule
Parâmetro & Valor \\
\midrule
Valorização anual (\%) & 4--6\% \\
Disponibilidade de financiamento & boa dentro da comunidade (cooperativas menonitas têm estrutura financeira sólida) \\
Bairros mais valorizados & Centro de Filadelfia, Loma Plata, Neuland \\
\bottomrule
\end{tabularx}


\textbf{Vale a pena comprar ou alugar?}
O Chaco paraguaio, e Filadelfia em particular, é um caso único: uma cidade próspera construída por menonitas no meio do Chaco árido. A infraestrutura é surpreendentemente boa --- energia elétrica, água tratada, hospitais modernos, supermercados, escolas de qualidade. Para quem investe em pecuária no Chaco, morar em Filadelfia é a opção mais confortável.

\textbf{Áreas recomendadas:}
- \textbf{Filadelfia:} Melhor infraestrutura do Chaco; hospital Menno de referência regional; supermercados, escolas bilíngues (alemão-espanhol-português).
- \textbf{Loma Plata:} Centro da Cooperativa Chortitzer; excelente organização comunitária.
- \textbf{Neuland:} Menor que as outras, mas comunidade coesa e bem organizada.

\textbf{Armadilhas comuns:}
- A comunidade menonita é fechada; não-membros podem ter dificuldade em adquirir imóveis dentro das colônias (algumas terras são coletivas).
- Clima extremo: verões de 40--45$^\circ$C, invernos frios à noite; casas precisam de ar condicionado de alto rendimento.
- Distância de Asunción: 480 km (5--6 horas de viagem pela Trans-Chaco).
- Serviços especializados médicos de alta complexidade requerem viagem a Asunción.
- Fuera das colônias menonitas, infraestrutura diminui rapidamente.


- [InfoCasas --- Filadelfia Boquerón](https://www.infocasas.com.py/venta/inmuebles/boqueron/filadelfia)
- [InfoCasas --- Imóveis Boquerón baratos](https://www.infocasas.com.py/venta/inmuebles/boqueron/baratos)
- [INDERT](https://www.indert.gov.py/)


\subsection{Cobertura Celular}

\textbf{Fonte:} CONATEL / Tigo / La Nación PY / Última Hora / CONATEL-FSU
\textbf{Ano de referência:} 2024
\textbf{Atualizado em:} 2026-03-20


\begin{tabularx}{\linewidth}{lXXXXX}
\toprule
Operadora & 2G & 3G & 4G & 5G & Cobertura Pop. \\
\midrule
Tigo & 45\% & 28\% & 15\% & --- & 38\% \\
Personal/Claro & 30\% & 18\% & 8\% & --- & 25\% \\
VOX & 12\% & 5\% & 2\% & --- & 10\% \\
\bottomrule
\end{tabularx}

> \textbf{Nota:} Boquerón é o maior departamento do Paraguai em área (~91.669 km$^2$) e tem a cobertura celular mais precária do país. A capital Filadelfia (~25k hab., capital do Mennonato) e comunidades menonitas próximas (Loma Plata, Neuland) têm cobertura razoável. O restante do vastíssimo território --- florestas, estâncias e missões indígenas --- tem cobertura muito limitada ou inexistente.


\begin{tabularx}{\linewidth}{lX}
\toprule
Parâmetro & Valor \\
\midrule
Melhor operadora rural & Tigo (único com cobertura FSU em pontos isolados) \\
Velocidade 4G média (Mbps) & 8--18 (Filadelfia/centros menonitas) / <2 (interior) \\
Áreas sem cobertura & 80\%+ do território; Chaco profundo sem nenhuma cobertura \\
Sinal em propriedade rural & muito ruim --- essencialmente ausente fora dos centros \\
\bottomrule
\end{tabularx}


\begin{tabularx}{\linewidth}{lXXX}
\toprule
Plano & Operadora & USD/mês & Dados \\
\midrule
Pré-pago básico (7 dias) & Tigo & ~2,50 & 3 GB \\
Pré-pago intermediário (30 dias) & Personal & ~5,00 & 8 GB \\
Pós-pago entrada & Tigo & ~12,00 & 15 GB + chamadas \\
\bottomrule
\end{tabularx}


- \textbf{Filadelfia, Loma Plata, Neuland (centros menonitas):} Cobertura 4G Tigo razoável; Personal com presença mais limitada. Comunidades menonitas têm infraestrutura própria bem organizada.
- \textbf{Projeto FSU Chaco:} Tigo instalou 12 radiobases via CONATEL em comunidades como Pozo Hondo, Pelícano, Cristo Rey, San Agustín, San José Estero (Boquerón) e Base Aérea Prats Gill --- cobertura 3G por energia solar, raio de 5--10 km.
- \textbf{Resto do departamento:} Sem cobertura celular. Radiotransmissores HF e Starlink são as únicas opções.
- \textbf{Para imigrantes no Chaco Boquerón:} Starlink é absolutamente essencial para qualquer propriedade fora dos centros menonitas.
- A rota Transchaco (Ruta 9) tem cobertura celular apenas nos trechos próximos a cidades.


- https://www.lanacion.com.py/negocios\_edicion\_impresa/2021/04/11/tigo-lleva-comunicacion-al-chaco-paraguayo-a-traves-de-la-conatel/
- https://www.ultimahora.com/tigo-expande-sus-servicios-el-chaco-n2935449
- https://www.con

\subsection{Internet}

\textbf{Fonte:} CONATEL / Tigo / INE / Starlink / Copaco
\textbf{Ano de referência:} 2024
\textbf{Atualizado em:} 2026-03-20


\begin{tabularx}{\linewidth}{lX}
\toprule
Indicador & Valor \\
\midrule
\% população com internet & ~35\% \\
\% com banda larga fixa & ~20\% dos domicílios (majoritariamente nos centros menonitas) \\
Velocidade média download (Mbps) & 15--50 (centros menonitas) / <3 (interior) \\
Velocidade média upload (Mbps) & 5--15 (centros menonitas) \\
\bottomrule
\end{tabularx}

> As comunidades menonitas (Filadelfia, Loma Plata, Neuland) têm, proporcionalmente, melhor conectividade do que a média do departamento, resultado do modelo cooperativo de gestão dessas comunidades. O restante do departamento tem acesso muito inferior.


\begin{tabularx}{\linewidth}{lXX}
\toprule
Tecnologia & Disponibilidade & Velocidade Típica \\
\midrule
Fibra ótica (FTTH) & Filadelfia, Loma Plata, Neuland & 30--200 Mbps \\
Cabo HFC & Filadelfia (limitado) & 20--80 Mbps \\
Rádio/Wireless & Centros menonitas + arredores & 5--40 Mbps \\
ADSL (Copaco) & Municípios com central telefônica & 1--6 Mbps \\
3G Solar (Tigo-FSU) & Comunidades indígenas específicas & 0,5--3 Mbps \\
Satélite (Starlink) & TODO o departamento & 50--200 Mbps \\
\bottomrule
\end{tabularx}


\begin{tabularx}{\linewidth}{lXXX}
\toprule
Provedor & Tecnologia & Área de cobertura & Preço ref. (USD/mês) \\
\midrule
Tigo & Fibra/Rádio & Centros menonitas + Ruta 9 (parcial) & 18--45 \\
Claro/Personal & Rádio & Filadelfia & 15--40 \\
Copaco & ADSL/Rádio & Centros com central + rota Transchaco & 10--25 \\
ISPs cooperativos & Fibra/Rádio & Comunidades menonitas & preços locais \\
Tigo FSU & 3G solar & Comunidades indígenas & subsidiado \\
Starlink & Satélite LEO & TODO o departamento & ~44 + equip. ~370 \\
\bottomrule
\end{tabularx}


Boquerón representa o \textbf{maior desafio de conectividade rural do Paraguai}:

\textbf{Centros menonitas (Filadelfia, Loma Plata, Neuland):}
- Surpreendentemente bem conectados para o interior chaquenho.
- As cooperativas menonitas (Fernheim, Menno, Neuland) investiram em infraestrutura própria.
- Fibra e rádio wireless com qualidade razoável dentro das colônias.
- Starlink foi adotado amplamente pelas comunidades menonitas desde 2023.

\textbf{Resto do departamento (80\%+ do território):}
- Sem infraestrutura de internet terrestre.
- Estâncias, missões indígenas e fazendas dependem de Starlink.
- Rádio HF (ondas curtas) ainda é usado em fazendas mais remotas como backup.
- A conect


\clearpage
\section{Alto Paraguay}
\label{sec:dept_17_Alto_Paraguay}

\subsection{Segurança Pública}

\textbf{Fonte principal:} Ministerio Público Paraguay / INE / Atlas da Violência CONACYT
\textbf{Data de referência:} 2019--2024
\textbf{Atualizado em:} 2026-03-20


\begin{tabularx}{\linewidth}{lX}
\toprule
Indicador & Valor \\
\midrule
Taxa de homicídios (por 100k hab) & ~2--4 (estimativa, abaixo da média nacional) \\
Crimes registrados (total/ano) & muito baixo (população total ~15 mil hab.) \\
Número de comisarías & ~3 (estimativa --- população mínima em território imenso) \\
Índice segurança relativo & alto \\
\bottomrule
\end{tabularx}


Alto Paraguay é o departamento menos populoso do Paraguai, com aproximadamente 15 mil habitantes distribuídos em um território de ~82.349 km$^2$ (maior que muitos países). O Ministerio Público registra apenas \textbf{1\% das causas de homicídio doloso} nacional --- em números absolutos, os homicídios são raríssimos.

O departamento faz fronteira com a Bolívia ao norte e noroeste, e com o Brasil (Mato Grosso do Sul) ao nordeste. Fuerte Olimpo, capital departamental, é uma pequena cidade ribeirinha às margens do Rio Paraguai. O Chaco Boreal nesta região é praticamente desabitado, com estradas precárias ou inexistentes.

O isolamento extremo é a característica dominante: as chuvas torrenciais tornam muitas áreas inacessíveis por meses. A infraestrutura policial é mínima dado o tamanho do território, mas isso não representa risco prático dado a ausência de população concentrada. Rotas de tráfico aéreo de drogas podem cruzar o território, mas sem impacto na segurança da população local.


- Alto Paraguay não é um destino viável para a maioria dos imigrantes --- extremo isolamento e ausência de infraestrutura
- Fuerte Olimpo é um destino de nicho para ecoturistas e pesquisadores do Pantanal Paraguaio (área do Alto Chaco)
- Potencial para turismo de natureza (fauna Pantaneira, pesca), mas sem infraestrutura comercial desenvolvida
- Brasileiros que considerem a região devem ter total autossuficiência logística
- Segurança em termos de criminalidade urbana é excelente --- mas os riscos são de natureza ambiental (isolamento, clima extremo, ausência de serviços médicos)
- Conexão aérea irregular com Asunción é a principal (às vezes única) opção de acesso


- https://www.ministeriopublico.gov.py/nota/estadisticas-del-ministerio-publico-la-mayoria-de-los-homicidios-dolosos-en-10564
- https://ods.ine.gov.py/ine-main/ods/paz-justicias-e-instituciones-solidas-16/meta-16.1/ind

\subsection{IDH e Desenvolvimento Humano}

\textbf{Fonte principal:} PNUD Paraguai / DGEEC / INE
\textbf{Ano de referência:} 2020 (IDH); 2022--2023 (pobreza e infraestrutura)
\textbf{Atualizado em:} 2026-03-20


\begin{tabularx}{\linewidth}{lX}
\toprule
Indicador & Valor \\
\midrule
IDH & 0,619 \\
Classificação IDH & baixo \\
Ranking nacional (1=melhor) & 18/18 (ou 17/18, disputado com Caazapá) \\
Esperança de vida & 68,4 anos \\
Anos de escolaridade média & 5,8 anos \\
RNB per capita (USD PPA) & 4.900 \\
\bottomrule
\end{tabularx}

> Nota: IDH estimado com base em dados da Global Data Lab e estudos de desenvolvimento do Chaco. Alto Paraguay é o departamento mais remoto e menos desenvolvido do Paraguai --- capital Fuerte Olimpo fica a mais de 500 km de Asunción. A pobreza mais alta do país (confirmada pelo INE 2025) e a falta de cobertura pela EPHC refletem o isolamento extremo.


\begin{tabularx}{\linewidth}{lX}
\toprule
Indicador & Valor \\
\midrule
Taxa de pobreza total (\%) & 38,7\% \\
Taxa de pobreza extrema (\%) & estimativa 14,2\% \\
Índice de Gini & estimativa 0,478 \\
\bottomrule
\end{tabularx}

> Fonte: Alto Paraguay foi confirmado pelo INE (setembro 2025) como o departamento com MAIOR incidência de pobreza do Paraguai (38,7\%). IMPORTANTE: A EPHC não inclui Alto Paraguay --- dados do Mapa de Pobreza 2022 (INE) e estimativas metodológicas. Pobreza extrema é estimativa.


\begin{tabularx}{\linewidth}{lX}
\toprule
Indicador & Valor \\
\midrule
Acesso a água potável (\%) & estimativa 34,5\% \\
Acesso a saneamento (\%) & estimativa 31,2\% \\
\bottomrule
\end{tabularx}

> Estimativas baseadas no Censo 2022 e perfil de remoticidade. Alto Paraguay tem o pior acesso a serviços básicos do país. O Rio Paraguai é a principal via de transporte para grande parte do departamento. Muitas comunidades indígenas (Ayoreo, Chamacoco/Ishir) dependem de fontes naturais não tratadas.


Alto Paraguay é o departamento mais extremo do Paraguai --- geograficamente, economicamente e em termos de desenvolvimento humano. Com IDH estimado em 0,619 (baixo), a maior taxa de pobreza do país (38,7\%) e infraestrutura rudimentar, representa o maior desafio logístico do Paraguai. Fuerte Olimpo, a capital, é uma cidade muito pequena acessível principalmente por via fluvial ou aérea. O departamento tem paisagens espetaculares (Pantanal paraguaio, Gran Chaco Húmedo), importante biodiversidade e potencial ecoturístico ainda inexplorado. Para imigrantes brasileiros, é completamente inadequado como destino de moradia padrão. Eventualmente pode ser considera

\subsection{Sistema Prisional}

\textbf{Fonte:} MJT (Ministerio de Justicia y Trabajo) / MNP Paraguay
\textbf{Ano de referência:} 2023--2024
\textbf{Atualizado em:} 2026-03-20


\begin{tabularx}{\linewidth}{lXXXX}
\toprule
Nome & Município & Capacidade & População & Superlotação \\
\midrule
Delegacia com celas provisórias --- Fuerte Olimpo & Fuerte Olimpo & N/D & ~20--30 & N/A \\
\bottomrule
\end{tabularx}

> \textbf{Nota:} Alto Paraguay não possui penitenciaría regional formal. Detentos com penas longas ou aguardando julgamento são transferidos para o presídio de referência de \textbf{Villa Hayes (Presidente Hayes)} ou para Tacumbú (Asunción). As celas em Fuerte Olimpo funcionam como custódia temporária pela Polícia Nacional.


\begin{tabularx}{\linewidth}{lX}
\toprule
Indicador & Valor \\
\midrule
Total de estabelecimentos penais formais & 0 \\
Estabelecimento de custódia provisória & 1 (delegacia) \\
Capacidade formal (vagas penitenciárias) & N/D \\
População em custódia local & ~20--30 \\
Presídio de referência & Penitenciaría de Villa Hayes \\
\% presos provisórios & ~80\% (estimativa) \\
\bottomrule
\end{tabularx}


Alto Paraguay é o departamento com menor população do Paraguai (estimativa de 13.000--18.000 habitantes) e menor infraestrutura penal. Não existe penitenciaría regional própria --- o departamento depende inteiramente do sistema de custódia provisória em delegacias e da transferência de detentos para outras regiões.

Fuerte Olimpo, a capital departamental, situa-se às margens do rio Paraguay, fronteira com o Brasil (Corumbá/MS) e Bolívia. Apesar do isolamento extremo, a região é rota de tráfico de drogas e contrabando de vida selvagem (animais silvestres) --- dada a escassez de infraestrutura de fiscalização.

O acesso ao departamento é majoritariamente fluvial (pelo rio Paraguay) ou por voo, já que as estradas são precárias e intransitáveis em período de chuvas. Isso torna o transporte de detentos para presídios formais extremamente oneroso e demorado.

\textbf{Para imigrantes:} Alto Paraguay é uma fronteira de extremo isolamento. Oportunidades econômicas existem no ecoturismo, pesca e pecuária extensiva, mas a ausência de infraestrutura básica (hospitais, internet, estradas asfaltadas) torna a instalação permanente muito desafiadora. O ambiente é de baixíssima violência urbana, mas o isolamento representa o principal risco para saúde e segurança.


- https://www.mjt.gov.py/index.php/penitenciaria
- https://www.mnp.gov.py/informes
- MNP --- Informe sobre Establecimiento

\subsection{Recursos Hídricos Subterrâneos}

\textbf{Fonte:} SENASA / MADES / Geología del Paraguay / Sistema Acuífero Yrenda (SAY) / IGRAC
\textbf{Atualizado em:} 2026-03-20


\begin{tabularx}{\linewidth}{lXXX}
\toprule
Aquífero & Cobertura no Dept. & Profundidade Média & Vazão Típica \\
\midrule
Yrenda-Toba-Tarijeño (SAYTT) --- confinado & ~70\% & 100--300 m & 20--80 m$^3$/h \\
Aluvial (rio Paraguay --- planície aluvial) & ~30\% (faixa leste) & 8--30 m & 5--15 m$^3$/h \\
Aquíferos locais (Pantanal) & Porção norte/leste & 10--40 m & 3--10 m$^3$/h \\
\bottomrule
\end{tabularx}


\begin{tabularx}{\linewidth}{lX}
\toprule
Parâmetro & Valor \\
\midrule
Profundidade média --- SAYTT (m) & 100--300 \\
Profundidade mínima para água menos salina (m) & >120--180 \\
Profundidade máxima recomendada (m) & 400 \\
Vazão típica (m$^3$/h) & 10--60 \\
Custo estimado perfuração (USD) & 4.500 -- 16.000 \\
Qualidade da água --- interior do Chaco & ruim a imprópria (salina/salobra) \\
Qualidade da água --- faixa ribeirinha Paraguay & moderada a boa \\
Condutividade elétrica típica --- interior (µS/cm) & 2.000--8.000+ \\
pH médio & 7,5--9,0 \\
Outorga necessária & sim \\
Órgão regulador & MADES --- DGPCRH \\
\bottomrule
\end{tabularx}


Alto Paraguay é o departamento mais setentrional e menos populoso do país. Compreende o extremo norte do Chaco e parte do Pantanal. O Sistema Acuífero Yrenda cobre a maior parte do território chaqueño, com problemas semelhantes a Boquerón: água subterrânea salobra a salina na maioria dos pontos do interior. A faixa ribeirinha do rio Paraguay oferece melhores condições, com captação aluvial de qualidade moderada a boa. Na área do Pantanal (extremo norte, fronteira com Brasil e Bolívia), o regime hidrológico é distinto, com lençol freático próximo à superfície durante as cheias e disponibilidade de água doce sazonal. A capital Fuerte Olimpo é abastecida principalmente por captação superficial do rio Paraguay. A presença indígena (comunidades Ayoreo, Ishir/Chamacoco) é significativa. O 47\% da população indígena do Chaco paraguaio reside nos três departamentos da Região Occidental (Presidente Hayes, Boquerón e Alto Paraguay), frequentemente dependendo de caravanas de distribuição de água.


Alto Paraguay é o departamento mais remoto do Paraguai, com infraestrutura precária e acesso difícil. Para fazendas no interior chaqueño, a questão hídrica é crítica e cara: perfuração (USD 4.500--16.000) mais, na maioria dos casos, sistema de dessaliniza

\subsection{Aptidão do Solo}

\textbf{Fonte:} MAG / SENAVE / Geología del Paraguay / FAO
\textbf{Atualizado em:} 2026-03-20


\begin{tabularx}{\linewidth}{lXX}
\toprule
Tipo de Solo & \% do Território & Características \\
\midrule
Arenossolos / Regossolos (dunas) & 30\% & Solos arenosos profundos do Chaco Seco Norte; baixíssima fertilidade e retenção hídrica \\
Fluvissolos (Aluviais) & 25\% & Solos aluviais das planícies do Rio Paraguay; fertilidade moderada; sazonalmente inundados \\
Cambissolos & 25\% & Solos de textura siltosa-argilosa; moderada fertilidade; camada de carbonatos variável \\
Solonetz / Salino-sódicos & 15\% & Solos com alto teor de sódio; estrutura desfavorável; muito limitados para uso agrícola \\
Histossolos (Orgânicos) & 5\% & Solos orgânicos em banhados permanentes; pantanais do norte \\
\bottomrule
\end{tabularx}


\begin{tabularx}{\linewidth}{lXX}
\toprule
Classe & \% Território & Uso Recomendado \\
\midrule
Lavoura intensiva & 5\% & Agricultura irrigada em áreas pontuais próximas a rios \\
Lavoura extensiva & 10\% & Sorgo, gergelim, amendoim em solos mais férteis \\
Pastagem natural/melhorada & 50\% & Pecuária extensiva de baixíssima lotação --- vocação central \\
Floresta / reserva & 30\% & Pantanal do Norte e Chaco Seco; alto valor ambiental \\
Sem aptidão agrícola & 5\% & Dunas ativas, salinas, banhados permanentes \\
\bottomrule
\end{tabularx}


\begin{tabularx}{\linewidth}{lX}
\toprule
Parâmetro & Valor \\
\midrule
pH médio & 7,0 -- 9,0 \\
Fertilidade natural & muito baixa a baixa \\
Teor de matéria orgânica & muito baixo (exceto Histossolos) \\
Necessidade de calcário & não (solos alcalinos) \\
\bottomrule
\end{tabularx}


- \textbf{Pecuária extensiva} --- única atividade em escala; bovinos de corte em densidades muito baixas (0,2--0,5 cab/ha)
- \textbf{Pesca artesanal} --- economicamente importante nas comunidades ribeirinhas
- \textbf{Extrativismo florestal} --- Quebracho colorado (Schinopsis balansae) para tanino e madeira
- \textbf{Sorgo} e \textbf{gergelim} em agriculturas de subsistência
- \textbf{Ecoturismo} --- Pantanal paraguaio tem potencial crescente


- \textbf{Déficit hídrico extremo:} precipitação de 400--700 mm/ano; estiagens de 6--8 meses frequentes
- \textbf{Salinidade generalizada:} pH alcalino e solos salino-sódicos em grande extensão
- \textbf{Água imprópria para irrigação:} lençol freático raso frequentemente salino; furos profundos necessários
- \textbf{Inundações e secas alternadas:} sistema hidrológico extremo; cheias do Rio Paraguai al

\subsection{Terra Rural}

\textbf{Fonte:} INDERT / DealStream / mercado local
\textbf{Ano de referência:} 2024
\textbf{Atualizado em:} 2026-03-20


\begin{tabularx}{\linewidth}{lXXX}
\toprule
Tipo & Preço (USD/ha) & Faixa & Tendência \\
\midrule
Agrícola (limitada, zonas mais elevadas) & 1.200 & 800 -- 2.000 & $\rightarrow$ \\
Pastagem para pecuária & 500 & 300 -- 900 & $\uparrow$ \\
Terra bruta / pantanosa & 200 & 100 -- 400 & $\rightarrow$ \\
\bottomrule
\end{tabularx}

> Referência confirmada: Alto Paraguay --- terra para pecuária USD 500/ha; terra bruta com potencial hídrico USD 500/ha.


\begin{tabularx}{\linewidth}{lX}
\toprule
Parâmetro & Valor \\
\midrule
Valorização anual média (\%) & 3--5\% \\
Lote mínimo rural (ha) & 100 (viável apenas em grandes escalas) \\
Imposto territorial rural (\%) & 1\% sobre valor fiscal \\
Liquidez do mercado & muito baixa \\
\bottomrule
\end{tabularx}


- \textbf{Fuerte Olimpo (entorno):} Capital departamental no Rio Paraguay; únicas terras com alguma infraestrutura básica.
- \textbf{Bahía Negra:} Fronteira com Bolivia e Brazil (Corumbá/Pantanal); zona remota com fazendas de gado extensivo.
- \textbf{Puerto Casado / Puerto Pinasco:} Cidades históricas às margens do Rio Paraguay; acesso fluvial; fazendas estabelecidas de gado.
- \textbf{Zona Pantanal paraguaio:} Terras alagáveis de baixíssimo custo; apenas para operações de pecuária extensiva com experiência no ambiente.


\begin{tabularx}{\linewidth}{lXXX}
\toprule
Indicador & Alto Paraguay (PY) & Mato Grosso --- Pantanal Norte (BR) & Pará --- Amazônia (BR) \\
\midrule
Pastagem (USD/ha) & 300--900 & 3.000--7.000 & 1.500--4.000 \\
Terra bruta (USD/ha) & 100--400 & 800--2.500 & 500--1.500 \\
Custo relativo & ~10\% do MT & referência MT & referência PA \\
\bottomrule
\end{tabularx}

Alto Paraguay é o departamento mais remoto e menos desenvolvido do Paraguai. Terra baratíssima, mas desafios logísticos imensos --- acesso quase exclusivamente fluvial, clima extremo, infraestrutura mínima. Para pecuaristas experientes com capital de giro, pode ser uma oportunidade especulativa de longo prazo.


- [INDERT --- Instituto Nacional de Desarrollo Rural y de la Tierra](https://www.indert.gov.py/)
- [DealStream --- Agricultural Land Alto Paraguay](https://dealstream.com/land-sale/agricultural/departamento-de-alto-paraguay-paraguay)
- [Plan B Paraguay --- Farmland Prices](https://x.com/planbparaguay/status/1834232169759211739)
- [Gateway to South America](https://www.gatewaytosouthamerica.com/es/farm-paraguay.php)


\subsection{Mercado Imobiliário}

\textbf{Capital:} Fuerte Olimpo
\textbf{Fonte:} INDERT / mercado local
\textbf{Ano de referência:} 2024
\textbf{Atualizado em:} 2026-03-20


\begin{tabularx}{\linewidth}{lXX}
\toprule
Tipo & Preço (USD/m$^2$) & Aluguel (USD/mês) \\
\midrule
Apartamento 2 quartos & 250 -- 450 & 100 -- 220 \\
Casa 3 quartos & 200 -- 380 & 90 -- 190 \\
\bottomrule
\end{tabularx}

> Nota: O mercado imobiliário formal em Fuerte Olimpo é praticamente inexistente. Pouquíssimas propriedades são ofertadas formalmente. Os valores acima são estimativas baseadas em padrão relativo ao resto do país e às poucas transações registradas.


\begin{tabularx}{\linewidth}{lX}
\toprule
Parâmetro & Valor \\
\midrule
Valorização anual (\%) & 2--4\% \\
Disponibilidade de financiamento & muito limitada (praticamente inexistente no local) \\
Bairros mais valorizados & Centro de Fuerte Olimpo (única zona urbana do departamento) \\
\bottomrule
\end{tabularx}


\textbf{Vale a pena comprar ou alugar?}
Alto Paraguay NÃO é recomendado para imigrantes em busca de qualidade de vida urbana. Fuerte Olimpo é uma cidade muito pequena (menos de 10.000 hab.) de difícil acesso, sem infraestrutura de saúde especializada, sem escolas de qualidade superior, com acesso quase exclusivamente fluvial. A única razão para morar aqui é operar grandes fazendas de pecuária extensiva na região.

\textbf{Áreas recomendadas (para quem precisa estar na região):}
- \textbf{Fuerte Olimpo:} única opção com serviços mínimos.
- \textbf{Puerto Casado:} cidade histórica com alguma infraestrutura remanescente; fazendas na região.

\textbf{Armadilhas comuns:}
- Acesso ao departamento é principalmente fluvial pelo Rio Paraguay; estradas são precárias ou inexistentes em grande parte do território.
- Calor extremo (40$^\circ$C+) grande parte do ano, com umidade elevada na época das chuvas.
- Mosquitos, cobras e fauna selvagem são realidade cotidiana.
- Serviços médicos básicos apenas; qualquer emergência exige evacuação por barco ou avião.
- Energia elétrica irregular; geradores são necessidade básica.
- Internet/telefonia extremamente limitados ou inexistentes em grande parte do departamento.


- [DealStream --- Alto Paraguay](https://dealstream.com/land-sale/agricultural/departamento-de-alto-paraguay-paraguay)
- [INDERT](https://www.indert.gov.py/)
- [InfoCasas --- Paraguay](https://www.infocasas.com.py/venta/inmuebles)


\subsection{Cobertura Celular}

\textbf{Fonte:} CONATEL / Tigo / La Nación PY / Copaco / CONATEL-FSU
\textbf{Ano de referência:} 2024
\textbf{Atualizado em:} 2026-03-20


\begin{tabularx}{\linewidth}{lXXXXX}
\toprule
Operadora & 2G & 3G & 4G & 5G & Cobertura Pop. \\
\midrule
Tigo & 35\% & 18\% & 8\% & --- & 28\% \\
Personal/Claro & 22\% & 10\% & 4\% & --- & 18\% \\
VOX & 8\% & 3\% & 1\% & --- & 6\% \\
\bottomrule
\end{tabularx}

> \textbf{Nota:} Alto Paraguay é o departamento mais isolado e com menor cobertura telecomunicacional do Paraguai --- e possivelmente um dos mais desconectados da América do Sul. Com apenas ~14.000 habitantes em 82.349 km$^2$ (menor densidade demográfica do país), Fuerte Olimpo (capital, ~5.000 hab.) está na margem do rio Paraguai e tem cobertura básica via rádio e torres Tigo. O resto do território --- Chaco do norte --- é praticamente sem cobertura.


\begin{tabularx}{\linewidth}{lX}
\toprule
Parâmetro & Valor \\
\midrule
Melhor operadora & Tigo (único com presença relevante) \\
Velocidade 4G média (Mbps) & 3--10 (Fuerte Olimpo) / sem dados (interior) \\
Áreas sem cobertura & 90\%+ do território --- Chaco Norte, Pantanal Norte \\
Sinal em propriedade rural & inexistente na maioria das áreas \\
\bottomrule
\end{tabularx}


\begin{tabularx}{\linewidth}{lXXX}
\toprule
Plano & Operadora & USD/mês & Dados \\
\midrule
Pré-pago básico (7 dias) & Tigo & ~2,50 & 3 GB \\
Pré-pago intermediário (30 dias) & Tigo & ~5,00 & 8 GB \\
\bottomrule
\end{tabularx}

> Presença de Personal e VOX é mínima e não confiável fora de Fuerte Olimpo.


- \textbf{Alto Paraguay não é destino típico de imigrantes} devido ao isolamento extremo, clima rigoroso e falta de infraestrutura.
- \textbf{Fuerte Olimpo:} Única cidade com alguma cobertura celular (Tigo 3G/4G precário); acessada principalmente por barco ou avião pequeno.
- \textbf{Bahía Negra (fronteira com Brasil/Bolívia):} Cobertura celular praticamente inexistente; chips brasileiros podem captar sinal de torres do MS ou Corumbá em condições específicas.
- \textbf{Projeto FSU:} CONATEL/Tigo instalaram torres em Alto Paraguay como parte do programa de conectividade universal, mas cobertura ainda é mínima.
- \textbf{Para qualquer propriedade em Alto Paraguay:} Starlink é a ÚNICA opção viável de comunicação e internet.
- Rádio HF (ondas curtas) e comunicação via satélite Inmarsat/Iridium ainda são usados em expedições e fazendas remotas.


- https://www.lanacion.com.py/negocios\_edicion\_impresa/2021/04/11/tigo-lleva-comunicacion-al-chaco-paraguayo-a-traves-de-la-conatel/


\subsection{Internet}

\textbf{Fonte:} CONATEL / Copaco / Starlink / INE
\textbf{Ano de referência:} 2024
\textbf{Atualizado em:} 2026-03-20


\begin{tabularx}{\linewidth}{lX}
\toprule
Indicador & Valor \\
\midrule
\% população com internet & ~20\% \\
\% com banda larga fixa & <5\% dos domicílios \\
Velocidade média download (Mbps) & 3--15 (Fuerte Olimpo) / N/A (interior) \\
Velocidade média upload (Mbps) & 1--5 (Fuerte Olimpo) \\
\bottomrule
\end{tabularx}

> Alto Paraguay tem a menor penetração de internet do Paraguai. Os ~20\% representam quase exclusivamente moradores urbanos de Fuerte Olimpo e Bahía Negra.


\begin{tabularx}{\linewidth}{lXX}
\toprule
Tecnologia & Disponibilidade & Velocidade Típica \\
\midrule
Fibra ótica (FTTH) & Não disponível & --- \\
Cabo HFC & Não disponível & --- \\
Rádio/Wireless & Fuerte Olimpo (limitado) & 2--15 Mbps \\
ADSL (Copaco) & Fuerte Olimpo (muito limitado) & 0,5--3 Mbps \\
3G Solar (Tigo-FSU) & Algumas comunidades isoladas & 0,5--2 Mbps \\
Satélite (Starlink) & TODO o departamento & 50--200 Mbps \\
Satélite convencional & Disponível (VSAT) & 5--25 Mbps (alto custo) \\
\bottomrule
\end{tabularx}


\begin{tabularx}{\linewidth}{lXXX}
\toprule
Provedor & Tecnologia & Área de cobertura & Preço ref. (USD/mês) \\
\midrule
Copaco & ADSL/Rádio & Fuerte Olimpo & 10--20 \\
Tigo (FSU) & 3G solar & Comunidades indígenas & subsidiado \\
Starlink & Satélite LEO & TODO o departamento & ~44 + equip. ~370 \\
VSAT (outros) & Satélite geoestacionário & Todo (com instalação técnica) & 80--200+ \\
\bottomrule
\end{tabularx}


Alto Paraguay é o \textbf{caso mais extremo de exclusão digital do Paraguai}:

- Praticamente não existe infraestrutura de internet fixa fora de Fuerte Olimpo.
- O projeto FSU da CONATEL instalou conectividade básica em comunidades indígenas selecionadas (Ayoreo, Chamacoco/Ishir), mas com velocidades mínimas.
- \textbf{Starlink transformou radicalmente a viabilidade de vida no Alto Paraguay} a partir de 2023 --- pela primeira vez, qualquer ponto do departamento pode ter 50--200 Mbps.
- Fazendas e estâncias na região dependem historicamente de satélite geoestacionário (VSAT), que custa 3--5x mais que o Starlink com velocidades menores.

\textbf{Contexto para imigrantes:}
Alto Paraguay não é recomendado como destino de imigração para uso geral. É destino para quem busca isolamento extremo, pecuária extensiva de escala muito grande ou ecoturismo de aventura. A conectividade com Starlink é suficiente para trabalho remoto básico, mas a logística de tud
